% Framework source integration v3.024 (audit v1.8): sharp nearest rooted Hopf source theorem; no physical selector promoted.
% Paper5B-TerminalSpectrum-v5_200: polyhedral full-2O working carrier and exact 2I meson/baryon geometry; physical occurrence remains conditional.
% Paper-series v2.381: complete algebraic exceptional torsion atlas and all-cocycle twist-stability; original claims retained.
% Paper-series v2.380: unified bibliography, exact exceptional double/length/gauge mathematics in I-B; previous theorem labels retained.
% Paper-series v2.379: full equivariant Dirac/Hodge/eta/torsion atlas in Paper I-B; paper-root references synchronized.
% Paper-series v2.378: 2026 exceptional literature consolidation and reference synchronization. The mathematical additions are located in Paper I-B.
% Paper-series v2.377: reference-only rebind to current Paper I-B v1.1; original mathematics retained except new Paper I-B comparison chapter.
% Framework v2.370: derived exceptional-root/Hurwitz comparison in I-A/I-C; other papers reference-only synchronized.
% Framework v2.368: preserve prior II-B convention/unital fixes and verified II-D corrections; current references synchronized.
% Framework v2.366: exact chamber--transporter/Hurwitz integration; physical action and hadron decay gates remain conditional.
% Framework v2.366: scoped Hurwitz/600-cell hadron-source integration; physical APBI/decay frontier remains open.
% Framework v2.358: current 30-root source reference synchronization; selectively updated III-E, III-F, VI-B, VI-C.
% Framework v2.356: bibliography/source-registry synchronization; mathematical body unchanged.
% Framework v2.348: Paper V audit, terminal observable scope, canonical Gram eigenvalues and exact weak-memory variance.
% Framework v2.344: Twisted Hurwitz C3/Petersen bridge, native baryon source and same-action curvature firewall.
% Framework v2.343: Reciprocal signed coupling, selected 15 and conditional physical boundary.
% Framework v2.342: Exact three-cycle Petersen reconstruction, class frames, typed nonactualization.
% Framework v2.341: Coherent coset--Hodge and parity cover integration. Previous v2.340 retained.
% Framework v2.340: Exact rank-two 600-cell incidence and typed operator/observable closure.
% Framework v2.339: Petersen--endomorphism isometry, incidence projector and typed physical-proof firewall.
% Framework v2.338: exact Petersen 120-carrier, oriented two-source and source-ownership results. No hadron mass prediction.
% Framework v2.321: Paper V Petersen rooted 3+12 induction, exact sector compression, conjugation orbit refinement and source-incidence firewall. Original source theorem bodies preserved.
% Framework v2.275: exact positive dense-twin word cubature, finite repeatability and eighth-moment classification, nested sampling frames, joint action and POVM-jet exports; preserves v2.274.
% Framework v2.274: two-scalar positive convolution certificates, first-leaked exceptional coset detection, branch-symmetric twin mixing, and finite endpoint moments; all v2.273 branches preserved.
% Framework v2.273: restored missing Problem 3 theorems from v2.268; exact exceptional frames, twin geometry, convolution repeatability and projector/occurrence firewalls; merged on v2.272 without overwriting independent branches.
% Framework v2.272: merged v2.270 covariant native-readout mathematics with v2.271 Problem 4 finite-action Hopf hypercharge/holonomy corrections; no universal bulk-gauge uniqueness claim.
% Independent branches: v2.270 readout/guard exports and v2.270 K16 action-consumed charge closure; all prior scoped assumptions retained.
% Framework v2.267: exact Paper II-A converse/observable certificates, III-F action tests, V-B right orbit and V-C native gauge scope; prior v2.266 preserved.
% Framework v2.266: exact Paper II-C outcome-count/readout converse results and scoped downstream imports; 30-root references synchronized.
% Framework v2.265: II-C coordinate-stage constitution and exact first-symbol host inheritance; prior mathematical domains retained.
% Framework v2.264: II-B reverse and action-bridge integration; preserves prior mathematical domains; active 30-root references synchronized.
% Framework v2.263: Paper I-D adds bridge-controlled root-carry parity and source-relative old-cap cancellation; other root mathematical bodies retained; active 30-root references synchronized.
% Framework v2.262: merges Paper II-F three-gap closure from v2.260 into v2.261, preserving newer II-B/II-D repairs; 30-root references synchronized.
% Framework v2.261: merged Paper II-B GP2-level/geometric repairs; preserves v2.260 II-D repairs; active references synchronized.
% Framework v2.260: Paper II-D signed RS, finite-symbol typing, and conditional curvature/Jordan repairs; references synchronized.
% Framework v2.259: conditional native SU(2) atlas and explicit later discharge; v2.258 mathematical bodies retained elsewhere.
% Framework v2.258: preserves parallel II-A/dynamics update and adds complete readout/frequency closures.
% Framework v2.257: exact minimal complete position readout and arbitrary-preparation frequency law.
% Framework v2.256: exact native inertia and sharp bounded ordered-source closure; preserves concurrent v2.255 work.
% Framework v2.255: preserves parallel detector integration and adds exact finite readout closures.
% Framework v2.254: shared bibliography/reference sync; substantive Pasterski detector-Ward additions only IV-B/IV-C/VII-C.
% Framework v2.253: typed raw coordinate defect in I-A and synchronized VI-A import; active references synchronized.
% Framework v2.252: optional shell-frame source-return calculation in III-G; active references synchronized.
% Framework v2.251: forward common-action execution, joint source jets and scoped coupled control.
% Framework v2.251: exact full-heat variations, fermion source jets and explicit physical completion frontier.
% Framework v2.251: certified interacting source jets, pole stability and scoped completion proofs.
% Framework v2.251: interacting finite-source evaluation and coupled local vacuum classification.
% Framework v2.251: exact finite source execution and native causal common-parent bridge.
% Framework v2.251: constructive parent and scoped source-incidence audit; see release report.
% Framework v2.251: common-action two-jet, corrected source/response closure and synchronized references.
% Framework v2.244: scoped rooted source-order Schur diagnostics and synchronized paper references.
% Framework v2.242: signed rooted phase lift, complete transverse return and exact family cubic.
% Framework v2.241: proved rooted conormal geometry, conditional cap realization and scoped action tests.
% Framework v2.240: observable-minimal action and response certificates integrated; inherited v2.239 geometry retained.
% Framework v2.234: restructured main physics owner, technical proofs in V-D.
\documentclass[12pt,a4paper]{article}

\usepackage{mathtools}
\usepackage{graphicx}
\usepackage{float}
\usepackage{caption}
\usepackage{microtype}
\PassOptionsToPackage{hidelinks}{hyperref}

% Canonical series preamble.
% Paper-series v2.381: new I-B torsion/modular results, bibliography and current source roots.
% Paper-series v2.380: new exceptional bibliography and synchronized paper references.
% Paper-series v2.379 current 31-paper macro/aux registry.
% Paper-series v2.378: current 31-paper AUX/bibliography registry updated for Paper I-B v1.2 literature additions.
% Framework v2.371: current-source reference registry refreshed; old theorem labels retained.
% Framework v2.370: current 30-paper source/AUX root rebinding; native mathematical macros unchanged.
% Framework v2.368: synchronized merged 30-root registry; mathematical notation unchanged.
% Framework v2.358: active 30-paper source registry synchronized; existing mathematical macros unchanged.
% Framework v2.357: selective new source-owned certificates; active reference roots updated for five revised papers.
% Framework v2.356: active 30-root registry synchronized to the 30 papers in this release.
% Framework v2.275: exact positive dense-twin word cubature, finite repeatability and eighth-moment classification, nested sampling frames, joint action and POVM-jet exports; preserves v2.274.
% Framework v2.274: filename and root-source synchronization.
% Framework v2.273: restored missing Problem 3 theorems from v2.268; exact exceptional frames, twin geometry, convolution repeatability and projector/occurrence firewalls; merged on v2.272 without overwriting independent branches.
% Framework v2.272: merged v2.270 covariant native-readout mathematics with v2.271 Problem 4 finite-action Hopf hypercharge/holonomy corrections; no universal bulk-gauge uniqueness claim.
% Independent branches: v2.270 readout/guard exports and v2.270 K16 action-consumed charge closure; all prior scoped assumptions retained.
% Framework v2.269: synchronized 30 active AUX roots; mathematics and bibliography preserved from Macros-v2_161.
% Framework v2.266: Paper II-C converse-derived outcome resources; II-E, III-C, IV-A/B, V-A and VI-B downstream mathematical closures; 30-root registry synchronized.
% Framework v2.265: II-C coordinate-stage constitution and exact first-symbol host inheritance; prior mathematical domains retained.
% Framework v2.264: paper II-B reverse/audit and IV-A observable-minimal exchange; synchronized external references to 30 active roots.
% Framework v2.263: Paper I-C adds bridge-controlled root-carry parity and source-relative old-cap cancellation; other root mathematical bodies retained; active 30-root references synchronized.
% Framework v2.262: integrates Paper II-F radius, Casimir and readout results; active 30-root references synchronized.
% Framework v2.261: merged Paper II-B GP2-level/geometric repairs; preserves v2.260 II-D repairs; active references synchronized.
% Framework v2.260: Paper II-D signed RS, finite-symbol typing, and conditional curvature/Jordan repairs; references synchronized.
% Framework v2.259: conditional native SU(2) atlas and explicit later discharge; v2.258 mathematical bodies retained elsewhere.
% Framework v2.258: preserves parallel II-A/dynamics update and adds complete readout/frequency closures.
% Framework v2.257: exact minimal complete position readout and arbitrary-preparation frequency law.
% Framework v2.256: exact native inertia and sharp bounded ordered-source closure; preserves concurrent v2.255 work.
% Framework v2.255: preserves parallel detector integration and adds exact finite readout closures.
% Framework v2.254: Pasterski bibliography and synchronized 30-paper external-reference registry.
% Framework v2.253: typed raw coordinate defect in I-A and synchronized VI-A import; active references synchronized.
% Framework v2.252: synchronized source registry; optional III-G shell-frame calculation.
% Framework v2.251: forward common-action execution, joint source jets and scoped coupled control.
% Framework v2.251: exact full-heat variations, fermion source jets and explicit physical completion frontier.
% Framework v2.251: certified interacting source jets, pole stability and scoped completion proofs.
% Framework v2.251: synchronized active source registry; mathematical notation unchanged.
% Framework v2.251: synchronized active source registry; mathematical notation unchanged.
% Framework v2.251: synchronized active source registry; mathematical notation unchanged.
% Framework v2.251: synchronized source registry; no mathematical notation changes.
% Framework v2.244: synchronized 30-paper AUX registry and rooted phase-source action-order tests.
% Framework v2.242: signed rooted phase lift, complete transverse return and exact family cubic.
% Framework v2.241: proved rooted conormal geometry, conditional cap realization and scoped action tests.
% Framework v2.239: all 30 active source roots synchronized; no mathematical macro changes.
% Framework v2.238: Paper IV A/B/C split and 30-root reference registry.
% Framework v2.237: Paper VII-D technical split and 28-root reference registry.
% Framework v2.236: Paper VI-C technical-volume reference added; all current root bindings advanced.
% Framework v2.235: 26-root AUX and bibliography bindings synchronized; mathematical macros unchanged.
% BEGIN INLINE SOURCE: Paper00-TeXBuildCompatibility-v1_0.tex
% Rendering-only repairs for inherited long appendices and bitmap text fonts.
% Integer counters through 26 retain the original labels; later values use AA,
% AB, ... . This file changes no action coefficient, state or theorem premise.
\RequirePackage{alphalph}
\makeatletter
\let\BHPrimitiveAlph\@Alph
\newalphalph{\BHExtendedAlph}{\BHPrimitiveAlph}{26}
\let\@Alph\BHExtendedAlph
\AtBeginDocument{%
  \@ifpackageloaded{microtype}{\microtypesetup{expansion=false}}{}%
}
\makeatother

% END INLINE SOURCE: Paper00-TeXBuildCompatibility-v1_0.tex

% Framework v2.164: current reference/bibliography profile; mathematical definitions retained.
% Framework v2.162: finite metric, functional Hilbert and scoped Ward integration.
% Framework v2.161: inherits Macros-v2_57.tex exactly except the current external-reference binding.
% Macros v2.57 / framework v2.156: current reference registry only; mathematical definitions unchanged.
% Framework v2.143 / Macros v2.55: current external-reference bindings; mathematical notation unchanged.
% Framework v2.110 / Macros v2.54: current external-reference registry and bibliography; mathematical notation inherited.
% Framework v2.67 / Macros-v2_53: adds Born-action-cell and joint-EH readout symbols; absolute unit realization is metrological; MATCH_{J/scal}^{EH} remains the sole G24 normalization bridge.
% Framework v2.35 / Macros-v2_52: grouped paper numbering, active AUX registry, and Born-v6_20 bibliography binding synchronized; notation unchanged.
% Framework v2.26: active-source AUX and bibliography bindings synchronized; notation unchanged.
% Framework v2.09: reviewed early-reciprocity integration; scoped physical imports.
% Framework v2.08: Born--Euler mathematics and conditional complete-action potential criterion.
% Framework v2.07: active AUX registry includes Paper I-B v1.2; shared bibliography Born-v6_18.
% Macros-v2_46 / framework v2.04: mathematical atlas and downstream realization correction.
% Macros-v2_46 / framework v2.03: scoped structure integration and current dependency bindings.
% Macros-v2_46 / framework v2.01: shared source bindings; prior revision comments are historical.
% v2.42 / framework v1.99: Synchronized RG/AS owner and bibliography bindings.
% v2.41 / framework v1.99: active manuscript AUX bindings and geometric-closure bibliography synchronized.
% Release-v1.80 centralized source synchronization for scalar-circle/comparison-band integration.
% Release v1.79: Hurwitz hierarchy and kinetic-normalized action response integration; physical mass matching remains conditional.
% Release v1.78: approved M8 closure-batch integration; physical and analytic domains retained.
% Current v1.78 registry: all active TeX roots are synchronized below.
% Earlier release comments are historical scope records, not current status.
% Release-v1.77 synchronization: native same-owner Noether closure, generator--flow/history typing, analytic-certification firewalls, scalaron/K16, LF/QCD, quantum-Ward, and RAR boundary-owner transport are registered; shared mathematical notation is unchanged.
% Release-v1.76 synchronization: registry points to the v1.76 active paper roots.
% Release-v1.76 synchronization: registers the local A+Phi+g Noether/constraint closure and dynamics-frontier propagation; shared mathematical notation is unchanged.
% Release-v1.75 synchronization: registry points to the v1.75 active paper roots.
% Release-v1.75 synchronization: integrates the joint-Dirac Einstein bulk real form and the local odd-scalar heat-coefficient action on explicit analytic first-jet/interior domains; v1.74 Yang--Mills closure is retained. Fermionic reality, pure-metric curvature-squared/exact-nonlocal terms, noncanonical boundaries, physical vacuum selection, and actual solution-history tokens remain open.
% Release-v1.74 synchronization: registry points to the v1.74 active paper roots and retains the v1.73 theorem labels.
% Release-v1.74 synchronization: canonical Macros-v2_34.tex; generated local nonlinear Yang--Mills action, constrained connection-field phase space, source-free Gauss momentum map, and Legendre closure are integrated on their proved domains. The remaining non-gauge ACT_cl sectors and actual solution-history selection remain open.
% Release-v1.71 synchronization: closes the former JOINT_x/JOINT_5
% connection-occurrence gap by canonical structure-group extension of the
% nonflat Born connection, faithful associated representation, and the local
% cotangent--Clifford operator constructor.  Actual fluctuation members,
% actions, vacua, solutions, calibrations, and analytic QFT domains remain
% separately typed.  The active-AUX registry is advanced once.
% Release-v1.71 canonical macro and active-AUX registry.
% ======================================================================
% Macros-v2_39.tex -- canonical publication macro file
% Release v1.71 retains the v1.70 null-ray/QCD boundaries and additionally
% removes the duplicate structural connection actuality premise.  It does not
% claim global bundle topology, select a gauge/Higgs member, or promote the
% finite construction to an interacting continuum theory.
% Release v1.68 constructs the finite light-front total-fraction operators and
% parton number/momentum operator measures, proves their simplex and Hausdorff-limit
% theorems, and separates positive cut measures from scheme-renormalized
% light-ray distributions.  It also separates finite real Schur poles, stable
% norm-resolvent/Riesz pole-residue limits, and continued resonance
% pole/residue limits.  The Paper-VII-C scattering package is claim-indexed:
% local QFT and spectral claims import no S-matrix or infrared premise, while
% asymptotic amplitudes, LSZ, S-matrix poles, completeness, and cross sections
% import only their named massive or massless/long-range clauses.
% Release v1.67 factorizes BND48_H into a same-action Hessian/Schur
% descendant and still-open typed boundary-trace, critical-lift, and
% refinement-limit conditions.  It proves the finite-collar Dirichlet-to-
% Neumann law sqrt(y)tanh(L sqrt(y)), its nonuniform critical limit, and the
% finite-return ultra-infrared turnover.  The Paper-I-A Hardy constant mode is
% an exact geometric near-hit, not yet the physical metric boundary mode.
% The current local f(R) Hessian remains massive and therefore still refutes
% a RAR prediction on the minimal action branch.
% Release v1.66 derives the transverse depth-six normalizer bridge from the
% already generated C6 Reynolds image and Paper-I-A binary-dihedral normalizer.
% The bridge is the unique least 2D3-stable scalar-containing realization of
% the faithful quotient 2D3/C6=C2, has trace defect six and hence K=48, but is
% transverse to the nested active carrier.  Exact K=44 and K=52 contraction,
% semigroup, and energy-midpoint identities are registered as structural
% readouts rather than projector events.  The active E48 carrier gap remains
% unchanged.  A same-action boundary-symbol/Calderon intertwiner is isolated
% as the necessary and still unproved owner bridge to the static DtN sort;
% representation equivalence alone neither derives RAR nor creates a K=48
% action junction, pole, occupation, or cosmological transition.
% Release v1.65 derives the native E_48-centered carrier mass gap from the
% GP1--GP3-generated active spectrum, distinguishes its carrier-free interval
% from a particle located at E_48, and withdraws the unsupported promotion to
% a positive continuum Yang--Mills mass gap.  The latter remains a valid
% conditional transfer theorem on an enlarged continuum branch, but it is not
% an active framework prediction or an unfinished native-gap proof.
% Release v1.64 synchronizes Derivation Batches 14--22.  It registers the
% GP1--GP3 heat-profile orbit, same-Hessian electroweak response and pole
% ownership, K16/K48 action-response corrections, perturbation state and
% non-Gaussianity typing, finite-vacuum/gap/QCD limit repairs, factorized
% position-area ownership, the actual static metric response, and the ordered
% Wilson/string-tension limits.  Strict continuum convergence, positive
% physical gap, surviving higher cumulants, and positive tension remain open.
% Release v1.61 closes TIMEWIDTH16 as an independent direct-production datum.
% JUNC_{2|7} continues to own the finite intrinsic microscopic profile/width,
% while the actual same-action background/production history determines its
% physical-time pullback, support, and duration.  The current Paper-I-A--VII-C AUX
% registry is advanced once; shared mathematical notation is unchanged.
% Release v1.60 derives the final-octet primitive-remnant depletion, canonical
% Hopf-current representative, exact Hopf descent to dOmega/(4pi), and
% closed-FLRW a^{-3} neutral-current shape.  It retypes K=48 as the onset
% of the final one-unit geometric transfer and records cDM/RAR as distinct
% physical projections of that common neutral parent while preserving
% occurrence, Hilbert-owner, stress-matching, and static-kernel premises.
% Release v1.59 transports the v1.58 exceptional-shell/K16 closure through
% Papers III-F and VI-B--VII-C.  It replaces Paper VI-B's obsolete mixed-Hessian
% output typing by the third-variation pair vertex and Bogoliubov occupation
% support, derives the post-seed Higgs-pair feedback vertex and its exact
% alpha--delta kernel factors, propagates the unique K16/K48 action-junction
% split, identifies the QCD delta count with the upstream 2O shell atom only
% as a common-scalar numerical compatibility, and narrows THERM16 to dynamical
% approach/beta selection because the finite Gibbs/KMS target is already
% unique on a supplied simple observable algebra once its Hamiltonian and
% inverse temperature are supplied; direct-sum central weights remain typed.
% Release v1.58 registers the exceptional nearest-shell closure, the derived
% transverse K=16 action event, the unique interior action junction, the
% type-correct ordered pair vertex, the thin-layer squeeze law, and the
% radiation-like/matter-like kinematic split.  REG_0, transfer magnitude,
% physical layer width, and thermalization remain explicitly conditional.
% Release v1.57 registers the active-depth-six obstruction, the native
% K=48 carrier desert, the typed 7|9 owner reduction, and the dynamical
% reheating handoff across Papers III-B--VI-B.
% Release v1.56 registers the least-support n_2 inflation dressing of the
% persistent n_7 Laurent response, the analytic global-convexity and positive
% endpoint-measure proofs, the derived inflation core and K=16 coefficient
% freeze, and the no-particle dark-sector split into static RAR plus the
% conditional homogeneous effective-curvature response at reflected K=48.
% There is no n_6 action, scalaron-DM mass, or second Einstein--Hilbert term.
% Release v1.55 registers the protected-octet resolved/residual reflection
% q_sp+Z_W=8, its j<->8-j and K<->64-K scale covariance, and the derived
% rank-one-exit quadratic grade C_1^2/U=C_2 in Papers III-C, III-E, and VIII.
% The reflection is a typed trace/scale theorem, not a time reversal or a
% physical event-occurrence theorem; reheating transfer remains conditional.
% Release v1.54 registers the exact positive-part Laurent indices, the
% active-domain/terminal-endpoint action split, the no-second-EH firewall, and
% the q_DM=6/K_DM=48 scalaron readout typing across Papers III-D/E/F, VI,
% VIII, and IX.
% Release v1.53 registers the common nonnegative Laurent obstacle projection
% for n_1^+ and n_7^+, discharges the three former exit/freeze assumptions,
% and advances the Paper-III-D/E/F, VIII, and IX consumers.  Shared notation
% and bibliography ownership are unchanged.
% Release v1.51 registers the PQTI derivation batch and its updated
% Paper-III-A/C/D/E/F, VIII, and IX consumers.  Shared notation and
% bibliography ownership are unchanged.
% Release v1.50 advances the Paper-III-B AUX target to v1.40 after the
% rigorous audit repair; shared notation and bibliography ownership are unchanged.
% Release v1.49 advances the Paper-II-E/F AUX target to v3.14 after the
% minimal audit repair; shared notation and bibliography ownership are unchanged.
% Release v1.46 synchronizes Paper II-A's minimal frozen constructor-origin
% theorem and its five dependent AUX targets.  Shared notation is unchanged.
% Release v1.45 updates the Paper-II-B AUX target to v2.23 after the
% GP3-rigidity and appendix-placement repair; shared notation is unchanged.
% Release v1.44 advances the centralized Paper-I-A--VII-C AUX registry after the
% Paper-I-A v5.16 protected-completion scope repair.  Every current
% paper imports this one file; no paper-level opt-in or \externaldocument list
% is permitted.  The bibliography remains centralized here through Born-v6_13.
% Shared bibliography, theorem infrastructure, colors, notation alignment,
% and current Paper-I-A--VII-C external-reference registration.
% This file centralizes common imports, bibliography sorting, colors, theorem
% infrastructure, and the notation shared by the complete series.
% v2.9 (release v1.43) made the Paper-I-A--VII-C AUX registry automatic and
% centralized the bibliography and all external-reference ownership.
% v2.8 release-registry refresh (v1.33) registers the rooted generated-parent
% cascade and the POL/REG/free-law/11+4/Born-group closure revisions.  No
% shared mathematical notation or Paper-II-A fundamental principle is changed.
% v2.8 release-registry refresh (v1.32) registers the universal generated-
% continuation revisions in Papers II-A, III-A, III-C, IV, and VIII.  No
% shared mathematical notation or Paper-II-A fundamental principle is changed.
% v2.8 release-registry refresh (v1.30) registers the CRP4/projective-state
% source revisions and their direct registry-only consumers.  No shared
% mathematical notation or Paper-II-A axiom is changed.
% v2.7 release-registry refresh (v1.27) updates only the current external
% Paper-I-A--VII-C basenames after the placement/assumption synchronization; no
% shared mathematical notation or Paper-II-A axiom is changed.
% v2.7 synchronizes the geometric defect-repair revision: Paper-II-A v3.42,
% Paper-III-C v1.26, Paper-III-D v1.27, Paper-III-E v1.32, and Paper-III-F
% v1.34.  It registers the common trace-defect/spectral-filtration theorem,
% its owner-typed response pullback, and the reduced GRED and TERM interfaces.
% No shared mathematical notation is changed.
% v2.6 synchronizes Paper-III-E v1.31's cone-type repair, minimal-response
% classification, and explicit classical action-role boundary, together with
% Paper-III-F v1.32's downstream response-atlas synchronization.  No shared
% mathematical notation is changed.
% v2.5 synchronizes the Paper-III-D v1.26 faithful-trace-duality proof,
% exact ordered-owner joint locus, and typed alternative-presentation
% bookkeeping repair, together with its Paper-III-E and closure-frontier
% references.  No shared mathematical notation is changed.
% v2.4 synchronizes the Paper-III-C v1.25 orbit/process/owner repair, its
% downstream Papers III-B--F, V, VI, VIII, and IX revisions, and Born-v6_13.
% v2.3 synchronizes the Paper-II-A v3.41 witness/support/principal inheritance
% layer, the Paper-III-A v1.24 joint-realizability/readout firewall, all dependent
% host-rate revisions, and the current Paper-I-A--VII-B source registry.
% v2.2 synchronizes the generated active-realization, determinant-depth,
% dimension-weighted hierarchy, and typed cosmological-event interfaces.
% v2.1 synchronizes the exceptional-support interface revisions: Paper I-A
% v5.9, Paper II-A v3.30, Paper II-E/F v3.7, and Paper IV v4.16.
% v2.0 introduces the stage-safe G0 occurrence notation, synchronizes the
% active Paper-I-A--IV source registry (current rigorous Paper-III closure:
% III-A v1.17, III-B v1.22, III-C--E v1.17, III-F v1.14, IV v4.14), and makes
% external-reference self-skip depend on logical paper role rather than a
% versioned job name.
% ======================================================================

% ======================================================================
% macros.tex — Shared macros for the Born Geometry series
% ======================================================================

% ======================================================================
% 1. PACKAGES, BIBLIOGRAPHY, COLORS, AND UNIT DEFINITIONS
% ======================================================================
% Core packages shared across the complete Paper-I-A--VII-C series.  Loading them here makes
% the macro file a genuine series preamble while remaining compatible with existing
% paper preambles that load the same packages before \input.
\usepackage{amsmath}
\usepackage{amssymb}
\usepackage{amsthm}
\usepackage{mathrsfs}
\usepackage{array}
\usepackage{booktabs}
\usepackage{xcolor}
\makeatletter
\@ifpackageloaded{geometry}{}{\usepackage{geometry}\geometry{margin=1in}}
\makeatother
\usepackage{mdframed}

% Shared support packages required by the macros below and by most papers.
\usepackage{ragged2e}
\usepackage{longtable}
\usepackage{tabularx}
\usepackage{enumitem}
\usepackage[most]{tcolorbox}

% Series-wide numeric bibliography convention using the portable BibTeX path.
% All Papers I-A--VII-C use the ordinary \cite interface.
\usepackage[numbers,sort&compress]{natbib}
\providecommand{\printbibliography}{%
  \bibliographystyle{plainnat}%
  \bibliography{Born-v6_59}%
}

% Shared colors.  \providecolor avoids clobbering a paper-specific override.
\providecolor{postulatecolor}{RGB}{230,242,255}
\providecolor{compare}{RGB}{230,245,230}
\providecolor{definitioncolor}{RGB}{230,242,255}

\providecolor{openriskcolor}{RGB}{255,235,235}
\providecolor{bridgecolor}{RGB}{240,240,255}
\providecolor{gold}{RGB}{255,215,0}

% ======================================================================
% 2. CUSTOM ENVIRONMENTS
% ======================================================================
% Common table column types used in all papers.  Guarded so that a paper may
% define its own versions before loading this file.
\makeatletter
\@ifundefined{NC@find@L}{%
  \newcolumntype{L}[1]{>{\RaggedRight\arraybackslash}p{#1}}%
}{}
\@ifundefined{NC@find@C}{%
  \newcolumntype{C}[1]{>{\Centering\arraybackslash}p{#1}}%
}{}
\@ifundefined{NC@find@Y}{%
  \newcolumntype{Y}{>{\RaggedRight\arraybackslash}X}%
}{}
\makeatother

% ======================================================================
% 3. THEOREM STYLES AND ENVIRONMENTS
% ======================================================================
% Custom underlined theorem style (useful for remarks or notes)
\newtheoremstyle{underlined}
  {3pt}   % Space above
  {3pt}   % Space below
  {}      % Body font (empty means normal/roman text, typical for remarks)
  {}      % Indent amount
  {}      % Theorem head font (can be \itshape, \bfseries, etc.)
  {.}     % Punctuation after theorem head
  { }     % Space after theorem head
  {\underline{\thmname{#1}\thmnumber{ #2}}\thmnote{ (#3)}} % Theorem head formatting

% Standard plain theorem style environments.  Each environment is guarded so
% the shared macro file can be loaded after a paper-specific theorem setup.
\theoremstyle{plain}
\makeatletter
\@ifundefined{theorem}{\newtheorem{theorem}{Theorem}[section]}{}
\@ifundefined{lemma}{\newtheorem{lemma}{Lemma}[section]}{}
\@ifundefined{proposition}{\newtheorem{proposition}{Proposition}[section]}{}
\@ifundefined{corollary}{\newtheorem{corollary}{Corollary}[section]}{}
\makeatother

% Definition style environments (typically non-italicized body)
\theoremstyle{definition}
\makeatletter
\@ifundefined{definition}{\newtheorem{definition}{Definition}[section]}{}
\@ifundefined{postulate}{\newtheorem{postulate}{Postulate}[section]}{}
\@ifundefined{closure}{\newtheorem{closure}{Closure}[section]}{}
\makeatother

% Applying the custom underlined style
\theoremstyle{underlined}
\makeatletter
\@ifundefined{remark}{\newtheorem{remark}{Remark}[section]}{}
\@ifundefined{convention}{\newtheorem{convention}{Convention}[section]}{}
\@ifundefined{principle}{\newtheorem{principle}{Principle}[section]}{}
\makeatother

% Framed proof environment.
% Every \begin{proof} ... \end{proof} is automatically framed/shaded.
% This avoids mdframed/trivlist nesting, which can drop the top border after
% theorem-like environments.
\makeatletter
\renewenvironment{proof}[1][\proofname]{%
  \par\addvspace{6pt}%
  \noindent
  \begin{tcolorbox}[%
    enhanced,
    breakable,
    colback=black!4,
    colframe=black!35,
    boxrule=0.35pt,
    arc=1pt,
    left=6pt,
    right=6pt,
    top=4pt,
    bottom=4pt,
    boxsep=0pt,
    before skip=0pt,
    after skip=6pt,
    parbox=false
  ]%
  \pushQED{\qed}%
  \normalfont\itshape #1\@addpunct{.}\enspace
  \normalfont\ignorespaces
}{%
  \unskip\nobreak\hfill\popQED
  \end{tcolorbox}%
  \par\addvspace{2pt}%
}
\makeatother

\mdfdefinestyle{postulatebox}{
  backgroundcolor=yellow!10,
  linecolor=black,
  linewidth=0.4pt,
  skipabove=\baselineskip,
  skipbelow=\baselineskip,
  nobreak=true
}

% ======================================================================
% PROOF-READABILITY HELPERS
% ======================================================================

\newcommand{\ProofPlan}[1]{%
  \par\smallskip
  \noindent\textbf{Proof plan. }#1%
  \par\smallskip
}

\newlist{proofsteps}{enumerate}{2}
\setlist[proofsteps,1]{
  label=\textbf{(\alph*)},
  ref=\alph*,
  leftmargin=2.2em,
  itemsep=0.45\baselineskip,
  topsep=0.35\baselineskip
}
\setlist[proofsteps,2]{
  label=\textbf{(\roman*)},
  ref=\roman*,
  leftmargin=2.0em,
  itemsep=0.25\baselineskip,
  topsep=0.2\baselineskip
}


% ======================================================================
% 4. NUMBER SYSTEMS AND ALPHABETS
% ======================================================================
\providecommand{\RR}{\mathbb{R}} % Real numbers
\providecommand{\PP}{\mathbb{P}} % Projective space / Probability
\providecommand{\CC}{\mathbb{C}} % Complex numbers
\providecommand{\ZZ}{\mathbb{Z}} % Integers
\providecommand{\NN}{\mathbb{N}} % Natural numbers
\providecommand{\HH}{\mathbb{H}} % Quaternions / Hyperbolic space
\providecommand{\LFRayGeom}{\mathsf{LF\!\mbox{-}RAY\!\mbox{-}GEOM}}
\providecommand{\LFRep}{\mathsf{LF\!\mbox{-}REP}}
\providecommand{\LFState}{\mathsf{LF\!\mbox{-}STATE}}
\providecommand{\LFLim}{\mathsf{LF\!\mbox{-}LIM}}
\providecommand{\LFScat}{\mathsf{LF\!\mbox{-}SCAT}}
% Number systems / projective spaces
\providecommand{\CP}{\mathbb{CP}}
\providecommand{\RP}{\mathbb{RP}}

% ======================================================================
% 5. GROUPS AND LIE THEORY
% ======================================================================
\providecommand{\SO}{\ensuremath{\mathrm{SO}}}     % Special Orthogonal group
\providecommand{\SU}{\ensuremath{\mathrm{SU}}}     % Special Unitary group
\providecommand{\Spin}{\ensuremath{\mathrm{Spin}}} % Spin group
\providecommand{\CS}{\ensuremath{\mathrm{CS}}}     % Chern-Simons (capitalized)
\providecommand{\cs}{\ensuremath{\mathrm{cs}}}     % Chern-Simons (lowercase)

% ======================================================================
% 6. COMMON MANIFOLDS AND TOPOLOGY
% ======================================================================
\providecommand{\Sone}{S^1}                 % Standard circle

% ======================================================================
% 7. MATH OPERATORS (Declared via \DeclareMathOperator)
% ======================================================================
\DeclareMathOperator{\End}{End}   % Endomorphisms
\DeclareMathOperator{\Aut}{Aut}   % Automorphisms
\DeclareMathOperator{\spec}{spec} % Lowercase spectrum
\DeclareMathOperator{\tr}{tr}     % Lowercase trace (if you already have \Tr)
\DeclareMathOperator{\supp}{supp} % Support
\DeclareMathOperator{\TrA}{Tr_{A}} % Partial trace over subsystem A
\DeclareMathOperator{\TrB}{Tr_{B}} % Partial trace over subsystem B
% Lie/algebra operators
\DeclareMathOperator{\Ad}{Ad}
\DeclareMathOperator{\Sym}{Sym}
\DeclareMathOperator{\su}{su}
% ======================================================================
% 8. STANDARD OPERATOR-LIKE SYMBOLS (Safe definitions via \providecommand)
% ======================================================================
\providecommand{\vol}{\operatorname{vol}}   % Volume form
\providecommand{\Vol}{\operatorname{Vol}}   % Total volume
\providecommand{\Ric}{\operatorname{Ric}}   % Ricci tensor
\providecommand{\Rm}{\operatorname{Rm}}     % Riemann tensor
\providecommand{\Tr}{\operatorname{Tr}}     % Uppercase trace
\providecommand{\Spec}{\operatorname{Spec}} % Uppercase spectrum
\providecommand{\diag}{\operatorname{diag}} % Diagonal matrix
\providecommand{\rank}{\operatorname{rank}} % Rank
\providecommand{\scal}{\mathcal R} % Scalar curvature
\providecommand{\Cl}{\operatorname{Cl}}     % Clifford algebra
\providecommand{\diam}{\operatorname{diam}} % Diameter

% ======================================================================
% 9. FUNCTIONS, TRIGONOMETRY, AND CALCULUS
% ======================================================================

% Codifferential naming conventions
% (Avoid \ddag: commonly defined by physics.sty and symbol packages)
\providecommand{\codiff}{\mathrm{d}^{\dagger}} % Codifferential operator

% Quantum mechanics bra-ket notation (needed for basis tables)
\providecommand{\ket}[1]{\lvert #1 \rangle}
\providecommand{\bra}[1]{\langle #1 \rvert}
\providecommand{\braket}[2]{\langle #1 \mid #2 \rangle}

% ======================================================================
% 10. PHYSICAL SCALES AND CONSTANTS (Paper 1c / Born Geometry)
% ======================================================================
\providecommand{\Rp}{\ensuremath{R_{\mathrm P}}}     % matched common EH curvature scale (not defined from observed G)
\providecommand{\hBorn}{\ensuremath{h_{\mathrm B}}} % physical representative 2\pi\hbar of one primitive Born scalar action cell
\providecommand{\UEHjoint}{\ensuremath{U_{J}}} % joint-Dirac EH curvature-coordinate scale 8 kappa_J
\providecommand{\Gread}{\ensuremath{G_{\rm scal}}} % scalar-action coordinate c_op^3/(h_B U); physical matching explicit
\providecommand{\GreadJoint}{\ensuremath{G_{J,\mathrm{raw}}}} % raw ordinary-trace coordinate = G_scal/96
\providecommand{\Rds}{R_{\mathrm{dS}}}    % de Sitter Curvature
\providecommand{\Mpl}{M_{\mathrm{Pl}}}    % Planck mass
\providecommand{\lp}{\ell_{\mathrm{Pl}}}  % Planck length
\providecommand{\Pmax}{P_{\mathrm{max}}}  % Maximum momentum
\providecommand{\Mplbar}{\bar{M}_{\mathrm{Pl}}} % Reduced Planck mass
\providecommand{\eff}{\mathrm{eff}}       % Subscript for effective quantities
\providecommand{\Cosmo}{\ensuremath{\Lambda}}                % Cosmological constant (use \Cosmo in source)
\providecommand{\LCDM}{\ensuremath{\Cosmo\mathrm{CDM}}}  % \LambdaCDM (typesets with \Cosmo)
\providecommand{\Lcut}{\ensuremath{\Lambda_{\mathrm{cut}}}}  % Cutoff scale (brace when used as script)
\providecommand{\Rcs}{\ensuremath{R_{\mathrm{cs}}}}      % Configuration-space radius
\providecommand{\CosmoEff}{\ensuremath{\Cosmo_{\eff}}}   % Effective cosmological constant

% Disambiguation: \Rmom is the momentum-space (S^3) radius. 
% Reserve plain R for the 4D curvature scalar in f(R).
\providecommand{\Rmom}{\ensuremath{R_{\mathrm{mom}}}} 
\providecommand{\Punit}{P_{0}} % Momentum unit set by \Rmom (= \hbar/\Rmom)
\providecommand{\Rvec}{\ensuremath{R_{\mathrm{vec}}}} % invariant vector-space radius in the reduced round sector
% Specific numeric values and expressions
\providecommand{\Vlogcap}{L_{\mathrm{Hopf}}^{(1)}}
\providecommand{\Hopf}{\operatorname{Hopf}}
% Scalar phase (S^1) sector + universal cover notation (Paper I-A)
\providecommand{\tper}{T_\star}      % Time period/scale
\providecommand{\Eper}{E_\star}      % Energy period/scale
\providecommand{\ttilde}{\tilde t}   % Universal cover time
\providecommand{\Etilde}{\tilde E}   % Universal cover energy

% ======================================================================
% 10A. SERIES-WIDE LOGICAL / NOTATIONAL DISAMBIGUATIONS
% ======================================================================
% Use these symbols to avoid collisions among: (i) the Paper--III
% no-insertion/geometric-closure class, (ii) Paper--I-A trimer diagonals,
% (iii) Paper--V-A/V-C family-space diagonals, (iv) scalar-index running,
% and (v) protected shell-clock labels.
\providecommand{\Rprim}{\ensuremath{\mathsf R_{\mathrm{prim}}^{\Sigma}}} % primitive finite structural readouts
\providecommand{\Rloop}{\ensuremath{\mathsf R_{\circ}^{\Sigma}}} % closed-loop readout role
\providecommand{\Rbdry}{\ensuremath{\mathsf R_{\partial}^{\Sigma}}} % boundary-conditioned readout role
\providecommand{\Rgen}{\ensuremath{\mathsf R_{\mathrm{fin}}^{\Sigma}}} % finite generated closure of primitive roles
\providecommand{\PhysOcc}{\operatorname{PhysOcc}}
\providecommand{\Actual}{\operatorname{Actual}}
\providecommand{\GeomLic}{\operatorname{GeomLic}}
\providecommand{\InvDist}{\operatorname{InvDist}}
\providecommand{\Dcl}{\ensuremath{\mathsf{D}_1}}      % unmarked positive Dirac-even closed spectral-readout form class
\providecommand{\Dheat}{\ensuremath{\mathsf D_{\mathrm{heat}}}} % one-parameter heat-semigroup subclass
\providecommand{\DtriI}{\ensuremath{D^{\mathrm{tri}}_1}} % Paper--I-A cyclic-trimer diagonal
\providecommand{\DtriII}{\ensuremath{D^{\mathrm{tri}}_2}}
\providecommand{\DfamI}{\ensuremath{D^{\mathrm{fam}}_1}} % Family-space diagonal selector
\providecommand{\DfamII}{\ensuremath{D^{\mathrm{fam}}_2}}
\providecommand{\alpharun}{\ensuremath{\alpha_{\mathrm{run}}}} % scalar spectral-index running
\providecommand{\Ksh}{\ensuremath{K_{\mathrm{sh}}}} % protected shell-clock coordinate
\providecommand{\ddef}{\ensuremath{\delta_{\mathrm{def}}}} % universal defect scalar after Paper--III realization

% ======================================================================
% 10B. SERIES-WIDE LOGICAL STATUS TAGS
% ======================================================================
% Upright small-caps wrapper for audit/status labels.  It prevents LaTeX from
% trying to use unavailable small-caps italic font shapes inside theorem bodies
% such as TU/lmr/m/scit.
\providecommand{\AuditLabel}[1]{}
\renewcommand{\AuditLabel}[1]{{\normalfont\scshape #1}}

% These tags are bookkeeping labels for consistency audits. They do not
% create new mathematical hypotheses; they record how a datum is being used
% in the proof chain.
\providecommand{\StatusPrimitive}{}
\renewcommand{\StatusPrimitive}{\AuditLabel{Primitive}}
\providecommand{\StatusDerived}{}
\renewcommand{\StatusDerived}{\AuditLabel{Derived}}
\providecommand{\StatusImported}{}
\renewcommand{\StatusImported}{\AuditLabel{Imported}}
\providecommand{\StatusConditional}{}
\renewcommand{\StatusConditional}{\AuditLabel{Conditional}}
\providecommand{\StatusBridge}{}
\renewcommand{\StatusBridge}{\AuditLabel{Bridge-derived}}
\providecommand{\StatusComparison}{}
\renewcommand{\StatusComparison}{\AuditLabel{Comparison-only}}

% ======================================================================
% 10C. SERIES-WIDE CLAIM-WORDING AUDIT LABELS
% ======================================================================
% These are textual audit labels used to keep modal language consistent
% across the series.  They introduce no new mathematical assumptions.
\providecommand{\ClaimForced}{}
\renewcommand{\ClaimForced}{\AuditLabel{forced}}
\providecommand{\ClaimUnique}{}
\renewcommand{\ClaimUnique}{\AuditLabel{unique}}
\providecommand{\ClaimCanonical}{}
\renewcommand{\ClaimCanonical}{\AuditLabel{canonical}}
\providecommand{\ClaimMinimal}{}
\renewcommand{\ClaimMinimal}{\AuditLabel{minimal}}
\providecommand{\ClaimClosed}{}
\renewcommand{\ClaimClosed}{\AuditLabel{closed}}
\providecommand{\ClaimPrediction}{}
\renewcommand{\ClaimPrediction}{\AuditLabel{prediction}}

% ======================================================================
% 10D. SERIES-WIDE THEOREM-CHAIN AUDIT LABELS
% ======================================================================
% These labels mark the direction of load-bearing theorem chains. They are
% audit annotations only; they introduce no new mathematical assumptions.
\providecommand{\ChainInput}{}
\renewcommand{\ChainInput}{\AuditLabel{Chain input}}
\providecommand{\ChainOutput}{}
\renewcommand{\ChainOutput}{\AuditLabel{Chain output}}
\providecommand{\ChainGuard}{}
\renewcommand{\ChainGuard}{\AuditLabel{Chain guard}}
\providecommand{\ChainNoBackflow}{}
\renewcommand{\ChainNoBackflow}{\AuditLabel{No backflow}}

% ======================================================================
% 10E. SERIES-WIDE LOCAL THEOREM-AUDIT LABELS
% ======================================================================
% These labels mark a statement-level dependency audit around individual
% theorem/definition nodes. They are bookkeeping labels only.
\providecommand{\LocalAudit}{}
\renewcommand{\LocalAudit}{\AuditLabel{Local audit}}
\providecommand{\UsesOnly}{}
\renewcommand{\UsesOnly}{\AuditLabel{Uses only}}
\providecommand{\Exports}{}
\renewcommand{\Exports}{\AuditLabel{Exports}}
\providecommand{\DoesNotUse}{}
\renewcommand{\DoesNotUse}{\AuditLabel{Does not use}}

% ======================================================================
% 10F. SERIES-WIDE NUMERICAL-ROW AUDIT LABELS
% ======================================================================
% These labels classify individual rows in numerical, benchmark, status,
% and comparison tables.  They are audit annotations only and do not change
% any theorem, value, formula, or comparison number.
\providecommand{\RowAudit}{}
\renewcommand{\RowAudit}{\AuditLabel{Numerical-row audit}}
\providecommand{\RowTheorem}{}
\renewcommand{\RowTheorem}{\AuditLabel{Theorem-fixed row}}
\providecommand{\RowBridge}{}
\renewcommand{\RowBridge}{\AuditLabel{Bridge-fixed row}}
\providecommand{\RowRepresentative}{}
\renewcommand{\RowRepresentative}{\AuditLabel{Representative row}}
\providecommand{\RowComparison}{}
\renewcommand{\RowComparison}{\AuditLabel{Comparison row}}
\providecommand{\RowCalibrated}{}
\renewcommand{\RowCalibrated}{\AuditLabel{Calibrated row}}
\providecommand{\RowEstimate}{}
\renewcommand{\RowEstimate}{\AuditLabel{Estimate row}}
\providecommand{\RowOpen}{}
\renewcommand{\RowOpen}{\AuditLabel{Open row}}
\providecommand{\RowNoSelector}{}
\renewcommand{\RowNoSelector}{\AuditLabel{No selector}}

% ======================================================================
% 10G. TABLE HELPERS
% ======================================================================
% Table helpers used by Paper V-C.  The unused readout-label preset has been
% removed; only active table helpers remain.
\providecommand{\gpsection}[1]{%
  \addlinespace[0.55em]
  \multicolumn{3}{@{}l}{\textsc{#1}}\\[-0.25em]
  \midrule
}
\providecommand{\mathcell}[1]{%
  \(\begin{gathered}#1\end{gathered}\)%
}

\newcommand{\DerBy}[1]{%
  \begin{tabular}[t]{@{}l@{}}#1\end{tabular}%
}

\newcommand{\TableNotes}[1]{%
  \par\medskip
  \begin{minipage}{0.94\linewidth}
  \footnotesize\emph{Notes.} #1
  \end{minipage}\par
}
% ======================================================================
% 10H. PAPER VI-A / QCD NOTATION
% ======================================================================
% Shared symbols used by Paper VI-A and downstream QCD/Yang--Mills handoffs.
\providecommand{\Pqcd}{\mathfrak P_{\rm VIII}^{\rm QCD}}
\providecommand{\Iqcd}{\mathfrak I_{\rm QCD}^{(7)}}
\providecommand{\muQCD}{\mu_{\rm QCD}^{\rm int}}
\providecommand{\LambdaQCD}{\Lambda_{\rm QCD}}
\providecommand{\LambdaBH}{\Lambda_{\rm BH}^{\rm QCD}}
\providecommand{\HQCD}{\mathcal H_{\rm QCD}^{\rm BH}}
\providecommand{\HQCDfin}{H_{\rm QCD}^{\rm fin}}
\providecommand{\Hraw}{\mathcal H_{\rm raw}}
\providecommand{\HphysQCD}{\mathcal H_{\rm phys}^{\rm QCD}}
\providecommand{\Vcol}{V_C}
\providecommand{\Ps}{P_{\rm sing}}
\providecommand{\PiOct}{\Pi_8}
\providecommand{\MSbar}{\overline{\rm MS}}
\providecommand{\YM}{\mathrm{YM}}
\providecommand{\Ktri}{K_{\triangle}}
\providecommand{\KYM}{K_{\YM}^{\rm BH}}
\providecommand{\aBH}{a_{\rm BH}}
\providecommand{\epsBH}{\varepsilon_{\rm BH}}
\newcommand{\RowConstructive}{\textnormal{\textsc{Constructive}}}
\newcommand{\RowObligation}{\textnormal{\textsc{Bridge obligation}}}
\newcommand{\RowTarget}{\textnormal{\textsc{Target row}}}
\newcommand{\RowValidation}{\textnormal{\textsc{Validation row}}}
\providecommand{\slashed}[1]{\not{#1}}

% ======================================================================
% 10I. PAPER VI-B / CONTINUOUS-QFT NOTATION
% ======================================================================
% Shared symbols used by Paper VI-B and any downstream continuous-QFT or
% Yang--Mills handoff paper.  All definitions are guarded so individual
% papers may override them locally before loading this file if needed.
\providecommand{\A}{\mathfrak A}
\providecommand{\Hphys}{\mathcal H_{\rm phys}}
\providecommand{\QCD}{\mathrm{QCD}}
\providecommand{\BH}{\mathrm{BH}}
\providecommand{\loc}{\mathrm{loc}}

\providecommand{\RBH}{\mathcal R_{\rm BH}}
\providecommand{\DG}{\mathfrak D_G}
\providecommand{\PW}{\mathcal{P\!W}} % Peter--Weyl spectral-symbol subspace
\providecommand{\Toe}{\operatorname{Toe}} % Toeplitz compression/representative
\providecommand{\Hol}{\operatorname{Hol}} % Holonomy

% ======================================================================
% 11. LATE-LOADED PACKAGES (Must be at the end of the preamble)
% ======================================================================
% Load late packages only if a paper has not already loaded them, avoiding
% option clashes in paper-specific preambles.
\makeatletter
\@ifpackageloaded{hyperref}{}{\usepackage{hyperref}}
\makeatother
% --- Series-wide disambiguations (Pass 3) ---
\providecommand{\diracdelta}{\delta_{\rm D}} % Dirac delta distribution (delta reserved for hierarchy)
\providecommand{\Hg}{H_{\mathrm g}} % Hodge harmonic projector (keep H for Hubble)
\providecommand{\Gr}{G_{\mathrm r}} % Green operator (keep G for Newton)
\providecommand{\FD}{\mathrm{D}} % functional/variational derivative symbol (delta reserved)
\providecommand{\Id}{\mathbf{1}} % identity/Kronecker tensor (avoid \delta_{ij})


% BEGIN INLINE SOURCE: Paper00-ActiveExternalReferences-v2_130.tex
% Framework v2.191: current reference binding; existing definitions retained.
% Framework v2.190: current22-root registry with native clock-evolution proofs.
% Framework v2.189: exact current22-root external registry, including retained v2.188 hardening owners.
% Framework v2.187: current registry synchronization only; source release v2.186 preserved.
% Framework v2.186: authoritative current 22-paper registry.
% Framework v2.182: all current 22 logical paper AUX bindings.
% Framework v2.170: current I-B v1.43 external AUX binding.
% Framework v2.167: exact current 22-root registry after CP--BL2 integration.
% Framework v2.166: exact current 22-root compact-release AUX registry.
% Framework v2.165: authoritative current root bindings.
% Framework v2.164: exact current 22 manuscript AUX bindings.
% Framework v2.162: finite metric, functional Hilbert and scoped Ward integration.
% Framework v2.161: current 22 manuscript AUX bindings; source only.
% Framework v2.156 / registry v2.113: synchronized current 22 manuscript roots; source-only registry.
% Framework v2.143 / registry v2.111: active manuscript AUX basenames; no generated AUX or PDF included.
% ======================================================================
% Central external-reference imports, numbers only, no PDF jumping
% ======================================================================

% xr-hyper is loaded centrally so that papers may retain their existing
% \usepackage{xr-hyper} lines without redefining \externaldocument after this
% registry is installed.  If plain xr was loaded earlier, retain it.
\makeatletter
\@ifpackageloaded{xr-hyper}{}{%
  \@ifpackageloaded{xr}{}{\usepackage{xr-hyper}}%
}
\makeatother

% Make external-document registration idempotent.  The registry identifies a
% logical paper by its Paper<number><optional letter> prefix.  Prefix-qualified
% imports remain independent.  Current paper roots contain no local
% \externaldocument calls; this guard also makes accidental duplicates harmless.
\ExplSyntaxOn
\prop_new:N \g_bh_external_refs_prop
\cs_new_eq:NN \BHOriginalExternalDocument \externaldocument

\cs_new_protected:Npn \bh_external_document:nnnn #1#2#3#4
 {
  \tl_set:Nn \l_tmpa_tl {#3}
  \regex_replace_once:nnN
   { \A (?: xref/ )? (Paper[0-9]+[A-F]?) .* \Z }
   { \1 }
   \l_tmpa_tl
  \tl_set:Nx \l_tmpb_tl { #1 | \tl_use:N \l_tmpa_tl }
  \prop_if_in:NVTF \g_bh_external_refs_prop \l_tmpb_tl
   {
    \PackageInfo{BornHopf}
     {Skipping~duplicate~external~refs~for~#3}
   }
   {
    \prop_gput:NVn \g_bh_external_refs_prop \l_tmpb_tl {#3}
    \IfNoValueTF{#4}
     { \BHOriginalExternalDocument[#1][#2]{#3} }
     { \BHOriginalExternalDocument[#1][#2]{#3}[#4] }
   }
 }

\RenewDocumentCommand{\externaldocument}{O{} O{} m o}
 { \bh_external_document:nnnn {#1}{#2}{#3}{#4} }

\NewDocumentCommand{\BHExternalRefs}{m}
 {
  % Do not import the aux file of the logical paper currently being compiled,
  % even when the job and registry use different version suffixes.
  \tl_set:Nx \l_tmpa_tl { \c_sys_jobname_str }
  \regex_replace_once:nnN
   { \A (Paper[0-9]+[A-F]?) .* \Z }
   { \1 }
   \l_tmpa_tl
  \tl_set:Nn \l_tmpb_tl { #1 }
  \regex_replace_once:nnN
   { \A (Paper[0-9]+[A-F]?) .* \Z }
   { \1 }
   \l_tmpb_tl
  \str_if_eq:eeTF { \tl_use:N \l_tmpa_tl } { \tl_use:N \l_tmpb_tl }
   {
    \PackageInfo{BornHopf}
     {Skipping~external~refs~for~#1~because~this~is~the~current~document}
   }
   {
    \IfFileExists{xref/#1.aux}
     { \externaldocument[][nocite]{xref/#1}[#1.pdf] }
     {
      \IfFileExists{#1.aux}
       { \externaldocument[][nocite]{#1} }
       {
        \PackageInfo{BornHopf}
         {External~aux~file~#1.aux~not~found;~no~labels~imported}
       }
     }
   }
 }
\ExplSyntaxOff

% External references are automatic.  This file is the sole owner of both the
% bibliography and the current Paper-I-A--VII-D AUX registry (30 roots including I-C).  Missing auxiliaries
% are informational on the first pass and are populated by the source-only
% two-pass series build.  Basenames deliberately omit temporary download/copy
% suffixes such as ``(1)'' or ``-1''.
  % Paper I-A: intrinsic S3 and exceptional group introduction.
  \BHExternalRefs{Paper1A-Geometry-v5_175}

  % Paper I-B: complete exceptional geometry and finite spectral structures.
  \BHExternalRefs{Paper1B-ExceptionalGeometry-v1_5}
  % Paper I-C: reciprocal product M8 geometry.
  \BHExternalRefs{Paper1D-TechnicalMaterial-v1_65}
  % Paper I-D: terminal technical proofs and calculations.
  \BHExternalRefs{Paper1C-M8Geometry-v1_137}

\ifdefined\BHMathOnly
  % Papers I-A and I-B import only the mathematical atlas labels.
\else
  % Paper II: five physical-mathematical main papers plus one technical proof volume.
  \BHExternalRefs{Paper2A-GeometricConstitution-v3_221}
  \BHExternalRefs{Paper2B-MomentumCarrier-v2_144}
  \BHExternalRefs{Paper2C-ReciprocalLocalization-v2_169}
  \BHExternalRefs{Paper2D-SymmetrySpectra-v2_141}
  % The former II-F export labels remain in II-E; current II-F is the technical volume.
  \BHExternalRefs{Paper2E-KinematicExports-v3_138}
  % Paper II-F: complete technical proof and calculation owner for II-A--II-E.
  \BHExternalRefs{Paper2F-TechnicalMaterial-v1_67}

  % Paper III: Dynamics and Intrinsic Geometric Constitution (six components).
  \BHExternalRefs{Paper3A-GeneratedDynamics-v1_176}
  \BHExternalRefs{Paper3B-SpectralAction-v1_187}
  \BHExternalRefs{Paper3C-ProtectedHierarchy-v1_187}
  \BHExternalRefs{Paper3D-RunningResponse-v1_165}
  \BHExternalRefs{Paper3E-GravityAction-v1_197}
  \BHExternalRefs{Paper3F-CompleteDynamics-v1_263}
  % Paper III-G: complete preserved technical proof owner.
  \BHExternalRefs{Paper3G-TechnicalMaterial-v1_53}

  % Paper IV-A: generated finite quantum mechanics.
  \BHExternalRefs{Paper4A-FiniteQuantumMechanics-v4_214}
  % Paper IV-B: complete finite quantum processes and physical response.
  \BHExternalRefs{Paper4B-QuantumResponse-v1_49}
  % Paper IV-C: extended proofs and finite technical calculations.
  \BHExternalRefs{Paper4C-TechnicalMaterial-v1_49}

  % Paper V-A: The Standard Model.
  \BHExternalRefs{Paper5A-FiniteMatterGeometry-v3_209}

  % Paper V-B: Finite Spectrum and Flavor Closure.
  \BHExternalRefs{Paper5B-TerminalSpectrum-v5_199}

  % Paper V-C: Finite Interaction Boundaries of the Standard Model.
  \BHExternalRefs{Paper5C-InteractionResponse-v2_177}
  % Paper V-D: preserved technical proof and calculation owner.
  \BHExternalRefs{Paper5D-TechnicalMaterial-v1_53}

  % Paper VI-A: Gravity and Closed Cosmology.
  \BHExternalRefs{Paper6A-GravityCosmology-v2_260}

  % Paper VI-B: Geometric Dark Sector and Cosmological Readouts.
  \BHExternalRefs{Paper6B-CosmologyReadouts-v1_206}

  % Paper VI-C: complete technical proof and comparison owner.
  \BHExternalRefs{Paper6C-TechnicalMaterial-v1_50}

  % Paper VII-A: Finite QCD Kernel and Nonperturbative Closure.
\BHExternalRefs{Paper7A-QCDKernel-v1_172}

  % Paper VII-B: Hadronic, Thermal, and Nuclear QCD Readouts.
  \BHExternalRefs{Paper7B-QCDReadouts-v1_157}

  % Paper VII-C: Continuous Quantum Field Theory.
  \BHExternalRefs{Paper7C-ContinuousQFT-v2_202}

  % Paper VII-D: complete technical material for VII-A through VII-C.
  \BHExternalRefs{Paper7D-TechnicalMaterial-v1_53}

\fi

% END INLINE SOURCE: Paper00-ActiveExternalReferences-v2_130.tex


\newcolumntype{R}[1]{>{\raggedleft\arraybackslash}p{#1}}

\newcommand{\PaperSixRevision}{v5.115 neutral geometry and source-scope integration}
\newcommand{\CoreSix}{\mathsf C_6}
\newcommand{\BridgeSix}{\mathsf B_6}
\newcommand{\BridgeLive}{\mathsf B_6^{\rm live}}
\newcommand{\BridgeRaw}{\mathsf B_6^{\rm hist}}
\newcommand{\BridgeObs}{\mathsf B_6^{\rm obs}}
\newcommand{\CLKrem}{\mathsf{CLK}_{\rm rem}}
\newcommand{\Gfive}{\mathsf G_5^{\rm fin}}
\newcommand{\Rrep}{\mathsf R_5^{\rm rep}}
\newcommand{\Rconn}{\mathsf J_5^{\rm conn,der}}
\newcommand{\Rssqft}{\mathsf R_5^{\rm SSQFT}}
\newcommand{\JOINTfiveocc}{\mathsf{JOINT}_5^{\rm der}}
\newcommand{\Ubase}{\mathcal U_6^{\rm str}}
\newcommand{\Urole}{\mathcal U_6^{\rm role}}
\newcommand{\Uabs}{\mathcal U_6^{\rm abs}}
\newcommand{\Uroleabs}{\mathcal U_6^{\rm role,abs}}
\newcommand{\CMPexp}{\mathsf{CMP}_{\rm exp}}
\newcommand{\UU}{\mathrm U}
\newcommand{\Ran}{\operatorname{Ran}}
\newcommand{\HomG}{\operatorname{Hom}}
\newcommand{\sgnline}{\mathbf 1_{\rm sgn}}
\newcommand{\etaH}{\eta_{\rm H}}
\newcommand{\PVIStatus}[1]{\textnormal{\textsc{#1}}}

\newcommand{\pvcite}[1]{\cite{#1}}

\title{\textbf{Born--Hopf Geometry in a Closed Universe V-B: Terminal Spectrum, Flavor and Physical Source Readouts}\\[0.45em]
\large Exact supports, finite mass returns, and conditional pole interfaces}

\author{Peter Kugelmann\thanks{\texttt{pkugelmann@proton.me}}}
\date{Version 5.202 --- 19 November 2026 \\ Framework v2.379}


% BEGIN FRAMEWORK v2.343: CURRENT EXTERNAL SOURCE-REFERENCE ROOTS
\BHExternalRefs{Paper1A-Geometry-v5_175}
\BHExternalRefs{Paper1C-M8Geometry-v1_137}
\BHExternalRefs{Paper1D-TechnicalMaterial-v1_65}
\BHExternalRefs{Paper2A-GeometricConstitution-v3_221}
\BHExternalRefs{Paper2E-KinematicExports-v3_138}
\BHExternalRefs{Paper2F-TechnicalMaterial-v1_67}
\BHExternalRefs{Paper3C-ProtectedHierarchy-v1_187}
\BHExternalRefs{Paper3E-GravityAction-v1_197}
\BHExternalRefs{Paper3F-CompleteDynamics-v1_263}
\BHExternalRefs{Paper3G-TechnicalMaterial-v1_53}
\BHExternalRefs{Paper4B-QuantumResponse-v1_49}
\BHExternalRefs{Paper4C-TechnicalMaterial-v1_49}
\BHExternalRefs{Paper5A-FiniteMatterGeometry-v3_209}
\BHExternalRefs{Paper5C-InteractionResponse-v2_177}
\BHExternalRefs{Paper5D-TechnicalMaterial-v1_53}
\BHExternalRefs{Paper7C-ContinuousQFT-v2_202}
% END FRAMEWORK v2.343: CURRENT EXTERNAL SOURCE-REFERENCE ROOTS


% Current release citation-only xref imports (no material DAG dependency).
\BHExternalRefs{Paper1A-Geometry-v5_175}
\begin{document}
% Publication line-breaking flexibility for technical prose (no math scaling).
\emergencystretch=2em
\maketitle
\begin{abstract}
Paper V-B derives native finite support, terminal coefficients,
family and flavor readouts of the constituted carrier. Its exact intrinsic
energy/mass grammar is not identified with a physical interacting pole unless
the declared source, clock, RG and matching certificates are present.
Rooted negative-Hopf transport and source-incidence statements remain in this
main manuscript; complete auxiliary spectral tables, countermodels and
comparison calculations are in Paper V-D.
\end{abstract}
% Allow the long AA--AK section numbering to fit the contents margin.
\begingroup
\makeatletter
\renewcommand*\l@section[2]{\@dottedtocline{1}{1.5em}{3.0em}{#1}{#2}}
\renewcommand*\l@subsection[2]{\@dottedtocline{2}{4.5em}{3.9em}{#1}{#2}}
\renewcommand*\l@subsubsection[2]{\@dottedtocline{3}{8.4em}{4.7em}{#1}{#2}}
\makeatother
\tableofcontents
\endgroup
\clearpage

\section{Input objects, spectral claims and comparison domains}
\label{p5b:sec:organized-scope}
\subsection{Framework status and logical contract}
\label{p6:sec:framework-status}

\subsubsection{Audited theorem domain}

The structural upstream domain used by the exact finite core is
\begin{equation}
\label{p6:eq:audited-upstream-domain}
\boxed{
\Ubase
=
\mathfrak F_{\le4}
+\Gfive.
}
\end{equation}
The minimal domain for named physical representation roles is the same
structural domain:
\begin{equation}
\label{p6:eq:physical-role-upstream-domain}
\boxed{
\Urole=\Ubase,
\qquad
\Rrep:=\varnothing.
}
\end{equation}
Thus \(\Ubase\) licenses the tensor carrier, structural projectors, generated
hypercharge, the generated even visible-chirality ancestry, and the
corresponding named representation supports.  Numerical Yukawa/Majorana
members, flavor orbits, vacua, schemes and calibrations remain separate data
when a theorem actually consumes them.

The local gauge datum is the derived certificate
\[
\Rconn:=\bigl(\JOINTfiveocc\bigr).
\]
It contains a canonical nonzero connection on the generated full candidate
bundle and the complete unmarked joint Dirac family.  Paper~III-F then derives the factor
connections, curvatures, profiles, and same-datum relations; there is no
second factorwise JOINT selector.  Common relative kinetic normalization is
derived from the uninserted trace on its typed kinetic domain.  The Lorentzian
\(\Rssqft\) package remains an analytic theorem domain rather than a
realization selector.  Absolute shell scales use
\begin{equation}
\label{p6:eq:absolute-upstream-domain}
\boxed{
\Uabs
=
\Ubase+
\bigl(\mathrm{REG}_{\sigma}\in\mathfrak A_{\rm REG}^{+}\bigr)
+(\mathrm{CRP4})+(\mathrm{MOMENERGY}_{0}^{\rm der}).
}
\end{equation}
and named physical absolute readouts use the corresponding role domain
\begin{equation}
\label{p6:eq:physical-role-absolute-domain}
\boxed{
\Uroleabs
=
\Urole+
\bigl(\mathrm{REG}_{\sigma}\in\mathfrak A_{\rm REG}^{+}\bigr)
+(\mathrm{CRP4})+(\mathrm{MOMENERGY}_{0}^{\rm der}).
}
\end{equation}
Here \((\mathrm{CRP4})\) is the one upstream physical parent of
Definition~\ref{p3e:def:crp4-cross-role-positive-parent}; the Paper-V-B name
\((\mathrm{ANC}_E)\) is a deprecated alias, not an additional domain item.
It is not a numerical identification of two positive scalars or an arbitrary
isomorphism of abstract lines: both transports descend from that one
occurring first-order family over their common base.  The derived
\((\mathrm{MOMENERGY}_{0}^{\rm der})\) import supplies
\(E_0=\hbar c_{\rm op}\sqrt U=\sqrt{18/\pi}\,\hbar c_{\rm op}/R_{\rm mom}\);
choice of numerical units is metrological.  Paper~III-C's
owner-normalization inheritance then transports the generated neutral ratio
without a second normalization.  The legacy \((\mathrm{HT})\) label can
certify this completed argument but proves neither ancestry nor calibration.
The positive REG family is the exact domain of the decreasing energy ordering.
A frozen \(\mathrm{REG}_{\sigma}\) member is required only for a
nonhorizontal numerical evaluation and is therefore a scheme boundary with no
natural affine origin, not a Paper-V-B structural selector.

\begin{table}[H]
\centering
\caption{Imported theorem, residual, and provenance contract.}
\label{p6:tab:imported-hypotheses}
\scriptsize
\setlength{\tabcolsep}{3pt}
\hbadness=10000
\renewcommand{\arraystretch}{1.12}
\begin{tabularx}{\textwidth}{@{}L{0.18\textwidth}L{0.49\textwidth}Y@{}}
\toprule
Clause & Content used in Paper--V-B & Owner/status \\
\midrule
\(\mathfrak F_{\le4}\)
& GP1--GP3 together with the derived finite geometry, protected shell/rank
data, conditional hierarchy/action interface, and finite spectral/Takagi
calculus established through Papers I-A--IV.
& Papers I-A--IV \pvcite{Kugelmann2026a,Kugelmann2026b,Kugelmann2026c,Kugelmann2026e};
imported prefix, not rederived wholesale.\\
\(\Gfive\)
& Derived tensor carrier, color/lepton incidence, three-family multiplicity,
scale-free coordinate-scalar type, Abelian candidate plane/kernel, and both
candidate relative-grading projectors.
& Paper V-A \pvcite{Kugelmann2026f}; structural finite theorem output.\\
\(\Rrep\)
& Empty structural representation residual: Paper~V-A generates hypercharge by
parent-relation PCAD, and visible chirality is generated upstream.  Particular
Yukawa/Majorana members are not part of this residual.
& Discharged upstream in Paper~V-A; \(\Rrep=\varnothing\).\\
\(\Rconn=(\JOINTfiveocc)\)
& Canonical nonflat full-bundle connection anchor and complete unmarked joint
Dirac family.  Factor connection, curvature, profile, and same-datum
relations follow functorially.
& Papers III-F and V-A; derived theorem import, with no G0 occurrence or
separate JOINT relation selector.\\
\(\Rssqft\)
& Lorentzian spin--statistics/local-QFT theorem domain invoked only at
explicitly marked interaction, anomaly, or loop interfaces.
& Papers IV--V-A \pvcite{Kugelmann2026e,Kugelmann2026f}; analytic theorem domain,
not a realization selector and never used to create a finite width.\\
\(({\rm HT})\)
& Legacy certification that the neutral determinant-depth character has been
generated and any physical owner ancestry used by the claim has been proved.
& Paper III \pvcite{Kugelmann2026c}; supplies neither ancestry nor a
normalization.\\
\((\mathrm{CRP4})\), historical alias \((\mathrm{ANC}_E)\)
& One positive common-base parent: fundamental determinant-dual transport on
the response summand and its unique positive fourth root on the rank-four
generator summand.  Consequently \(\lambda_{ED}=1\).
& Generated on the occurring rooted action branch of Papers III-B/III-E;
Theorem~\ref{p6:thm:crp4-c13-discharge} derives C13.  The alias is not a
separate Paper-V-B live family.\\
\((\mathrm{MOMENERGY}_{0}^{\rm der})\)
& Derived native shell anchor \(E_0=\hbar c_{\rm op}\sqrt U\) used for absolute shell energies.
& G24 theorem import; not an independent premise. Ratios do not depend on numerical unit realization.\\
\(\mathrm{REG}_{\sigma}\)
& Frozen member of the fixed-shell family
\(\delta_{\sigma}=1/(6\sqrt2)+\sigma_{\rm fin}\),
\(\sigma_{\rm fin}\in\RR\).  The source member \(\mathrm{REG}_0\) has
\(\sigma_{\rm fin}=0\).
& Paper III; unresolved scheme family with covariance tracked in
Section~\ref{p6:sec:hierarchy-scheme}.\\
\bottomrule
\end{tabularx}
\end{table}

\begin{definition}[Corrected finite core]
\label{p6:def:corrected-finite-core}
\(\CoreSix\) is the collection of group, incidence, trace, projector,
normalizer, representation-modulus, and operator-classification theorems
whose hypotheses lie entirely in \(\Ubase\).  It includes the structural
theorems of Part~\ref{p6:part:exact-core}, the bridge-obstruction and
classification theorems of Part~\ref{p6:part:conditional-model}, and the
exact heavy-carrier/no-width boundaries of
Part~\ref{p6:part:interfaces}.  It includes the derived active anchors and
excludes every residual clock, comparison exponent assignment,
normalization, flavor matrix, finite-part value, and pole/rate datum.
An algebraic identity evaluated \emph{after} such a choice is exact
conditional mathematics but is not thereby promoted into \(\CoreSix\).
The logical direction is
\[
\begin{aligned}
\Ubase&\vdash\CoreSix,\\
\Ubase+(\mathrm{CRP4})+\BridgeLive
&\vdash\mathfrak P_{5B}^{\rm premise\text{-}minimal},\\
\Urole+(\mathrm{CRP4})+\BridgeLive+b_{\rm hist}
&\vdash\mathfrak P_{5B}^{\rm historical\ comparison},\\
b_{\rm hist}&\in\BridgeRaw.
\end{aligned}
\]
\(\CoreSix\) is an output of the first judgment, not a primitive premise of
the second.
\end{definition}

\begin{definition}[Paper--V-B live-owner and comparison split]
\label{p6:def:bridge-split}
The premise-minimal live relational package is
\begin{equation}
\label{p6:eq:live-bridge-split}
\boxed{
\BridgeLive
=
\bigl(
\mathrm{REG}_{0},
\texttt{G16}_{\rm claim}
\bigr).
}
\end{equation}
On the occurring anchored Hurwitz/Hopf branch Paper~III-C fixes
$\mathrm{REG}_0$; the unanchored affine $\mathrm{REG}_\sigma$ family remains
only a comparison branch. The complete rooted charged mass/incidence packet
and its finite source realization are generated by
Theorem~\ref{p6:thm:v590-charged-channel-representatives} and Paper~V-A,
Theorem~\ref{p5:thm:v397-charged-preterminal-srcval-realization}; no SVD
assembly interface remains. The claim-relative \texttt{G16} entry contains
only an actually consumed source-cyclic/pole-rest, reporting or limiting
domain. It is empty for a native terminal coefficient or an already covered
internal finite source response on its stated actual owner/state/source domain,
and is not an independent dynamical law.  Here CRP4 is an already-counted upstream occurrence,
not a Paper-V-B-local family; \(\mathrm{ANC}_{E,\lambda}\) survives only in the countermodel
appendix, with \(\mathrm{ANC}_E\equiv\mathrm{CRP4}\) as a deprecated
weight-one alias.  The former finite/raw packet is retained only as the
historical optional comparison atlas
\begin{equation}
\label{p6:eq:bridge-split}
\BridgeRaw
=
\bigl(
\mathsf{SLOT}_{\rm hist},
\mathsf{TMF}_{\rm hist},
\CLKrem^{\rm opt},
\CMPexp^{\rm opt},
\mathsf H_{\rm comp},
\mathsf F_{\rm comp}
\bigr),
\end{equation}
\begin{equation}
\label{p6:eq:observable-bridge-split}
\BridgeObs
=
\bigl(
\mathsf{PV},
\mathsf{QRG},
\mathsf{DEC},
\mathsf{FLRW},
\mathsf N_{\rm th}
\bigr),
\qquad
\BridgeSix=\BridgeLive\times\BridgeRaw\times\BridgeObs .
\end{equation}
The products in \eqref{p6:eq:bridge-split}--\eqref{p6:eq:observable-bridge-split}
are optional enumeration atlases, not the premise-minimal structural package.
A historical numerical member must still satisfy the joint compatibility
conditions of Definition~\ref{p6:def:compatible-raw-bridge}; a formal Cartesian
tuple is not automatically a physical candidate.  In particular, the former
\(\mathsf H_{\rm norm}\) is not live: the invariant embedding projector and
state are derived, and a positive kinetic coefficient is removed by canonical
field rescaling.  A conventional scalar comparison value is a downstream readout coordinate; no native scalar magnitude is required to constitute the electroweak stabilizer.
The optional package \(\mathsf{SLOT}_{\rm phys}\), when supplied, embeds
fermionic lookup rows into Paper--V-A physical role blocks on \(\Urole\) and is
needed only for a finite Dirac realization; it is not part of the external
lookup construction.
Scalar stationary/effective-action dynamics belong to the downstream claim-scoped
\(\mathsf H_{\rm dyn}\) branch when such a claim is requested; native electroweak reduction does not consume that branch.  Paper~V-B records the finite matrix part of
\(\mathsf{PV}\) and the \(\mathsf{QRG}\) dictionary only as comparison
interfaces.  Physical pole selection, vertices, rates, cosmology, and
thermal-scale assignments are downstream contracts and are never premises of
the native finite theorem.
\end{definition}

\begin{definition}[Selector-status partition]
\label{p6:def:selector-status-partition}
The exact finite structural core has no active realization selector.  The
current premise-minimal physical interface is
\[
 \BridgeLive=
 \bigl(\mathrm{REG}_{0},\texttt{G16}_{\rm claim}\bigr).
\]
These entries have different types.  On the occurring anchored Hurwitz/Hopf
branch Paper~III-C fixes \(\mathrm{REG}_0\); the affine
\(\mathrm{REG}_{\sigma}\) family is retained only for unanchored comparison
and scheme studies. The rooted charged terminal coefficients and their
noncentral source/incidence siblings are already generated; their synthesis
does not reconstruct masses as singular values. The claim-relative
\texttt{G16} entry retains only the actual source/support evaluation,
particle/rest identification, named-scheme comparison or limiting transport
that a particular output consumes. Neither primitive source occurrence nor
charged preterminal assembly is an independent interface.  None is a
Paper--V-B-local structural selector.
Pole certification is added only for claims that require a pole mass.

The former \(\mathsf{G5Q\mbox{-}LOOP}\) item is retired as a separate physical
operator/member selector.  Its useful weighted-quiver and spectral algebra is
retained as conditional owner/readout machinery and comparison diagnostics for the generated packet.  The
single upstream CRP4 occurrence, the derived G24 shell anchor, and continuum or
observable interfaces keep their previously typed roles and do not add local
selector roots.
\end{definition}

\begingroup
\footnotesize
\renewcommand{\arraystretch}{1.12}
\setlength{\tabcolsep}{3pt}
\begin{longtable}{@{}>{\raggedright\arraybackslash}p{0.25\linewidth}
                    >{\raggedright\arraybackslash}p{0.24\linewidth}
                    >{\raggedright\arraybackslash}p{0.44\linewidth}@{}}
\caption{Reconstructed Paper--V-B status ledger.}
\label{p6:tab:status-ledger}\\
\toprule
Block & Status & Exact boundary \\
\midrule
\endfirsthead
\caption[]{Reconstructed Paper--V-B status ledger (continued).}\\
\toprule
Block & Status & Exact boundary \\
\midrule
\endhead
\bottomrule
\endfoot

Paper--V-A tensor projectors
& \PVIStatus{branched exact}
& Six structural projectors, refined to eighteen family blocks, exhaust the
derived \(48\)-dimensional tensor carrier under \(\Gfive\); physical role
names use the structurally generated hypercharge role. \\
Nearest \(2I\) shell
& \PVIStatus{corrected}
& The twelve nearest elements generate \(2I\).  The occurring \(q=4\) active
stage supplies a conjugate \(2T\) representative; invariant statements may
equivalently average over all five conjugates. \\
Twenty-five cells
& \PVIStatus{exact}
& Incidence and left--right symmetry are exact; \(24/25\) is label-carrier
self-exclusion, not an intrinsic mass amplitude. \\
Terminal radical grammar
& \PVIStatus{mixed}
& On the atomic integral-transfer branch, a transitive node radical is the
first minimum of its invariant Hom-lattice and is exact up to sign.
Atomicity, common-domain assembly, additive-path
phase/coherence, pseudoinverse metrics, and mass insertion remain bridge data. \\
Physical clock map
& \PVIStatus{mixed}
& \((A^{\rm even},2T)\mapsto32\),
\([C_6\curvearrowright E_3]\mapsto40\), and scalar reclosure \(\mapsto56\)
are derived active anchors.  \(K=36\) uses declared antipodal covariance,
\(K=16\) is a separate heavy rule, and particle/no-particle assignments are
residual bridge data. \\
Scalar normalization/readout chart
& \PVIStatus{derived/readout}
& The rank-fifteen projector, invariant embedding projector, pure state, and
canonical field coordinate are exact.  The conventional scalar comparison
coordinate is a readout chart on the native \(E_{32}\) owner; scalar stationary/effective-action dynamics are claim-scoped in Paper~V-C and are not upstream electroweak constitution. \\
Charged and neutral spectra
& \PVIStatus{conditional}
& The charged-root and light-neutrino operator families are explicitly
retained inside \(\mathsf F\); the displayed source member chooses one root
and one radical neutral operator. \\
CKM/Dirac-PMNS/weak CP
& \PVIStatus{dependent/readout}
& Once the physical quark or Dirac-lepton simultaneous-conjugacy orbit member
is fixed, its spectra, relative eigenframe class, and polynomial weak-CP
invariants are deterministic readouts.  The orbit member itself remains
occurrence/action data; Majorana-specific occurrence is separately typed. \\
Heavy neutral carrier
& \PVIStatus{mixed}
& The quotient-odd rank-three carrier is exact under \(\Gfive\);
\(E_N,\Delta_M\) remain inputs. \\
Widths and reheating
& \PVIStatus{downstream}
& Exact finite Hermitian dynamics has zero width.  Continuum rates and FLRW
histories are not Paper--V-B results. \\

\end{longtable}
\endgroup

\begin{definition}[Physical predictivity]
\label{p6:def:physical-predictivity}
For \(\mathcal U\in\{\Ubase,\Urole\}\), let
\(\mathfrak D_{\rm adm}(\mathcal U)\) be the set of all finite actions and
readout operators compatible with the declared carriers, gradings, reality
structure, gauge action, incidence data, and hypotheses, before a residual
bridge normalization is fixed.
Let \(\mathcal G_{\rm phys}\) contain only gauge and basis changes that leave
physical observables invariant.  Full canonical closure requires
\begin{equation}
\label{p6:eq:singleton-criterion}
\mathfrak D_{\rm adm}(\Urole)/\mathcal G_{\rm phys}
\cong\{\mathrm{point}\}.
\end{equation}
A particular observable is predicted precisely when its evaluation map is
constant on this quotient, including all inequivalent discrete branches.
\end{definition}

\begin{theorem}[Positive-dimensional closure obstruction]
\label{p6:thm:positive-dimensional-obstruction}
The unrestricted equivariant finite terminal-map \emph{comparison class}
contains a non-singleton half-line modulo isometric source/target changes.
The same non-singleton conclusion holds for the structural or physical-role
\emph{admissible quotient} only when at least two members of this family
satisfy all its fixed compatibility, action and readout conditions and extend
to that quotient. Neither admissible-domain membership nor such a role
extension follows solely from equivariance.
\end{theorem}

\begin{proof}
Fix any nonempty transitive terminal support \(S\), put
\(u_S=|S|^{-1/2}\sum_{s\in S}|s\rangle\), and consider
\[
T_\lambda(1)=\lambda u_S,
\qquad \lambda>0.
\]
Every \(T_\lambda\) is an equivariant finite comparison map and has
singular value \(\lambda\). Isometric changes of basis preserve that
value; distinct \(\lambda\) therefore define distinct comparison
classes. But a compatible admissible action/readout tuple may impose
additional constraints on \(\lambda\), and a role-specific extension
can further restrict it. Restricting to those \(\lambda\) with verified
admissible representatives gives exactly the conditional quotient statement
in the theorem. No unverified physical half-line of admissible operators is
inferred merely from this algebraic comparison family.
\end{proof}

\begin{corollary}[Additional nonminimal obstruction witnesses]
\label{p6:cor:additional-obstruction-witnesses}
The shorter proof above does not remove the other demonstrated freedoms.  The
unresolved role/evaluation placement
\(K_a=32+8a\), \(0<a<1\), the scalar coordinate \(\kappa>0\), and a relative
family rotation \(R_{12}(\theta)\), \(0<\theta<\pi/2\), give additional typed
counterfamilies through the derived shell law,
Theorem~\ref{p6:thm:scalar-counterfamily}, and
Theorem~\ref{p6:thm:eigenframe-counterfamily}, respectively.  They strengthen
but are not premises of Theorem~\ref{p6:thm:positive-dimensional-obstruction}.
\end{corollary}

\begin{corollary}[Absolute and finite-comparison extensions]
\label{p6:cor:absolute-comparison-obstruction}
On the enlarged absolute domain \(\Uabs\), the finite-part deformation
\(c\in\RR\) of Theorem~\ref{p6:thm:finite-part-covariance} changes different
shells by different factors.  If the finite pole-matrix comparison is
appended, the coherent-boundary deformation \(b\in\CC\) of
Theorem~\ref{p6:thm:coherent-boundary-family} changes its displacement by
\(|b|\sqrt R\).  These are additional obstructions in their respective
typed extensions; neither is needed for the native structural or physical-role
non-singleton theorem.
\end{corollary}

\begin{corollary}[Continuum-observable obstruction]
\label{p6:cor:extended-observable-obstruction}
If a continuum interaction is appended downstream, the rescaling
\(V\mapsto gV\) preserves its channel and covariance data while changing a
weak-coupling rate by \(\Gamma\mapsto g^2\Gamma\).  This is an additional
obstruction in the extended observable domain; it is not a modulus of the
native finite Paper--V-B quotient.
\end{corollary}

\begin{corollary}[Correct status of ``no free coefficient'']
\label{p6:cor:no-free-coefficient-status}
Checking that every symbol in a chosen formula was previously defined is a
valid provenance audit, but it does not prove predictivity.  For one fixed
integer \(n>1\), the coefficient-free rules
\[
1,\qquad\sqrt n,\qquad n,\qquad\sqrt{n+1}
\]
are distinct.  The source's former no-free-coefficient theorem is therefore
replaced by Definition~\ref{p6:def:physical-predictivity}.
\end{corollary}

\subsection{Acyclic ownership and interfaces}
\label{p6:sec:ownership}

\begin{proposition}[Required acyclic target]
\label{p6:thm:acyclic-ownership}
To remove the audited cross-paper cycles, the companion manuscripts must
adopt the following one-way ownership order:
\[
\begin{aligned}
&\text{Papers I-A--II: typed geometry and role signatures}\\
&\longrightarrow
\text{Paper III: derived dual determinant-depth hierarchy and REG scheme}\\
&\longrightarrow
\text{Paper IV: finite operator and limit interfaces}\\
&\longrightarrow
\text{Paper V-A}:\Gfive\text{ structural carrier and physical representation roles}\\
&\longrightarrow
\text{Paper V-B: exact finite core and }\BridgeRaw\\
&\longrightarrow
\text{Paper V-C: interactions, scalar potential, EWSB, QCD interface,
and }g_*\\
&\longrightarrow
\text{Paper VI-A: gravity, }R+R^2\text{ matching, and conditional reheating}.
\end{aligned}
\]
Post-reheating cosmology and any thermal heavy-scale assignment occur after
Paper~VI-A.  Finite QCD kernels may be developed independently downstream,
and a later continuum reconstruction may validate earlier interfaces but is
not an earlier premise.  This proposition is the target dependency contract
exported by the revised Paper~V-B; it does not assert that the currently
supplied downstream sources have already been synchronized.
\end{proposition}

\begin{proof}
Under the stated target, every arrow points from the owner of a proved or
explicitly conditional object to a consumer.  Paper~V-B leaves the scalar
potential, \(E_N\), and \(\Delta_M\) symbolic where required.  Enforcing the
target forbids Paper~V-C from using a Paper--VI-A reheating result to prove a
Paper--V-B or Paper--V-C input; permits Paper~VI-A to consume \(g_*\) without
being cited back into its owner; and removes any later continuum theorem from
the finite carrier's premise set.
\end{proof}

\begin{table}[H]
\centering
\caption{Corrected ownership of formerly entangled results.}
\label{p6:tab:ownership}
\footnotesize
\hbadness=10000
\renewcommand{\arraystretch}{1.12}
\begin{tabularx}{\textwidth}{@{}L{0.27\textwidth}L{0.22\textwidth}Y@{}}
\toprule
Object & Owner & Paper--V-B contract \\
\midrule
Dual-character/REG shell family & Paper III &
Import the dimensionless neutral hierarchy ratios as theorems.  A physical
energy ladder additionally uses \((\mathrm{CRP4})\); the dimensionful
reference and REG member remain declared calibration/scheme data, and
decreasing order uses \(\mathfrak A_{\rm REG}^{+}\).\\
Finite spectral calculus & Paper IV &
Use for finite eigenvalues and Takagi/singular-value constructions; do not
infer nonzero widths.\\
Tensor carrier and physical roles & Paper V-A, \(\Gfive\) &
The tensor carrier/projectors are structural \(\Gfive\) outputs; named
hypercharge representation slots are generated structurally, while relative chirality is inherited from the generated upstream ancestry theorem.  Numerical Yukawa--Majorana members remain downstream.  Numerical
family matrices remain downstream.\\
Terminal, clock, cell, and flavor readouts & Paper V-B &
Exact core plus explicit \(\BridgeRaw\); theorem status retained row by row.\\
Native EW owner/stabilizer & Papers III-C + V-A + III-F &
Paper V-B derives the invariant scalar embedding/state; the former
\(\mathsf H_{\rm norm}\) is not live. Optional historical scalar comparison
coordinates lie in \(\mathsf H_{\rm comp}\); a physical scalar magnitude
requires its downstream owner.\\
Scalar stationary/effective-action branch & Paper V-C &
Claim-scoped only for equilibrium, scalar-Hessian, global energetic-preference, or actual scalar-history claims; not an upstream electroweak owner.\\
Self-energies, vertices, and physical poles & Paper V-C/QFT owner &
The finite \(\mathsf{PV}_{\rm mat}\) packet is only a comparison matrix.\\
Quark common-slice values & QCD/QFT owner &
Require \(\mathsf{QRG}\); native finite coordinates are boundary data.\\
Scalaron mass & Paper VI-A &
Obtained from its \(R+R^2\) coefficient match, not from Paper V-B.\\
Scalaron partial/total widths & continuum interaction owner &
Require \(\mathsf{DEC}\); channel labels and totals must be explicit.\\
Reheating history and temperature & Paper VI-A/cosmology owner &
Require \(\mathsf{FLRW}\); equality-time and history-level temperatures are
distinct functionals.\\
Thermal heavy scale & post-reheating owner &
\(\mathsf N_{\rm th}\) may assign \(E_N,\Delta_M\) only downstream and may
not be cited back into the carrier theorem.\\
\bottomrule
\end{tabularx}
\end{table}

\subsection{Complete typed bridge package}
\label{p6:sec:complete-bridge}

\begin{definition}[Compatible historical comparison member]
\label{p6:def:compatible-raw-bridge}
A member \(b_{\rm hist}\in\BridgeRaw\) is a compatible optional comparison
member when the following historical data obey their stated typing,
covariance, positivity, self-adjointness, and cross-sector conditions.  None
of these rows belongs to \(\BridgeLive\).
\begin{description}[style=nextline,leftmargin=2.8em]
\item[\(\mathsf{SLOT}\)]
An explicit external lookup-row specification, including every row domain
and codomain, multiplicity, reused support, and whether the construction is
a direct sum, tensor product, common-domain column, or non-projective ledger.
It does not assert an embedding into Paper--V-A role projectors.  On \(\Urole\),
such embeddings, if desired, form the separate optional package
\(\mathsf{SLOT}_{\rm phys}\) of
Corollary~\ref{p6:cor:fermionic-dirac-restriction}.

\item[\(\mathsf{TMF}\)]
One retained member of the residual terminal-mass realization family:
atomic/primitive transfer in the canonical integral Hom-lattice; common-domain
incidence for Pythagorean sums; positive codiagonal and phase alignment for
additive paths; fixed Hilbert metrics for pseudoinversion; and one dimensionful
shell insertion.  The transitive node magnitude and sign-insensitivity follow
from atomicity and are not separately selected coordinates.

\item[\(\CLKrem\)]
The optional antipodal-orbit role assignment placing the interface packet at
\(K=36\), the historical comparison association of the \(2O/2T\) quotient
\(K=16\), and the declaration that \(K=8,24\) carry no role assignment in
this paper.  Native \(K=48\) carrier absence is derived in
Theorem~\ref{p6:thm:native-k48-carrier-desert}; it is not optional clock data.
The continuous
dimension-weighted shell law, all seven scalar depth levels, and the anchors
\((A^{\rm even},2T)\mapsto32\),
\([C_6\curvearrowright E_3]\mapsto40\), and scalar reclosure \(\mapsto56\)
are upstream theorems and are not members of this residual package.

\item[\(\CMPexp\)]
The transitive orthogonal twenty-five-cell label carrier and self-exclusion
map together with the charged-derived / otherwise retained exponent assignment
\(n_{\etaH}(p)\in\tfrac12\ZZ\), and the operator that inserts
\(\etaH^{n_{\etaH}(p)}\) in the selected finite readout.  Factor count, factor
names, and a word lift are quotient-redundant for all current outputs.

\item[\(\mathsf H_{\rm comp}\)]
One historical scalar comparison coordinate.  The invariant embedding
projector and state are derived, a positive kinetic coefficient is quotiented
by field rescaling, and the physical scalar coordinate is a downstream vacuum
readout.  This row contains no live normalization, potential, vacuum selector,
or EWSB theorem.

\item[\(\mathsf F\)]
The product-valued conditional spectrum/flavor construction and readout
family on \(\Urole\): charged norm-duality and retained-root member, the charged absolute
map, one retained positive light-neutrino operator, the relative CKM/PMNS
eigenframe invariants and their standard representatives, and the heavy-family
package \((E_N,\Delta_M,\widehat D_N)\).  Positive Majorana parity is derived
after the positive neutral operator is frozen.  This raw Paper--V-B atlas is
not the physical weak-basis quotient of
Definition~\ref{p6:def:weak-basis-flavor-quotient}: presentation-equivalent
operator representatives are not distinct physical flavor members, and the
high-energy neutrino lift remains a later Paper--V-C fiber.
\end{description}
\end{definition}

\begin{definition}[Observable/downstream package]
\label{p6:def:observable-package}
The downstream interfaces are:
\begin{description}[style=nextline,leftmargin=2.8em]
\item[\(\mathsf{PV}\)]
\(\mathsf{PV}_{\rm mat}\) together with normalized self-energies, branch
selection, vertex maps, Ward identities, and a pole/renormalization scheme.

\item[\(\mathsf{QRG}\)]
The complete package of Definition~\ref{p6:def:qrg-package}.

\item[\(\mathsf{DEC}\)]
A normalized interaction, channel multiplicities and curvature couplings,
a continuum or reservoir limit, and an explicit partial-versus-total label.

\item[\(\mathsf{FLRW}\)]
Initial scalaron/radiation data, \(g_*(T)\), thermalization and collision
terms, expansion equations, and a precise temperature functional.

\item[\(\mathsf N_{\rm th}\)]
A post-reheating heavy-scale rule that either leaves
\(E_N,\Delta_M\) symbolic or assigns them without being cited back into
Paper~V-B.
\end{description}
\end{definition}

\begin{remark}[Downstream-interface status]
\label{p6:rem:bridgeobs-downstream-status}
The package \(\BridgeObs\) is not part of the native Paper--V-B selector count.
Its members are downstream continuum/observable realization interfaces.  A
later paper may consume or independently discharge one of them, but such a
later use does not retroactively select the finite/raw bridge.
\end{remark}

\begin{theorem}[Conditional finite sufficiency]
\label{p6:thm:conditional-finite-sufficiency}
Fix a compatible \(b_{\rm raw}\in\BridgeRaw\), and let
\[
b_{\rm str}
=
(\mathsf{SLOT},\mathsf{TMF},\CLKrem,\CMPexp,\mathsf H_{\rm comp})
\]
be its historical structural/comparison subtuple, with the components
understood as the corresponding \(\mathsf{SLOT}_{\rm hist}\),
\(\mathsf{TMF}_{\rm hist}\), optional clock and exponent data
of \(\BridgeRaw\).  Then the domains split minimally as
follows.
\begin{enumerate}[leftmargin=2em]
\item On \(\Ubase+b_{\rm str}\), enlarged to \(\Uabs+b_{\rm str}\) for
absolute coordinates, there is a single-valued evaluation to the six-block
and family-refined structural projectors, one external typed lookup ledger,
one finite lookup operator and self-adjoint dilation, and the declared
historical scalar comparison coordinate without a physical vacuum or native
scalar-magnitude claim.
\item On \(\Urole+b_{\rm raw}\), enlarged to
\(\Uroleabs+b_{\rm raw}\) for named absolute coordinates, the representation
roles additionally type one retained-order charged spectrum, one positive
light-neutrino operator, one symbolic or declared heavy-family operator,
exact CKM/PMNS representatives from the declared sector operators, and the
positive Majorana-parity consequence when the optional Majorana operator
occurs.
\item If \(\mathsf{SLOT}_{\rm phys}\) is also supplied on \(\Urole\), the
fermionic lookup rows transport to the conditional finite Dirac realization
of Corollary~\ref{p6:cor:fermionic-dirac-restriction}.
\end{enumerate}
No numerical member, flavor orbit, pole, vacuum, scheme or calibration from
item~2 is licensed by \(\Ubase\) alone; the displayed compatible member data
remain required.
\end{theorem}

\begin{proof}
\(\mathsf{SLOT}\) fixes all domains and codomains.
The first-minimum theorem fixes transitive node magnitudes from the atomicity
clause in \(\mathsf{TMF}\), while \(\mathsf{TMF}\), \(\CLKrem\), and
\(\CMPexp\) fix the remaining finite
transfer on every declared lookup row.  On the external row-sum carrier,
\[
M_{{\rm lookup},b}
=
\sum_{p\in\mathcal P_{\rm mass}}
E_{K(p)}\widehat J_{p,R}
\bigl(w_pF_{\rm term}(t_p)\bigr)
\widehat J_{p,L}^* .
\]
The doubling
\[
D_{{\rm lookup},b}
=
\begin{pmatrix}
0&M_{{\rm lookup},b}^*\\
M_{{\rm lookup},b}&0
\end{pmatrix}
\]
is self-adjoint and has a finite singular spectrum.
Theorem~\ref{p6:thm:canonical-invariant-scalar-state} fixes the invariant
scalar embedding/state, whereas \(\mathsf H_{\rm comp}\) is merely the
historical comparison-coordinate entry of the raw tuple. It does not fix a
new native scalar normalization or vacuum.  These statements consume only \(\Ubase\) and
\(b_{\rm str}\).

On \(\Urole\), \(\mathsf F\) fixes the charged branch, the exclusive neutrino
operator, sector matrices, phases, and heavy affine operator before
diagonalization; finite spectral and Takagi calculus make the named
evaluations single valued.  If the optional fermionic
\(\mathsf{SLOT}_{\rm phys}\) embeddings of
Corollary~\ref{p6:cor:fermionic-dirac-restriction} are supplied, transport
gives the finite Dirac operator; bosonic readouts are not included in that
interpretation.
\end{proof}

\begin{corollary}[Downstream extension]
\label{p6:cor:downstream-extension}
If compatible owners additionally supply \(\BridgeObs\), the finite export
has single-valued comparison maps to physical poles and vertices, a
common-scheme quark table, declared continuum partial or total rates, an FLRW
history and temperature functional, and a downstream heavy-scale placement.
This is a composition theorem for the extended model, not a Paper--V-B
derivation of \(\BridgeObs\).
\end{corollary}

\begin{theorem}[Relative minimality by unresolved type]
\label{p6:thm:relative-minimality}
Relative to the outputs retained in
Theorem~\ref{p6:thm:conditional-finite-sufficiency} and
Corollary~\ref{p6:cor:downstream-extension}, omitting a package component
reintroduces the corresponding demonstrated ambiguity:
\begin{center}
\footnotesize
\setlength{\tabcolsep}{4pt}
\renewcommand{\arraystretch}{1.15}
\begin{tabularx}{\linewidth}{@{}>{\raggedright\arraybackslash}p{0.21\linewidth}
                          >{\raggedright\arraybackslash}X@{}}
\toprule
Omitted component & Unresolved type\\
\midrule
\(\mathsf{SLOT}\)& Lookup domains, multiplicities, reuse, and row construction\\
\(\mathsf{TMF}\)& Atomicity, common/additive assembly, phase, metric, and mass unit\\
\(\CLKrem\)& \(K=16\) carrier placement and \(K=36\) orbit-role assignment;
\(K=8,24\) no-role declarations (not the derived \(K=48\) desert)\\
\(\CMPexp\)& Charged grade derived in v5.91; noncharged assignment where still retained\\
\(\mathsf H_{\rm comp}\)& Optional historical scalar comparison coordinate; no native vacuum\\
\(\mathsf F\)& Charged branch, neutral choice, frames, parities, heavy matrix\\
\(\mathsf{PV}\)& Pole normalization/sign and vertex definition\\
\(\mathsf{QRG}\)& Quark scheme, scale, and thresholds\\
\(\mathsf{DEC}\)& Nonzero rate, multiplicity, and total-width meaning\\
\(\mathsf{FLRW}\)& History, thermalization, and temperature definition\\
\(\mathsf N_{\rm th}\)& Thermal heavy-scale placement\\
\bottomrule
\end{tabularx}
\end{center}
This is relative minimality of unresolved \emph{types}; one future normalized
action may derive and subsume several components.
\end{theorem}

\begin{proof}
The counterfamilies are, respectively: the non-projective packet lookup;
nonprimitive node multiples, independent-domain versus common-domain assembly,
and relative path phases;
logarithmic role placement and the clock/normalizer mismatch; distinct
exponent functions such as \(n_p=0\) and \(n_p=1/2\); the scalar
\(\kappa\)-family; the two charged roots,
three neutral spectra, arbitrary eigenframes, parity choices, and general
\(H_N\); the coherent \(b\)-family and independent vertex shifts; ordinary
RG scheme freedom; interaction rescaling; alternative cosmological initial
data and temperature functionals; and the free thermal heavy anchor.
\end{proof}

\begin{corollary}[Correct meaning of closure]
\label{p6:cor:correct-closure-meaning}
Theorem~\ref{p6:thm:conditional-finite-sufficiency} proves that the repaired
model is mathematically coherent and evaluable once its typed finite bridge
is fixed.  Theorem~\ref{p6:thm:positive-dimensional-obstruction} exhibits a
continuous comparison-map modulus, but does not by itself prove two
compatible physical admissible representatives. No unconditional uniqueness
or nonuniqueness of that narrower physical quotient is inferred. Hence the
accurate title-level claim is
\[
\boxed{\text{rigorous finite core plus conditional spectrum closure},}
\]
not unconditional parameter-free prediction.
\end{corollary}

\label{p6:part:integrated}

\label{p6:part:exact-core}

\paragraph{Technical proof owner and claim-status guard.}
Detailed auxiliary matrix, representation, finite-spectral and comparison
calculations formerly embedded in this manuscript now belong to
Paper~V-D. The original named theorem/proof labels remain
cross-referenced, and no optional finite comparison is promoted
to physical source occurrence, interacting mass pole or measured response.

\section{Exact finite supports, projectors and shell geometry}
\label{p5b:sec:organized-support}
\subsection{Common projector exhaustion and family moduli}
\label{p6:sec:projector-exhaustion}

\begin{theorem}[Exhaustion of the derived Paper--V-A tensor carrier]
\label{p6:thm:matter-exhaustion}
On the structural domain \(\Ubase\), let
\[
W=W_L\oplus W_{R+}\oplus W_{R-},\qquad
Q=\ell\oplus C,\qquad G\cong\CC^3,
\]
and
\[
F=W\otimes Q\otimes G.
\]
For
\(\sigma\in\{L,R+,R-\}\) and \(q\in\{\ell,C\}\), define
\[
Q_{\sigma q}=P_\sigma\otimes P_q\otimes I_G.
\]
Then
\[
Q_{\sigma q}Q_{\sigma' q'}
=\delta_{\sigma\sigma'}\delta_{qq'}Q_{\sigma q},
\qquad
\sum_{\sigma,q}Q_{\sigma q}=I_F.
\]
Their ranks are
\[
\begin{array}{c|rrrrrr}
(\sigma,q)&(L,C)&(L,\ell)&(R+,C)&(R-,C)&(R+,\ell)&(R-,\ell)\\
\hline
\rank Q_{\sigma q}&18&6&9&9&3&3
\end{array}
\]
Thus the six sectors exhaust the derived \(48\)-dimensional structural tensor
carrier.  No physical hypercharge or chirality name is used in this statement.
\end{theorem}

\begin{proof}
Factorwise orthogonality gives the product relation, and
\[
\sum_{\sigma,q}Q_{\sigma q}
=
(P_L+P_{R+}+P_{R-})
\otimes(P_\ell+P_C)\otimes I_G=I_F.
\]
The ranks use
\(\dim W_L=2\), \(\dim W_{R\pm}=1\),
\(\dim C=3\), \(\dim\ell=1\), and \(\dim G=3\).
\end{proof}

\begin{corollary}[Physical role-identification domain]
\label{p6:cor:physical-role-identification-domain}
On \(\Urole=\Ubase\), Paper~V-A's parent-derived Abelian
kernel and even graded bridge identify the ordered structural blocks as
\[
\begin{array}{c|rrrrrr}
(\sigma,q)&(L,C)&(L,\ell)&(R+,C)&(R-,C)&(R+,\ell)&(R-,\ell)\\
\hline
\text{physical role}&Q_L&L_L&u_R&d_R&\nu_R&e_R
\end{array}
\]
The former \(\Rrep\) distinction is discharged upstream; the projector/rank theorem and these
physical names.  Hence the corollary is domain-minimal relative to Paper~V-A's
representation theorem.
\end{corollary}

\begin{corollary}[Eighteen structural family blocks]
\label{p6:cor:eighteen-family-blocks}
If \(P_0,P_1,P_2\) are the primitive orthogonal family projectors, then
\[
Q_{\sigma qi}:=P_\sigma\otimes P_q\otimes P_i
\]
are eighteen orthogonal projectors and
\(
\sum_{\sigma,q,i}Q_{\sigma qi}=I_F
\).
They are structural blocks, not particle-mass eigenprojectors.
\end{corollary}


\begin{corollary}[Structural family projectors are not a common mass eigenframe]
\label{p6:cor:structural-projector-mass-firewall}
Let \(H_s\) and \(H_t\) be simple self-adjoint sector operators on the
family carrier.  If the same rank-one projector resolution diagonalizes both,
then \([H_s,H_t]=0\), and their relative diagonalizing matrix is monomial up
to diagonal rephasings.  In particular, the corresponding Jarlskog invariant
vanishes.  Therefore the structural projectors
\(P_0,P_1,P_2\) cannot be identified with the mass-eigenprojector resolution
of every physical sector without an additional relation theorem.
\end{corollary}

\begin{proof}
A simple self-adjoint operator is a linear combination of its unique rank-one
spectral projectors.  If two such operators use the same projector resolution,
they commute and are simultaneously diagonal.  Their relative eigenframe can
then differ only by permutation and phase matrices, which has no nontrivial
Jarlskog invariant.
\end{proof}

\begin{theorem}[Typed orthogonal-exhaustion criterion]
\label{p6:thm:typed-exhaustion}
Finite role carriers \(\mathcal H_r\) form a complete mutually exclusive
decomposition of one previously constituted carrier \(\mathcal U\) if and
only if there are isometries
\[
\iota_r:\mathcal H_r\longrightarrow\mathcal U
\]
such that
\[
\iota_r^*\iota_s=\delta_{rs}I_{\mathcal H_r},
\qquad
\sum_r\iota_r\iota_r^*=I_{\mathcal U}.
\]
Without \(\mathcal U\) and the \(\iota_r\), a projector sum across different
role carriers is ill typed.
\end{theorem}

\begin{proof}
The conditions are exactly the range-isometry form of an orthogonal
projector resolution.  Defining an abstract direct sum after the fact gives a
tautological resolution; it does not prove that the roles were sectors of one
pre-existing physical carrier.
\end{proof}

\begin{theorem}[Rank-five packet obstruction]
\label{p6:thm:rank-five-packet-obstruction}
Let \(X=2I/2T\), so \(|X|=5\).  The natural \(A_5\) permutation module is
\[
\CC[X]\cong\mathbf1\oplus\mathbf4
\]
with irreducible four-dimensional complement.  The only
\(A_5\)-equivariant orthogonal projectors are
\[
0,\qquad P_{\mathbf1},\qquad P_{\mathbf4},\qquad I.
\]
In particular, there is no equivariant four-block resolution of ranks
\(1,1,1,2\).
\end{theorem}

\begin{proof}
The commutant of \(\mathbf1\oplus\mathbf4\) is
\(\CC\oplus\CC\).  A self-adjoint idempotent independently chooses zero or
identity on the two irreducible summands.
\end{proof}

\begin{corollary}[Status of \(1+1+1+2=5\)]
\label{p6:cor:rank-five-accounting}
The identity
\(
1+1+1+2=[2I:2T]=5
\)
is normalized accounting only.  The Paper--V-B packet table may be a complete
lookup table after \(\mathsf{SLOT}\) is declared, but it is not one common
orthogonal projector exhaustion.
\end{corollary}

\begin{theorem}[Admissible Yukawa--Majorana space on the role domain]
\label{p6:thm:yukawa-majorana-space}
On \(\Urole\), the four named Dirac slots and the unique optional neutral
Majorana support are well typed.  Their coefficient atlas is
\[
\mathfrak M_5
=
M_3(\CC)_u\oplus M_3(\CC)_d\oplus M_3(\CC)_e
\oplus M_3(\CC)_\nu\oplus\operatorname{Sym}_3(\CC)_M,
\]
and
\[
\dim_\CC\mathfrak M_5=4\cdot9+6=42.
\]
Thus core-slot multiplicity one does not imply family-map uniqueness.  The
summand \(\operatorname{Sym}_3(\CC)_M\) records allowed coefficients if a
Majorana operator occurs; the theorem does not assert that occurrence.
\end{theorem}

\begin{proof}
Each Dirac coefficient is an arbitrary element of
\(\End(G)\cong M_3(\CC)\).  A Majorana bilinear has an arbitrary complex
symmetric \(3\times3\) coefficient.  Their direct-sum dimensions give the
claim.
\end{proof}

\begin{definition}[Classical weak-basis flavor quotient]
\label{p6:def:weak-basis-flavor-quotient}
This definition is made on \(\Urole=\Ubase\); the same generated
hypercharge/chirality role typing therefore applies throughout this structural
domain.  The weak-basis quotient removes family-basis presentation freedom but
does not select a normalized flavor/source member.  Let
\[
\mathcal G_{\rm wb}
=
U(3)_Q\times U(3)_u\times U(3)_d\times
U(3)_L\times U(3)_e\times U(3)_N
\]
act on the Paper--V-A family coefficients by the induced kinetic-term-preserving
changes of family bases.  On the Majorana branch define the classical
weak-basis quotient
\[
\mathfrak A_{\rm flav}^{M}
:=
\mathfrak M_5/\mathcal G_{\rm wb}.
\]
On the Dirac-neutrino branch, where no Majorana operator occurs, replace
\(\mathfrak M_5\) by
\[
\mathfrak M_5^{D}
:=
M_3(\CC)_u\oplus M_3(\CC)_d\oplus
M_3(\CC)_e\oplus M_3(\CC)_\nu
\]
and write
\(\mathfrak A_{\rm flav}^{D}:=
\mathfrak M_5^{D}/\mathcal G_{\rm wb}\).
These are presentation quotients of the flavor coefficients, not claims that
the upstream carrier selects one orbit.  When the quark quotient is embedded
in full QCD, an anomalous axial rephasing must be accompanied by the
compensating shift of the topological/fermion-measure coordinate owned by the
downstream strong-CP interface.  The present quotient therefore does not
remove the physical invariant \(\bar\theta\), which is not one of its
coordinates.
\end{definition}

\begin{theorem}[Generic dimensions of the weak-basis flavor quotient]
\label{p6:thm:weak-basis-flavor-dimensions}
On the generic full-rank, nondegenerate strata,
\[
\boxed{
\dim_{\RR}\mathfrak A_{\rm flav}^{M}=31,
\qquad
\dim_{\RR}\mathfrak A_{\rm flav}^{D}=20.
}
\]
Equivalently, the generic quark quotient contributes ten real moduli; the
generic Majorana/type-I lepton quotient contributes twenty-one; and the
generic Dirac-neutrino lepton quotient contributes ten.  Degenerate,
rank-deficient, or symmetry-enhanced strata have different stabilizers and
must be counted separately.
\end{theorem}

\begin{proof}
The quark pair \((Y_u,Y_d)\) has \(36\) real coefficients.  The group
\(U(3)_Q\times U(3)_u\times U(3)_d\) has real dimension \(27\), and on a
generic pair its continuous stabilizer is the common baryon phase
\(U(1)_B\).  Hence the generic quark quotient has dimension
\[
36-(27-1)=10.
\]
For a Majorana/type-I lepton tuple
\((Y_e,Y_\nu,M_R)\), the raw real dimension is
\(18+18+12=48\).  The corresponding group
\(U(3)_L\times U(3)_e\times U(3)_N\) has dimension \(27\), and a generic
Majorana tuple has no continuous lepton-number stabilizer.  Thus its quotient
has dimension \(48-27=21\), giving \(10+21=31\) in total.  Without a
Majorana operator, the lepton tuple \((Y_e,Y_\nu)\) has \(36\) real
coefficients and a generic common lepton phase \(U(1)_L\) stabilizer, so its
quotient has dimension \(36-(27-1)=10\).  Together with the quark sector this
gives \(20\).  Discrete stabilizers do not change these generic dimensions.
\end{proof}

\begin{corollary}[Raw coefficient atlas versus physical flavor moduli]
\label{p6:cor:raw-versus-physical-flavor}
The equality
\(\dim_\CC\mathfrak M_5=42\) is a raw coefficient-space statement and is not
the physical flavor-assumption count.  Weak-basis presentation directions are
removed by Definition~\ref{p6:def:weak-basis-flavor-quotient}, while the
remaining generic quotient is still positive dimensional.  Consequently the
family carrier and its Weyl algebra do not select a physical flavor member.
The conditional package \(\mathsf F\) below is a finite construction/readout
atlas that maps into selected spectra and relative eigenframe data; it is not
identical to the complete quotient \(\mathfrak A_{\rm flav}^{M,D}\).
\end{corollary}

% BEGIN FLAVOR-QUOTIENT v5.48
\subsection{\texorpdfstring{The \(A_2\) family plane}{The A2 family plane}}
\label{p6:sec:a2-plane}

\begin{theorem}[Trace-separated \(A_2\) identities]
\label{p6:thm:a2-trace}
Let \(P_0,P_1,P_2\) be orthogonal rank-one projectors on \(\CC^3\) with
\(\sum_iP_i=I\), and define
\[
\tau_G(A)=\frac13\Tr_G(A),
\qquad
\Delta_i=P_i-\frac13I.
\]
Then
\[
\Tr_G(\Delta_i\Delta_j)=\delta_{ij}-\frac13
=
\begin{cases}
2/3,&i=j,\\
-1/3,&i\ne j,
\end{cases}
\]
whereas
\[
\tau_G(\Delta_i\Delta_j)
=\frac{\delta_{ij}}3-\frac19
=
\begin{cases}
2/9,&i=j,\\
-1/9,&i\ne j.
\end{cases}
\]
\end{theorem}

\begin{proof}
Expand \(\Delta_i\Delta_j\), use
\(P_iP_j=\delta_{ij}P_i\), \(\Tr P_i=1\), and \(\Tr I=3\).  Division by
three gives the normalized-trace formula.
\end{proof}

\begin{corollary}[Retained Hilbert--Schmidt invariants]
\label{p6:cor:a2-invariants}
For
\[
D_I=\diag(2/3,-1/3,-1/3),
\qquad
D_{II}=\diag(0,1,-1),
\]
ordinary trace gives
\[
\tau_D:=\Tr(D_I^2)=\frac23,
\qquad
\rho_D:=\Tr((D_I+D_{II})^2)=\frac83.
\]
The values survive after the manuscript explicitly separates normalized
trace from the Hilbert--Schmidt metric.
\end{corollary}

\subsection{Binary-icosahedral shell and embedded cells}
\label{p6:sec:icosahedral-cells}

Let \(I=2I\subset\SU(2)\), let \(\pi:I\to A_5\) be the central quotient,
and let
\[
S=\{g\in I\setminus\{1\}:\Re(g)\text{ is maximal}\}.
\]
The inherited \(H_4\) realization gives \(|S|=12\),
\(\Re(s)=\cos(\pi/5)\), and \(s^5=-1\).

\begin{theorem}[Nearest-shell no-go and generated subgroup]
\label{p6:thm:nearest-shell-no-go}
Put \(S^\pm=S\cup(-S)\).  Then
\begin{enumerate}[label=(\roman*),leftmargin=2em]
\item \(|S^\pm|=24\), but \(S^\pm\) is not a subgroup;
\item no element of \(S^\pm\) belongs to a binary tetrahedral subgroup of
\(I\); and
\item \(\langle S\rangle=I\).
\end{enumerate}
Thus central-sign closure of the nearest twelve directions does not produce
\(2T\).
\end{theorem}

\begin{proof}
The real parts on \(S\) and \(-S\) have opposite sign, so the two sets are
disjoint; neither contains \(1\).  Each \(s\in S\) has order ten and its
image in \(A_5\) has order five.  A binary tetrahedral subgroup projects to
\(A_4\), which has no element of order five.

Since \(S\) is conjugacy invariant, \(H=\langle S\rangle\) is normal in
\(I\).  Its nontrivial image is a normal subgroup of simple \(A_5\), hence
is all of \(A_5\).  Also \(s^5=-1\in H\), so the central kernel lies in
\(H\), and \(H=I\).
\end{proof}

\begin{theorem}[All binary tetrahedral subgroups]
\label{p6:thm:five-binary-tetrahedral}
For the natural five-point action, define
\[
H_j=\pi^{-1}(\operatorname{Stab}_{A_5}(j)),
\qquad j=1,\ldots,5.
\]
These are exactly the five binary tetrahedral subgroups of \(I\).  They are
conjugate, have order \(24\), and are self-normalizing.  For distinct
indices,
\[
H_i\cap H_j\cong C_6,
\qquad
H_i\cap H_j\cap H_k=\{\pm1\}
\quad(i,j,k\text{ distinct}).
\]
\end{theorem}

\begin{proof}
The five point stabilizers are the five \(A_4\) subgroups of \(A_5\);
their central preimages are the \(2T\) subgroups.  Stabilizer intersections
fix two or three points, giving respectively an order-three subgroup or the
identity in \(A_5\); central lifting gives \(C_6\) and \(\{\pm1\}\).
\end{proof}

% BEGIN HM v1.79 EX01
\begin{corollary}[The four cyclic Hopf axes as tetrahedral intersections]
\label{p6:cor:v179-cell-intersection-hopf-axes}
For a fixed \(H_i\) in Theorem~\ref{p6:thm:five-binary-tetrahedral},
the four intersections \(H_i\cap H_j\), \(j\ne i\), are exactly its four
order-six axis subgroups.  Hence there is a canonical incidence bijection
\[
 \{H_j:j\ne i\}
 \longleftrightarrow\{H_i\cap H_j:j\ne i\}
 \longleftrightarrow
 \{\text{four unoriented Hopf-axis lines of }H_i\}.
\]
The octahedral normalizer acts on these axes through \(S_4\), its
tetrahedral subgroup through \(A_4\), and \(2I\) moves the five \(H_i\)
globally.
\end{corollary}
\begin{proof}
If two intersections coincided, their order-six subgroup would lie in a
triple intersection, which is only \(\{\pm1\}\).  They are therefore
distinct.  The quotient \(H_i/\{\pm1\}\simeq A_4\) has exactly four
order-three subgroups, one for each unoriented tetrahedral axis.
Their binary preimages exhaust the order-six axis subgroups.
The normalizer and tetrahedral actions on the four body diagonals give
the stated permutation groups; conjugation gives the global \(2I\) action.
\end{proof}

This constructs the \(2I\)-to-\(2T\)-to-\(C_6\) incidence used by the
cyclic contact-stage support.  The subgroup alone does not select its
generator or orientation.  The exceptional relation remains the fork of
Paper I-A, Theorem~\ref{p1:thm:exceptional-normalizer-fork}:
\(2T\subset2O\), \(2T\subset2I\), \(2O\cap2I=2T\), and
\(N_{\mathrm{SU}(2)}(2T)=2O\).  In particular, \(2I/2T\) is a
homogeneous space, not a quotient group.
Each tetrahedral cell has heavy-stage type; the midpoint response of
Theorem~\ref{p3c:thm:v179-balanced-hierarchy-response} instead joins the
\(q=4\) and \(q=5\) Clifford blocks.
% END HM v1.79 EX01



\begin{corollary}[Correct onset boundary]
\label{p6:cor:correct-onset-boundary}
The count equalities
\[
|S|=12,\qquad 2|S|=24=|2T|,\qquad56-24=32
\]
do not by themselves define a selected \(2T\).  Instead, Paper~III-C's
derived \(q=4\) active realization fixes the \emph{type} \(2T\) and the
anchor \(K=8q=32\).  Its G0 occurrence supplies one conjugate representative
\(H_j\); alternatively one may use the full conjugacy orbit and prove
\(j\)-independence.  Thus \(H_j\mapsto K=32\) is no longer a
\(\CLKrem\) clause.  The displayed count equalities are compatibility checks,
not the proof of that map.
\end{corollary}

For \(i,j\in\{1,\ldots,5\}\), set
\[
\mathcal C_{ij}=\{g\in I:\pi(g)(j)=i\}.
\]

\begin{theorem}[Twenty-five-cell incidence system]
\label{p6:thm:cell-incidence}
The \(\mathcal C_{ij}\) are twenty-five distinct \(24\)-element translates
of binary tetrahedral subgroups.  Every vertex belongs to five cells and
\[
25\cdot24=120\cdot5.
\]
Moreover,
\[
|\mathcal C_{ij}\cap\mathcal C_{kl}|
=
\begin{cases}
24,&(i,j)=(k,l),\\
0,&i=k,\ j\ne l,\text{ or }i\ne k,\ j=l,\\
6,&i\ne k,\ j\ne l.
\end{cases}
\]
\end{theorem}

\begin{proof}
For fixed \(j\), the permutations sending \(j\) to \(i\) form a coset of
the order-twelve point stabilizer; its central preimage has order \(24\).
For every \(g\) and every \(j\), precisely one cell
\(\mathcal C_{\pi(g)(j),j}\) contains \(g\).  Two distinct source and target
requirements leave three even permutations, each with two central lifts.
Coincidence in only one coordinate contradicts bijectivity.
\end{proof}

\begin{theorem}[Diagonal and left--right symmetry]
\label{p6:thm:cell-actions}
The diagonal \(A_5\) action on labels has two orbits,
\[
\{(i,i)\},\qquad\{(i,j):i\ne j\},
\]
of sizes \(5\) and \(20\).  Its invariant label space is two-dimensional.
By contrast,
\[
(a,b)\cdot\mathcal C_{ij}
=a\mathcal C_{ij}b^{-1}
=\mathcal C_{\pi(a)(i),\pi(b)(j)}
\]
defines a transitive \(A_5\times A_5\) action with one invariant line.
\end{theorem}

\begin{proof}
Diagonal action preserves the relation \(i=j\) and is transitive on both
relation classes.  Independent transitivity in the two coordinates proves
the left--right statement.
\end{proof}

\begin{theorem}[Incidence-operator spectrum]
\label{p6:thm:incidence-spectrum}
For
\[
B:\CC^{25}\to\CC^{120},
\qquad
B|\mathcal C_{ij}\rangle
=\sum_{g\in\mathcal C_{ij}}|g\rangle,
\]
one has
\[
\Spec(B^*B)=120^{(1)}\oplus30^{(16)}\oplus0^{(8)}.
\]
Hence \(\rank B=17\), and the orthogonal label carrier is not unitarily
identical to the span of actual incidence vectors.
\end{theorem}

\begin{proof}
The intersection theorem gives
\[
B^*B=24I_{25}+6(J_5-I_5)\otimes(J_5-I_5).
\]
Since \(J_5-I_5\) has eigenvalues \(4^{(1)}\) and \((-1)^{(4)}\), its
tensor square has eigenvalues \(16^{(1)},(-4)^{(8)},1^{(16)}\).  Substitution
gives the result.
\end{proof}

% BEGIN HM v1.79 EX02
\begin{theorem}[Canonical metric on the effective incidence carrier]
\label{p6:thm:v179-incidence-polar-normalization}
For the cell-incidence map \(B\) of
Theorem~\ref{p6:thm:incidence-spectrum}, put
\[
 P=J_5/5,\quad Q=I_5-P,\quad
 P_{\rm glob}=P\otimes P,\quad P_{\rm con}=Q\otimes Q,\quad
 P_{\rm eff}=P_{\rm glob}+P_{\rm con}.
\]
Then
\[
 G_B=B^\dagger B=120P_{\rm glob}+30P_{\rm con},
 \qquad \operatorname{rank}P_{\rm glob}=1,\quad
 \operatorname{rank}P_{\rm con}=16.
\]
The polar normalization on its effective carrier is
\begin{equation}
\label{p6:eq:v179-incidence-polar-isometry}
 U_B=B\left(\frac{P_{\rm glob}}{\sqrt{120}}
           +\frac{P_{\rm con}}{\sqrt{30}}\right),\qquad
 U_B^\dagger U_B=P_{\rm eff}.
\end{equation}
Thus the global and contrast modes acquire their prescribed kinetic
normalization before any scalar response is extracted.
\end{theorem}
\begin{proof}
Theorem~\ref{p6:thm:incidence-spectrum} supplies the actual Gram
$120P_{\rm glob}+30P_{\rm con}$; its kernel is
$(P\otimes Q)\oplus(Q\otimes P)$. On its rank-seventeen support,
apply the positive inverse-square-root normalization of Paper~1B,
Proposition~\ref{p1b:prop:kinetic}~\cite{Kugelmann2026Paper1B}.
The two positive eigenvalues give the coefficients $1/\sqrt{120}$ and
$1/\sqrt{30}$ in \eqref{p6:eq:v179-incidence-polar-isometry}, with unit
$P_{\rm eff}$ on that support.
\end{proof}

This is a second rank-\(17\) construction, involving the \(2I\) incidence.
If an admitted action consumes \(B\) as its field source, this is its
normalized source, up to the action's fixed common prefactor.
It is not yet a field identification with the Hurwitz unital carrier of
Theorem~\ref{p6:thm:v179-hurwitz-unital-completion}.
The natural restriction to a selected Hurwitz cell has rank \(10\), and
orthogonal single-label exclusion does not descend through \(B\).
Appendix~\ref{p6:app:v179-incidence-descent} proves both statements and
distinguishes the corresponding normalized exclusion readouts.
% END HM v1.79 EX02



\begin{corollary}[Carrier dependence of self-exclusion]
\label{p6:cor:self-exclusion}
In the orthogonal label carrier,
\[
\Omega_{25}=\frac15\sum_{i,j}|\mathcal C_{ij}\rangle,
\qquad
\langle\Omega_{25},
(I-|\mathcal C_0\rangle\langle\mathcal C_0|)
\Omega_{25}\rangle=\frac{24}{25}.
\]
In the normalized vertex-incidence carrier,
\[
\Omega_{120}=\frac1{\sqrt{120}}\sum_{g\in I}|g\rangle,
\qquad
c_0=\frac1{\sqrt{24}}\sum_{g\in\mathcal C_0}|g\rangle,
\]
and
\[
|\langle\Omega_{120},c_0\rangle|^2=\frac15,
\qquad
\langle\Omega_{120},(I-|c_0\rangle\langle c_0|)
\Omega_{120}\rangle=\frac45.
\]
Thus \(24/25\) is exact only after the label representation is declared.
\end{corollary}


% BEGIN M8C v1.81 EXCLUSION-OWNER
\begin{theorem}[Claim-complete owner of the declared self-exclusion]
\label{p6:thm:v181-claim-complete-self-exclusion-owner}
Let
\[
 \Pi_H:=|\mathcal C_0\rangle\langle\mathcal C_0|,
 \qquad Q_{\rm exc}:=I-\Pi_H,
\]
and retain the typed owner
\[
 \mathfrak E_H=(\mathbb C^{25},\Omega_{25},\Pi_H,Q_{\rm exc}).
\]
Any comparison claim that consumes this already constituted owner has the
unique exclusion readout
\[
 \boxed{\eta_H(\mathfrak E_H)
 =\langle\Omega_{25},Q_{\rm exc}\Omega_{25}\rangle=\frac{24}{25}.}
\]
It is not permissible to replace this consumed owner by the incidence or
three-family compression while calling the result the same operation,
because Proposition~\ref{p6:prop:v179-label-exclusion-nondescent} proves that
$Q_{\rm exc}$ does not descend through the incidence map.  Appendix~\ref{p6:app:v181-exclusion-owner-completion}
records the least faithful completion when the normalized three-family
comparison is retained simultaneously.
\end{theorem}

\begin{corollary}[Conditional Higgs exclusion value]
\label{p6:cor:v181-conditional-higgs-exclusion-value}
If the physical Higgs/Yukawa same-action source word is proved to consume
$\mathfrak E_H$, then
\[
 \boxed{\eta_H=24/25}
\]
is forced for that word.  The rooted q=5 transport selector is already
closed upstream by the Paper-V-A $C_3$ source-cell theorem; this conditional
exclusion statement neither selects nor reopens matter/antimatter orientation.
Its remaining scope is the licensed physical Higgs/Yukawa source word itself.
\end{corollary}
% END M8C v1.81 EXCLUSION-OWNER

\label{p6:part:conditional-model}

\subsection{Binary-octahedral quotient-odd carrier}
\label{p6:sec:heavy-carrier}

\begin{theorem}[Normalizer and quotient]
\label{p6:thm:binary-octahedral-normalizer}
For a binary tetrahedral subgroup \(2T\subset\SU(2)\),
\[
N_{\SU(2)}(2T)=2O,
\qquad
2O/2T\cong C_2.
\]
Consequently the quotient permutation module
\[
\mathcal Q_O:=\CC[2O/2T]
\]
has the canonical character decomposition
\[
\mathcal Q_O
=
\Ran P_O^+\oplus\Ran P_O^-,
\qquad
P_O^\pm=\frac12(I\pm T_O),
\]
where \(T_O\) exchanges the two cosets.  Both summands have rank one, and
\(\Ran P_O^-\) is the sign line.
\end{theorem}

\begin{proof}
Projection to \(\SO(3)\) sends \(2T\) to the rotational tetrahedral group
\(A_4\).  Its rotational normalizer is the octahedral group \(S_4\), because
normalization permutes the four tetrahedral vertex axes and every such
orientation-preserving permutation is octahedral.  Taking the full inverse
image in \(\SU(2)\) gives \(2O\).  Since \(2T\) has index two in \(2O\), it is
normal and the quotient is \(C_2\).  The last statements are the spectral
decomposition of the regular \(C_2\)-module.
\end{proof}

\begin{corollary}[M8 defect provenance of the quotient-odd line]
\label{p6:cor:v188-heavy-sign-defect-provenance}
The quotient in Theorem~\ref{p6:thm:binary-octahedral-normalizer} is also the
abelianization of the normalizer:
\[
 \boxed{2O/[2O,2O]=2O/2T\cong C_2.}
\]
Therefore $\Ran P_O^-$ is the unique nontrivial real sign residue obtained by
restricting the parent $M_8$ finite order-swap defect to the generated
exceptional supports.  This supplies a common geometric provenance for the
quotient-odd sign.  It does not place that line in the protected hierarchy
owner, assign $K=16$, select a physical heavy packet, or supply a clock or
mass map.
\end{corollary}

\begin{proof}
Paper~I-A and Paper~III-B give $[2O,2O]=2T$, while
$\mathcal Q_O=\mathbb C[2O/2T]$ is already the regular module of the quotient.
The alternating summand is therefore exactly the nontrivial character of the
abelianized defect residue.  The final firewall follows because parent
ancestry does not identify independently typed hierarchy and physical-member
owners.
\end{proof}

\begin{corollary}[Conditional rank-three odd carrier]
\label{p6:cor:rank-three-heavy-carrier}
Under the Paper--V-A neutral-line and family-trimer clauses,
\[
\mathcal H_N^{\rm odd}
=
\Ran P_O^-\otimes\CC^3_{\rm fam}
\]
is a canonical rank-three quotient-odd carrier.  This proves the existence
and multiplicity of one minimal odd family packet.  It does not assign a
clock, a mass, or a family matrix.
\end{corollary}

\begin{theorem}[Classification of positive quotient operators]
\label{p6:thm:quotient-operator-classification}
Let \(X\) be self-adjoint, nonnegative, \(C_2\)-equivariant on
\(\mathcal Q_O\), and suppose \(X|_{\Ran P_O^+}=0\).  Then
\[
X=cP_O^-
\qquad(c\ge0).
\]
If \(X\ne0\) is additionally required to be an orthogonal projector, then
\(c=1\).  Thus
\[
T_O|_{\Ran P_O^-}=-1,\qquad
(I-T_O)|_{\Ran P_O^-}=2,\qquad
P_O^-|_{\Ran P_O^-}=1
\]
are three equally intrinsic normalizations of the same odd line.
\end{theorem}

\begin{proof}
Equivariance diagonalizes \(X\) on the two inequivalent one-dimensional
characters.  The even eigenvalue vanishes by hypothesis and positivity makes
the odd eigenvalue nonnegative.  Idempotence gives \(c^2=c\).
\end{proof}

\begin{corollary}[Meaning of \(x_N=1\)]
\label{p6:cor:heavy-projector-normalization}
The former number \(x_N=1\) is exactly the eigenvalue of the selected
orthogonal projector \(P_O^-\).  It is a coordinate normalization, not a
dimensionless mass theorem and not a threshold prediction.
\end{corollary}

\begin{definition}[Conditional heavy-family operator]
\label{p6:def:heavy-family-operator}
The flavor component of \(\mathsf F\) may select
\[
\widehat D_N
=
\diag\left(0,-\frac{55}{56},-\frac{57}{56}\right)
\]
and the affine operator
\[
\boxed{
M_N
=
P_O^-\otimes
\left(E_NI_3+\Delta_M\widehat D_N\right).
}
\]
Its three odd-sector eigenvalues are
\[
E_N,\qquad
E_N-\frac{55}{56}\Delta_M,\qquad
E_N-\frac{57}{56}\Delta_M.
\]
For \(E_N>57\Delta_M/56\ge0\) it is positive.  The anchor \(E_N\), splitting
\(\Delta_M\), diagonal family orientation, and the displayed affine
restriction are bridge data.  The most general self-adjoint
quotient-equivariant family operator is \(P_O^-\otimes H_N\) with arbitrary
\(H_N=H_N^*\), a nine-real-parameter family.
\end{definition}

\begin{corollary}[Heavy carrier versus heavy spectrum]
\label{p6:cor:heavy-carrier-spectrum-firewall}
The carrier
\[
\boxed{
\mathcal H_N^{\rm odd}=\Ran P_O^-\otimes\CC^3_{\rm fam}
}
\]
is theorem-level.  The affine spectrum data
\[
\boxed{(E_N,\Delta_M,\widehat D_N)}
\]
form a separate conditional heavy-spectrum realization family.  Quotient
oddness and the projector normalization do not determine a dimensionful
heavy mass or a family splitting.
\end{corollary}

\begin{proposition}[Normalizer multiplicity versus hierarchy-depth ancestry]
\label{p6:prop:heavy-clock-separation}
The normalizer quotient has the canonical real permutation multiplicity
\[
 Q_O^{\mathbb R}:=\mathbb R[2O/2T],
 \qquad
 \dim_{\mathbb R}Q_O^{\mathbb R}=2,
 \qquad
 Q_O^{\mathbb R}\otimes_{\mathbb R}\mathbb C\cong\mathbb C[2O/2T].
\]
This object must not be confused with the one-dimensional odd character line
\(\Ran P_O^-\).  The index \([2O:2T]=2\), the odd rank
\(\rank P_O^-=1\), and the protected upstream hierarchy depth are typed in
different categories.  In particular, the current normalizer construction
does not generate a morphism
\[
 \boxed{
 \mathsf{ANC}_O:
 \mathsf{QuotMult}_{2O/2T}\longrightarrow\mathsf{Mult}_{\rm hier}}
\]
placing \(Q_O^{\mathbb R}\) in Paper~III-C's physical protected
determinant-depth owner.

Hence the retained placement
\[
 \kappa_{\rm phys}(2O/2T):=16
\]
remains \(\mathsf{CLK}_O\), but its unresolved content is now isolated as the
typed ancestry/owner bridge \(\mathsf{ANC}_O\), not as a second determinant
normalization.  If a later theorem constructs \(\mathsf{ANC}_O\) and preserves
real rank, Paper~III-C's determinant-owner inheritance immediately gives
\(q_O=2\) and hence \(K_O=8q_O=16\) with no additional normalization choice.
Until that ancestry theorem exists, neither equal dimension nor the odd
projector rank closes the placement.  Even after the placement, \(E_N\) and
\(\Delta_M\) remain symbolic.

Moreover, Paper~II-A joint-family exhaustion does not reduce this residual:
in the present frozen grammar there is no generated cross-role relation or
common constructor between the quotient-multiplicity owner and the protected
depth owner.  Their joint zero-defect locus is therefore the relation-free
product, not the graph \(q_O=2\).  Inserting that graph without
\(\mathsf{ANC}_O\) would itself be a relation defect.
\end{proposition}

\begin{proof}
The quotient elements form a canonical real permutation basis, giving the
rank-two real multiplicity and its complexification displayed above.  The odd
projector instead selects one character line from the complex quotient module.
Paper~III-C's protected depth coordinate is generated from its own protected
resolution role; its determinant-rank theorem applies physically to an object
only after that object is typed in the corresponding multiplicity owner.
Paper~II-A's owner-typing firewall forbids inferring that ancestry from the
coincidence \(\dim_{\mathbb R}Q_O^{\mathbb R}=2\).  Thus the missing datum is
precisely the cross-role ancestry map \(\mathsf{ANC}_O\).  Once supplied, the
same-owner determinant theorem fixes the exponent and no second scale factor is
available.  Without it, using quotient rank to declare \(q=2\), or odd rank to
declare \(q=1\), is type-incorrect.

For the final joint-locus statement, apply Paper~II-A's relation-free
composition clause: the frozen presentation contains neither an ancestry map
nor any other generated relation joining the two owners, so every separately
admissible pair remains jointly admissible.  The proposed diagonal
\(q_O=\dim Q_O^{\mathbb R}\) is an ungenerated cross-role relation and has
nonzero GP2 relation defect.  This proves that no further reduction is
available from the current data.
\end{proof}

\subsection{Canonical terminal supports and the left-family quotient}
\label{p6:sec:v585-terminal-supports}
The scalar magnitudes of the charged raw compiler have an operator
realization on the nonzero supports of their actual terminal maps.
Throughout this section the common color identity factor is suppressed;
its multiplicity is not a family multiplicity.

\begin{theorem}[Canonical three-dimensional terminal support operator]
\label{p6:thm:v585-terminal-support-operator}
Let $F_r:\mathcal H_r^{\rm in}\to\mathcal H_r^{\rm out}$ be the three
nonzero rank-one terminal maps of a charged sector in
Definitions~\ref{p6:def:canonical-terminal-primitives}
and~\ref{p6:def:selected-typed-rows}, with unique singular magnitude
$\mu_r>0$. Their Moore--Penrose inverses define
\[
 P_r^{\rm in}=F_r^+F_r,\qquad P_r^{\rm out}=F_rF_r^+,
 \qquad V_r=F_r/\mu_r.
\]
For the role-tagged direct sum $T_s^{\rm term}=\bigoplus_rF_r$, put
\[
 \mathcal L_s=\bigoplus_r\operatorname{ran}P_r^{\rm in},\qquad
 \mathcal O_s=\bigoplus_r\operatorname{ran}P_r^{\rm out}.
\]
Both spaces have dimension three, and the support restriction has the
canonical polar decomposition
\[
 T_s^{\rm term}|_{\mathcal L_s}=V_s^{\rm term}\Sigma_s^{\rm term},\qquad
 V_s^{\rm term}=\bigoplus_rV_r,\quad
 \Sigma_s^{\rm term}=\sum_r\mu_rP_r^{\rm in}.
\]
Here $V_s^{\rm term}:\mathcal L_s\to\mathcal O_s$ is unitary, and
$(T_s^{\rm term})^\dagger T_s^{\rm term}|_{\mathcal L_s}
=\sum_r\mu_r^2P_r^{\rm in}$.
\end{theorem}
\begin{proof}
On its supports a rank-one map has the form
$F_r=\mu_r|l_r\rangle\langle r_r|$ with unit vectors. Hence
$V_r^\dagger V_r=P_r^{\rm in}$ and $V_rV_r^\dagger=P_r^{\rm out}$.
The three tagged summands are orthogonal, so their direct sum gives the
claimed unitary and positive diagonal factor, without an extra support choice.
\end{proof}

For the raw charged packets of
Theorem~\ref{p6:thm:v566-derived-charged-raw-compiler}, this construction is
explicit:
\begin{align*}
 T_u^{\rm term}&=N_2\oplus6I\oplus(\sqrt7+3)I,&
 &\mathbb C^3\longrightarrow\mathbb C^4,\\
 T_d^{\rm term}&=(\sqrt2+\sqrt3)I\oplus N_6^+
                 \oplus(\sqrt8+5)^{-1}I,&
 &\mathbb C^8\longrightarrow\mathbb C^3,\\
 T_e^{\rm term}&=(1+\sqrt5)^{-1}I\oplus N_4^+
                 \oplus(2\sqrt6+\sqrt{12})I,&
 &\mathbb C^6\longrightarrow\mathbb C^3.
\end{align*}
The support of $N_n$ on the output side is $\mathbb C\Omega_n$ with
normalized direction $\Omega_n/\sqrt n$; for $N_n^+$ it is the input side.
Every compressed packet is therefore exactly $3\to3$.
This establishes \texttt{CHARGED\_TERMINAL\_SUPPORT\_OPERATORIZATION}
as \texttt{DERIVED\_STRUCTURAL}. The tagged summands are mathematical role
summands, not asserted disjoint subsets of one Hurwitz carrier.

\begin{corollary}[Only left support incidence enters the charged quotient]
\label{p6:cor:v585-left-support-quotient}
If a source constructor supplies an isometry
$J_s:\mathcal O_s\to G_L$, the associated left return is
\[
 H_s=J_sT_s^{\rm term}(T_s^{\rm term})^\dagger J_s^\dagger.
\]
Coordinates on $\mathcal L_s$ cancel. Equivalently,
$U_s\Sigma_s^2U_s^\dagger$ is invariant under the simultaneous relabeling
$\Sigma_s\mapsto P\Sigma_sP^\dagger$, $U_s\mapsto U_sP^\dagger$.
\end{corollary}
\begin{proof}
A unitary change on the input multiplies $T_s^{\rm term}$ on the right
and cancels against its adjoint. Substitution proves the permutation formula.
\end{proof}

This is an exact role-tagged support/polar presentation, not native mass
genesis by singular values. The actual rooted charged source and incidence
packet is supplied by
Corollary~\ref{p6:cor:v590-preterminal-contraction-explicit} and
Corollary~\ref{p5:cor:v397-charged-preterminal-assembly-closure}.
The superseded directed matching ledger is placed in
Appendix~\ref{p6:app:v2187-historical-fermion-assembly}; its conditional
algebra does not reopen charged source assembly or identify a physical pole.
% END TERMINAL-SUPPORT-OPERATORIZATION v5.85


% BEGIN N56-ROLE-SPACE-RECLOSURE v5.87

\section{Protected hierarchy, Hurwitz scale and terminal energy}
\label{p5b:sec:organized-hierarchy}
\subsection{Hierarchy finite part and absolute-scale covariance}
\label{p6:sec:hierarchy-scheme}

\begin{definition}[Separated geometric defect and REG realization family]
\label{p6:def:separated-defects}
Write
\[
\delta_{\rm geom}:=\frac1{\sqrt{72}},
\qquad
\mathfrak A_{\rm REG}
=
\{\mathrm{REG}_{\sigma}:\sigma_{\rm fin}\in\RR\},
\]
\[
\delta_{\sigma}
:=
\frac1{6\sqrt2}+\sigma_{\rm fin},
\qquad
\log\alpha_{\sigma}:=4\pi^2+\delta_{\sigma}.
\]
The exact decreasing-energy domain is the inherited positive subfamily
\[
 \mathfrak A_{\rm REG}^{+}
 :=\left\{\mathrm{REG}_{\sigma}\in\mathfrak A_{\rm REG}:
 4\pi^2+\frac1{6\sqrt2}+\sigma_{\rm fin}>0\right\}.
\]
The geometric coordinate \(\delta_{\rm geom}\) is independent of the
fixed-shell finite part.  The source member is
\[
\boxed{
\mathrm{REG}_0:\quad\sigma_{\rm fin}=0,
\qquad
\delta_0=\frac1{6\sqrt2}.
}
\]
On the occurring Hurwitz/Hopf branch, Theorem~\ref{p3c:thm:occurring-hurwitz-reg-origin} selects \(\mathrm{REG}_0\). The unanchored family remains \(\mathfrak A_{\rm REG}^{+}\). Throughout
the source numerical ledger \(\alpha_{\rm hier}:=\alpha_0\); generic
scheme-labelled statements use \(\alpha_{\sigma}\).
\end{definition}

% BEGIN M8 v1.78 VI_IMPORT
\begin{remark}[Updated spectrum input partition]
\label{p6:rem:m8-batch-spectrum-inputs}
The occurring Hurwitz/Hopf hierarchy now imports the anchored
REG theorem of III-C.  The flavor member and its kinetic
normalization remain separate from this finite-part origin.
Running/matching and physical pole extraction retain their
own maps; the absolute ledger retains its calibration.
The geometric comparison of I-A and the conditional physical
criterion of V-A add no mass eigenvalue or family bilinear here.
Action derivatives consume the complete owner of the particular
spectrum claim, rather than an additional universal missing-action
premise.
\end{remark}
% END M8 v1.78 VI_IMPORT



\begin{theorem}[Exact finite-part covariance]
\label{p6:thm:finite-part-covariance}
If \(X(\alpha)=X_0\alpha^{-p}\), then for a shift
\(\sigma_{\rm fin}\mapsto\sigma_{\rm fin}+c\),
\[
X(\alpha_{\sigma+c})=e^{-pc}X(\alpha_{\sigma}).
\]
On the active G24 branch the upstream curvature scale is
\[
\boxed{K_p=\frac6{R_{\rm mom}^2}},
\qquad
\boxed{U=\frac3\pi K_p=\frac{18}{\pi R_{\rm mom}^2}}
\quad\text{on }(\mathrm{CELL}_{K\to U}),
\]
\[
\boxed{M_P:=8\hbar c_{\rm op}\sqrt U}
\quad\left(=8\sqrt{18/\pi}\,\frac{\hbar c_{\rm op}}{R_{\rm mom}}
\text{ on }(\mathrm{CELL}_{K\to U})\right),
\]
where \(K_p\) is the intrinsic momentum-sphere curvature,
\(U\) is the distinct scalarized cell curvature, and their equality up to
\(3/\pi\) is licensed only on the existing upstream action/geometry
comparison \((\mathrm{CELL}_{K\to U})\). Without that certificate the
\(U\)-owned Planck-unit definition remains meaningful, but its
\(R_{\rm mom}\)-rewrite is not an unconditional identity. Here
\(c_{\rm op}\) is the already inherited host metric rate.  Paper~V-B
does not require a photon-channel identification; after Paper~V-C proves
vacuum-photon principal ancestry one may equivalently write
\(c_{\rm light}^{\rm th}=c_{\rm op}\).
For the physical energy specialization, set
\(\alpha_{\rm hier}=\alpha_{\sigma}\) on a frozen
\(\mathrm{REG}_{\sigma}\in\mathfrak A_{\rm REG}^{+}\) member and work on
\((\mathrm{CRP4})+(\mathrm{MOMENERGY}_{0}^{\rm der})\).  The Paper--V-B
terminal unit is
\[
E_{32}:=\frac{M_P}{8\alpha_{\rm hier}}.
\]
For the owner-inherited shell law with the derived native reference
\[
E_K
=
E_{32}\alpha_{\rm hier}^{-(K-32)/32}
=
\frac{M_P}{8}\alpha_{\rm hier}^{-K/32},
\]
one has
\[
\boxed{E_K(c)=e^{-Kc/32}E_K(0).}
\]
In particular,
\[
\begin{array}{c|ccccc}
K&16&32&36&40&56\\
\hline
E_K(c)/E_K(0)
&e^{-c/2}&e^{-c}&e^{-9c/8}&e^{-5c/4}&e^{-7c/4}.
\end{array}
\]
The \(K=36\) entry assumes the declared logarithmic-midpoint bridge.
After the first dimensional definition, numerical energy formulas use
\(\hbar=c_{\rm op}=1\).  On the gravitational convention
\[
R_p=4\overline M_{\rm Pl}^{\,2},
\qquad
M_P=16\overline M_{\rm Pl}.
\]
Thus \(M_P\) is the series shell coordinate and is not the reduced Planck
mass appearing in the downstream scalaron coupling.
\end{theorem}

\begin{proof}
Substitute \(\alpha_{\sigma+c}=e^c\alpha_{\sigma}\).  The shell exponents are \(K/32\).
\end{proof}

\begin{theorem}[Momentum-rooted shell normalization]
\label{p6:thm:momentum-rooted-shell-normalization}
On \((\mathrm{CRP4})+(\mathrm{G24})\),
\[
 \boxed{E_0=\hbar c_{\rm op}\sqrt U
 =\sqrt{18/\pi}\,\frac{\hbar c_{\rm op}}{R_{\rm mom}},\qquad
 E_K=\sqrt{18/\pi}\,\frac{\hbar c_{\rm op}}{R_{\rm mom}}
 \alpha_{\rm hier}^{-K/32}.}
\]
Thus the absolute shell anchor is derived and the former independent
\(\mathrm{CAL}_{E_0}\) is retired on the native branch.  Pole extraction,
RG matching, flavor realization, and scheme selection remain separately typed.
\end{theorem}

\begin{corollary}[Exact energy--curvature/action ladder normalization]
\label{p6:cor:exact-energy-action-ladder-normalization}
For the same owner-homogeneous hierarchy,
\[
 \boxed{I_K=\left(\frac{E_K}{\hbar c_{\rm op}}\right)^4,\qquad
 \frac{\mathcal A_K}{U}=\left(\frac{E_K}{E_0}\right)^4,\qquad
 \frac{\mathcal R_K}{\sqrt U}=\left(\frac{E_K}{E_0}\right)^2.}
\]
The previous fourth-power exponent correspondence is therefore an absolute
normalization identity on the G24 branch, not merely a shared exponent law.
\end{corollary}

\begin{corollary}[Status of the absolute ledger]
\label{p6:cor:absolute-ledger-status}
Integer carrier dimensions, orbit counts, projector ranks, and ratios sharing
one shell exponent are invariant under \(c\).  Absolute masses and ratios
using different shells are not.  Setting the finite part to zero corresponds, relative to the source member,
to \(\sigma_{\rm fin}=-1/(6\sqrt2)\); this raises \(E_{32}\) by about
\(12.51\%\) and \(E_{56}\) by about \(22.90\%\).  Every native absolute
numerical ledger in this paper therefore carries the tag
\[
[\bigl(\mathrm{REG}_{\sigma}\in\mathfrak A_{\rm REG}^{+}\bigr)
 +(\mathrm{CRP4})+(\mathrm{MOMENERGY}_{0}^{\rm der})+\BridgeRaw].
\]
The legacy \((\mathrm{HT})\) label certifies the already derived neutral law
and proved owner inheritance; it adds no ancestry or normalization.
The finite pole-coordinate packet additionally carries
\([\mathsf{PV}_{\rm mat}]\); a physical pole comparison requires the full
\(\mathsf{PV}\) package.
\end{corollary}

\begin{corollary}[Charged-map covariance]
\label{p6:cor:charged-map-covariance}
For the selected charged absolute map and fixed shell coordinate
\(\widetilde n\),
\[
m_{\widetilde n}(c)
=
\exp\left[
-\frac{(14+1/12)c}{8(1+\widetilde n)}
\right]m_{\widetilde n}(0).
\]
Thus even charged-family ratios obtained from distinct shell coordinates
depend on the declared finite part.
\end{corollary}

\subsection{the native \texorpdfstring{\(E_{48}\)}{E48} carrier mass gap}
\label{p6:sec:v165-native-e48-mass-gap}

\begin{definition}[Native carrier mass-energy operator]
\label{p6:def:v165-native-carrier-mass-energy}
Fix the occurring active branch, a positive
\(\mathrm{REG}_{\sigma}\in\mathfrak A_{\rm REG}^{+}\), and the derived
dimension-weighted shell law.  On the finite active carrier define
\begin{equation}
 \label{p6:eq:v165-native-carrier-mass-energy}
 \boxed{
 H_{\rm car}(u)
 :=E_0\alpha_{\rm hier}^{-\mathcal Q_{\rm act}(u)/4}.}
\end{equation}
This is the functional-calculus grading constructed in the main-text
Theorem~\ref{p3c:thm:v179-balanced-hierarchy-response} and attached to
this energy character in
Corollary~\ref{p6:cor:v179-midpoint-action-attachment}.
Its absolute energy presentation uses the derived
\((\mathrm{MOMENERGY}_{0}^{\rm der})\) anchor; its spectral support and all
ratios are independent of the numerical unit realization.  Calling a
particular endpoint a particle rest mass still requires the pole/role bridge
of Definition~\ref{p6:def:typed-terminal-row-energy}.
\end{definition}

\begin{theorem}[Native \(E_{48}\)-centered carrier mass gap]
\label{p6:thm:v165-native-e48-carrier-mass-gap}
The GP1--GP3 carrier theorem of Paper~III-C and the strictly decreasing shell
character give
\begin{equation}
 \label{p6:eq:v165-native-carrier-energy-spectrum}
 \boxed{
 \operatorname{spec}H_{\rm car}
 =\{E_0,E_{32},E_{40},E_{56}\},
 \qquad
 \mathbf1_{(E_{56},E_{40})}(H_{\rm car})=0.}
\end{equation}
The open interval
\begin{equation}
 \label{p6:eq:v165-native-mass-gap-interval}
 \boxed{\mathcal I_{\rm gap}^{\rm car}:=(E_{56},E_{40})}
\end{equation}
is therefore a positive native carrier mass gap.  Its unique intermediate
allowed shell-lattice coordinate is the carrier-free point
\begin{equation}
 \label{p6:eq:v165-e48-gap-center}
 \boxed{
 E_{48}=E_0\alpha_{\rm hier}^{-3/2}
 =\sqrt{E_{40}E_{56}},
 \qquad P_{48}^{\rm act}=0.}
\end{equation}
Equivalently, the gap has exact additive width and multiplicative span
\begin{equation}
 \label{p6:eq:v165-native-mass-gap-width}
 \boxed{
 \Delta_{\rm car}^{\rm nat}
 :=E_{40}-E_{56}
 =E_{48}\!\left(\alpha_{\rm hier}^{1/4}
                 -\alpha_{\rm hier}^{-1/4}\right)>0,
 \qquad
 \frac{E_{40}}{E_{56}}=\alpha_{\rm hier}^{1/2}.}
\end{equation}
The framework predicts the gap because the carrier projector vanishes, not
because a particle of mass \(E_{48}/c_{\rm op}^2\) is placed at its centre.
\end{theorem}

\begin{proof}
Paper~III-C gives
\(\operatorname{spec}\mathcal Q_{\rm act}=\{0,4,5,7\}\) and
\(\mathbf1_{(5,7)}(\mathcal Q_{\rm act})=0\).  Applying the decreasing
positive function \(q\mapsto E_0\alpha_{\rm hier}^{-q/4}\) by finite
functional calculus gives
\(E(0),E(32),E(40),E(56)\) and no spectral support between the final two
values.  Since \(48/32=3/2\), direct evaluation gives \(E_{48}\); the
geometric-mean, width, and ratio identities follow from the shell exponents.
Positivity follows from \(\alpha_{\rm hier}>1\).
\end{proof}

\begin{corollary}[Profile invariance and source-member presentation]
\label{p6:cor:v165-native-gap-profile-invariance}
The gap projector and its shell labels are constant on the GP1--GP3-generated
heat-profile orbit
\((a,\Lambda)\sim(\rho a,\sqrt\rho\Lambda)\).  The profile action leaves the
exact process operator and \(\mathcal Q_{\rm act}\) unchanged, while the
energy ratios above depend only on the independently generated hierarchy
character.  On the displayed \(\mathrm{REG}_0\) and absolute-calibration
member,
\begin{equation}
 \label{p6:eq:v165-e48-numerical-coordinate}
 \boxed{E_{48}=78.1493227888\ {\rm eV}.}
\end{equation}
This decimal is a unit-presented gap centre, not an occupied particle row.
\end{corollary}

\begin{theorem}[No promotion to a continuous mass gap]
\label{p6:thm:v165-no-continuum-gap-promotion}
The native conclusion
\(\mathbf1_{(E_{56},E_{40})}(H_{\rm car})=0\) does not imply a positive
gap for a reconstructed continuum Hamiltonian.  Indeed, retaining the same
finite carrier factor and adjoining auxiliary refinement Hamiltonians
\[
 B_L^{(m)}=\operatorname{diag}(0,m),
 \qquad
 B_L^{(0)}=\operatorname{diag}(0,L^{-1})
\]
preserves \(P_{48}^{\rm act}=0\) at every level, while the auxiliary physical
gap has respectively positive infimum \(m\) and zero infimum.  Hence the
carrier-desert premises underdetermine every continuous-gap assertion.

The active framework therefore claims the finite structural carrier mass gap
of Theorem~\ref{p6:thm:v165-native-e48-carrier-mass-gap}, but no continuous
Yang--Mills mass-gap theorem.  A statement in Paper~VII-C transferring a
separately proved refinement-uniform lower bound remains a correct
conditional implication on an enlarged continuum branch; it is not an open
step in the proof of the native gap.
\end{theorem}

\begin{proof}
The two displayed families have identical native-carrier data and different
uniform-gap behavior.  Thus no function of the carrier projector, its missing
depth-six eigenspace, or the coordinate \(E_{48}\) can decide the continuum
infimum.  A continuum result would require additional refinement dynamics,
state convergence, scaling, and reconstruction data not contained in the
finite carrier theorem.
\end{proof}

\subsection{transverse bridge and adjacent half-step energies}
\label{p6:sec:v166-transverse-bridge-halfstep-energies}

\begin{theorem}[Exact \(E_{44}\), \(E_{48}\), and \(E_{52}\) structural
coordinate theorem]
\label{p6:thm:v166-e44-e48-e52-structural-coordinates}
On the occurring active branch, Paper~III-C derives
\[
 B_{48}(u)=P_0+\Pi_u,
 \qquad
 \delta_W(B_{48})=6,
 \qquad
 B_{48}\not\le P_{40}(u),
 \qquad
 P_{48}^{\rm act}=0.
\]
The dimension-weighted shell law then gives the exact adjacent half-step
coordinates
\begin{equation}
 \label{p6:eq:v166-e44-e52-geometric-midpoints}
 \boxed{
 E_{44}=E_0\alpha_{\rm hier}^{-11/8}
       =\sqrt{E_{40}E_{48}},
 \qquad
 E_{52}=E_0\alpha_{\rm hier}^{-13/8}
       =\sqrt{E_{48}E_{56}}.}
\end{equation}
On the displayed \(\mathrm{REG}_0\) and absolute-calibration member these
evaluate to
\begin{equation}
 \label{p6:eq:v166-e44-e52-numerical-coordinates}
 E_{44}\simeq11.0275833\ {\rm keV},
 \qquad
 E_{52}\simeq0.553821858\ {\rm eV}.
\end{equation}
Together with \(E_{48}\simeq78.1493228\,{\rm eV}\), these numbers are
coordinate readouts of the continuous shell character.  None is added to
\(\operatorname{spec}H_{\rm car}\) by this theorem.
\end{theorem}

\begin{proof}
The structural statements are
Theorems~\ref{p3c:thm:v166-unique-transverse-k48-bridge} and
\ref{p3c:thm:v166-transverse-diamond-halfsteps}.  Applying
\(E(K)=E_0\alpha_{\rm hier}^{-K/32}\) gives
\[
 \frac{44}{32}=\frac12\left(\frac{40}{32}+\frac{48}{32}\right),
 \qquad
 \frac{52}{32}=\frac12\left(\frac{48}{32}+\frac{56}{32}\right),
\]
which proves the geometric means.  The decimals are direct evaluations of
the already displayed \(\mathrm{REG}_0\) anchors.  The spectrum statement
follows from
\(\operatorname{spec}H_{\rm car}=\{E_0,E_{32},E_{40},E_{56}\}\).
\end{proof}

\begin{corollary}[No pole or mass-gap collapse from the transverse bridge]
\label{p6:cor:v166-transverse-bridge-no-pole}
The notation \(B_{48}\) records the trace valuation
\(8\delta_W(B_{48})=48\).  It does not make \(B_{48}\) a spectral projector
of \(\mathcal Q_{\rm act}\) at eigenvalue six or of \(H_{\rm car}\) at
energy \(E_{48}\).  Consequently
\[
 \mathbf1_{(E_{56},E_{40})}(H_{\rm car})=0
\]
and the native carrier mass gap are unchanged.  A physical pole at any of
\(E_{44},E_{48},E_{52}\) would still require a nonzero same-action propagator
pole and residue in the relevant target sector.
\end{corollary}




% BEGIN HM v1.79 HU02
\subsection{carrier-desert and endpoint-readout scope}
\label{p5b:sec:v565-carrier-desert-scope}

The terminal hierarchy should be read with the hard-carrier skeleton rather
than as seven consecutive particle masses.  In particular, $E_{40}$ is the
last relevant charged-carrier anchor and $E_{56}$ anchors the terminal neutral
sector.  The electron is a dressed descendant below $E_{40}$ rather than a new
hierarchy carrier.  Between $E_{40}$ and $E_{56}$ the native carrier projector
has no $K=48$ particle support:
\[
 \boxed{P_{48}^{\rm act}=0,\qquad E_{48}=\sqrt{E_{40}E_{56}}.}
\]
Thus $E_{48}$ is a balanced response/transition scale inside the native
particle-empty interval, not a rest mass.  Any logarithmic placement of
observed electron or neutrino readouts inside this corridor is diagnostic only
unless a separate carrier theorem is supplied.


% BEGIN CHARGED-PRIMITIVE-OWNER v5.66
\subsection{rank-preserving fermion corridor and gap typing}
\label{p6:sec:v581-fermion-rank-corridor}

\begin{proposition}[Rank-preserving middle corridor]
\label{p6:prop:v581-rank-preserving-corridor}
After an actual finite family source $T_s$ is supplied, each ordinary
fixed-support step used here---positive hierarchy dressing, canonical kinetic
normalization, finite regular RG transport and basis/scheme change---is
owner-equivalent transport in the sense of Paper~II-A.  Therefore
\[
 \boxed{\operatorname{rank}T_s(\mu)=\operatorname{rank}T_s(\mu_0)}
\]
on every finite regular interval.  These operations neither create a missing
family rank nor erase an existing one.
\end{proposition}

\begin{remark}[Four gap notions]
\label{p6:rem:v581-gap-taxonomy}
Keep distinct: (i) a finite source gap
$\sigma_{\min}(T)>0$; (ii) a local/operator gap
$0\notin\sigma(D_{A,\Phi})$; (iii) a whole Hamiltonian-sector gap; and
(iv) a physical two-point/pole mass certificate.  None is silently identified
with another.  In particular a massive charged fermion can coexist with a
gapless photon sector.
\end{remark}

\begin{remark}[Fermion frontier after rank inheritance]
\label{p6:rem:v581-fermion-frontier}
The structural question is localized to source/family compilation: if the
actual family source is rank deficient, ordinary normalization/hierarchy/RG
cannot repair it; if it is full rank, those regular transports preserve that
rank.  The physical pole remains the separately typed G16 analytic readout.
\end{remark}
% END FERMION-RANK-CORRIDOR v5.81


% BEGIN INVARIANT-FLAVOR-COMPILER v5.82
\subsubsection*{invariant flavor compiler and comparison-normalization quotient}
\label{p6:sec:v582-invariant-flavor-compiler}

\begin{theorem}[Exact spectral/frame factorization at the primitive-trimer witness]
\label{p6:thm:v582-spectral-frame-factorization}
At the witness $(A_*,B_*)$ of
Theorem~\ref{p6:thm:v548-trimer-local-submersion}, the two single-sector
spectral maps
\[
 A\mapsto(\tr A,\tr A^2,\tr A^3),\qquad
 B\mapsto(\tr B,\tr B^2,\tr B^3)
\]
have rank three.  In the ordered trimer basis
$\{I,\DtriI,\DtriII,C_1,C_2\}$ the minors on
$\{I,\DtriI,\DtriII\}$ are
\[
 \boxed{\det J_A=144},\qquad \boxed{\det J_B=72}.
\]
The corresponding isospectral tangent kernels admit the bases
\[
\begin{aligned}
 \delta A_1&=2\DtriI-\frac{11}{6}\DtriII+C_1,&
 \delta A_2&=-4\DtriI-\frac56\DtriII+C_2,\\
 \delta B_1&=7\DtriI+\frac{32}{3}\DtriII+C_1,&
 \delta B_2&=24\DtriI+25\DtriII+C_2.
\end{aligned}
\]
On these four fixed-spectrum directions the Jacobian of
\[
 (\tr AB,\tr AB^2,\tr A^2B,\tr A^2B^2)
\]
has determinant
\[
 \boxed{-15402240\ne0}.
\]
Consequently the full determinant factorizes exactly as
\[
 \boxed{144\cdot72\cdot(-15402240)=-159690424320},
\]
recovering the determinant in
Theorem~\ref{p6:thm:v548-trimer-local-submersion}.
\end{theorem}

\begin{corollary}[Six spectra plus four frame coordinates]
\label{p6:cor:v582-six-plus-four-flavor}
On the generic simple-spectrum stratum, fixing the three spectral invariants
in each quark sector leaves a smooth four-real-dimensional relative-family
fiber.  The four mixed moments
\[
 M_{11}=\tr(H_uH_d),\quad M_{12}=\tr(H_uH_d^2),\quad
 M_{21}=\tr(H_u^2H_d),\quad M_{22}=\tr(H_u^2H_d^2)
\]
are local coordinates on that fiber near the witness.  Hence generic quark
flavor decomposes locally as
\[
 \boxed{3_u+3_d+4_{\rm frame}=10.}
\]
In particular, two continuous $Q_0/Q_\perp$ response ratios plus discrete
cyclic offsets cannot be locally sufficient for generic CKM data after the
spectra are fixed.
\end{corollary}

\begin{proposition}[Vandermonde reconstruction of the CKM modulus matrix]
\label{p6:prop:v582-vandermonde-ckm-reconstruction}
Let
\[
 H_u=U_u\operatorname{diag}(x_1,x_2,x_3)U_u^\dagger,\qquad
 H_d=U_d\operatorname{diag}(y_1,y_2,y_3)U_d^\dagger,
\]
with simple spectra, and $V=U_u^\dagger U_d$.  Put
$P_{ij}=|V_{ij}|^2$ and
\[
 M_{mn}:=\tr(H_u^mH_d^n),\qquad 0\le m,n\le2.
\]
Then
\[
 M_{mn}=\sum_{i,j}x_i^m y_j^nP_{ij}.
\]
Writing $S_x=(x_j^m)_{m=0,1,2}$ and
$S_y=(y_j^n)_{n=0,1,2}$, one has
\[
 \boxed{P=S_x^{-1}MS_y^{-T}}.
\]
The entries of $M$ with $m=0$ or $n=0$ are already fixed by the two spectra,
so the four mixed moments in
Corollary~\ref{p6:cor:v582-six-plus-four-flavor} determine the CKM modulus
matrix.  The sign of
\[
 \chi:=\frac1{6i}\tr[H_u,H_d]^3
\]
retains the rooted CP orientation once the eigenvalue ordering is fixed.
\end{proposition}

\begin{theorem}[Common comparison-exponent shift is a normalization redundancy]
\label{p6:thm:v582-eta-g16-normalization-quotient}
On the optional comparison branch write a massive-row continuum coordinate as
\[
 m_p^S(\mu)=\frac{A_p\etaH^{n_p}}{C_S(\mu)},
 \qquad n_p\in\tfrac12\mathbb Z.
\]
For every common $k\in\tfrac12\mathbb Z$ the transformation
\[
 \boxed{n_p\mapsto n_p+k,\qquad C_S(\mu)\mapsto\etaH^kC_S(\mu)}
\]
leaves all such coordinates unchanged.  Therefore the old comparison atlas
contains only the relative class
\[
 [n]\in(\tfrac12\mathbb Z)^N/\Delta(\tfrac12\mathbb Z)
\]
as flavor data.  The common exponent baseline is not an additional physical
mass normalization.
\end{theorem}

\begin{corollary}[Near-unity reporting-point matching is not invariant before source normalization]
\label{p6:cor:v582-near-unity-not-invariant}
Before the same-action family compiler fixes the absolute native source
normalization, a statement such as $C_S(2\,{\rm GeV})\simeq1$ on the optional
historical ledger is representative-dependent under
Theorem~\ref{p6:thm:v582-eta-g16-normalization-quotient}.  Consequently the
numerical relation $E_{32}/\sqrt{240}\simeq2\,{\rm GeV}$ may be retained only
as a later matching candidate; it is not promoted here to a theorem or to
independent evidence for the historical exponent baseline.
\end{corollary}

\begin{remark}[Historical flavor frontier]
\label{p6:rem:v582-sharper-live-flavor-frontier}
\noindent\textbf{Historical scope.} The member-selection frontier below
belongs to the former eigenframe comparison target. It is superseded for
native rooted charged mass/source genesis by main-text Release v5.90,
Theorem~\ref{p6:thm:v592-terminal-mass-no-backgeneration}, and Paper~V-A
Corollary~\ref{p5:cor:v397-charged-preterminal-assembly-closure}.
Theorem~\ref{p6:thm:v548-trimer-local-submersion} and
Theorem~\ref{p6:thm:v582-spectral-frame-factorization} close the algebraic
capacity question.  The live physical question is same-action selection:
derive an occurring positive rooted-trimer pair $(H_u,H_d)$, or equivalently
its six spectral invariants and four mixed moments, from the actual fermion
source/Feshbach compiler.  No diagnostic terminal word, response ratio or
comparison exponent table is promoted to that member by this release.
\end{remark}
% END INVARIANT-FLAVOR-COMPILER v5.82

% BEGIN ANTIPRISM-SVD-ASSEMBLY v5.83
\subsubsection*{antiprism frame capacity and same-action SVD assembly frontier}
\label{p6:sec:v583-antiprism-svd-assembly}

\begin{definition}[Magnitude, frame and same-action assembly layers]
\label{p6:def:v583-magnitude-frame-assembly}
Separate the charged-family compiler into three typed layers:
\[
 \begin{gathered}
 \boxed{\text{structural magnitude grammar}}
 +\boxed{\text{noncentral frame grammar}}\\[3pt]
 +\boxed{\text{same-action co-ownership/assembly}}.
 \end{gathered}
\]
The existing terminal words provide structural magnitude data only.  Paper~I-C
v1.33 provides a coefficient-free full-rank antiprism frame primitive.  A
physical source member is obtained only after the action proves that the
relevant magnitude and frame outputs are co-owned by one sector source.
\end{definition}

\begin{proposition}[Positive-return target after co-ownership]
\label{p6:prop:v583-positive-return-target}
If a sector source co-owns a singular-value matrix $\Sigma_s$ and a left frame
$U_s$, then all charged mass eigenvalues and left-handed mixing readouts are
fixed by
\[
 \boxed{H_s=U_s\Sigma_s^2U_s^\dagger.}
\]
No independent right-handed unitary is required for these readouts.  This
proposition is a uniqueness statement after co-ownership and is not an
occurrence theorem.
\end{proposition}

\begin{theorem}[G04 exchange-odd slice has full local frame capacity]
\label{p6:thm:v583-g04-odd-frame-capacity}
Let $\mathcal T_{\rm loc}^{-}$ be the four-real-dimensional exchange-odd local
product-trimer slice of Paper~I-C-G04, with canonical basis represented by
\[
 D_{1,x}-D_{1,p},\quad D_{2,x}-D_{2,p},\quad
 C_{1,x}-C_{1,p},\quad C_{2,x}-C_{2,p}.
\]
At the witness of Theorem~\ref{p6:thm:v582-spectral-frame-factorization},
orthogonally project these four directions, using the inherited
Hilbert--Schmidt metric, onto the direct sum of the two fixed-spectrum tangent
kernels.  In the tangent bases displayed in that theorem the projection matrix
is
\[
 P_{\rm odd\to iso}=\begin{pmatrix}
 27/409&-105/409&325/818&41/818\\
 -57/409&-51/409&41/818&277/818\\
 501/3778&-255/1889&-1845/3778&363/1889\\
 -117/1889&171/7556&363/1889&-599/7556
 \end{pmatrix},
\]
with
\[
 \boxed{\det P_{\rm odd\to iso}=-\frac{11925}{6180808}\ne0.}
\]
Composing with the four-moment Jacobian of
Theorem~\ref{p6:thm:v582-spectral-frame-factorization} gives
\[
 \boxed{
 \det\frac{\partial(M_{11},M_{12},M_{21},M_{22})}
 {\partial(\text{G04 odd coordinates})}
 =\frac{22958964000}{772601}\ne0.}
\]
Hence the G04 exchange-odd local slice has full local capacity for the
generic fixed-spectrum quark-frame fiber at the witness.
\end{theorem}

\begin{remark}[Born-exchange/flavor typing firewall]
\label{p6:rem:v583-bornodd-flavor-firewall}
Theorem~\ref{p6:thm:v583-g04-odd-frame-capacity} is a capacity/comparison
theorem only.  The two reciprocal factors of G04 are not identified with the
up/down quark sectors by dimension or rank.  A physical use requires an
explicit same-parent action map
\[
 \boxed{\Xi_{BF}:\mathcal T_{\rm loc}^{-}
 \longrightarrow T_{(H_u,H_d)}\mathcal F_{\rm iso}.}
\]
No such typed Born-odd-to-flavor bridge is asserted in this release.
\end{remark}

\begin{remark}[Live family frontier after v5.83]
\label{p6:rem:v583-live-family-frontier}
The algebraic magnitude grammar, noncentral frame capacity and generic
$3+3+4$ flavor capacity are no longer the missing objects.  The physical
debt is
\[
 \boxed{\texttt{FERMION\_SAME\_ACTION\_SVD\_ASSEMBLY}},
\]
meaning the theorem that the physical finite action co-owns the appropriate
sector magnitude and frame outputs.  The historical $\eta_H$ baseline remains
quotiented against common G16 normalization as in Release v5.82 and is not
used to set the absolute native source scale.
\end{remark}
% END ANTIPRISM-SVD-ASSEMBLY v5.83


% BEGIN ROOTED-SOURCE-HANDOFF v5.84
\subsubsection*{primitive source-incidence contraction and family-assembly frontier}
\label{p6:sec:v584-rooted-source-handoff}

Paper~V-A v3.92 closes primitive fermion-source existence on the rooted
triality branch by generating a specific nonzero finite Dirac member and
placing it in the local same-action fermion owner.  Therefore the historical
v5.78 token
\[
 \texttt{G0\_FERMION\_SOURCE\_MEMBER}
\]
is no longer a live prerequisite for primitive source existence on that
branch.  Its generic member-occurrence logic is retained only for source
members not independently generated by a licensed constructor.

The live charged-family frontier is
\[
 \boxed{\texttt{FERMION\_SAME\_ACTION\_SVD\_ASSEMBLY}},
\]
namely proof that one sector source/action owner co-realizes the structural
magnitude data and the physical noncentral left frame, giving
\[
 \boxed{H_s=U_s\Sigma_s^2U_s^\dagger.}
\]
Spectra, relative CKM/Dirac-PMNS frames and weak-CP invariants remain dependent
readouts of the assembled member.  RG/threshold transport remains downstream;
pole conditions remain claim-scoped.

No terminal diagnostic word, antiprism capacity result, Born-odd comparison
slice, scalar VEV or RG normalization is promoted into the missing same-action
family assembly by this retyping.
% END ROOTED-SOURCE-HANDOFF v5.84


% Continue beyond the existing A--Z appendix series without renumbering it.
\setcounter{section}{0}
\renewcommand{\thesection}{A\Alph{section}}
\renewcommand{\theHsection}{extended.\arabic{section}}
\subsection{Physical clock map and comparison tensor}
\label{p6:sec:clock-comparison}

\subsubsection{Derived active anchors and the residual clock}

Paper~III-C derives the equivariant active Clifford rank-chain orbit
\[
8\longrightarrow4\longrightarrow3\longrightarrow1,
\qquad q\in\{0,4,5,7\},
\qquad K_{\rm sh}=8q\in\{0,32,40,56\},
\]
and identifies its carrier/action types as
\[
q=4:\ A^{\rm even}\text{ with }2T,\qquad
q=5:\ E_3(u)\text{ with }C_6(u)\curvearrowright E_3(u),\qquad
q=7:\ \mathbb R\mathbf1.
\]
The final identification is orbit-invariant and operator-level:
\[
 \mathbf1_{\{7\}}(\mathcal Q_{\rm act}(u))
 =P_0=\mathbb E_{\rm IR}.
\]
Thus the same G0-occurring Hurwitz/Hopf token that generates the active orbit
co-realizes its common infrared constituent; no second active or infrared
selector is present.  A particular \(u\) is not selected by G0.  On the
Paper--V-A oriented-link branch, its joint zero-defect graph supplies the
matching active member without a second diagonal mark.
These are active-support theorems, not particle-name assignments.  Paper~III
also supplies the basis-neutral trace interpolation
\[
q_{\rm sp}(\tau)=7(1-e^{-\tau}),
\qquad
K_{\rm sh}(\tau)=8q_{\rm sp}(\tau).
\]
Its integer crossings are the ordered neutral depth coordinate.  On
\(\mathrm{REG}_{\sigma}\in\mathfrak A_{\rm REG}^{+}\), and after the
imported \((\mathrm{CRP4})\) occurrence supplies the common-parent DETFRAME
identity of Theorem~\ref{p6:thm:crp4-c13-discharge}, the positive first-order
energy scale is the unique fourth root of the weight-one density transport:
\[
 \frac{E(K)}{E(0)}
 =\alpha_{\rm hier}^{-K/32}.
\]
No \((\mathrm{FLAG})\) is needed for these owner-typed scalar scale levels; a full
\(q=1,\ldots,7\) rank-one projector chronology still requires that
refinement.

\begin{theorem}[Canonical anchor export and residual non-export]
\label{p6:thm:clock-non-export}
The active-realization theorem defines the representative-independent partial
map
\begin{equation}
\label{p6:eq:derived-active-clock-map}
 \boxed{(A^{\rm even},2T)\mapsto32,\qquad
        [C_6\curvearrowright E_3]\mapsto40,
        \qquad\mathbb R\mathbf1\mapsto56.}
\end{equation}
It does not define a complete map
\[
\kappa_{\rm phys}:
\{\text{finite carriers or physical roles}\}\longrightarrow[0,56].
\]
In particular it does not place the quotient-odd \(2O/2T\) carrier at
\(K=16\) or assign particle names to any anchor.  The continuous
dimension-weighted neutral law is generated by Paper~III-C's dual character;
the physical energy law and scalar \(K=36\) midpoint additionally use
\((\mathrm{CRP4})\).  A subgroup index, an orbit size,
and a rank loss remain differently typed integers unless a theorem supplies
the connecting map.
\end{theorem}

\begin{proof}
Paper~III-C constructs the projector orbit from the grade involution, the
Hopf contact line, and the order-six Hurwitz-unit average, so its ranks and
finite support/action types are invariant under conjugating a calculation
representative.
Finite functional calculus gives the terminal projector \(P_0\), proving both
the partial map and the derived \((\mathrm{IR}_{O(3)})\) certification.  None
of those operators contains a
\(2O/2T\) quotient, an antipodal midpoint law, a particle slot, or an energy
interpolation; those conclusions therefore require separately typed data.
\end{proof}

\begin{theorem}[Exact finite split grammars at the derived anchors]
\label{p6:thm:derived-anchor-split-grammars}
At the derived \(q=4\) Hurwitz anchor, Paper~I-A's radial shell theorem gives
\[
 24=1+8+6+8+1,
 \qquad
 \mathcal G_{16}=\mathcal G_7\sqcup\mathcal G_9,
 \qquad
 \boxed{16=7+9}.
\]
The fixed unital completion retaining the global complement and its
scalar unit gives the exact count \(17=16+1\); its actual same-carrier
projector and inherited metric are constructed in
Theorem~\ref{p6:thm:v179-hurwitz-unital-completion}.  At the derived \(q=5\) hexagonal anchor, let \(R\) be the
regular \(C_6\) shift and \(J=R^3\).  The projectors
\[
 P_1=\frac16\sum_{k=0}^{5}R^k,
 \qquad P_3=\frac{I-J}{2},
 \qquad P_2=\frac{I+J}{2}-P_1
\]
have ranks \(1,3,2\) and sum to \(I\).  Hence
\[
 \boxed{6=5+1,\qquad5=3+2}.
\]
The radicands \(7,9,17\) used by the heavy packet and \(1,5,3,2\) used by
the light packet are therefore finite-support outputs rather than fitted
integer choices.  The fixed-word projector response and its native polar
operations are constructed in
Theorem~\ref{p6:thm:v179-projector-response-grammar}; the former
independent \(\mathsf{TMF}\) is superseded as explained in
Section~\ref{p6:sec:post-v125-supersession}.
The heavy and light split geometry fixes the raw terminal words for
\(H,t,b,e,u,d\) once the owner type is specified.  What remains is not a
radicand/row lookup choice but the generated physical same-action source and
family-member occurrence, with kinetic normalization retained.
\end{theorem}

\begin{proof}
The heavy identities are the radial-shell and global-complement decomposition
of Paper~I-A, Theorem~\ref{p1:thm:hurwitz-radial-shells}.  The scalar line is
unique by the Paper-III-B finite Hurwitz projection.  For the light identities,
\(J\) is an orthogonal involution with three \(+1\) and three \(-1\)
eigenvalues.  Thus \(P_3\) has rank three, \((I+J)/2\) has rank three, and
removing its constant line leaves \(P_2\) of rank two.  Orthogonality and the
sum identity prove the two-stage split.
\end{proof}

% BEGIN HM v1.79 HU01
\subsubsection{The Hurwitz scalar-adjoining word on its actual carrier}
\label{p6:sec:v179-hurwitz-unital-response}

\begin{theorem}[Minimal unital Hurwitz completion and inherited metric]
\label{p6:thm:v179-hurwitz-unital-completion}
Use the radial partition of
Theorem~\ref{p6:thm:derived-anchor-split-grammars} on
\(\mathcal H_{24}=\mathbb C^{2T}\), with its counting metric:
\[
 2T=\mathcal G_7(b)\sqcup\mathcal G_9(b)\sqcup\mathcal N_8^+(b),
 \qquad P_{16}=P_7+P_9=I-P_8.
\]
Let \(\Omega_{24}=\sum_{h\in2T}|h\rangle\),
\(\Pi_{\rm c}=|\Omega_{24}\rangle\langle\Omega_{24}|/24\), and
\(\Omega_8=P_8\Omega_{24}\).  The least unital pointwise
\(*\)-subalgebra of the already finite function algebra retaining every
function supported on \(\mathcal G_{16}\) is
\[
 \mathcal A_{17}
 =\mathbb C^{\mathcal G_{16}}+\mathbb C1_{2T}
 =\{f:f|_{\mathcal N_8^+}\text{ is constant}\}.
\]
Its orthogonal projector is the same-carrier join
\begin{equation}
\label{p6:eq:v179-hurwitz-unital-join}
 P_{17}=P_{16}\vee\Pi_{\rm c}
       =P_{16}+\tfrac18|\Omega_8\rangle\langle\Omega_8|,
 \qquad \operatorname{rank}P_{17}=17.
\end{equation}
As a function operator, \(P_{17}\) fixes the sixteen singleton values and
averages the eight values on \(\mathcal N_8^+\).  It is a positive unital
counting-preserving conditional expectation onto \(\mathcal A_{17}\).
The seventeen atom squared norms are \(1,\ldots,1,8\); the unit still has
squared norm \(24\).
\end{theorem}
\begin{proof}
Every unital algebra retaining the sixteen singleton indicators contains
\(1_{\mathcal N_8^+}=1_{2T}-\sum_{h\in\mathcal G_{16}}1_{\{h\}}\).
These seventeen disjoint indicators form a closed pointwise algebra and
are independent, proving leastness and dimension.  The new orthogonal
line is \(\mathbb C\Omega_8\), of squared norm eight, giving the join.
The displayed averaging is positive, preserves the total sum and satisfies
\(P_{17}(af)=aP_{17}f\) for \(a\in\mathcal A_{17}\); the atom metric follows.
\end{proof}

The scalar-unit map from a point to the finite function carrier is retained
explicitly.  It is not an identification of \(\Pi_{\rm c}\) with a Clifford
projector on another carrier.  GP2 retains this generated unit/complement
diagram, GP1 retains its finite envelope, and GP3 proves leastness for this
fixed unital-completion problem.  The full parent and its other outputs
remain.  No unitality rule is imposed on arbitrary subspaces, and no new
pointwise multiplication is inserted into the momentum-seed grammar.
Left translation transports \(b\) and the complete partition, giving
\(U_gP_{17}(b)U_g^\dagger=P_{17}(gb)\).

\begin{theorem}[Projector response and the cyclic split]
\label{p6:thm:v179-projector-response-grammar}
For a nonzero generated orthogonal projector \(P\), with inherited
Hilbert--Schmidt pairing on its operator space, define
\(\mathcal N_P(z)=zP\).  If \(d=\operatorname{rank}P\), then
\[
 \mathcal N_P^\dagger\mathcal N_P=dI,\qquad
 |\mathcal N_P|=\sqrt d\,I,\qquad
 \mathcal N_P^+(X)=\operatorname{Tr}(PX)/d,\qquad
 \sigma(\mathcal N_P^+)=d^{-1/2}.
\]
The mixed Gram is \(\operatorname{Tr}(PQ)\).
Denote by \(C_1,C_3,C_2\) the rank-\(1,3,2\) cyclic projectors already
constructed in Theorem~\ref{p6:thm:derived-anchor-split-grammars}.
Their operator-source magnitudes are \(1,\sqrt3,\sqrt2\), respectively.
\end{theorem}
\begin{proof}
The source adjoint is \(\mathcal N_P^\dagger X=\operatorname{Tr}(PX)\).
Idempotence gives its squared norm \(d\), and the positive square root and
Moore--Penrose formula give the claims.  Apply that identity to the
already proved orthogonal cyclic resolution.
\end{proof}

An orthogonal common-source column has magnitude \(\sqrt{d+e}\);
the positive scalar codiagonal of two magnitudes is \(\sqrt d+\sqrt e\).
These are different licensed words.  Inversion of a dimensionless
magnitude precedes attachment of its energy unit.  With the already
derived heavy, cyclic and scalar clock maps, these words attach to
\(E_{32},E_{40},E_{56}\), respectively.  In particular,
\[
 \|E_{32}P_{17}\|_{\rm HS}^2=17E_{32}^2,\qquad
 \|E_{40}C_j\|_{\rm HS}^2=E_{40}^2\operatorname{rank}C_j .
\]
The rank \(17\) is the projector norm, not the norm of the algebra unit.
Its normalized source is \(P_{17}/\sqrt{17}\), with normalized fourth
moment \(1/17\).  Physical mass use consumes the kinetic and potential
forms together through
Theorem~\ref{p3f:thm:v179-normalized-scalar-hessian}.
Appendix~\ref{p6:app:v179-projector-diagnostics} records why the direct
flat-projector Higgs identification fails, without discarding the
positive unital construction.
% END HM v1.79 HU01



\begin{corollary}[Canonical counting measure on the \(16=7+9\) shell complement]
\label{p6:cor:G16-canonical-counting-measure}
At the derived \(q=4\) Hurwitz anchor,
\(\mathcal G_{16}=\mathcal G_7\sqcup\mathcal G_9\).  Paper~I-A's
tetrahedral shell theorem gives
\begin{equation}
\label{p6:eq:G16-counting-atom}
 \boxed{\omega_{2T}=\frac1{16}=\frac1{|\mathcal G_{16}|}.}
\end{equation}
Therefore normalized counting measure has block weights
\begin{equation}
\label{p6:eq:G16-counting-block-weights}
 \boxed{w_7^{\rm count}=\frac7{16},\qquad
 w_9^{\rm count}=\frac9{16}.}
\end{equation}
This corollary does not turn counting measure into a thermal preparation,
response amplitude, branching fraction, or mass law.
\end{corollary}

\begin{theorem}[Derived dimension-weighted shell law]
\label{p6:thm:derived-dimension-weighted-shell-law}
Fix \(\mathrm{REG}_{\sigma}\in\mathfrak A_{\rm REG}^{+}\), and assume
\((\mathrm{CRP4})+(\mathrm{MOMENERGY}_{0}^{\rm der})\).  Let the master
determinant-depth readout have weight one.  Paper~III-C's generated neutral
dual character, its owner-typed inheritance theorem at weight \(1/4\), and
\(K=8q\) then give
\begin{equation}
\label{p6:eq:continuous-shell-law}
E(K)
=
E_0\alpha_{\rm hier}^{-K/32}
=E_{32}\alpha_{\rm hier}^{-(K-32)/32}.
\end{equation}
For each such fixed upstream determinant-shell scheme, all physical energy
ratios are derived and
there is no shell-by-shell interpolation freedom.  Their numerical value
changes only if the upstream REG finite part is changed.  Every absolute value
of \(E(K)\) additionally remains conditional on the reference calibration in
\(\Uabs\).
\end{theorem}

\begin{proof}
Paper~III-C generates the neutral ratio
\(\chi_{\sigma,w}(q)=\alpha_{\rm hier}^{-wq}\).
The common-parent DETFRAME identity and
Theorem~\ref{p6:thm:physical-cone-fourth-power-density} force
\(E(q)/E(0)=\chi_{\sigma,1}(q)^{1/4}\).  Substitute \(q=K/8\).
Re-expressing the same law relative to the \(K=32\) anchor gives the second
equality.  The reference value cancels from every ratio, so no shellwise
normalization remains.
\end{proof}

\begin{theorem}[Native carrier desert at \(K=48\)]
\label{p6:thm:native-k48-carrier-desert}
At the generated smooth crossing
\begin{equation}
 \label{p6:eq:k48-smooth-crossing}
 q_{\rm sp}=6,
 \qquad K_{\rm sh}=48,
 \qquad N_W=1,
\end{equation}
Paper~III-C supplies the normalized residual density
\(D_*=\frac17P_W\), with \(\Tr_WD_*=1\) and
\(D_*^2=\frac17D_*\).  It is not an orthogonal projector.  On the generated
active orbit,
\begin{equation}
 \label{p6:eq:k48-active-projector-zero}
 \boxed{P_{48}^{\rm act}(u)
 :=\mathbf1_{\{6\}}(\mathcal Q_{\rm act}(u))=0.}
\end{equation}
Paper~III-C additionally derives the transverse structural projector
\[
 B_{48}(u)=P_0+\Pi_u,
 \qquad
 \delta_W(B_{48})=6,
 \qquad
 B_{48}\not\le\Pi_3(u).
\]
The subscript records its common trace-defect placement; \(B_{48}\) is not
\(\mathbf1_{\{6\}}(\mathcal Q_{\rm act})\), not an eigenspace of
\(H_{\rm car}\) at \(E(48)\), and not a one-particle carrier.
Thus \(E(48)=E_0\alpha_{\rm hier}^{-3/2}\) is a valid continuous shell-energy
coordinate, but the native construction supplies no active one-particle
carrier at that coordinate.  By Paper~IV's absent-target theorem, quadratic
evolution internal to the constituted carrier has zero production into
\(P_{48}^{\rm act}\).  Any later effective-curvature or cosmic \(K=48\)
response is a separately typed realization whose occurrence and matching
must be proved; it is not a native particle sector.
\end{theorem}

\begin{proof}
Equation~\eqref{p6:eq:k48-smooth-crossing} and the properties of \(D_*\) are
Corollary~\ref{p3c:cor:primitive-final-octet-residual}.  The active projector
identity is Theorem~\ref{p3c:thm:active-depth-six-obstruction}.  Evaluating
\eqref{p6:eq:continuous-shell-law} at \(K=48\) gives the energy coordinate,
and Theorems~\ref{p3c:thm:v166-unique-transverse-k48-bridge} and
\ref{p3c:thm:v166-transverse-diamond-halfsteps} give the transverse bridge
and its nonnesting.  Corollary~\ref{p4:cor:absent-target-production} gives the
internal production statement.
\end{proof}

\begin{corollary}[Carrier-free \(K=48\) scale]
\label{p6:cor:carrier-free-k48-scale}
The data \(q_{\rm sp}=6\), \(N_W=1\), \(D_*\), \(B_{48}\), \(K=48\), and
\(E(48)\) are mutually compatible neutral structural/trace/scale readouts.
None defines a rest mass, occupation number, field, or particle species
because the required native active target projector is zero.  The derived
\(B_{48}\) role is transverse and therefore does not count as an externally
adjoined carrier.  Any actual external particle carrier would still require
new ancestry and occurrence data.
\end{corollary}

\begin{definition}[Residual Paper--V-B clock-to-role realization family]
\label{p6:def:complete-clock-bridge}
After Theorem~\ref{p6:thm:derived-dimension-weighted-shell-law},
\(\mathfrak A_{\rm CLK}\) is the family of compatible maps from retained
carrier/interface/particle roles to coordinates of the already-derived shell
law.  The retained member \(\CLKrem\) supplies the \(K=16\) quotient-carrier
placement and the interface-packet role at the already-defined value \(K=36\).
Particle packet placements at the active anchors belong to \(\mathsf{SLOT}\),
while the no-role entries at \(K=8,24\) remain Paper--V-B declarations.  The
native \(K=48\) carrier absence is
Theorem~\ref{p6:thm:native-k48-carrier-desert}, not a member of \(\CLKrem\).
Neither the continuous scale law nor its logarithmic midpoint is a member of
\(\CLKrem\).
\end{definition}

\begin{table}[H]
\centering
\caption{Derived anchors and residual conditional clock-to-role map.}
\label{p6:tab:clock-map}
\small
\renewcommand{\arraystretch}{1.10}
\begin{tabularx}{\textwidth}{@{}L{0.18\textwidth}L{0.27\textwidth}Y@{}}
\toprule
\(K_{\rm sh}\) & Retained Paper--V-B role & Logical status \\
\midrule
\(8\) & none & derived neutral depth level; physical energy uses \((\mathrm{CRP4})\); no-role status is a Paper--V-B declaration \\
\(16\) & quotient-odd \(2O/2T\) carrier & independent heavy clock bridge;
the neutral depth level exists independently of that carrier assignment \\
\(24\) & none & derived neutral depth level; no-role status is a Paper--V-B declaration \\
\(32\) & occurring \(2T\) active anchor; \(W,Z,H,t,b\) support packet
& \((A^{\rm even},2T)\mapsto32\) derived; particle packet is \(\mathsf{SLOT}\) data \\
\(36\) & \(\mu,s,c,\tau\) interface packet
& neutral ratio is the derived logarithmic midpoint; its physical energy uses
\((\mathrm{CRP4})\); packet role is \(\mathsf{SLOT}\) data \\
\(40\) & matching \(C_6\curvearrowright E_3\) active/link stage; \(e,u,d\) packet
& \([C_6\curvearrowright E_3]\mapsto40\) derived; particle packet is \(\mathsf{SLOT}\) data \\
\(48\) & no active carrier; transverse \(2D_3/C_6\) normalizer bridge
& \(q_{\rm sp}=6\), \(N_W=1\), \(B_{48}=P_0+\Pi_u\), and \(E(48)\) are
derived, while \(P_{48}^{\rm act}=0\).  The bridge label is a trace-defect
placement, not an \(H_{\rm car}\) eigenvalue.  A downstream reflected cosmic response is separately
typed and conditional; it is not an \(n_6\) action, scalaron particle, or
rank-48 species assignment \\
\(56\) & scalar reclosure anchor; light-neutral radical packet
& scalar \(\mapsto56\) derived; particle packet is \(\mathsf{SLOT}\) data \\
\bottomrule
\end{tabularx}
\end{table}

\begin{remark}[The \(K=48\) carrier-desert/cosmic-response firewall]
\label{p6:rem:k48-cosmological-readout-firewall}
The zero projector in
Theorem~\ref{p6:thm:native-k48-carrier-desert} is a native generated-carrier
statement.  The later cosmological label
\[
 j_{\rm cDM}=6,\qquad K_{\rm cDM}=8j_{\rm cDM}=48,
 \qquad J(2)=8-2=6,
\]
is the protected-octet reflection of the depth-two dressing label.  If its
separate occurrence and matching premises are discharged, it indexes an
effective curvature-fluid response; it introduces neither a sixth Laurent
action channel \(n_6\), nor a rest mass, occupation number, nor particle
species.  This does not prohibit an externally adjoined dark sector, but no
such sector is generated by the present native carrier chain.
\end{remark}

\begin{theorem}[Logarithmic interpolation family]
\label{p6:thm:logarithmic-interpolation}
For every real \(a\), let \(K_a=32+8a\).  Then
\[
E(K_a)=E_{32}^{1-a}E_{40}^{a}.
\]
In particular,
\[
E(36)=\sqrt{E_{32}E_{40}}.
\]
The identity is an exact consequence of the generated dual character.  It
does not assign a particle carrier or particle name to an intermediate value.
\end{theorem}

\begin{proof}
Since \(E_{40}=E_{32}\alpha_{\rm hier}^{-1/4}\),
\[
E_{32}^{1-a}E_{40}^a
=E_{32}\alpha_{\rm hier}^{-a/4}
=E(32+8a).
\]
\end{proof}

% BEGIN HM v1.79 EN01
\begin{corollary}[The midpoint as a represented Clifford response]
\label{p6:cor:v179-midpoint-action-attachment}
The quarter-character in
Theorem~\ref{p6:thm:derived-dimension-weighted-shell-law} is represented on
the active Clifford depths by \(H_{\rm car}\) in
Theorem~\ref{p3c:thm:v179-balanced-hierarchy-response}.
Its multiplication source \(T_u(v)=Q_5L_vQ_4\) and self-adjoint completion
\(S_u(v)\) give
\[
 H_{\rm car}^{1/2}S_u(v)H_{\rm car}^{1/2}=E_{36}S_u(v),\qquad
 |\,v\mapsto E_{36}S_u(v)\,|=\tfrac12E_{36}I_V .
\]
This realizes the logarithmic midpoint as a specified operator word with
the inherited octet trace.  The complete even-operator envelope also
contains the \(E_{44}\) and \(E_{48}\) weights of
Corollary~\ref{p3c:cor:v179-complete-balanced-weights}.
\end{corollary}
\begin{proof}
The character has the same endpoint values on the same active depth
projections.  Its functional calculus and the positive two-sided square
root therefore coincide with the cited Clifford construction.
\end{proof}

The actual carrier-orbit derivative is
\((E_{32}-E_{40})S_u(v)\), with the exact conversion in
\eqref{p3c:eq:v179-carrier-orbit-conversion}.
The normalized source is \(2S_u\).  A charged-lepton interpretation must
reach the chiral operator of
Corollary~\ref{p5:cor:v179-canonical-fermion-comparison} with its kinetic
pairings.  The bare midpoint supplies neither that map nor a new active
projector event.
% END HM v1.79 EN01



Let
\[
2O\setminus2T=X_+\sqcup X_0\sqcup X_-,
\qquad
(|X_+|,|X_0|,|X_-|)=(6,12,6),
\]
where the antipodal involution exchanges \(X_+\) and \(X_-\) and preserves
\(X_0\).

\begin{theorem}[Conditional antipodal midpoint]
\label{p6:thm:antipodal-midpoint}
Assume the clock bridge is constant on the three signed-axis orbits and
satisfies
\[
k(-x)=72-k(x),
\qquad
k(X_+)=32.
\]
Then
\[
k(X_-)=40,\qquad k(X_0)=36.
\]
\end{theorem}

\begin{proof}
The first equality is the reflection law.  Transitivity makes \(k\) constant
on \(X_0\), while antipodal preservation of that orbit gives
\(k(X_0)=72-k(X_0)\).
\end{proof}

\begin{remark}[Midpoint status]
The orbit decomposition is exact finite mathematics.  The affine reflection
law that assigns the neutral orbit \(X_0\) to \(K=36\) is \(\CLKrem\) data;
the endpoint anchors \(32,40\), the continuous scale law, and its value
\(E_{36}=\sqrt{E_{32}E_{40}}\) are derived upstream.  Thus \(E_{36}\) is a
generated continuous hierarchy scale, but \(K=36\) is not one of the seven
integer protected hitting levels.  Its particle-interface role remains
conditional.
\end{remark}


\section{Native mass grammar and exact terminal coefficients}
\label{p5b:sec:organized-mass}
\subsection{Terminal transfers and the finite mass grammar}
\label{p6:sec:terminal-transfers}

\begin{theorem}[Classification of equivariant terminal transfers]
\label{p6:thm:terminal-transfer-classification}
Let a finite group \(K\) act on
\(
S=S_1\sqcup\cdots\sqcup S_r
\)
with orbit sums
\[
u_a=|S_a|^{-1/2}\sum_{s\in S_a}|s\rangle.
\]
Then
\[
\HomG_K(\CC,\CC[S])
=\CC[S]^K
=\bigoplus_{a=1}^r\CC u_a.
\]
Every equivariant \(T:\CC\to\CC[S]\) has
\[
T(1)=\sum_a\lambda_au_a,
\qquad
\sigma(T)=\left(\sum_a|\lambda_a|^2\right)^{1/2}.
\]
For a transitive support, symmetry fixes the invariant line but not its
magnitude.
\end{theorem}

\begin{proof}
An invariant vector is constant on each orbit.  The normalized orbit sums
are an orthonormal basis of the invariant space, and the sole singular value
of a map from \(\CC\) is the norm of its image of \(1\).
\end{proof}

\begin{corollary}[Three natural but inequivalent normalizations]
\label{p6:cor:terminal-normalizations}
For a transitive support of size \(n\),
\[
\begin{array}{c|c}
T(1)&\sigma(T)\\
\hline
\sum_{s\in S}|s\rangle&\sqrt n\\
n^{-1/2}\sum_{s\in S}|s\rangle&1\\
n^{-1}\sum_{s\in S}|s\rangle&n^{-1/2}.
\end{array}
\]
All three maps are permutation-equivariant and basis-free relative to the
counting carrier.
\end{corollary}

\begin{definition}[Terminal-mass realization family]
\label{p6:def:tmf-realization-family}
Before a terminal-mass functor is frozen, the compatible realization family
contains independent-domain versus common-domain assembly, relative path
phases \(e^{i\theta_i}\), compatible Hilbert metrics and pseudoinverses, and
dimensionful shell-insertion rules.  For transitive counting carriers,
Theorem~\ref{p6:thm:lattice-transfer-radical} reduces the node-normalization
problem: on the atomic/primitive integral branch, the unsigned magnitude is
forced to be the first minimum of the equivariant Hom-lattice and no positivity
clause remains.  The retained \(\mathsf{TMF}\) member fixes this atomicity
condition together with the remaining common-domain/codiagonal semantics,
additive-path phase law, and metric/insertion rule.
\end{definition}

\begin{definition}[Invariant integral terminal Hom-lattice]
\label{p6:def:primitive-lattice-transfer}
For a transitive finite \(K\)-set \(S\), equip \(\ZZ[S]\) with its counting
inner product and define
\[
\mathcal L_S
:=
\operatorname{Hom}_{\ZZ K}(\ZZ,\ZZ[S]).
\]
This lattice is fixed by the counting carrier and the group action; it is not
a terminal-normalization selector.
\end{definition}

\begin{theorem}[First-minimum lattice-transfer radical]
\label{p6:thm:lattice-transfer-radical}
Let \(N_S(1)=\Omega_S:=\sum_{s\in S}|s\rangle\).  Then
\[
\mathcal L_S=\ZZ N_S,
\qquad
\min_{0\ne T\in\mathcal L_S}\sigma(T_\CC)=\sqrt{|S|}.
\]
The minimizers are exactly \(\{N_S,-N_S\}\).  Equivalently, a nonzero
equivariant lattice map is primitive if and only if it has minimal nonzero
singular value, and for either primitive orientation
\[
\qquad
\sigma((N_S)_\CC)=\sqrt{|S|}.
\]
\end{theorem}

\begin{proof}
Transitivity gives
\(
\ZZ[S]^K=\ZZ\sum_s|s\rangle
\).
Thus every \(T\in\mathcal L_S\) has
\(T(1)=m\Omega_S\) for a unique \(m\in\ZZ\), proving
\(\mathcal L_S=\ZZ N_S\).  Its sole complex singular value is
\(|m|\|\Omega_S\|=|m|\sqrt{|S|}\).  Among nonzero integers the minimum of
\(|m|\) is one, attained exactly at \(m=\pm1\); these are also exactly the
primitive elements of the rank-one lattice.
\end{proof}

\begin{corollary}[Atomic node radical is independent of positivity]
\label{p6:cor:tmf-radical-firewall}
For a transitive support,
\[
\boxed{
T\in\mathcal L_S\text{ primitive}
\quad\Longrightarrow\quad
\sigma(T_\CC)=\sqrt{|S|}.
}
\]
The atomic/primitive condition is the residual geometric clause; no
coordinatewise-positivity or separately chosen magnitude is needed.  Without
atomicity, the maps \(mN_S\) give the distinct magnitudes
\(|m|\sqrt{|S|}\), and complex equivariant maps retain the full rescaling
family.  Positivity or phase alignment remains relevant only when several
oriented maps are coherently added, as in
Proposition~\ref{p6:prop:path-phase-obstruction}.
\end{corollary}

\begin{proposition}[Common-domain correction]
\label{p6:prop:common-domain-correction}
For
\(
\varepsilon_i(1)=\sum_{s\in S_i}|s\rangle
\),
the operator direct sum
\[
\varepsilon_{\rm ind}=\bigoplus_i\varepsilon_i:
\CC^r\longrightarrow\bigoplus_i\CC[S_i]
\]
has singular values \(\sqrt{|S_i|}\) separately.  The different column map
\[
\varepsilon_{\rm col}:\CC\longrightarrow\bigoplus_i\CC[S_i],
\qquad
\varepsilon_{\rm col}(1)
=\bigoplus_i\sum_{s\in S_i}|s\rangle
\]
has
\[
\sigma(\varepsilon_{\rm col})
=\sqrt{\sum_i|S_i|}.
\]
\end{proposition}

\begin{proof}
The two Gram operators are respectively
\[
\diag(|S_1|,\ldots,|S_r|)
\quad\text{and}\quad
\sum_i|S_i|.
\]
\end{proof}

\begin{proposition}[Completed-path phase obstruction]
\label{p6:prop:path-phase-obstruction}
For normalized orbit sums \(u_i\), the most general coherent scalar
codiagonal on the same supports includes
\[
P_{\boldsymbol\theta}
=\sum_i e^{i\theta_i}u_i^*\varepsilon_i,
\qquad
\sigma(P_{\boldsymbol\theta})
=\left|\sum_i e^{i\theta_i}\sqrt{|S_i|}\right|.
\]
Support stabilizers do not fix the relative phases.  Moreover, serial
composition multiplies scalar amplitudes; it does not add them.  The rule
\(\sum_i\sqrt{|S_i|}\) therefore requires a declared positive codiagonal,
once-per-segment semantics, and phase alignment.
\end{proposition}

\begin{definition}[Finite terminal mass functor]
\label{p6:def:terminal-mass-functor}
The retained member \(\mathsf{TMF}\) uses
Theorem~\ref{p6:thm:lattice-transfer-radical} and consists of the following
minimal residual data:
\begin{enumerate}[label=(\roman*),leftmargin=2em]
\item the atomic/primitive condition in the canonical integral terminal
Hom-lattice, with either sign orientation;
\item common-domain, once-per-block columns for Pythagorean sums;
\item a positive once-per-segment codiagonal and phase alignment for
additive paths;
\item fixed Hilbert metrics before Moore--Penrose inversion;
\item the explicit slot assignment \(\mathsf{SLOT}\); and
\item one dimensionful shell insertion \(E_{K(p)}\).
\end{enumerate}
These clauses make multi-support and dimensionful grammar deterministic.  The
node magnitude and its sign-insensitivity are theorem consequences of item~(i),
not separately selected real parameters.
\end{definition}

\begin{definition}[Canonical terminal primitives]
\label{p6:def:canonical-terminal-primitives}
For every \(n\ge1\) and copy label \(a\), choose pairwise disjoint abstract
sets \(S_n^{(a)}\) of cardinality \(n\), put
\(\mathcal H_n^{(a)}=\CC[S_n^{(a)}]\), and define
\[
N_n^{(a)}:\CC\longrightarrow\mathcal H_n^{(a)},
\qquad
N_n^{(a)}(1)=\sum_{s\in S_n^{(a)}}|s\rangle.
\]
For a list \(\boldsymbol n=(n_1,\ldots,n_r)\), define the common-domain
column and positive scalar path
\[
C_{\boldsymbol n}(1)
=
\bigoplus_{a=1}^rN_{n_a}^{(a)}(1),
\qquad
P_{\boldsymbol n}
=
\left(\sum_{a=1}^r\sqrt{n_a}\right)I_{\CC}.
\]
Their nonzero singular values are
\[
\sigma(N_n)=\sqrt n,\qquad
\sigma(C_{\boldsymbol n})=\sqrt{\sum_an_a},\qquad
\sigma(P_{\boldsymbol n})=\sum_a\sqrt{n_a}.
\]
The superscript \(+\) denotes the Moore--Penrose inverse with respect to
these counting metrics.  Thus \(N_n^+\), \(C_{\boldsymbol n}^+\), and
\(P_{\boldsymbol n}^+\) have reciprocal nonzero singular values.
\end{definition}

\begin{corollary}[Primitive \(7\mid9\) Gram and Hessian theorem]
\label{p6:cor:primitive-seven-nine-gram}
For the canonical primitives on the two shell summands,
\begin{equation}
 \label{p6:eq:primitive-seven-nine-gram}
 N_7^*N_7=7I_{\mathbb C},
 \qquad
 N_9^*N_9=9I_{\mathbb C}.
\end{equation}
Consequently the quadratic functional
\[
 \Phi_{79}(w,z)
 :=\frac12\bigl(\|N_7w\|^2+\|N_9z\|^2\bigr)
\]
has complex Hermitian Hessian \(7I_{\mathbb C}\oplus9I_{\mathbb C}\).
The ratio \(7:9\) is therefore an exact primitive counting-metric result.
It is not yet a ratio of physical gauge-boson masses: that interpretation
also requires a common physical owner, ancestry for the two response lines,
and an equal-amplitude statement.
\end{corollary}

\begin{proof}
The vector \(N_n(1)\) is the sum of \(n\) orthonormal basis vectors, so its
squared norm is \(n\).  Since the source is one-dimensional this is exactly
\(N_n^*N_n=nI_{\mathbb C}\).  The Hessian statement follows by taking the
orthogonal direct sum.  No role or amplitude identification enters this
calculation.
\end{proof}

\begin{definition}[Unified lookup rows with derived charged subset]
\label{p6:def:selected-typed-rows}
The following table is retained as the unified terminal lookup atlas.  In
Release v5.66 the charged rows \(t,b,e,u,d,\mu,s,c,\tau\) are no longer
independent \(\mathsf{SLOT}\) choices: their raw terminal words are outputs of
the charged split/primitive-owner compiler below.  The neutral rows retain
their previous status, and actual physical family/source occurrence remains a
separate Paper--V-A obligation.
The multiplicity \(m_p\) is an explicit row degeneracy; for \(m_p>1\), the
listed map means \(I_{m_p}\otimes F_p\).  The support sets are the abstract
counting copies of
Definition~\ref{p6:def:canonical-terminal-primitives}, not claims that
different particle rows are disjoint subsets of one Hurwitz carrier.
\[
\begin{array}{c|c|c|c}
p&m_p&F_p&\text{source}\longrightarrow\text{target}\\
\hline
g&8&0&\CC^8\to\CC^8\\
\gamma&1&0&\CC\to\CC\\
W^\pm&2&N_7&\CC^2\to\CC^2\otimes\mathcal H_7\\
Z&1&N_9&\CC\to\mathcal H_9\\
H&1&C_{(16,1)}&\CC\to\mathcal H_{16}\oplus\mathcal H_1\\
t&3&P_{(7,9)}&\CC^3\to\CC^3\\
b&3&P_{(8,16,1)}^+&\CC^3\to\CC^3\\
e&1&P_{(1,5)}^+&\CC\to\CC\\
u&3&N_2&\CC^3\to\CC^3\otimes\mathcal H_2\\
d&3&P_{(2,3)}&\CC^3\to\CC^3\\
\mu&1&N_4^+&\mathcal H_4\to\CC\\
s&3&N_6^+&\CC^3\otimes\mathcal H_6\to\CC^3\\
c&3&P_{(4,4,4)}&\CC^3\to\CC^3\\
\tau&1&P_{(6,12,6)}&\CC\to\CC\\
\nu_1&1&N_8^+&\mathcal H_8\to\CC\\
\nu_2&1&C_{(4,1)}&\CC\to\mathcal H_4\oplus\mathcal H_1\\
\nu_3&1&P_{(1,\ldots,1)}\ (13\text{ entries})&\CC\to\CC
\end{array}
\]
Clock assignments and comparison exponents are supplied by
Table~\ref{p6:tab:conditional-particle-ledger}.  The physical
three-family operators are not inferred from this lookup direct sum; they
are declared separately in \(\mathsf F\) on the Paper--V-A family spaces.
\end{definition}

\begin{theorem}[Conditional terminal operator and dilation]
\label{p6:thm:conditional-terminal-operator}
Fix \(\mathsf{SLOT}\), \(\mathsf{TMF}\), shell assignments \(K(p)\), and
the retained exponent function \(n_{\etaH}(p)\).  Let
\(\mathcal P_{\rm mass}\) be the
massive rows of the declared packet ledger.  For each
\(p\in\mathcal P_{\rm mass}\), let \(F_{\rm term}(t_p)\) be its node,
common-column, positive-path, or pseudoinverse transfer, with typed row
spaces \(\mathcal H_{p,L},\mathcal H_{p,R}\).  Insert the declared positive
active-line scalar
\[
w_p:=\etaH^{n_{\etaH}(p)}.
\]
Form the external lookup
carriers
\[
\widehat{\mathcal H}_L
=
\bigoplus_{p\in\mathcal P_{\rm mass}}\mathcal H_{p,L},
\qquad
\widehat{\mathcal H}_R
=
\bigoplus_{p\in\mathcal P_{\rm mass}}\mathcal H_{p,R},
\]
and let \(\widehat J_{p,L},\widehat J_{p,R}\) be their canonical injections.
Then
\[
M_{\rm lookup}
=
\sum_{p\in\mathcal P_{\rm mass}}
E_{K(p)}\widehat J_{p,R}
\bigl(w_pF_{\rm term}(t_p)\bigr)\widehat J_{p,L}^*
\]
is a well-defined finite lookup map.  Extend it by the zero map on the
gauge-null \(g,\gamma\) roles.  The dilation
\[
D_{\rm lookup}
=
\begin{pmatrix}
0&M_{\rm lookup}^*\\
M_{\rm lookup}&0
\end{pmatrix}
\]
is self-adjoint.  Its nonzero spectrum is the signed singular spectrum of
\(M_{\rm lookup}\).  This dilation is an operator-theoretic device; it is not
called a fermionic Dirac operator on bosonic lookup slots.
\end{theorem}

\begin{proof}
Every summand has a fixed source, target, metric, and scalar shell insertion.
The external row copies are orthogonal by construction; this does not assert
that reused underlying supports are disjoint physical subspaces.
Self-adjointness is immediate, and the standard singular-value decomposition
of an off-diagonal self-adjoint doubling gives eigenvalues
\(\pm\sigma_j(M_{\rm lookup})\).
\end{proof}

\begin{definition}[Typed terminal row-energy and physical mass bridge]
\label{p6:def:typed-terminal-row-energy}
For a massive lookup row \(p\), let
\(s_p:=\sigma_{\rm nz}(F_{\rm term}(t_p))\) denote the specified nonzero
singular-value readout (with its common-column/path/pseudoinverse convention).
Its native terminal scalar is the energy coordinate
\begin{equation}
 \label{p6:eq:typed-terminal-row-energy}
 \boxed{\mathcal E_p^{\rm term}
 :=E_{K(p)}\etaH^{n_{\etaH}(p)}s_p.}
\end{equation}
The multiplicity symbol \(m_p\) in the lookup ledger remains a row
degeneracy and is never reused as a mass.  A physical rest mass is introduced
only if a downstream pole/role bridge identifies this row with a physical
one-particle pole and supplies
\begin{equation}
 \label{p6:eq:physical-rest-mass-bridge}
 \mathfrak m_p^{\rm rest}c_{\rm op}^{2}=\mathcal E_p^{\rm term},
 \qquad c_{\rm op}=\frac{2\pi R_x}{T_t}.
\end{equation}
Without that bridge, \(\mathcal E_p^{\rm term}\) is not a particle mass.
\end{definition}

\begin{corollary}[One-\(\alpha\) regrading as a unit conversion]
\label{p6:cor:one-alpha-mass-regrading}
For every fixed \(\mathrm{REG}_\sigma\) member, the derived shell law obeys
\begin{equation}
 \label{p6:eq:one-alpha-shell-shift}
 \boxed{E(K+32)=\frac{E(K)}{\alpha_{\rm hier}}.}
\end{equation}
Define the auxiliary conversion speed
\begin{equation}
 \label{p6:eq:auxiliary-conversion-speed}
 v_\sigma:=\frac{c_{\rm op}}{\sqrt{\alpha_{\rm hier}}}.
\end{equation}
If the physical bridge \eqref{p6:eq:physical-rest-mass-bridge} holds at
coordinate \(K\), then the same row scalar can be regraded as
\begin{equation}
 \label{p6:eq:one-alpha-mass-regrading}
 \mathfrak m_p^{\rm rest}v_\sigma^2
 =\frac{\mathcal E_p^{\rm term}}{\alpha_{\rm hier}},
\end{equation}
which is the corresponding \(K+32\) energy coordinate with the unchanged
dimensionless row factor.  The quantity \(v_\sigma\) is a unit-conversion
coordinate.  It is not a second propagation speed and does not modify the
operational characteristic cone fixed by \(c_{\rm op}\).
\end{corollary}

\begin{proof}
Equation~\eqref{p6:eq:one-alpha-shell-shift} is immediate from
\eqref{p6:eq:continuous-shell-law}.  Divide
\eqref{p6:eq:physical-rest-mass-bridge} by \(\alpha_{\rm hier}\) and use
\eqref{p6:eq:auxiliary-conversion-speed}.  No dynamical or causal statement
enters this algebraic regrading.
\end{proof}

\begin{corollary}[Conditional fermionic realization]
\label{p6:cor:fermionic-dirac-restriction}
On \(\Urole\), suppose a separate optional package
\(\mathsf{SLOT}_{\rm phys}\) supplies
isometric source and target
embeddings of every fermionic row into mutually orthogonal Paper--V-A matter
blocks, with all reused-support functionals already evaluated as row
scalars.  Transporting the fermionic part of \(M_{\rm lookup}\) through
those embeddings gives a finite map \(M_F\).  Only on that graded matter
subspace do we write
\[
D_{F,6}
=
\begin{pmatrix}0&M_F^*\\M_F&0\end{pmatrix}
\]
and interpret the self-adjoint doubling as a conditional finite Dirac
operator.  The \(W,Z,H\) entries remain separately typed bosonic terminal or
matrix readouts.  Without these physical embeddings, the row ledger remains
a lookup construction and is not a common matter operator.
\end{corollary}

\subsection{owner-lifted charged compiler and terminal readout}
\label{p6:sec:v588-owner-lifted-charged-compiler}
The charged terminal words remain exact magnitude readouts.  Their tagged
three-role direct sums, however, are downstream support presentations; they
must not be promoted to the physical family carrier when richer owner data
exists upstream.

\begin{definition}[Owner-resolved charged packet]
\label{p6:def:v588-owner-resolved-charged-packet}
For every terminal role $r$ in sector $s$, retain the tuple
\[
 \widehat{\mathcal T}_{s,r}
 =(E_{s,r},A_{s,r},K_{s,r},\rho_{s,r},F_{s,r}),
\]
where $E_{s,r}$ is the actual generated owner when known,
$A_{s,r}$ its rooted owner arrow, $K_{s,r}$ the inherited metric,
$\rho_{s,r}$ its retained root/orientation/Hodge/triality/shell provenance,
and $F_{s,r}$ the already derived terminal transfer.  Put
\[
 \widehat{\mathscr C}_s
 =\{\widehat{\mathcal T}_{s,r}:r\in\mathcal I_s\}.
\]
The existing terminal packet is the downstream support/singular-value readout
\[
 \boxed{T_s^{\rm term}=\operatorname{TermReadout}(\widehat{\mathscr C}_s),}
\]
not a declaration that its abstract support summands are physical family
subspaces.
\end{definition}

\begin{proposition}[Primitive terminal transfers lift to the actual owners]
\label{p6:prop:v588-primitive-owner-lift}
For $n=2,4,6$, let $E_n$ be the generated primitive owner of Paper I-C.
The counting transfer $N_n$ and its inverse have a canonical support-preserving
partial isometry to
\[
 \iota_{E_n}:\mathbb C\to\operatorname{End}(E_n),\qquad z\mapsto zI_{E_n},
\]
so their nonzero magnitudes remain
\[
 \sqrt n,\qquad \frac1{\sqrt n}
\]
on the actual owner.  Therefore the primitive entries of the charged raw
compiler need not pass through an anonymous counting carrier before action
contraction.
\end{proposition}

\begin{proposition}[$q=5$ light split retains family incidence before rank readout]
\label{p6:prop:v588-q5-split-family-incidence}
Let $R$ be the rooted regular $C_6$ shift of the derived $q=5$ anchor, with
projectors $P_1,P_2,P_3$ of ranks $1,2,3$ as in the existing light split, and
let $\iota_+:G\hookrightarrow\mathbb C^{C_6}$ include one alternating rooted
trimer.  With $Q_0=J_G/3$ and $Q_\perp=I-Q_0$,
\[
 \boxed{
 \iota_+^\dagger P_1\iota_+=\frac12Q_0,\qquad
 \iota_+^\dagger P_2\iota_+=\frac12Q_\perp,\qquad
 \iota_+^\dagger P_3\iota_+=\frac12I_G.}
\]
Hence the later numerical split $6=5+1$, $5=3+2$ is already accompanied by
operator-valued singlet/augmentation incidence before its ranks are converted
into radicals.
\end{proposition}

\begin{proof}
On the alternating trimer,
$\iota_+^\dagger R^{2m}\iota_+=X_G^m$ and odd powers map to the opposite
trimer, hence have zero compression.  Substituting these identities into the
standard formulas for $P_1,P_2,P_3$ gives the three displayed expressions.
\end{proof}

\begin{remark}[Direct-sum firewall]
\label{p6:rem:v588-direct-sum-firewall}
The role-tagged direct sum used in
Theorem~\ref{p6:thm:v585-terminal-support-operator} remains a correct exact
support/spectral readout.  Because distinct tagged roles are orthogonal there
by construction, it does not retain cross-role Grams of a possible common
upstream receiver.  Such cross returns must be formed before this readout if a
physical action claim depends on them.
\end{remark}

\begin{corollary}[Charged compiler chronology]
\label{p6:cor:v588-charged-compiler-chronology}
The charged mass values derived by
Theorem~\ref{p6:thm:v566-derived-charged-raw-compiler} are unchanged.  Their
physical family incidence, however, must descend from the occurring
owner-resolved parent contraction of Paper V-A before the terminal support
packet is formed.  That charged contraction is explicitly supplied by
Corollary~\ref{p6:cor:v590-preterminal-contraction-explicit} and physically
co-realized by Paper~V-A,
Corollary~\ref{p5:cor:v397-charged-preterminal-assembly-closure}.
No independent terminal role-to-$G$ isometry is required. Linear
assembly is explicit on a certified generated packet; the hierarchy-arrow
origin certificate is retained as in
Definition~\ref{p5:def:v228-hierarchy-origin}.
\end{corollary}
% END OWNER-LIFTED-CHARGED-COMPILER v5.88

% BEGIN CHARGED-HIERARCHY-INCIDENCE-EXPORT v5.89
\subsection*{hierarchy/incidence export of the already-massed charged rows}
\label{p6:sec:v589-hierarchy-incidence-export}
The charged terminal mass ledger remains unchanged.  The new input from Paper
I-B is the explicit rooted pullback $\Phi_{\rm HG}$ identifying the
$E_{32}/E_{36}/E_{40}$ Clifford corridor with the $1\oplus2$ block
decomposition of $\operatorname{End}(G)$.

\begin{definition}[Hierarchy-level family-incidence block]
\label{p6:def:v589-hierarchy-incidence-block}
For a charged row $p$ with hierarchy level $K(p)\in\{32,36,40\}$ define
\[
 \mathscr I_{K(p)}:=
 \begin{cases}
 Q_{\rm fam,0}\operatorname{End}(G)Q_{\rm fam,0},&K(p)=32,\\
 Q_{\rm fam,0}\operatorname{End}(G)Q_{\rm fam,\perp}
 \oplus Q_{\rm fam,\perp}\operatorname{End}(G)Q_{\rm fam,0},&K(p)=36,\\
 Q_{\rm fam,\perp}\operatorname{End}(G)Q_{\rm fam,\perp},&K(p)=40.
 \end{cases}
\]
This is an owner block, not a mass-generating matrix and not yet a particular
physical member of that block.
\end{definition}

\begin{proposition}[Mass/incidence sibling export]
\label{p6:prop:v589-mass-incidence-sibling-export}
Let $\widehat{\mathcal T}_{s,r}$ be an owner-resolved charged role from
Definition~\ref{p6:def:v588-owner-resolved-charged-packet}, with its already
derived terminal mass readout $m_{s,r}$ and a hierarchy-level owner arrow
$A_{s,r}^{\rm hier}$ when that arrow has been independently generated by the
source owner.  Membership in $\mathscr I_{K(s,r)}$ alone does not construct
this role arrow.  Conditional on such an arrow, before terminal scalarization
its correct typed export is
\[
 \boxed{
 \widehat{\mathcal T}_{s,r}
 \longmapsto
 \Bigl(m_{s,r},\;\iota_{s,r}\Bigr),\qquad
 \iota_{s,r}:=\Phi_{\rm HG}(A_{s,r}^{\rm hier})\in\mathscr I_{K(s,r)}.}
\]
The first component is the existing Clifford/split terminal mass readout.  The
second is its same-root family-incidence/orientation datum.  No spectral
relation between the two components is asserted or required.
\end{proposition}

\begin{corollary}[Charged hierarchy host ambiguity retired]
\label{p6:cor:v589-hierarchy-host-retired}
The question of where the charged $E_{32}/E_{36}/E_{40}$ owner information
lives relative to $G$ is closed: their complete corridor is canonically
identified with $\operatorname{End}(G)$ by the rooted pullback.  The particular rooted charged contraction is supplied by
Corollary~\ref{p6:cor:v590-preterminal-contraction-explicit} and Paper~V-A's
charged source-value realization. This host theorem alone still does not
select arbitrary additional source members.  The charged mass values and their named row
assignments are unchanged.
\end{corollary}

\begin{remark}[No SVD mass reinterpretation]
\label{p6:rem:v589-no-svd-mass-reinterpretation}
Neither $\iota_{s,r}$ nor any later family return is used to regenerate
$m_{s,r}$ as an eigenvalue or singular value.  Mixing/current calculations may
consume the incidence data only after the mass-labelled owner relation has
been preserved.
\end{remark}
% END CHARGED-HIERARCHY-INCIDENCE-EXPORT v5.89

% BEGIN CHARGED-MASS-INCIDENCE-SOURCE-EXPORT v5.90
\subsection*{charged finite mass/incidence source export}
\label{p6:sec:v590-charged-mass-incidence-source-export}
The charged rows of Theorem~\ref{p6:thm:v566-derived-charged-raw-compiler}
are now consumed in their direct finite-mass role.  This does not identify a
finite source coefficient with an interacting pole mass and does not alter any
numerical charged row.

\begin{definition}[Finite charged mass coefficient]
\label{p6:def:v590-finite-charged-mass-coefficient}
For a charged row $p=p(s,r)$ define
\[
 \boxed{m^{\rm fin}_{s,r}:=\mathcal E_{p}^{\rm term}.}
\]
This is the already derived Clifford/hierarchy plus split/terminal coefficient
of the finite charged zero-order mass source.  Equation
\eqref{p6:eq:physical-rest-mass-bridge} remains a separate pole/rest-mass
identification: $m^{\rm fin}_{s,r}$ need not equal a renormalized pole mass
without that downstream bridge.
\end{definition}

\begin{definition}[Complete charged mass/incidence sibling packet]
\label{p6:def:v590-charged-mass-incidence-packet}
For each charged sector whose hierarchy-arrow origins have been
independently certified, retain the complete generated tuple
\[
 \boxed{
 \mathfrak P_s^{\rm ch}
 :=\bigl((m^{\rm fin}_{s,r},\iota_{s,r})\bigr)_{r\in\mathcal I_s},}
\]
where $\iota_{s,r}$ is the hierarchy/rooted family-incidence export of
Proposition~\ref{p6:prop:v589-mass-incidence-sibling-export}, with all inherited
orientation, Hodge, triality and shell marks retained in that operator-valued
entry.  The mass coefficient and incidence are siblings of one owner-resolved
role; neither is reconstructed from the other.
\end{definition}

\begin{theorem}[Canonical channel representatives of the charged incidence packet]
\label{p6:thm:v590-charged-channel-representatives}
Let $\chi(W)=D_\eta W D_\eta$ be the signed rooted triality
channel--family isomorphism of Paper~I-A.  Every incidence entry has the unique
channel representative
\[
 \boxed{W_{s,r}:=\chi^{-1}(\iota_{s,r})=D_\eta\iota_{s,r}D_\eta.}
\]
The packet therefore admits the frozen role-linear preterminal synthesis
\[
 \boxed{
 W_s^{\rm pre}:=\sum_{r\in\mathcal I_s}m^{\rm fin}_{s,r}W_{s,r},
 \qquad
 R_s^{\rm pre}:=\chi(W_s^{\rm pre})
 =\sum_{r\in\mathcal I_s}m^{\rm fin}_{s,r}\iota_{s,r}.}
\]
The formula is covariant under any simultaneous presentation relabelling of
the generated role pairs.  It contains no fitted coefficient, no independent
three-space isometry and no spectral reconstruction of the masses.
\end{theorem}

\begin{proof}
Paper~I-A makes $\chi$ a unital $*$-isomorphism, so every $\iota_{s,r}$ has
one channel representative.  The finite odd source grammar is linear in the
independently generated charged row contributions; generated co-consumption
requires every nonredundant pair $(m^{\rm fin}_{s,r},\iota_{s,r})$ to be
retained.  Hence the displayed sum is the result-independent linear synthesis
of the complete packet.  A simultaneous relabelling permutes the summands and
therefore leaves the synthesized operator unchanged as an abstract rooted
output.
\end{proof}

\begin{corollary}[Charged preterminal parent contraction is explicit]
\label{p6:cor:v590-preterminal-contraction-explicit}
The residual map of Paper~V-A v3.96 is now
\[
 \boxed{
 \mathfrak K_s(\widehat{\mathscr C}_s)=R_s^{\rm pre}
 =\sum_r m^{\rm fin}_{s,r}\iota_{s,r}.}
\]
Thus \texttt{FERMION\_PRETERMINAL\_PARENT\_CONTRACTION} is an explicit
linear assembly on the independently generated-arrow domain. Its input
origins must be certified before the conditional Paper~V-A source-value
handoff is used; the present sum does not construct those input arrows.
\end{corollary}

\begin{remark}[No mass-spectrum reinterpretation]
\label{p6:rem:v590-no-mass-spectrum-reinterpretation}
The numbers $m^{\rm fin}_{s,r}$ remain the direct charged mass coefficients
from the Clifford/split compiler.  They are not defined as eigenvalues or
singular values of $R_s^{\rm pre}$, $R_s^{\rm pre}(R_s^{\rm pre})^\dagger$,
a Yukawa matrix, an SVD or a Feshbach return.  The latter operators encode the
co-consumed family incidence/current geometry of the already-massed roles.
\end{remark}
% END CHARGED-MASS-INCIDENCE-SOURCE-EXPORT v5.90

% BEGIN CHARGED-COMPARISON-GRADE-CLOSURE v5.91
\subsection*{derived charged exclusion comparison and polar grade}
\label{p6:sec:v591-charged-comparison-grade-closure}
The numerical charged finite coefficients are unchanged.  This release
separates their raw Clifford/split magnitude from the finite rooted comparison
character and derives the latter from the same owner provenance already
retained in $\widehat{\mathscr C}_s$.

\begin{definition}[Raw charged geometric coefficient]
\label{p6:def:v591-raw-charged-geometric-coefficient}
For every charged row $p$ put
\[
 \boxed{a_p:=E_{K(p)}s_p,}
\]
where $s_p$ is the already derived nonzero terminal singular-value/radical
readout before the orthogonal-label comparison factor is applied.  This is a
bookkeeping separation inside the existing finite coefficient; it introduces
no new scale or mass law.
\end{definition}

\begin{definition}[Canonical polar comparison classes]
\label{p6:def:v591-polar-comparison-classes}
On the physically consumed rooted exclusion owner of Paper~V-A v3.98, let
$A=Q_{\rm cmp}|\Omega_{25}\rangle\langle\Omega_{25}|$.  Its active singular
value is $\eta_H^{1/2}$ with $\eta_H=24/25$.  Define the four scalar-output
classes
\[
 \mathsf S=[A],\qquad
 \mathsf R=[A^\dagger A],\qquad
 \mathsf R^\vee=[(A^\dagger A)^+],\qquad
 \mathsf O^\vee=[(A(A^\dagger A))^+].
\]
Their polar degrees and active scalar outputs are
\[
\begin{array}{c|c|c}
\text{class}&\deg_{\rm pol}&w\\\hline
\mathsf S&1&\eta_H^{1/2}\\
\mathsf R&2&\eta_H\\
\mathsf R^\vee&-2&\eta_H^{-1}\\
\mathsf O^\vee&-3&\eta_H^{-3/2}.
\end{array}
\]
Thus $n_{\eta_H}=\deg_{\rm pol}/2$ on every role for which one of these
owner classes is generated.
\end{definition}

\begin{theorem}[Owner-provenance classification of the nine charged grades]
\label{p6:thm:v591-owner-provenance-charged-grades}
After quotienting owner-equivalent isometries, scalar-only internal refinements
and common finite normalization, the retained charged owner provenance of
Definition~\ref{p6:def:v588-owner-resolved-charged-packet} has the following
polar classes:
\[
\boxed{
\begin{array}{c|c|c}
p&\operatorname{PolType}(p)&n_{\eta_H}(p)\\\hline
 u&\mathsf S&+\tfrac12\\
d&\mathsf R&+1\\
e&\mathsf O^\vee&-\tfrac32\\
t&\mathsf S&+\tfrac12\\
b&\mathsf O^\vee&-\tfrac32\\
\mu&\mathsf R&+1\\
s&\mathsf R^\vee&-1\\
c&\mathsf R&+1\\
\tau&\mathsf R&+1
\end{array}}
\]
where the first row is the up quark $u$.
\end{theorem}
\begin{proof}
Theorem~\ref{p6:thm:v566-derived-charged-raw-compiler} supplies the
owner-independent terminal provenance.  The unique non-scalar $q=5$
primitive owner $u$ is a one-endpoint source, hence $\mathsf S$; the inner
$5=3+2$ down-quark split retains a source--target owner relation before rank
readout, hence $\mathsf R$; and the outer $6=5+1$ electron pseudoinverse is the
reciprocal completion of the primitive source plus its closed relation, hence
$\mathsf O^\vee$.  At $E_{32}$ the $7+9$ top refinement stays inside the
one-dimensional hierarchy source block and therefore remains $\mathsf S$,
whereas the bottom word is the outer heavy split plus unital reciprocal
completion and is $\mathsf O^\vee$.

For the $E_{36}$ pair, Paper~I-A proves
\[
 H_{\Lambda^2}=\tfrac14H_V^{-1}.
\]
The common scalar $1/4$ is normalization degree zero, so the vector-owner
return is $\mathsf R$ while its exterior-square partner is the reciprocal
class $\mathsf R^\vee$, giving $n_\mu=+1$ and $n_s=-1$.  The charm
$K_{4,4,4}$/triality path and the tau primal--dual $6,12,6$ incidence path are
two-ended generated relations with no reciprocal owner inversion, so both are
$\mathsf R$.
\end{proof}

\begin{corollary}[Derived charged comparison character]
\label{p6:cor:v591-derived-charged-comparison-character}
For every charged row,
\[
 \boxed{m_p^{\rm fin}
 =a_p\left(\frac{24}{25}\right)^{\deg_{\rm pol}(p)/2}.}
\]
Equivalently the historical charged exponent vector is now a theorem output,
\[
 \boxed{\left(\tfrac12,1,-\tfrac32,\tfrac12,-\tfrac32,1,-1,1,1\right)}
\]
in the order $(u,d,e,t,b,\mu,s,c,\tau)$.  No particle-specific exponent is an
independent input.
\end{corollary}

\begin{corollary}[Additive charged comparison cocycle]
\label{p6:cor:v591-additive-charged-comparison-cocycle}
On the reduced charged comparison category, define
\[
 \chi(\gamma):=\frac12\deg_{\rm pol}(\gamma).
\]
For licensed compositions and Moore--Penrose reciprocity,
\[
 \boxed{\chi(\gamma_2\circ\gamma_1)=\chi(\gamma_2)+\chi(\gamma_1),
 \qquad \chi(\gamma^+)=-\chi(\gamma).}
\]
Thus the exact additive $\tfrac12\mathbb Z$ cocycle required by the earlier
conditional discussion is generated on the charged branch.
\end{corollary}

\begin{remark}[No conflict with the factor-count no-go]
\label{p6:rem:v591-no-conflict-factor-count-no-go}
Theorem~\ref{p6:thm:comparison-factor-no-go} remains valid: the $25$-cell
design by itself does not choose an arbitrary tensor word or factor count.
Release v5.91 instead uses independently generated owner provenance to select
the scalar-output class $\mathsf S$, $\mathsf R$, $\mathsf R^\vee$ or
$\mathsf O^\vee$.  The word lift remains nonphysical whenever only the scalar
comparison output is consumed.
\end{remark}

\begin{corollary}[Charged comparison closure status]
\label{p6:cor:v591-charged-comparison-closure-status}
On the rooted same-action charged comparison branch,
\[
 \boxed{\begin{gathered}
 \texttt{CHARGED\_FINITE\_COMPARISON}\\
 =\texttt{CLOSED\_DEPENDENT\_ON\_ROOTED}\\
 \texttt{\_CHARGED\_COMPARISON\_BRANCH}.
 \end{gathered}}
\]
The downstream distinction among $m_p^{\rm fin}$, a named-scheme running
coordinate $m_p^S(\mu)$ and a pole/rest mass remains unchanged.
\end{corollary}
% END CHARGED-COMPARISON-GRADE-CLOSURE v5.91


% BEGIN FRAMEWORK v2.221: EXTERIOR-POWER CHARGED COMPILER AUDIT
\subsection{terminal-mass ontology and current-based flavor readout}
\label{p6:sec:v592-mass-ontology-flavor-readout}

\noindent\textbf{Framework v2.208 branch scope.}
The neutral inverse/direct/closed-Gram entries in this retained section
belong to the historical $D_4$ response constructor. They are obsolete
on the harmonic rank branch of \ref{p6:sec:v2208-neutral-rank-word}.
Their algebraic theorems neither force nor constrain that new mass word.
The charged results and branch-independent geometric/source rules keep
their stated domains.


\begin{theorem}[Terminal masses are not eigenvalues of the later family/source operator]
\label{p6:thm:v592-terminal-mass-no-backgeneration}
On the active rooted branch the charged coefficients $m_{s,r}^{\rm fin}$ and
the neutral $E_{56}\{x^{-1},x,x^2\}$ role energies are generated by their
terminal geometric/hierarchy constructors before the later finite
family/source synthesis.  Consequently
\[
 R_s^{\rm ch}=\sum_r m_{s,r}^{\rm fin}\iota_{s,r}
\]
is an incidence-weighted source coefficient, not a mass matrix whose spectrum
defines the $m_{s,r}^{\rm fin}$.  Likewise the N56 response-jet coefficients
are terminal role readouts, not Takagi eigenvalues of a neutral family matrix.
No eigenvalue, singular-value or Takagi reconstruction of a later operator is
part of native mass genesis.
\end{theorem}

\begin{definition}[Framework-native weak-current incidence]
\label{p6:def:v592-current-incidence}
Let $\mathscr L_u,\mathscr L_d,\mathscr L_e$ denote the charged terminal role
spaces and $\mathscr L_{56}$ the neutral inverse/direct/closed-Gram role space.
For a generated weak-current morphism $\mathscr J_W$ and canonically normalized
terminal-support evaluations $W_s$, define the native current readouts
\[
 C_{ud}:=W_u^\dagger\mathscr J_W W_d,
 \qquad
 C_{e\nu}:=W_e^\dagger\mathscr J_W W_\nu .
\]
These maps compare already generated role spaces.  They do not require masses
to be eigenvalues of a common family operator.
\end{definition}

\begin{definition}[Neutral propagation readout]
\label{p6:def:v592-neutral-propagation-readout}
If downstream continuum propagation supplies phases $\Phi_r(\gamma)$ for the
three neutral terminal roles along a physical path $\gamma$, define
\[
 \mathcal P_\nu(\gamma)=\sum_{r\in\{\rm inv,dir,Gram\}}
 e^{-i\Phi_r(\gamma)}\Pi_r^{\rm role}.
\]
For equal normalized production and detection current maps the corresponding
flavor-transition kernel is
\[
 \mathcal A_\ell(\gamma)=C_{e\nu}\,\mathcal P_\nu(\gamma)\,C_{e\nu}^\dagger.
\]
The propagation law fixing the $\Phi_r$ is downstream and claim-scoped; the
formula does not assume a neutrino mass matrix or a Takagi mass basis.
\end{definition}

\begin{corollary}[CKM/PMNS are optional conventional coordinate dictionaries]
\label{p6:cor:v592-effective-ckm-pmns-dictionary}
When a conventional unitary coordinate section is useful, one may define
\[
 V_{\rm CKM}^{\rm eff}:=[C_{ud}]_{\rm std},\qquad
 U_{\rm PMNS}^{\rm eff}:=[C_{e\nu}]_{\rm std}
\]
after the relevant normalization and presentation quotients.  These symbols
are downstream comparison representatives of generated current maps.  The
former native targets \texttt{CKM\_RELATIVE\_FRAME} and
\texttt{PMNS\_DIRAC\_RELATIVE\_FRAME} are retired as framework obligations.
They are replaced, where the corresponding observables are claimed, by
\texttt{QUARK\_WEAK\_CURRENT\_INCIDENCE},
\texttt{LEPTON\_WEAK\_CURRENT\_INCIDENCE},
\texttt{NEUTRAL\_PROPAGATION\_PHASE} and
\texttt{OSCILLATION\_CURRENT\_READOUT}.
\end{corollary}

\begin{remark}[Scope of the older PMNS/eigenframe sections]
The older simultaneous-conjugacy, eigenframe-counterfamily, PMNS-packet,
Takagi-compatibility and Majorana-frame passages below/above remain useful as
Standard--Model coordinate dictionaries, countermodels and comparison
constructions.  They no longer define the framework-native mass or mixing
ontology and may not be used as premises for terminal mass generation.
\end{remark}
% END MASS-ONTOLOGY-FLAVOR-READOUT v5.92

% BEGIN N56-NEUTRAL-INCIDENCE-CLOSURE v5.93
\subsection{Complete charged transfer certificate and its actual-source antecedent}
\label{p6:sec:v2168-charged-transfer-certificate}
The positive terminal coefficients, already generated by the direct
Clifford/split compiler, are fixed before the source response is calculated:
\[
 M=\diag(\varepsilon_e,\varepsilon_\mu,\varepsilon_\tau)>0,\qquad
 \begin{aligned}
 \varepsilon_e&=E_{40}(1+\sqrt5)^{-1}(24/25)^{-3/2},\\
 \varepsilon_\mu&=(E_{36}/2)(24/25),\\
 \varepsilon_\tau&=E_{36}(2\sqrt6+\sqrt{12})(24/25).
 \end{aligned}
\]
The occurring source is $R_{\rm act}=\sum_r\varepsilon_r\iota_r$,
with $\iota_r=\Phi_{\rm HG}(A_r^{\rm hier})$. Its physical occurrence and
assembly are closed by V-A; no independent source choice is introduced.
The current frame $U=(g_e,g_\mu,c)$ is the generated incidence frame,
not a spectrum-generated mass frame. The primitive momentum level remains
$m_p=1$. A local rest-comparison symbol requires its admitted physical
dictionary; it creates neither a primitive zero-momentum state nor a
position lattice.
\subsubsection{A right-role completion from the existing negative-Hopf parent}
The current source supplies
\[
 c=\frac1{\sqrt3}(1,1,1)^T,\quad Q_0=cc^\dagger,\quad Q_\perp=I-Q_0,
 \quad R=\diag(1,i,-i),
\]
\[
 g_\mu=\frac1{2\sqrt6}(2,-1+3i,-1-3i)^T,\quad
 h_\tau=\overline{g_\mu},\quad
 g_e=\frac12(-1-i,i,1)^T.
\]
The existing left frame is $U=(g_e,g_\mu,c)$ and $P_e=g_eg_e^\dagger$.
The two already generated mixed relation blocks are
$A_\mu=Q_\perp RQ_0$ and $A_\tau=Q_0RQ_\perp$.

\begin{theorem}[Generated right completion and three orthogonal role arrows]
\label{p6:thm:v2168-right-completion}
Define the geometric relation descendant
\[
 A_e=P_eRQ_\perp,\qquad
 h_e=R^\dagger g_e=\frac12(-1-i,1,i)^T,
 \qquad V=(h_e,c,h_\tau).
\]
Then $U$ and $V$ are unitary, and with $s=\sqrt{8/9}$,
\[
 A_e=g_eh_e^\dagger,\qquad
 A_\mu=s\,g_\mu c^\dagger,\qquad
 A_\tau=s\,c h_\tau^\dagger.
\]
Consequently the normalized arrows
\[
 B_e=g_eh_e^\dagger,\quad B_\mu=g_\mu c^\dagger,\quad
 B_\tau=c h_\tau^\dagger
\]
have pairwise orthogonal source lines and pairwise orthogonal target lines.
They contain no free angle, amplitude or right-frame selector.
\end{theorem}
\begin{proof}
The existing left-frame theorem gives $U^\dagger U=I$.
Since $g_e$ is orthogonal to both $c$ and $g_\mu$, and
$Rc=(c/3)+s g_\mu$, one has
$c^\dagger R^\dagger g_e=(Rc)^\dagger g_e=0$.
Thus $h_e\in Q_\perp G$, and $R$ being unitary gives $\|h_e\|=1$.
Moreover
\[
 h_\tau=s^{-1}Q_\perp R^\dagger c,
 \qquad Rh_\tau=s^{-1}(c-Rc/3),
\]
so $h_e^\dagger h_\tau=g_e^\dagger Rh_\tau=0$.
Both are orthogonal to $c$, proving $V^\dagger V=I$.
The electron identity follows from
$g_e^\dagger RQ_\perp=h_e^\dagger$; the two mixed identities follow from
the displayed compressions and their norms $s$.
Orthogonality of the three columns on each side proves the last statement.
\end{proof}

The map $A\mapsto P_eAQ_\perp$ is an explicit generated relation projection.
II-A permits retention and co-consumption of this output in its geometric
relation sort. This does not itself identify $B_e$ with the physical
$\mathsf{SrcVal}$ arrow $\iota_{e,e}$. The mass/clock test below states
that additional equality visibly.

\subsubsection{Unit angular arrows from the retained relation}
The known rooted relation is $R_-=\diag(1,i,-i)$, and
\[
 c=\frac1{\sqrt3}(1,1,1)^T,\quad Q_0=cc^\dagger,\quad Q_\perp=I-Q_0,
\]
\[
 g_e=\tfrac12(-1-i,i,1)^T,\quad
 g_\mu=\tfrac1{2\sqrt6}(2,-1+3i,-1-3i)^T,\quad
 h_\tau=\overline{g_\mu},\quad h_e=R_-^\dagger g_e.
\]
The frames $U=(g_e,g_\mu,c)$ and $V=(h_e,c,h_\tau)$ are unitary.
Define the unit relation arrows
\[
 B_e=g_eh_e^\dagger=P_eR_-Q_\perp,
 \quad B_\mu=g_\mu c^\dagger,
 \quad B_\tau=c h_\tau^\dagger.
\]
The raw mixed relation compressions have norm $s=\sqrt{8/9}$, with
\[
 Q_\perp R_-Q_0=sB_\mu,\qquad
 Q_0R_-Q_\perp=sB_\tau.
\]
This $s$ is a relation-compression norm. Theorem~\ref{p6:thm:v2168-complete-transfer-lift} uses
unit terminal ports together with unit angular arrows; it does not rescale an
already fixed physical mass source.

\begin{proposition}[Explicit hierarchy lifts of the unit angular arrows]
\label{p6:prop:v2168-angular-hierarchy}
Let $J_{\rm HG}$ be the existing hierarchy/family unitary and put
\[
 \widetilde B_r=J_{\rm HG}B_rJ_{\rm HG}^\dagger.
\]
Then $\widetilde B_e$ lies in the $Q_5$--$Q_5$ hierarchy corner,
$\widetilde B_\mu$ in the $Q_5$--$Q_4$ corner, and
$\widetilde B_\tau$ in the $Q_4$--$Q_5$ corner.
For $w_\mu=J_{\rm HG}g_\mu$ and $w_\tau=J_{\rm HG}h_\tau$, set
$v_\mu=-w_\mu u$ and $v_\tau=-w_\tau u$ in the complexified augmentation
plane. Both have unit inherited norm. The existing quaternionic edge map gives
\[
 \boxed{\widetilde B_\mu=T(v_\mu),\qquad
 \widetilde B_\tau=T(v_\tau)^\dagger,\qquad
 T(v)=|vu\rangle\langle u|.}
\]
Moreover
\[
 \boxed{B_rB_s^\dagger=\delta_{rs}P_r,
 \qquad (P_e,P_\mu,P_\tau)=(P_e,P_\mu,Q_0).}
\]
\end{proposition}
\begin{proof}
The rooted unitary sends $c$ to $u$ and $Q_\perp$ to $Q_5$.
Right multiplication by $u$ preserves the Hermitian norm of the
complexified coefficient space. Since $u^2=-1$,
$(-w u)u=w$. Thus $T(v_\mu)=|w_\mu\rangle\langle u|$ and
$T(v_\tau)^\dagger=|u\rangle\langle w_\tau|$.
The corner statements follow by transporting the corresponding projectors.
Finally $B_r=|Ue_r\rangle\langle Ve_r|$ and the columns of both frames are
orthonormal, proving the return certificate including the electron--tau zero.
\end{proof}

These are explicit geometric descendants in the relation/operator sort.
The proposition does not identify $\widetilde B_r$ with the actual retained
$A_r^{\rm hier}$ of the physical source packet.

\subsubsection{The complete owner-transfer map and its source lift}
Define the already fixed positive hierarchy/comparison multipliers
\[
 \kappa_e=E_{40}\eta_H^{-3/2},\qquad
 \kappa_\mu=\kappa_\tau=E_{36}\eta_H,
\]
and the positive path coefficients
$p_e=(1+\sqrt5)^{-1}$ and $p_\tau=2\sqrt6+\sqrt{12}$.
Use the geometric transfer spaces
\[
 \mathcal H_{\rm in}=\mathbb C_e\oplus\End(E_4)_\mu\oplus\mathbb C_\tau,
 \qquad \mathcal H_{\rm out}=\mathbb C^3_{e,\mu,\tau}.
\]
The $\End(E_4)$ presentation is the native owner lift of the primitive muon
port. These spaces are transfer-owner spaces; their entire dimensions are
not asserted additional physical fermion multiplicities.
The electron and tau terminal maps are already $p_eI$ and $p_\tau I$ in
\path{p6:thm:v585-terminal-support-operator}; the muon counting input
is displayed through the equivalent full native owner presentation here. These terminal maps
do not replace the richer incidence siblings of the physical source.
In particular this tagged transfer presentation does not prove that the
actual common physical receiver has orthogonal cross-role returns.
The full primitive-support normalization is I-B, Theorem~\ref{p1b:thm:v2168-primitive-owner-support}. The complete weighted transfer is
\[
 \boxed{\mathbb F(z_e,X,z_\tau)=
 \begin{pmatrix}
 \kappa_ep_e z_e\\
 \kappa_\mu\Tr(X)/4\\
 \kappa_\tau p_\tau z_\tau
 \end{pmatrix}.}
\]
Its active input embedding is
\[
 \Gamma:\mathbb C^3\longrightarrow\mathcal H_{\rm in},\qquad
 \Gamma(a_e,a_\mu,a_\tau)=(a_e,a_\mu I_4/2,a_\tau).
\]

\begin{theorem}[Complete weighted transfer and unit-support contraction]
\label{p6:thm:v2168-complete-transfer-lift}
Let $\Pi_{\rm in}=\Gamma\Gamma^\dagger$. Then
\[
 \boxed{\Gamma^\dagger\Gamma=I,\qquad
 \mathbb F\Gamma=M,\qquad
 \mathbb F=M\Gamma^\dagger,\qquad
 \mathbb F\mathbb F^\dagger=M^2.}
\]
In particular the full transfer annihilates precisely the unused traceless
muon-port directions. Define generated comparison maps
\[
 \mathcal E_L=U:\mathcal H_{\rm out}\to G_L,\qquad
 \mathcal E_R=\Gamma V^\dagger:G_R\to\mathcal H_{\rm in}.
\]
Both are isometric, and $\mathcal E_R\mathcal E_R^\dagger=\Pi_{\rm in}$.
Their operator contraction is
\[
 \boxed{R_{\rm lift}:=\mathcal E_L\mathbb F\mathcal E_R
       =UMV^\dagger=\sum_r\varepsilon_rB_r.}
\]
It satisfies all three the free-return criterion return identities exactly.
\end{theorem}
\begin{proof}
The three tagged input directions are orthogonal and $\|I_4/2\|_{\rm HS}=1$.
The electron and tau coordinates give their fixed positive path coefficients;
the native primitive retraction gives $\kappa_\mu/2=\varepsilon_\mu$.
For arbitrary $X$, the muon component of $\Gamma^\dagger$ is
$\langle I_4/2,X\rangle_{\rm HS}=\Tr(X)/2$.
Multiplying it by $\varepsilon_\mu=\kappa_\mu/2$ gives
$\kappa_\mu\Tr(X)/4$, proving the full-domain identity
$\mathbb F=M\Gamma^\dagger$. Positivity of all three coefficients shows
that its kernel is exactly the traceless muon-port space.
The isometry and range identities follow from unitarity of $U,V$ and
$\Gamma^\dagger\Gamma=I$. Multiplication gives $R_{\rm lift}$.
The unit angular return certificate gives
$R_{\rm lift}R_{\rm lift}^\dagger=UM^2U^\dagger$ and each of the three
owner-block tests separately.
\end{proof}

\paragraph{The combined muon port explicitly.}
The contraction separates the primitive transfer from its angular incidence:
\[
 \boxed{\kappa_\mu\epsilon_{E_4}(I_4/2)\,
       \Phi_{\rm HG}(T(v_\mu))
       =\frac{E_{36}\eta_H}{2}\,g_\mu c^\dagger.}
\]
Its $1/2$ comes from the native inverse transfer, not from a fitted mass or
a chosen singular value of the final family operator.
The electron and tau contributions are analogously
$\kappa_ep_e B_e$ and $\kappa_\tau p_\tau B_\tau$.
The map consumes scalar-transfer and angular-relation siblings jointly;
tracing an off-diagonal angular arrow first is a different, zero word
(Appendix~\ref{p6:app:v2168-transfer-diagnostics}).

All embeddings and angular arrows in this displayed construction depend on
primitive owner metrics, hierarchy labels, and the rooted Hopf relation.
They are independent of the evaluated terminal energy values. The positive
path and comparison scalars retain the existing compiler values. Thus this
is an explicit forward construction from the complete port packet; it is
not a spectral reconstruction of the native masses.

\subsubsection{A certificate that does not select this particular right frame}
\begin{theorem}[Exact equivalence to unit-support transport]
\label{p6:thm:v2168-unit-support-equivalence}
Let $R_{\rm act}:G_R\to G_L$ be the actual normalized family source on the
uninserted branch. With the already fixed $U,M,\mathbb F,\Gamma$, the following
are equivalent:
\begin{enumerate}
\item $R_{\rm act}R_{\rm act}^\dagger=UM^2U^\dagger$.
\item There exists $\mathcal J_R:G_R\to\mathcal H_{\rm in}$ such that
\[
 \boxed{R_{\rm act}=U\mathbb F\mathcal J_R,\qquad
 \mathcal J_R^\dagger\mathcal J_R=I,
 \qquad\mathcal J_R\mathcal J_R^\dagger=\Pi_{\rm in}.}
\]
\end{enumerate}
The concrete lift uses $\mathcal J_R=\Gamma V^\dagger$, but this choice is
not required for the free left-current claim.
\end{theorem}
\begin{proof}
For (2), $\mathbb F\Pi_{\rm in}=\mathbb F$ gives
$R_{\rm act}R_{\rm act}^\dagger=U\mathbb F\mathbb F^\dagger U^\dagger$.
For (1), let $K=M^{-1}U^\dagger R_{\rm act}$. Then $KK^\dagger=I$, hence
$K$ is unitary. Define $\mathcal J_R=\Gamma K$. All three identities in (2)
follow from $\mathbb F\Gamma=M$ and $\Gamma^\dagger\Gamma=I$.
\end{proof}

The converse formula is a comparison readout of an already given $R_{\rm act}$;
it cannot prove its own unitarity or serve as a mass-dependent source selector.
A native proof supplies the actual source-to-port graph from its owner word,
or evaluates the weaker return identity directly. For inserted kinetic metrics
apply the theorem to $Z_L^{-1/2}D_eZ_R^{-1/2}$ and the same action's normalized
current frame. Energy-dependent physical jets retain the complete inverse
propagator and derivative normalization.


\paragraph{Actual source application.}
The complete transfer certificate is proved for the displayed geometric
lift. To apply it to the occurring $R_{\rm act}$, the physical exporter must
supply its own unit-support graph, or the equivalent actual return identity
of III-F, Theorem~\ref{p3f:thm:v2168-left-return-test}. The converse
comparison construction does not generate that graph. Paper II-A's
owner-normalization inheritance consumes proved ancestry; it does not
establish ancestry from numerical agreement. The direct dynamical route
of IV does not require this stronger free comparison for a pole-only claim.

% END INLINE SOURCE: Paper5B-ChargedTransferCertificate-Main-v1_0.tex



% v2.169: generated action/support constructors and consuming scope.

% BEGIN INLINE SOURCE: Paper5B-ChargedPoleComparison-Main-v1_1.tex
% Framework v2.169: Batches 7--13, integrated at the owning source.
\subsection{Fixed terminal coefficients and the downstream physical root}
\label{p6:sec:v2169-terminal-root-comparison}

The terminal compiler fixes $\varepsilon_e,\varepsilon_\mu,\varepsilon_\tau$
and the source incidence before the response calculation. IV's complete
physical-support polynomial and V-C's evaluated weak occupation returns
are downstream outputs of that source/action word. Their first unproved
finite calculation is the actual source-visible root set of $p_{r,P}$ on the
certified support.  If the terminal coefficient is claimed to be the physical
pole, the comparison certificate is $p_{r,P}(\varepsilon_r)=0$; the root is
not selected by that target.  A shifted dynamical root is retained as the
forward output and does not back-generate a replacement coefficient or fit.
The explicit Hopf transfer lift is still a comparison lift until its actual
exporter equality is proved for a whole-free-response target. That stronger
route is unnecessary for the direct physical pole construction of
Theorem~\ref{p4:v2169:b10:thm:support}.

% END INLINE SOURCE: Paper5B-ChargedPoleComparison-Main-v1_1.tex


% v2.189: complete-word extraction and actual marked hierarchy moments.

% BEGIN INLINE SOURCE: Paper5B-ElectronWordMoments-Main-v1_0.tex
% Framework v2.189: native charged complete-word coefficient extraction.
% Positive source-moment chain; no actual C01/A05 pole closure is asserted.
\subsection{Native electron source moments from the inherited action word}
\label{p6:sec:v2189-electron-word-moments}

The rooted charged source assembly is already closed by
Corollary~\ref{p5:cor:v397-charged-preterminal-assembly-closure}.
Its fixed coefficients and actual incidence entries are consumed together
in Theorem~\ref{p6:thm:v590-charged-channel-representatives}.
The following calculations extract finite source moments from that word;
they neither change a terminal coefficient nor require uniqueness of the
whole action. The complete canonical generator is the inherited output of
Theorem~\ref{p3f:thm:v1148-finite-metric-constructor}, with its actual field
and dual-field embeddings, ordering, occupation and constraint data.

Work on its admitted finite tensor-CAR chart, with $H=H^\dagger$,
$\rho\succeq0$, $\operatorname{Tr}\rho=1$ and
$\omega(A)=\operatorname{Tr}(\rho A)$. The native internal tag is
the current ray of Corollary~\ref{p6:cor:v597-electron-plane-ray},
\[
 g_e=\tfrac12(-1-i,i,1)^T,\qquad
 j=j_{\rm ext}\otimes g_e,\quad \|j_{\rm ext}\|=1,\quad
 q_e=a^\dagger(j).
\]
The external source belongs to the actual licensed chart; normalization
does not select its mode or replace its inherited external slot factors.
If a physical constraint or source compression changes this CAR algebra,
its actual compressed products and boundary returns are used instead.

\begin{theorem}[Unit source weight and complete-word fermionic Hessian]
\label{p6:thm:v2189-complete-word-hessian}
For this normalized tensor-CAR source,
\[
 \{q_e^\dagger,q_e\}=I,\qquad
 \boxed{\mu_0:=\omega(\{q_e^\dagger,q_e\})=1.}
\]
Write the complete even normal word, including all inherited contacts, as
\[
 H=I\otimes H_B+
   \sum_{I,J}a_I^\dagger a_J\otimes B_{I,J},\qquad
 |I|+|J|\equiv0\pmod2.
\]
Here $I=(i_1<\cdots<i_p)$ and $J=(j_1<\cdots<j_s)$,
$a_I^\dagger=a_{i_1}^\dagger\cdots a_{i_p}^\dagger$ and
$a_J=a_{j_s}\cdots a_{j_1}$. The possibly noncommuting $B_{I,J}$ commute
with these CAR slots and are the actual inherited coefficient operators.
Then
\begin{align}
 [H,q_e]
 &=\sum_{I,J}\sum_{r=1}^{|J|}(-1)^{r-1}j_{j_r}\,
   a_I^\dagger a_{J\setminus j_r}\otimes B_{I,J},
   \label{p6:eq:v2189-complete-source-column}\\
 \mathcal D_jH:=\{q_e^\dagger,[H,q_e]\}
 &=\sum_{I,J}\sum_{r=1}^{|J|}\sum_{t=1}^{|I|}
   (-1)^{r+t-2}\overline{j_{i_t}}j_{j_r}\,
   a_{I\setminus i_t}^\dagger a_{J\setminus j_r}\otimes B_{I,J}.
   \label{p6:eq:v2189-directional-hessian}
\end{align}
Consequently the complete first moment is the exact coefficient extraction
\[
 \boxed{\mu_1=\omega(\mathcal D_jH).}
\]
Pure creation or pure annihilation terms, and $I\otimes H_B$, have zero
first-moment contribution. This does not assert that their source columns
vanish or that the complete first moment is zero.
\end{theorem}
\begin{proof}
Normalized CAR give the first identity before choosing a state or generator.
Move $q_e$ through the ordered annihilation slots. Contraction in slot $r$
has sign $(-1)^{r-1}$. The uncontracted terms cancel in the commutator
because the parent word is even. Move $q_e^\dagger$ through the remaining
creation slots; contraction in slot $t$ contributes $(-1)^{t-1}$.
The uncontracted terms cancel in the anticommutator because the first
source column is odd. Their product gives
\eqref{p6:eq:v2189-directional-hessian}. Empty $I$ or $J$ leaves no double
contraction. Taking the actual density trace proves the stated moments.
\end{proof}

\begin{corollary}[Exact quartic occupation contraction]
\label{p6:cor:v2189-quartic-first-moment}
Where the parent actually inherits a quartic fermion vertex, III-F supplies
$\mathsf V_F=(\Lambda^2Z)^\dagger\mathsf V_p(\Lambda^2Z)$ and
\[
 H_4=\sum_{i<k,l<m}\mathsf V_{ik;lm}\,
                 a_i^\dagger a_k^\dagger a_m a_l.
\]
Its first-moment operator is
\begin{align*}
 \mathcal D_jH_4=\sum_{i<k,l<m}\mathsf V_{ik;lm}\bigl(&
  \overline{j_i}j_l a_k^\dagger a_m
 -\overline{j_i}j_m a_k^\dagger a_l\\
 &-\overline{j_k}j_l a_i^\dagger a_m
 +\overline{j_k}j_m a_i^\dagger a_l\bigr).
\end{align*}
The complete contribution is its actual coefficient-weighted occupation
trace, without Wick factorization. This formula adds no vertex when the
parent has none.
\end{corollary}
\begin{proof}
Apply \eqref{p6:eq:v2189-directional-hessian} to the two creation and two
annihilation slots. The four ordered contractions have signs $+,-,-,+$.
The inherited slot transport fixes $\mathsf V$, so no new coupling is used.
\end{proof}

\begin{theorem}[Complete second moment and native marked cancellations]
\label{p6:thm:v2189-marked-moment-cancellations}
Suppose now that $[H,\rho]=0$ for the complete licensed generator.
For its whole first column $\mathcal C_e=[H,q_e]$,
\[
 \boxed{\mu_2:=\omega(\{q_e^\dagger,[H,[H,q_e]]\})
       =\omega(\{\mathcal C_e^\dagger,\mathcal C_e\})\ge0.}
\]
Every inherited cross term is retained. In particular, for a bilinear
column $Y=\sum_i a_i^\dagger A_i$ with finite noncommuting coefficient
operators,
\[
 \boxed{\{Y^\dagger,Y\}=\sum_i A_i^\dagger A_i
       +\sum_{i,k}a_k^\dagger a_i[A_k,A_i^\dagger].}
\]
This is the coefficient-pairing identity of
Proposition~\ref{p4:prop:v2168-coefficient-pairing}, including the exact
occupation correction on an entangled state.

On the admitted local constant zero-order chiral mass mark
$h_R=\bigl(\begin{smallmatrix}0&R\\R^\dagger&0\end{smallmatrix}\bigr)$,
with the normalized external tensor factor retained, its source column is
$C_M=a_{e,R}^\dagger(R^\dagger g_e)$. The actual charged weak column of
Proposition~\ref{p5c:prop:v2168-weak-receiver-columns} is
\[
 C_W=\tfrac\pi2\bigl(Y_1+iY_2-(1+i)Y_3\bigr),
\]
where the $Y_j$ retain every actual clock, spin, spatial and finite Bose
factor. Exact port CAR therefore give
\begin{align*}
 \{q_e^\dagger,C_M\}&=\{q_e^\dagger,C_W\}=0,\\
 \{C_M^\dagger,C_M\}&=\|R^\dagger g_e\|^2I,\\
 \{C_M^\dagger,C_W\}+\{C_W^\dagger,C_M\}&=0.
\end{align*}
The last identity is state independent: the constant mass column has
charged-right ports, while the weak column has neutral-left ports.
The weak internal squared weights are $1/4,1/4,1/2$; their full norm still
consumes the actual external/Bose coefficients and density.

Pure pair marks have zero first-moment Hessian. The native $B-L=2$ pair
uses neutral-right ports, so its direct commutator with this charged
electron source is zero. These marked identities do not delete the pair
from the complete stationary density or from later neutral/scalar/weak
returns, and do not prove $\mu_1=0$ for the complete word.
\end{theorem}
\begin{proof}
Stationarity implies $\omega([H,A])=0$. Apply this to
$q_e^\dagger\mathcal C_e$ and $\mathcal C_eq_e^\dagger$, using
$[H,q_e^\dagger]=-\mathcal C_e^\dagger$, to obtain the two terms
$\omega(\mathcal C_e^\dagger\mathcal C_e)$ and
$\omega(\mathcal C_e\mathcal C_e^\dagger)$.
For the bilinear identity, CAR reorder $a_i a_k^\dagger$; the two
coefficient orders differ by $[A_k,A_i^\dagger]$.
The native weak coefficient is the conjugate electron row of the existing
current matrix $C=U^\dagger W_\nu$; inserting that row gives the displayed
column and weights. Orthogonal charged-right and neutral-left ports have
zero mixed anticommutator. The mass coefficient is a c-number matrix and
commutes with the finite Bose coefficient, so the two mixed products
cancel as operators. The pair statements follow from the empty-slot
case of Theorem~\ref{p6:thm:v2189-complete-word-hessian} and its actual
neutral-right support.
\end{proof}

Quantum scalar/Yukawa coefficients, source-dependent kinetic pullbacks,
contacts and constraints need not obey that constant-mass cross
cancellation. Their complete interference terms remain. Thus the full
$\mu_2$ is not bounded below by simply adding isolated positive marked
norms. A field-independent central energy drops out of these commutators
on the fixed chart; its consumed metric/source derivatives and absolute
energy allocation remain in the parent action.

\begin{theorem}[Hierarchy-reduced mass moments and two-dimensional return]
\label{p6:thm:v2189-hierarchy-mass-moments}
Use the actual output coordinates of
Theorem~\ref{p3f:thm:v2168-hierarchy-source-return}:
\[
 L=(g_e,g_\mu),\quad W=(h_e,h_\tau),\qquad
 \iota_e=LAW^\dagger,\quad
 \iota_\mu=Lbc^\dagger,\quad
 \iota_\tau=cd^\dagger W^\dagger.
\]
Here $A\in M_2(\mathbb C)$ and $b,d\in\mathbb C^2$ are evaluations of the
existing source arrows, not adjustable inputs. For
$R=\sum_r\varepsilon_r\iota_r$, put $U=(g_e,g_\mu,c)$ and
\[
 B=U^\dagger RR^\dagger U=
 \begin{pmatrix}
 \varepsilon_e^2AA^\dagger+\varepsilon_\mu^2bb^\dagger
       &\varepsilon_e\varepsilon_\tau Ad\\
 \varepsilon_e\varepsilon_\tau d^\dagger A^\dagger
       &\varepsilon_\tau^2d^\dagger d
 \end{pmatrix}.
\]
For the constant mass-only marked generator $H_M=d\Gamma(h_R)$, put
$v=(g_e,0)^T$ and $G_M(z)=\langle v,(zI-h_R)^{-1}v\rangle$;
let $p_M$ be the first monic annihilator of this marked source vector.
The source moments
$\mu_n^M=\omega(\{q_e^\dagger,\operatorname{ad}_{H_M}^nq_e\})$
are state independent and satisfy
\begin{align}
 \mu_0^M&=1,&\mu_1^M&=\mu_3^M=0,\\
 \mu_2^M&=B_{11}
  =\varepsilon_e^2(|A_{11}|^2+|A_{12}|^2)
                    +\varepsilon_\mu^2|b_1|^2,
 \label{p6:eq:v2189-mass-second-moment}\\
 \mu_4^M-(\mu_2^M)^2
 &=\left|\varepsilon_e^2(AA^\dagger)_{21}
                 +\varepsilon_\mu^2b_2\overline{b_1}\right|^2
       +\varepsilon_e^2\varepsilon_\tau^2|(Ad)_1|^2.
 \label{p6:eq:v2189-mass-return-defect}
\end{align}
There is no direct tau term in $\mu_2^M$; the shared right plane can return
tau incidence in $\mu_4^M$.

The two return scalars in \eqref{p6:eq:v2189-mass-return-defect} vanish
exactly when $Be_1=B_{11}e_1$. On that actual marked certificate, for
$B_{11}>0$ the electron mass-source cyclic space has dimension two and
\[
 G_M(z)=\frac{z}{z^2-B_{11}},\qquad p_M(z)=z^2-B_{11}.
\]
For $B_{11}=0$, positivity gives $h_R(g_e,0)^T=0$, and instead
$G_M(z)=1/z$ and $p_M(z)=z$. This reduced certificate does not require
the full three-role equality $RR^\dagger=UM^2U^\dagger$ or one chosen
right frame.
\end{theorem}
\begin{proof}
The hierarchy supports eliminate the four mixed muon products already
proved by III-F. In the retained frames, multiplication gives the
displayed $B$. Exact bilinear CAR give
$\operatorname{ad}_{H_M}^n a^\dagger(v)=a^\dagger(h_R^n v)$, so these
moments are $\langle v,h_R^n v\rangle$ independently of the density.
For $v=(g_e,0)^T$, odd powers lie in the opposite chiral port and even
powers have left return $B$. Thus $\mu_2^M=B_{11}$ and
$\mu_4^M=(B^2)_{11}$. Since $B=B^\dagger$,
$(B^2)_{11}-B_{11}^2=|B_{21}|^2+|B_{31}|^2$; substitute the actual
blocks to obtain \eqref{p6:eq:v2189-mass-return-defect}.
Their vanishing makes $h_R^2v=B_{11}v$. If $B_{11}>0$, $v$ and $h_Rv$
are nonzero orthogonal chiral vectors and their span is exactly invariant,
giving the displayed degree-two resolvent and first monic annihilator.
If $B_{11}=0$, $\|h_Rv\|^2=0$, so the source is a nonnull zero-frequency
vector and the first monic annihilator is $z$.
\end{proof}

\paragraph{Actual execution handoff.}
The actual contractions in \eqref{p6:eq:v2189-mass-second-moment} and
\eqref{p6:eq:v2189-mass-return-defect} have not been evaluated here.
The mass-only certificate does not close the complete interacting
electron source polynomial. The latter consumes the whole column
\eqref{p6:eq:v2189-complete-source-column}, its actual stationary support
and the exact closing relation of
Theorem~\ref{p4:thm:v2188-moment-pencil}. Source-visible pole/rest
identification remains the separately typed comparison in
\eqref{p6:eq:physical-rest-mass-bridge}.
Appendix~\ref{p6:app:v2189-electron-execution-lineage} records why primitive
scalar transfers, host transport and normalized current comparisons do
not substitute for evaluation of those existing actual source arrows.

% END INLINE SOURCE: Paper5B-ElectronWordMoments-Main-v1_0.tex



% BEGIN INLINE SOURCE: Paper5B-ElectronMassRoot-Main-v1_0.tex
% Framework v2.191: narrow corollaries of the existing actual mass moments.
\subsubsection{Exact electron mass-source root tests}
\label{p6:sec:v2191-electron-mass-root}
Keep the actual constant mass word, current frame and coordinates
$B,A,b,d$ of Theorem~\ref{p6:thm:v2189-hierarchy-mass-moments}, and write
$(e,m,t)=(\varepsilon_e,\varepsilon_\mu,\varepsilon_\tau)$.
The coefficient $e>0$ is fixed by the existing native compiler before any
source-visible root comparison. The following results sharpen its finite
execution without replacing the complete interacting electron generator.

\begin{corollary}[Electron-only two-pole defect]
\label{p6:cor:v2191-electron-defect}
For the inherited normalized left electron source, define
\[
 \mathcal D_e:=\|(B-e^2I)e_1\|^2
       =\mu_4^M-2e^2\mu_2^M+e^4.
\]
Then
\[
 \boxed{\mathcal D_e=0\quad\Longleftrightarrow\quad
   Be_1=e^2e_1\quad\Longleftrightarrow\quad
   G_M(z)=\frac{z}{z^2-e^2}.}
\]
Equivalently the defect is the following sum of three actual contractions:
\[
 \boxed{\begin{aligned}
 \mathcal D_e={}&
  [e^2((AA^\dagger)_{11}-1)+m^2|b_1|^2]^2\\
 &+|e^2(AA^\dagger)_{21}+m^2b_2\overline{b_1}|^2
   +e^2t^2|(Ad)_1|^2.
 \end{aligned}}
\]
When it vanishes, $p_M(z)=z^2-e^2$, with poles $\pm e$ and residues
$1/2$. This electron-only statement needs neither the full three-role
left-return equality nor a particular right-arrow completion.
\end{corollary}
\begin{proof}
Hermiticity gives
$\|(B-e^2I)e_1\|^2=(B^2)_{11}-2e^2B_{11}+e^4$.
Substitute the existing mass moments and three entries of the first column
of $B$. The nonnegative norm vanishes exactly at the displayed vector
identity, which yields the two-dimensional mass response by
Theorem~\ref{p6:thm:v2189-hierarchy-mass-moments}.
Conversely equality of the rational responses fixes the second and fourth
moments to $e^2$ and $e^4$, so the norm is zero.
\end{proof}


% BEGIN FRAMEWORK v2.239: PROTECTED ELECTRON CORNER FAMILY
\begin{corollary}[Exact protected electron on a non-orthonormal muon--tau family]
\label{p6:cor:v239-protected-electron-corner-family}
Retain the inherited three-corner weighted source of
Theorem~\ref{p6:thm:v2189-hierarchy-mass-moments}, and consider the
explicit output-coordinate \emph{family} (not an action-source assertion)
\[
 A=\begin{pmatrix}1&0\\0&\gamma\end{pmatrix},\qquad
 b=\begin{pmatrix}0\\\beta\end{pmatrix},\qquad
 d=\begin{pmatrix}0\\\delta\end{pmatrix},
 \quad \gamma,\beta,\delta\in\mathbb C.
\]
For $e=\varepsilon_e,m=\varepsilon_\mu,t=\varepsilon_\tau$,
\[
 B=\begin{pmatrix}
 e^2&0&0\\
 0&e^2|\gamma|^2+m^2|\beta|^2&et\gamma\delta\\
 0&et\overline{\gamma\delta}&t^2|\delta|^2
 \end{pmatrix}.
\]
Every member has the same electron-visible mass-only response,
\begin{equation}\label{p6:eq:v239-electron-protected-family}
 \boxed{Be_1=e^2e_1,\qquad\mathcal D_e=0,\qquad
 G_M(z)=\frac{z}{z^2-e^2},\quad p_M(z)=z^2-e^2.}
\end{equation}
The remaining $2\times2$ return has
\[
 \boxed{\Tr B_{\mu\tau}=e^2|\gamma|^2+m^2|\beta|^2
  +t^2|\delta|^2,\quad
 \det B_{\mu\tau}=m^2t^2|\beta\delta|^2.}
\]
For $\gamma=-1/3$, $\beta=\delta=1$ the latter has
$\operatorname{tr}B_{\mu\tau}=m^2+t^2+e^2/9$ and
$\det B_{\mu\tau}=m^2t^2$, even though the complete left return need
not be diagonal. Thus full orthonormal three-role closure is strictly
stronger than the electron two-pole condition.
\end{corollary}
\begin{proof}
Insert the displayed coordinates into the exact return block of
Theorem~\ref{p6:thm:v2189-hierarchy-mass-moments}. Its first row
and column are $(e^2,0,0)$. Corollary~\ref{p6:cor:v2191-electron-defect}
proves the mass-only response and residues. Direct calculation of the
remaining trace and determinant cancels the two $e^2t^2|\gamma\delta|^2$
terms. The statement is constant on this explicitly parameterized
\emph{mathematical} family; membership of the actual mixed action
coefficient remains the separate certificate of
Definition~\ref{p2a:def:v229-action-incidence-certificate}.
\end{proof}
% END FRAMEWORK v2.239: PROTECTED ELECTRON CORNER FAMILY

\begin{corollary}[Complete mass-only source-visible polynomial]
\label{p6:cor:v2191-mass-cofactor}
For the actual Hermitian positive return $B$, set
\[
 \chi_B(s)=\det(sI_3-B),\qquad
 N_B(s)=\operatorname{cof}_{11}(sI_3-B).
\]
The exact mass-only electron response and its first monic source polynomial
are
\[
 \boxed{G_M(z)=\frac{zN_B(z^2)}{\chi_B(z^2)},\qquad
 p_M(z)=\frac{\chi_B(z^2)}
              {\gcd(\chi_B(z^2),zN_B(z^2))}.}
\]
This includes invisible, degenerate and zero roots; no root value is imposed.
\end{corollary}
\begin{proof}
The existing chiral block identity gives
$G_M(z)=z e_1^\dagger(z^2I_3-B)^{-1}e_1$.
The adjugate formula gives its numerator. Also
$\det(zI_6-h_R)=\det(z^2I_3-RR^\dagger)=\chi_B(z^2)$;
the block computation away from $z=0$ extends as a polynomial identity.
Apply Theorem~\ref{p4:thm:v2191-source-cofactor} to the complete Hermitian
mass matrix $h_R$ and its specified electron vector. The reduced monic
denominator is the exact source-vector annihilator, including at zero.
\end{proof}

\paragraph{Actual execution and scope.}
The three contractions have not been evaluated here. They may be proved
directly on the existing source owner; evaluating every directed matrix
entry or choosing a unique right completion is not a prerequisite for this
weaker claim. Alternatively an independently evaluated occurring matrix
can be read through Corollary~\ref{p3f:cor:v2191-action-corner-inverse}.
A nonzero $\mathcal D_e$ excludes this exact two-pole response but does
not exclude a visible $e$ pole together with other modes; compute $p_M$
for that question. The complete interacting support polynomial still
consumes the entire source column, complete action and licensed stationary
support. Particle/rest interpretation and consumed limits remain separate.
The current mass coefficients and historical fourth-order comparisons
retain their existing scale domains; no new matching factor is inserted.
Appendix~\ref{p6:app:v2191-electron-arrow-limits} explains why mathematical
arrow availability and free masses alone do not select a physical source word.

% END INLINE SOURCE: Paper5B-ElectronMassRoot-Main-v1_0.tex


% BEGIN NEUTRAL-RANK-WORD-SCOPE v2.208
\subsection{Charged transfer tests with the source origin retained}
\label{p5b:sec:v228-source-character-consumption}
\noindent\textbf{Current scope of earlier assembly labels.}
The complete charged packet is a generated input when its hierarchy arrows
have the certificate of Definition~\ref{p5:def:v228-hierarchy-origin}.
Proposition~\ref{p5:prop:v228-charged-origin-scope} states the current
conditional source handoff. The inverse channel map and linear synthesis
are explicit; the conditional arrow export does not independently construct
its inputs. A missing origin certificate must be distinguished from an
unprinted entry table of a fully specified constructor.
\begin{corollary}[Complete source-cell electron test]
\label{p5b:cor:v228-source-character-test}
If the actual certified row expressions have the form
$T=eF_\gamma C_e+mB_\mu+tB_\tau C_\tau$ of
Theorem~\ref{p1b:thm:v228-character-defect}, its positive defect is a
necessary and sufficient mass-only electron certificate. For the full plane
corner $\gamma=-1/3$, the exact requirements are unit modulus on both
electron plane characters and equal relative electron--tau action there.
For the rank-one corner $\gamma=0$, only the weighted electron norm and
relative equality are required. The muon coefficient is absent from this
marked test. These character values are outputs of the actual word, not
physical parameters to be chosen by a desired spectrum.
\end{corollary}
\begin{proof}
Substitute the evaluated character values in the geometric defect identity
and apply Corollary~\ref{p3f:cor:v228-derived-electron-pair}.
\end{proof}
\begin{corollary}[Mass-only closure from a proved complete output fibre]
\label{p5b:cor:v228-source-fibre-readout}
If the actual rank-one exports lie in the normalized opposite-odd family of
Proposition~\ref{p1b:prop:v228-opposite-odd-fibre}, their complete left
return is $\sum_r\varepsilon_r^2P_r$ throughout that family. This readout
requires no selection of its member, by
Corollary~\ref{p2a:cor:v228-constant-source-fibre}.
\end{corollary}
\begin{proof}
The geometric proposition proves the return constant on the complete
specified family; the hypothesis identifies it with the actual output
domain. Generated common-readout consumption then applies.
\end{proof}
The terminal coefficients and pure structural current rays keep their
existing meanings. Neither this conditional test nor the source-cell tuple
defines a physical pole, chooses a row function, or equates terminal and
spectral masses without the displayed relation certificate.

% END FRAMEWORK v2.228: CHARGED ORIGIN AND CHARACTER TESTS

% BEGIN INLINE SOURCE: Paper5B-HistoricalFermionAssembly-Appendix-v1_0.tex
% Framework v2.187: superseded charged comparison/incidence frontiers.
\subsection{derived charged terminal compiler}
\label{p6:sec:v566-derived-charged-terminal-compiler}

\begin{definition}[Primitive charged finite node]
\label{p6:def:v566-primitive-charged-node}
A charged node is evaluated only after all mandatory finite refinements
already generated by its owner.  It is eligible as an independent charged
node only when the remaining block carries a primitive non-scalar owner type
in the sense of Paper~I-C.  This is not the rule ``even ranks are particles.''
\end{definition}

\begin{theorem}[Primitive charged node set]
\label{p6:thm:v566-primitive-charged-node-set}
For the charged grammar one has
\[
 \boxed{\mathcal N_{\rm ch}^{\rm prim}=\{2,4,6\}}
\]
with owner identifications
\[
 \boxed{2=\dim u^\perp,\qquad4=\dim V,\qquad6=\dim\Lambda^2V.}
\]
The exclusions are geometric:
\[
 1:\text{ scalar},\qquad
 3:\text{ killed by the }C_6\to C_3\text{ owner quotient},
\]
\[
 8=1+7':\text{ nonprimitive complement},\qquad
 12=6_+\oplus6_-:\text{ composite incidence layer},
\]
\[
 16=7+9:\text{ refinement parent}.
\]
Hence the canonical node words are \(N_2,N_4,N_6\), with reverse readouts
where the already typed terminal orientation requires the Moore--Penrose
inverse.
\end{theorem}

\begin{proposition}[Heavy hinge refinement]
\label{p6:prop:v566-heavy-hinge-refinement}
With the same-carrier projectors of
Theorem~\ref{p6:thm:v179-hurwitz-unital-completion}, set
\[
 \Pi_1^{(8)}:=\frac18|\Omega_8\rangle\langle\Omega_8|,
 \qquad W_7:=P_8-\Pi_1^{(8)}.
\]
Then
\[
 P_8=\Pi_1^{(8)}+W_7,\qquad \operatorname{rank}W_7=7,
\]
so the rank-eight complement has the compulsory refinement
\[
 \boxed{8=1+7'}.
\]
Together with
\[
 \boxed{16=7+9},\qquad \boxed{16+1\to17},
\]
the rank-sixteen block is a hinge parent: it is simultaneously an internal
refinement parent and the recipient of the minimal unital completion.  Thus
rank eight does not define an additional charged node.
\end{proposition}
\begin{proof}
The line \(\mathbb C\Omega_8\) is the orthogonal line adjoined to
\(P_{16}\) in Theorem~\ref{p6:thm:v179-hurwitz-unital-completion}.  Its
orthogonal complement inside the rank-eight block has rank seven, giving the
displayed decomposition.
\end{proof}

\begin{theorem}[Derived charged raw terminal-word compiler]
\label{p6:thm:v566-derived-charged-raw-compiler}
The charged finite grammar produces the following raw structural words:
\[
\begin{array}{c|c|c}
p&F_p&\text{geometric provenance}\\ \hline
 t&P_{(7,9)}&\text{inner heavy }16=7+9\text{ split}\\
 b&P_{(8,16,1)}^+&\text{outer heavy split plus unital completion}\\
 e&P_{(1,5)}^+&\text{outer }6=5+1\text{ split}\\
 u&N_2&\text{unique non-scalar }q=5\text{ primitive owner}\\
 d&P_{(2,3)}&\text{inner }5=3+2\text{ split}\\
 \mu&N_4^+&V\text{-owner node}\\
 s&N_6^+&\Lambda^2V\text{-owner node}\\
 c&P_{(4,4,4)}&K_{4,4,4}\text{/triality split path}\\
 \tau&P_{(6,12,6)}&\text{primal--dual incidence split path}
\end{array}
\]
Thus the raw singular factors are
\[
 \boxed{
 \frac1{1+\sqrt5},\ \sqrt2,\ \sqrt2+\sqrt3,\
 \frac12,\ \frac1{\sqrt6},\ 6,\ 2\sqrt6+\sqrt{12},\
 \sqrt7+3,\ \frac1{\sqrt8+4+1}}
\]
in the row order \(e,u,d,\mu,s,c,\tau,t,b\).
These are structural terminal words, not pole-mass predictions.  Their
promotion to the normalized physical Yukawa/Dirac member remains the G15
same-action source-occurrence problem of Paper~V-A.
\end{theorem}

\begin{remark}[Neutral sector unchanged]
\label{p6:rem:v566-neutral-unchanged}
No statement in this release changes the \(K=56\) light-neutrino lookup or
its separate split/realization grammar.
\end{remark}
% END CHARGED-PRIMITIVE-OWNER v5.66

% BEGIN D4-TERMINAL-NEUTRINO v5.67

\section{Charged source incidence and rooted family transport}
\label{p5b:sec:organized-source}
\subsection{Charged-family branch and absolute normalization}
\label{p6:sec:charged-family}

Let
\[
\RR^3=\RR e_{\mathbf1}\oplus A_2,
\qquad
e_{\mathbf1}=\frac1{\sqrt3}(1,1,1),
\]
and decompose a charged root vector as
\(
r=r_{\mathbf1}+r_\triangle
\).

\begin{proposition}[Conditional norm duality and Koide implication]
\label{p6:prop:koide-implication}
For \(h=\|r_{\mathbf1}\|>0\) and \(d=\|r_\triangle\|>0\), the map
\[
\mathcal B_L(r)
=
\frac dhP_{\mathbf1}r+\frac hdP_{A_2}r
\]
is a nonlinear, homogeneous, \(S_3\)-equivariant involution exchanging
\((h,d)\).  Its fixed cone is the locus \(h=d\), and that locus implies
\[
\frac{m_e+m_\mu+m_\tau}
{(\sqrt{m_e}+\sqrt{m_\mu}+\sqrt{m_\tau})^2}
=\frac23.
\]
The implication is exact; the norm-exchange map is a charged-family bridge,
not an \(S_3\)-module automorphism.
\end{proposition}

\begin{proof}
The projectors commute with \(S_3\), and the two norms are invariant.
Applying the map twice restores each component.  On the fixed cone,
\[
\|r\|^2=2h^2,
\qquad
\left(\sum_i r_i\right)^2=3h^2,
\]
which gives the ratio.  Schur's lemma shows that a linear equivariant map
preserves the nonisomorphic summands \(\mathbf1\) and \(A_2\); it cannot
exchange them.
\end{proof}

\begin{lemma}[Ordered Koide root]
\label{p6:lem:ordered-koide-root}
With
\[
u=\sqrt{m_e/m_\tau},
\qquad
v=\sqrt{m_\mu/m_\tau},
\]
the Koide equation gives
\[
v_\pm(u)=2(1+u)\pm\sqrt{3(u^2+4u+1)}.
\]
The ordered chamber \(0<u<v<1\) selects \(v_-(u)\), but it does not by
itself select a unique intersection with a separate shell curve.
This is the algebraic form of the charged-lepton Koide relation
\pvcite{Koide1983}.
\end{lemma}

\begin{definition}[Retained charged shell template]
\label{p6:def:charged-shell-template}
On a frozen \(\mathrm{REG}_{\sigma}\) member set
\[
\delta_{\rm geom}=\frac1{\sqrt{72}},
\qquad
\delta_{\sigma}=\frac1{6\sqrt2}+\sigma_{\rm fin},
\]
\[
\log\alpha_{\rm hier}=4\pi^2+\delta_{\sigma},
\qquad
A_8=\frac{14+1/12}{8},
\qquad
B_L=A_8\log\alpha_{\rm hier}.
\]
The retained charged center, denominator, and bow are
\[
\widetilde c_L
=
\frac12-2\delta_{\rm geom}^2
-\frac23\delta_{\rm geom}^3
+\frac14\delta_{\rm geom}^4
+\frac18\delta_{\rm geom}^5,
\]
\[
C_L=1+\widetilde c_L,
\qquad
q_L=\frac95+\frac32\delta_{\rm geom}^3.
\]
For \(0<\Delta<\widetilde c_L\), define
\[
u(\Delta)
=
\exp\left[
-\frac{B_L}{2}
\left(
\frac1{C_L-\Delta}-\frac1{C_L+\Delta}
\right)
\right],
\]
\[
v_{\rm shell}(\Delta)
=
\exp\left[
-\frac{B_L}{2}
\left(
\frac1{C_L+q_L\Delta^2}-\frac1{C_L+\Delta}
\right)
\right].
\]
The associated retained shell coordinates are
\[
\widetilde n_{L,1}=\widetilde c_L-\Delta,
\qquad
\widetilde n_{L,2}=\widetilde c_L+q_L\Delta^2,
\qquad
\widetilde n_{L,3}=\widetilde c_L+\Delta.
\]
Finally, the displayed fourth-order width/amplitude germ is
\[
\begin{aligned}
\Delta_L^{[4]}(p)
&=
\frac{C_L^2}{B_L}p
+\frac{C_L^3(C_Lq_L+1)}{3B_L^2}p^2\\
&\quad
+\frac{C_L^4(2C_L^2q_L^2+4C_Lq_L-7)}
{9B_L^3}p^3\\
&\quad
+\frac{C_L^5(C_Lq_L+1)
(5C_L^2q_L^2+C_Lq_L-31)}
{27B_L^4}p^4 .
\end{aligned}
\]
All displayed expressions are retained-order bridge definitions; no omitted
term is asserted to vanish.
\end{definition}

\begin{theorem}[At least two charged shell/Koide intersections]
\label{p6:thm:two-koide-intersections}
For the retained fourth-order shell template \emph{at the anchored
\(\mathrm{REG}_0\) member}, with
\(\log\alpha_{\rm hier}=4\pi^2+1/(6\sqrt2)\), the equation
\[
v_{\rm shell}(\Delta)=v_-(u(\Delta))
\]
has at least two positive solutions. The same two transverse branches
persist in a sufficiently small open neighborhood of this member by the
implicit-function theorem; no two-root assertion is made for all positive
\(\mathrm{REG}_\sigma\) schemes.  Their numerical isolations and
coordinates are
\[
\begin{array}{c|cc}
\Delta&u(\Delta)&v_{\rm shell}(\Delta)\\
\hline
0.125314767527813&
0.0171563829158774&
0.243570009483205\\
0.377358234148415&
2.23985349328759\times10^{-6}&
0.267945913071040 .
\end{array}
\]
Both satisfy
\[
0<\Delta<\widetilde c_L<q_L^{-1}<C_L,
\qquad
0<u<v<1,
\]
and both zeros are transverse.  Their positive displayed amplitudes are
\[
p_1=3.64964007950150,
\qquad
p_2=9.22071322448997.
\]
\end{theorem}

\begin{proof}
Put \(F=v_{\rm shell}-v_-\circ u\).  Direct outward-rounded evaluation gives
\[
\begin{array}{c|r}
x&F(x)\\
\hline
0.125314766&>3.67\times10^{-9}\\
0.125314769&<-3.53\times10^{-9}\\
0.377358233&<-1.74\times10^{-9}\\
0.377358235&>1.29\times10^{-9}.
\end{array}
\]
Continuity therefore gives one zero in each of the two disjoint intervals.
Substitution gives the displayed coordinates and inequalities.
Outward-rounded evaluation of the differentiated elementary expression gives
\[
-2.405<F'(\Delta)<-2.402
\quad\text{on the first isolating interval},
\]
\[
1.520<F'(\Delta)<1.524
\quad\text{on the second isolating interval}.
\]
In particular both roots are transverse.  Their point values are
\[
-2.40395978348\ldots,
\qquad
1.52216690983\ldots,
\]
so both roots persist under small perturbations.  The fourth-order
width-amplitude coefficients on this slice are approximately
\[
0.031050108,\quad
0.000797698,\quad
0.000027156,\quad
0.000000269,
\]
all positive.  Hence \(\Delta_L^{[4]}\) is strictly increasing for \(p>0\)
and assigns one positive amplitude to each root.
\end{proof}

\begin{definition}[Retained-order charged-root realization family]
\label{p6:def:principal-charged-branch}
At the anchored \(\mathrm{REG}_0\) member, the explicit fourth-order
retained equation has the two transverse positive members
\[
\boxed{
\mathfrak A_{L,[4]}=\{\Delta_{L,1},\Delta_{L,2}\}.
}
\]
Write \(\Delta_{L,1}\) for the root isolated in
\[
(0.125314766,\ 0.125314769).
\]
\(\mathsf F_{L,[4]}^{(1)}\) retains this root as the source representative of the
displayed retained-order model and defines \(p_{L,1}\) as the unique positive
solution of
\[
\Delta_L^{[4]}(p_{L,1})=\Delta_{L,1}.
\]
Numerically,
\[
p_{L,1}=3.64964007950150\ldots .
\]
The nearby bare carrier amplitude is
\[
p_L^{(0)}=\frac72+\frac1{56}.
\]
Proximity motivates the name ``principal'' but is not an upstream selector.
This definition chooses the fourth-order polynomial as a finite
representative; it does not assert that the unknown exact completion has
zero higher coefficients.  An exact charged theorem would require
\(\mathsf F\) to supply the complete analytic function and a chamber
selecting its root.
\end{definition}

\begin{definition}[Sole retained-order charged absolute map]
\label{p6:def:charged-absolute-map}
The charged-lepton coordinates in the conditional retained-order branch use
only
\[
\boxed{
m(\widetilde n)
=
M_P\alpha_{\rm hier}^{-\frac{14+1/12}
{8(1+\widetilde n)}}.
}
\]
The three shell coordinates are obtained from the selected charged center,
the principal width, and the retained middle bow.  The direct
\(E_{36}\) radical rows in Table~\ref{p6:tab:conditional-particle-ledger}
are support diagnostics only and are not a second absolute mass
normalization.
Explicitly, for the selected first intersection \(\Delta=\Delta_{L,1}\),
\[
\boxed{
(m_e,m_\mu,m_\tau)
=
\left(
m(\widetilde n_{L,1}),
m(\widetilde n_{L,2}),
m(\widetilde n_{L,3})
\right),
}
\]
with the three \(\widetilde n_{L,i}\) of
Definition~\ref{p6:def:charged-shell-template}.
\end{definition}

\begin{corollary}[Actual charged output of the selected branch]
\label{p6:cor:actual-charged-output}
On the source \(\mathrm{REG}_0\) member, \(\sigma_{\rm fin}=0\), with
\[
E_{32}=30.9845880638852\,{\rm GeV},
\qquad
M_P=8E_{32}\alpha_{\rm hier},
\]
the selected root gives
\[
(\widetilde n_{L,1},\widetilde n_{L,2},\widetilde n_{L,3})
\simeq
(0.345867308,\ 0.499487455,\ 0.596496843)
\]
and hence
\[
\boxed{
(m_e,m_\mu,m_\tau)_{[4]}
\simeq
(1.251549\,{\rm MeV},\
252.257466\,{\rm MeV},\
4.252031\,{\rm GeV}).
}
\]
These values obey the selected root ratios and Koide relation.  They are not
the \(E_{40}/E_{36}\) diagnostic coordinates in
Table~\ref{p6:tab:selected-finite-coordinates}.  The present bridge cannot
retain both absolute charged ledgers simultaneously.
\end{corollary}

\begin{remark}[Exact-completion status]
The local shell-width inversion has a nonzero linear coefficient and hence a
formal local inverse.  No global convergence radius or numerical remainder
bound is presently derived.  The values in
Corollary~\ref{p6:cor:actual-charged-output} are therefore coordinates of the
explicit \([4]\)-model, not exact predictions of an unspecified analytic
completion.
\end{remark}

\subsection{canonical charged current-incidence evaluations}
\label{p6:sec:v594-charged-current-incidence-closure}

\begin{definition}[Charged terminal-role incidence evaluation]
\label{p6:def:v594-charged-role-incidence-evaluation}
For each charged sector $s\in\{u,d,e\}$ let
$\mathscr L_s^{\rm role}$ be its generated three-role terminal space.  The
closed mass/incidence packet of Release v5.90 supplies the actual role-indexed
left-current incidence evaluation
\[
 E_s:\mathscr L_s^{\rm role}\longrightarrow G_L
\]
as a descendant of the same rooted charged source word.  On the current
normalization domain where that actual evaluation has full rank, put
\[
 Z_s:=E_s^\dagger E_s>0,
\]
define the canonical normalized incidence
\[
 \boxed{W_s:=E_sZ_s^{-1/2}.}
\]
Then $W_s^\dagger W_s=I_{\mathscr L_s^{\rm role}}$.  No independent family
isometry is introduced: $E_s$ is the current-owner evaluation of the already
generated incidence siblings $\iota_{s,r}$.
\end{definition}

\begin{theorem}[Charged current evaluations are generated siblings]
\label{p6:thm:v594-charged-current-evaluations-generated}
On the rooted charged branch, the complete source closure and same-action
current owner co-consume the packet
\[
 \bigl((m_{s,r}^{\rm fin},\iota_{s,r})\bigr)_r
\]
before terminal scalarization.  Consequently $E_s$, $Z_s$ and, on the
displayed full-rank domain, $W_s$ above are
dependent readouts of that already generated packet and require no new
selector.  In particular $W_u,W_d,W_e$ are fixed up to simultaneous role
presentation/rephasing transformations carried together with their role
labels.
\end{theorem}

\begin{corollary}[Quark weak-current incidence is ready for same-action pullback]
\label{p6:cor:v594-quark-current-ready}
The two quark terminal-role spaces are now equipped with generated normalized
maps on that same current-normalization domain,
\[
 W_u:\mathscr L_u^{\rm role}\xrightarrow{\cong}G_L,
 \qquad
 W_d:\mathscr L_d^{\rm role}\xrightarrow{\cong}G_L.
\]
Thus no mass-eigenframe, Yukawa diagonalization or additional $J_u,J_d$ is
required to form the physical charged-current comparison.  Paper~V-C performs
the weak-current contraction.
\end{corollary}

\begin{proposition}[Actual quark-current output certificate]
\label{p6:prop:v2211-quark-current-output-certificate}
On the full-rank current domain of
Definition~\ref{p6:def:v594-charged-role-incidence-evaluation}, the complete
native quark current is the dependent output
\[
 \boxed{C_{ud}=Z_u^{-1/2}E_u^\dagger E_dZ_d^{-1/2},\qquad
 C_{ud}^\dagger C_{ud}=C_{ud}C_{ud}^\dagger=I.}
\]
The positivity/rank certificate is
\(\det Z_s=|\det E_s|^2>0\), \(s=u,d\), on their inherited metrics.
With the already bound terminal-role projectors, its observable coefficient
outputs include
\[
 P_{ij}=\operatorname{tr}
 (\Pi_{u,i}C_{ud}\Pi_{d,j}C_{ud}^\dagger)=|(C_{ud})_{ij}|^2,
\]
\[
 J_{ik;jl}=\operatorname{Im}\operatorname{tr}
 (\Pi_{u,i}C_{ud}\Pi_{d,j}C_{ud}^\dagger
  \Pi_{u,k}C_{ud}\Pi_{d,l}C_{ud}^\dagger).
\]
The actual entries, rank determinants and these contractions must be
evaluated on the retained source before a numerical quark-flavor claim.
No entry table is supplied by the four legacy CKM comparison inputs in
Definition~\ref{p6:def:selected-ckm-branch}.  Identifying that benchmark
with the native current requires an actual rephasing-invariant comparison.
This is dependent execution of an existing source/current, not a new CKM
selector or a reconstruction of terminal masses as eigenvalues.
\end{proposition}
\begin{proof}
Insert \(W_s=E_sZ_s^{-1/2}\) into \(W_u^\dagger W_d\).
Each full-rank map between the two three-dimensional Hilbert carriers has
a unitary polar factor, proving both unitarity identities.  The determinant
identity is multiplicativity of determinant under dagger.  In role bases,
the first trace is the squared modulus of one entry and the second is the
imaginary part of
\((C_{ud})_{ij}\overline{(C_{ud})_{kj}}
 (C_{ud})_{kl}\overline{(C_{ud})_{il}}\).
Role rephasings cancel in these products.  An arbitrary physical
normalization or distinct actual operator is not inserted by this algebra.
\end{proof}


% BEGIN FRAMEWORK v2.239: QUARK JOINT GRAM EXECUTION
\begin{corollary}[Complete joint Gram and no independent relative-frame selector]
\label{p6:cor:v239-quark-joint-gram}
On the already stated full-rank common three-dimensional weak-current
receiver, retain the \emph{action-evaluated} cross packet
\[
 Z_u=E_u^\dagger E_u,\quad Z_d=E_d^\dagger E_d,
 \quad X_{ud}=E_u^\dagger E_d.
\]
Then its block Gram is positive semidefinite of rank three and obeys
\begin{equation}\label{p6:eq:v239-quark-gram-schur}
 \boxed{\begin{pmatrix}Z_u&X_{ud}\\X_{ud}^\dagger&Z_d\end{pmatrix}
  =\begin{pmatrix}E_u^\dagger\\E_d^\dagger\end{pmatrix}
    \begin{pmatrix}E_u&E_d\end{pmatrix},\qquad
 Z_d-X_{ud}^\dagger Z_u^{-1}X_{ud}=0.}
\end{equation}
The physical normalized current is then the dependent unitary
$C_{ud}=Z_u^{-1/2}X_{ud}Z_d^{-1/2}$. The exact raw directed-colour
$a=2,t=1$ candidate has invariant magnitude
$|\mathcal J|=316/58681$ by
Theorem~\ref{p6:thm:v585-cp-reweighting}; all 36 spectral
assignment permutations preserve this \emph{candidate} magnitude.
It is not thereby the physical $C_{ud}$, and no choice of role
permutation substitutes for evaluating $X_{ud}$ on the common action
receiver.
\end{corollary}
\begin{proof}
The Gram factorization is immediate. Since $E_u$ is an isomorphism
onto the three-dimensional receiver,
$E_u(E_u^\dagger E_u)^{-1}E_u^\dagger=I$, hence the Schur complement
vanishes. The polar-normalization formula follows; the candidate
invariance is the separately stated Vandermonde theorem. None of
these algebraic identities evaluates an unevaluated incidence vector
or proves the action-current ownership of a candidate colour frame.
\end{proof}
% END FRAMEWORK v2.239: QUARK JOINT GRAM EXECUTION

% BEGIN FRAMEWORK v2.245: TRANSVERSE GRAM RECEIVER DIMENSION AUDIT
\begin{proposition}[The rooted transverse comparison cannot be the complete three-dimensional CKM Gram]
\label{p6:prop:v245-transverse-gram-rank-six}
For the explicit candidate substitution
$Z_u=Z_d=3I_3$ and
$X_{ud}=\rho=\operatorname{diag}(1,1-2i,1+2i)$,
its joint Hermitian Gram is strictly positive with spectrum
\[
 \boxed{\operatorname{spec}\begin{pmatrix}3I_3&\rho\\
  \rho^\dagger&3I_3\end{pmatrix}
 =\{4,2,3+\sqrt5,3-\sqrt5,3+\sqrt5,3-\sqrt5\}.}
\]
It has rank six and Schur complement
\[
 \boxed{Z_d-X_{ud}^\dagger Z_u^{-1}X_{ud}
 =\operatorname{diag}(8/3,4/3,4/3)>0.}
\]
It therefore cannot equal the complete cross Gram of two full-rank
maps $E_u,E_d:\mathbb C^3\to G_L\cong\mathbb C^3$ on the native
common weak-current receiver. Its normalized transverse return
$\rho/3$ is not unitary and its independently calculated orbit
$|\mathcal J|=316/58681$ is not thereby the physical CKM invariant.
\end{proposition}
\begin{proof}
Each complex diagonal channel produces a Hermitian $2\times2$
block with eigenvalues $3\pm|\rho_j|$, where
$|\rho_j|=1,\sqrt5,\sqrt5$; all eigenvalues are positive.
Direct Schur elimination gives $3I_3-\rho^\dagger\rho/3$.
A $6\times6$ joint Gram of six columns in a three-dimensional
receiver has rank at most three, contradicting the calculated rank six.
The actual current Gram instead satisfies
Eq.~\eqref{p6:eq:v239-quark-gram-schur} once its $E_s$ are
independently evaluated in that receiver.
\end{proof}
% END FRAMEWORK v2.245: TRANSVERSE GRAM RECEIVER DIMENSION AUDIT

\begin{remark}[Hard scope of legacy CKM/PMNS/Takagi sections]
\label{p6:rem:v594-legacy-sm-flavor-dictionary-scope}
Every earlier theorem in this paper formulated with $Y_fY_f^\dagger$ mass
Grams, relative mass eigenframes, Takagi masses, selected CKM/PMNS matrices,
Majorana column phases, $m_\beta$ or $m_{\beta\beta}$ is retained only as an
external Standard--Model comparison dictionary or historical diagnostic.
None may be used to generate or select terminal masses, $W_s$, $C_{ud}$,
$C_{e\nu}$, or a native family member.  A physical claim that explicitly
consumes the external Majorana/seesaw dictionary keeps that additional
hypothesis visible.
\end{remark}
% END CHARGED-CURRENT-INCIDENCE-CLOSURE v5.94

\subsection{role-space reclosure of the terminal \texorpdfstring{$N56$}{N56} response}
\label{p6:sec:v587-n56-role-space-reclosure}

\noindent\textbf{Framework v2.208 branch scope.}
The neutral inverse/direct/closed-Gram entries in this retained section
belong to the historical $D_4$ response constructor. They are obsolete
on the harmonic rank branch of \ref{p6:sec:v2208-neutral-rank-word}.
Their algebraic theorems neither force nor constrain that new mass word.
The charged results and branch-independent geometric/source rules keep
their stated domains.


The historical v5.74 projector-first proof is superseded.  The terminal
inverse/\allowbreak direct/\allowbreak closed-Gram packet is closed before any neutral family
projectors are introduced.  This avoids promoting the central response to a
particular noncentral family operator merely by co-realization.

\begin{definition}[Terminal response-role space]
\label{p6:def:v587-terminal-role-space}
Let
\[
 \mathscr L_{56}
 =\operatorname{span}_{\mathbb C}
   \{e_{\rm inv},e_{\rm dir},e_{\rm Gram}\}
\]
be the formal three-dimensional role-index space for the already licensed
one-pass response ports.  Write
\(\Pi_{\rm inv},\Pi_{\rm dir},\Pi_{\rm Gram}\) for its canonical rank-one
coordinate projectors.  This is a role packet, not an additional physical
family carrier or mass basis.
\end{definition}

\begin{definition}[Complete terminal depth-one response jet]
\label{p6:def:v587-terminal-response-jet}
For a generated positive terminal response $R$, declare the typed finite
response constructor
\[
 \mathfrak J_{56}^{(1)}(R)
 :=(R^+,R,R^\dagger R)
\]
with target the product of the existing inverse, direct and closed-Gram
response sorts.  The constructor packages the complete declared depth-one
output; it does not insert those ports into a fermion action or identify their
target sorts.
\end{definition}

\begin{theorem}[Generated closure of the terminal $K=56$ response jet]
\label{p6:thm:v587-n56-response-jet-closure}
On the rooted terminal $D_4$ response branch let
\[
 R=xI_G,\qquad x^3=192,\qquad x>1.
\]
Then
\[
 \mathfrak J_{56}^{(1)}(R)
 =(x^{-1}I_G,\,xI_G,\,x^2I_G),
\]
and the canonical role-packet encoding is
\[
 \boxed{
 \Sigma_{56}^{\rm role}
 =x^{-1}\Pi_{\rm inv}+x\Pi_{\rm dir}+x^2\Pi_{\rm Gram}.}
\]
All three constituents are nonredundant: their frozen operation signatures are
distinct and their positive scalar readouts are pairwise distinct.  Therefore
Paper~II-A generated co-consumption retains the complete three-role output and
no independent \texttt{N56\_ROLE\_COMPLETION} occurrence token is licensed.
Equivalently,
\[
 \boxed{\texttt{N56\_ROLE\_COMPLETION
 = CLOSED\_DEPENDENT\_ON\_ROOTED\_TERMINAL\_RESPONSE}.}
\]
\end{theorem}

\begin{proof}
For the central positive response, the Moore--Penrose inverse is
$x^{-1}I_G$, the direct response is $xI_G$, and the closed Gram return is
$x^2I_G$.  Since $x=192^{1/3}>1$, the three scalar readouts are nonzero and
pairwise distinct.  The response-jet constructor has these three ports as its
complete frozen output, so Paper~II-A
Theorem~\ref{p2a:thm:v3101-generated-coconsumption} retains them jointly.
No family-space projector or mass eigenframe is used in this argument.
\end{proof}

\begin{corollary}[Terminal role-energy spectrum]
\label{p6:cor:v587-terminal-role-energy-spectrum}
The native $K=56$ row-energy functor gives
\[
 \boxed{\mathcal E_{56}^{\rm role}=E_{56}\Sigma_{56}^{\rm role},}
\]
with ordered role-energy multiset
\[
 \boxed{E_{56}\{x^{-1},x,x^2\},\qquad x^3=192.}
\]
The scale-free terminal role-energy squared-gap quotient remains
\[
 \frac{x^6-1}{x^4-1}=33.3105115561\ldots .
\]
Identification of these terminal row energies with renormalized pole/rest
masses remains a downstream readout problem.
\end{corollary}

\begin{proposition}[Family embedding is downstream of N56 closure]
\label{p6:prop:v587-family-embedding-firewall}
A physical neutral family realization requires a generated isometric map
\[
 J_\nu:\mathscr L_{56}\longrightarrow G_L.
\]
Only after such a map is generated may one form
\[
 T_\nu=J_\nu\Sigma_{56}^{\rm role}J_\nu^\dagger
 =x^{-1}P_{\rm inv}+xP_{\rm dir}+x^2P_{\rm Gram},
 \qquad P_r=J_\nu\Pi_rJ_\nu^\dagger.
\]
The rooted $D_3\cong S_3$ presentation action may simultaneously relabel the
role basis and the columns of $J_\nu$; this covariance does not construct
$J_\nu$ and does not determine a charged/neutral relative frame.  PMNS-type
readouts therefore remain downstream of the N56 role closure.
\end{proposition}

\begin{remark}[Supersession of the v5.74 projector-first proof]
\label{p6:rem:v587-n56-supersession}
Theorem~\ref{p6:thm:v587-n56-response-jet-closure} supersedes the proof route
in historical Release v5.74 that first asserted
$T_{\nu,D_4}=(R^+)P_1+RP_2+(R^\dagger R)P_3$ and then invoked constituent
co-realization.  Co-consumption cannot create that projector-resolved
composite.  The corrected proof closes the response roles first and leaves the
family embedding separately typed.  The numerical terminal magnitude packet
is unchanged; only its logical provenance is repaired.
\end{remark}
% END N56-ROLE-SPACE-RECLOSURE v5.87



% BEGIN TAU-PRIMAL-DUAL-INCIDENCE-CLOSURE v5.96
\subsection{rooted tau primal--dual incidence and $E_{36}$ direction closure}
\label{p6:sec:v596-tau-primal-dual-incidence-closure}
\noindent\textbf{Current scope.} The in-plane open status and the no-consumer remark below are superseded by main-text Release v5.97, which uses the already physical $R^-$ slot relation as an incidence parent rather than as a new mass/source coefficient.  The tau-singlet and direction proofs remain active inputs.
The terminal charged masses are unchanged.  This release evaluates one piece
of the physical charged-lepton current incidence upstream of terminal
scalarization: the singlet endpoint carried by the tau primal--dual owner.

\begin{theorem}[Tau primal--dual owner has a canonical singlet family return]
\label{p6:thm:v596-tau-singlet-return}
For the tau role, retain the rooted primal--dual incidence packet whose layer
multiplicities are $(6,12,6)$ in
Theorem~\ref{p6:thm:v566-derived-charged-raw-compiler}.  Before its positive
codiagonal is terminalized to
\[
 P_{(6,12,6)}=(2\sqrt6+\sqrt{12})I,
\]
form the same-parent family cross return of Paper~I-A,
Corollary~\texttt{p1:cor:v570-\allowbreak rooted-\allowbreak primal-\allowbreak dual-\allowbreak singlet-\allowbreak return}.  Its
normalized left-family image is
\[
 \boxed{v_\tau=c,\qquad P_\tau=Q_0.}
\]
This evaluates the tau family endpoint without reconstructing it from the
terminal scalar and without changing the tau mass owner $E_{36}$.
\end{theorem}
\begin{proof}
The tau compiler already identifies the terminal word's geometric provenance
as the rooted primal--dual $6,12,6$ incidence split, and the rooted family
trimer is a frozen same-root receiver of that parent.  Thus the displayed
$S^\dagger s_{u_+}$ is the exact relation-valued sibling before the positive
codiagonal, not an optional composite available merely because the objects
share a host.  Paper~II-A preterminal relation preservation requires this
common-parent relation to be retained before the noninjective terminal
readout.  The cited Paper~I-A corollary evaluates it as $6\sqrt3\,c$.
Normalization therefore gives the unique line $\mathbb Cc$, hence
$P_\tau=Q_0$.  No inverse of $2\sqrt6+\sqrt{12}$ is used.
\end{proof}

\begin{corollary}[The two charged-lepton $E_{36}$ directions are fixed]
\label{p6:cor:v596-mu-tau-e36-directions}
Under the hierarchy/family bridge,
\[
 E_{36}\longleftrightarrow
 Q_0\End(G)Q_\perp\oplus Q_\perp\End(G)Q_0.
\]
Since the tau left image is $Q_0G$,
\[
 \boxed{\iota_\tau\in Q_0\End(G)Q_\perp,\qquad
 \Delta_{36}(\iota_\tau)=-\iota_\tau.}
\]
The electron remains an $E_{40}$ role and therefore has left image in
$Q_\perp G$.  On the full-rank charged-current branch required by the
normalized evaluation, the muon cannot also terminate on $Q_0G$.  Hence
\[
 \boxed{\iota_\mu\in Q_\perp\End(G)Q_0,\qquad
 \Delta_{36}(\iota_\mu)=+\iota_\mu,}
\]
and after canonical current normalization
\[
 \boxed{P_\tau=Q_0,\qquad P_e+P_\mu=Q_\perp,}
\]
provided the raw current evaluation respects the generated
\(Q_0\oplus Q_\perp\) block decomposition. The individual in-plane
rays are not thereby preserved: identifying each normalized role ray with
a particular preterminal incidence ray additionally requires the role-column
Gram matrix to be diagonal in those lines (or an equivalent source-matched
polar-frame certificate).
Thus the previous discrete $\mu\leftrightarrow\tau$ singlet swap is closed.
\end{corollary}
\begin{proof}
The first statement follows from the $E_{36}$ block typing and the derived tau
range.  If the muon also terminated on $Q_0G$, then the tau and muon columns of
the role evaluation would be collinear, while the electron column lies in the
two-plane $Q_\perp G$; the three-column evaluation would have rank at most two.
The positive metric $Z_e=E_e^\dagger E_e>0$ required for the normalized current
therefore forces the opposite muon direction.  Orthogonality of $Q_0G$ and
$Q_\perp G$ keeps the tau line fixed under the polar normalization, while the
two remaining normalized role lines resolve $Q_\perp G$.
\end{proof}




The retained tau/direction inputs are used by the current in-plane
evaluation below. The superseded open-plane and no-consumer statements are
archived in Appendix~\ref{p6:app:v2187-historical-fermion-assembly}.
% END TAU-PRIMAL-DUAL-INCIDENCE-CLOSURE v5.96



% BEGIN LEPTON-INPLANE-HOPF-PCAD-CLOSURE v5.97
\subsection{charged-lepton in-plane incidence and orientation closure}
\label{p6:sec:v597-lepton-inplane-closure}
Release v5.96 fixed $P_\tau=Q_0$ and the opposite muon/tau $E_{36}$
directions.  Paper~V-A v3.108 now identifies the structural negative-Hopf
family relation as the preterminal parent of the charged-current incidence
sibling.  The remaining plane rays can therefore be evaluated exactly.

\begin{theorem}[Negative-Hopf compression fixes the muon and tau relation rays]
\label{p6:thm:v597-hopf-pcad-lepton-rays}
Put
\[
 c=\frac1{\sqrt3}(1,1,1)^T,
 \qquad Q_0=|c\rangle\langle c|,
 \qquad Q_\perp=I-Q_0,
 \qquad R_-:=R_G^\dagger=\operatorname{diag}(1,i,-i).
\]
The retained relation blocks are
\[
 A_\mu:=Q_\perp R_-Q_0,
 \qquad A_\tau:=Q_0R_-Q_\perp.
\]
They have the normalized rank-one forms
\[
 \boxed{\widehat\iota_\mu=|g_\mu\rangle\langle c|,
 \qquad
 g_\mu=\frac1{2\sqrt6}
 \begin{pmatrix}2\\-1+3i\\-1-3i\end{pmatrix},}
\]
\[
 \boxed{\widehat\iota_\tau=|c\rangle\langle h_\tau|,
 \qquad
 h_\tau=\frac1{2\sqrt6}
 \begin{pmatrix}2\\-1-3i\\-1+3i\end{pmatrix}.}
\]
Equivalently,
\[
 \boxed{P_\mu=\frac98A_\mu A_\mu^\dagger,
 \qquad P_\tau=Q_0.}
\]
The $E_{36}$ direction eigenvalues remain
$\Delta_{36}(\widehat\iota_\mu)=+\widehat\iota_\mu$ and
$\Delta_{36}(\widehat\iota_\tau)=-\widehat\iota_\tau$.
\end{theorem}
\begin{proof}
Directly,
\[
 Q_\perp R_-c
 =\frac1{3\sqrt3}
 \begin{pmatrix}2\\-1+3i\\-1-3i\end{pmatrix},
 \qquad
 \|Q_\perp R_-c\|^2=\frac89.
\]
Thus $A_\mu=(Q_\perp R_-c)\langle c|$ is nonzero rank one and normalizes to
the displayed $g_\mu$.  Taking the adjoint-side support of
$A_\tau=|c\rangle\langle c|R_-Q_\perp$ gives the conjugate plane vector
$h_\tau$.  The direction statements are the already derived block-commutator
eigenvalues and are unchanged by normalization.
\end{proof}

\begin{corollary}[The electron line is the unique remaining plane ray]
\label{p6:cor:v597-electron-plane-ray}
Assume the positive raw current evaluation has the three nonzero incidence
columns along the generated $g_e$, $g_\mu$, and $c$ rays, with their
role-column Gram matrix diagonal. Then polar normalization preserves each
ray, the normalized charged-current role map is an isometry, and the
Release v5.96 block identity gives $P_e+P_\mu=Q_\perp$. Under this
explicit polar-frame hypothesis,
\[
 \boxed{P_e=Q_\perp-P_\mu.}
\]
This is rank one.  One normalized representative is
\[
 \boxed{g_e=\frac12
 \begin{pmatrix}-1-i\\ i\\1\end{pmatrix},
 \qquad P_e=|g_e\rangle\langle g_e|.}
\]
Together with $g_\mu$ and $c$ it gives the exact geometric incidence
frame
\[
 \boxed{W_e=(g_e,g_\mu,c)}
\]
up to phases. It is a physical post-polar charged-current frame only on
the explicit role-column Gram/polar-frame hypothesis above.
\end{corollary}
\begin{proof}
The three normalized left role images are orthonormal.  The tau line is $Q_0G$
and the muon line has just been fixed inside $Q_\perp G$, so the electron is
the unique orthogonal line completing the two-plane.  Direct substitution
checks $\langle c,g_e\rangle=\langle g_\mu,g_e\rangle=0$ and
$\|g_e\|=1$.
\end{proof}

\begin{corollary}[Charged-lepton orientation moment]
\label{p6:cor:v597-zeta-e-closure}
For
\[
 C_{2,G}=-i(X_G-X_G^\dagger),
 \qquad X_Ge_j=e_{j+1},
\]
the derived rays satisfy
\[
 \frac1{\sqrt3}\operatorname{Tr}(P_\mu C_{2,G})=-\frac{\sqrt3}{2},
 \qquad
 \boxed{\zeta_e:=\frac1{\sqrt3}\operatorname{Tr}(P_eC_{2,G})
 =\frac{\sqrt3}{2}.}
\]
Therefore, as a geometric relation-ray statement (with its physical-current
interpretation subject to the preceding polar-frame guard),
\[
 \boxed{\texttt{LEPTON\_QPERP\_ORIENTATION}=\texttt{CLOSED},}
\]
\[
 \boxed{\texttt{CHARGED\_LEPTON\_ORIENTATION\_MOMENT}
 =\texttt{CLOSED}[\sqrt3/2].}
\]
Under complete rooted reversal the projector and $C_{2,G}$ are conjugated
together, so the scalar $\zeta_e$ is unchanged.  The internal Hopf switch by
itself is not a presentation symmetry of the fixed charged $R^-$ sector.
\end{corollary}
\begin{proof}
Evaluation on the displayed unit vectors gives the two expectations
$-3/2$ and $+3/2$.  Alternatively, $P_e+P_\mu=Q_\perp$ and
$\operatorname{Tr}(Q_\perp C_{2,G})=0$ show that they are opposite once one
is evaluated.  If $S$ denotes rooted reversal, then
$SR_-S^{-1}=R_+$ and $SC_{2,G}S^{-1}=-C_{2,G}$ while the derived projectors
transform by $P\mapsto SPS^{-1}$; cyclicity of trace leaves the paired
orientation readout invariant.
\end{proof}

\begin{remark}[Mass/source-value firewall retained]
\label{p6:rem:v597-mass-incidence-firewall}
The closure above uses $R_-$ only as the generated structural parent relation
of the charged-current incidence sibling.  It does not set the charged mass
operator equal to $R_-$, replace any terminal coefficient, or evaluate a
physical pole/rest claim. The rooted charged source contraction is already
closed by Paper~V-A,
Corollary~\ref{p5:cor:v397-charged-preterminal-assembly-closure}; a genuinely
different generic fermion target must declare its own source sort.  The direct
charged mass compiler and its terminal scalars remain unchanged.
\end{remark}
% END LEPTON-INPLANE-HOPF-PCAD-CLOSURE v5.97



% BEGIN INLINE SOURCE: Paper5B-ChargedTransferCertificate-Main-v1_0.tex
% Framework v2.168: integrated finite charged-source derivations.
\subsection{dependent native closure of the terminal three-role response}
\label{p6:sec:v574-n56-role-completion}

\noindent\textbf{Framework v2.208 branch scope.}
The neutral inverse/direct/closed-Gram entries in this retained section
belong to the historical $D_4$ response constructor. They are obsolete
on the harmonic rank branch of \ref{p6:sec:v2208-neutral-rank-word}.
Their algebraic theorems neither force nor constrain that new mass word.
The charged results and branch-independent geometric/source rules keep
their stated domains.


Release v5.70 correctly distinguished mathematical availability of
\(R^+,R,R^\dagger R\) from physical use of an arbitrary licensed operation.
Its reopening of the terminal three-role gate was nevertheless too strong once
the already frozen Paper-V-A source-consumption word is composed with the
Paper-II-A generated-constituent theorem.  The present release performs that
composition and introduces no new constructor or physical premise.

\begin{theorem}[Dependent native closure of the terminal three-role gate]
\label{p6:thm:v574-n56-role-dependent-closure}
Assume G0 actuality of a nonzero terminal-neutral source coordinate in the
already generated Yukawa--Dirac source role.  Then
\[
 \boxed{\mathsf{N56}_{\rm role}}
\]
is a dependent theorem: the occurring claim-complete terminal neutral source
consumes all three one-pass response constituents
\[
 \boxed{R^+,\qquad R,\qquad R^\dagger R.}
\]
No independent neutral-role occurrence token is required.
\end{theorem}

\begin{proof}
Definition v3.74 of Paper~V-A defines
\(\mathsf{N56}_{D_4}^{\rm consume}\) to mean that the actual same-action
\(K=56\) neutral source uses the generated central response through the
terminal neutral response word already declared in
Definition~\ref{p6:def:v567-inverse-direct-gram-branch}.  Corollary v3.75 of
Paper~V-A proves that, on a nonzero terminal-neutral source coordinate, this
consumption relation is a descendant of the same D4 triality-source theorem
and is not an independent assumption.  Hence the occurring neutral source has
the frozen composite presentation
\[
 T_{\nu,D_4}=(R^+)P_1+RP_2+(R^\dagger R)P_3.
\]
Paper~II-A generated-constituent co-realization then retains every
nonredundant generated constituent of that occurring composite in its native
sort.  The three constituents are nonredundant: the orthogonal rank-one
compressions are
\[
 P_1T_{\nu,D_4}P_1=x^{-1}P_1,\qquad
 P_2T_{\nu,D_4}P_2=xP_2,\qquad
 P_3T_{\nu,D_4}P_3=x^2P_3,
\]
with \(x=192^{1/3}>1\).  Deleting any one constituent changes the
corresponding nonzero compression to zero and therefore changes the frozen
operator presentation.  Thus all three response roles are simultaneously
consumed.  This is exactly \(\mathsf{N56}_{\rm role}\).
\end{proof}

\begin{proposition}[Rooted role matching is presentation-only]
\label{p6:prop:v574-rooted-role-matching-orbit}
On the rooted axial branch, let \(P_a,P_b,P_c\) be the generated endpoint
projectors of Proposition~\ref{p6:prop:v570-rooted-axial-neutral-frame-narrowing}.
The generated rooted three-cycle and any affine rooted reflection act on this
projector triple as a 3-cycle and a reflection.  They generate
\(D_3\cong S_3\), hence act transitively on the six bijections between
\(\{R^+,R,R^\dagger R\}\) and \(\{P_a,P_b,P_c\}\).  Therefore the
choice of one printed matching is presentation data, not an additional
physical selector.
\end{proposition}

\begin{proof}
The rooted cycle is the cyclic transport \(X_G\), while every affine rooted
reflection conjugates it to \(X_G^{-1}\).  A 3-cycle together with a
reflection generates the six-element dihedral group \(D_3\), canonically
isomorphic to the full permutation group of three objects.  Postcomposition
therefore connects any two role-to-projector bijections.
\end{proof}

\begin{corollary}[Terminal row-energy spectrum after neutral-source actuality]
\label{p6:cor:v574-e56-terminal-spectrum}
Under the theorem's G0 antecedent,
\[
 \boxed{
 \mathcal E_\nu^{\rm term}
 =E_{56}\bigl(x^{-1}P_{\sigma(1)}+xP_{\sigma(2)}+x^2P_{\sigma(3)}\bigr),
 \qquad x^3=192,
 }
\]
where \(\sigma\in S_3\) is presentation-only modulo the generated rooted
\(D_3\) action.  Normal ordering labels the three eigenvalues by
\(x^{-1}<x<x^2\).  The relative charged/neutral eigenframe and the
terminal-row-energy to renormalized pole/RG map remain separate downstream
obligations.
\end{corollary}

\begin{remark}[Corrected finite-neutral frontier]
\label{p6:rem:v574-corrected-neutral-frontier}
The open gate introduced in Definition~\ref{p6:def:v570-terminal-three-role-gate}
is superseded by Theorem~\ref{p6:thm:v574-n56-role-dependent-closure}.
Accordingly
\[
 \boxed{\texttt{N56\_ROLE\_COMPLETION = CLOSED\_DEPENDENT\_ON\_G0\_NEUTRAL\_ACTUALITY}.}
\]
The surviving neutral obligations are only G0 actuality of a nonzero neutral
source member where a nonzero spectrum is claimed, the full charged/neutral
family-orbit member from which PMNS-type readouts descend, and the downstream
pole/RG analytic certificate.  No new constitutional principle is introduced.
\end{remark}
% END N56-ROLE-COMPLETION v5.74


% BEGIN G15-C3-SYNC v5.75
\subsection{family-orbit readouts after the G0 fermion-source boundary}
\label{p6:sec:v578-family-member-g0-retyping}

Paper~V-A v3.85 closes the finite G15 source geometry structurally while
retaining nonzero physical source/member actuality as a G0 boundary.  The
simultaneous-conjugacy quotient used in this paper is therefore a
presentation-free readout of that occurring member, not a structural object
that GP1--GP3 still have to select.

\begin{theorem}[Actual family/source member is not a residual structural selector]
\label{p6:thm:v578-family-member-g0-retyping}
Let an actual nonzero quark source member determine the positive left-return
pair $(H_u,H_d)$, and write
\[
 \mathcal O_q=[(H_u,H_d)]_{U(3)_{Q_L}}.
\]
The generated structural family grammar does not select one inequivalent
$\mathcal O_q$ from the admissible member family when the already proved
counterfamilies preserve the frozen structural presentation.  The actual
member is therefore typed as
\[
 \boxed{\texttt{G0\_FERMION\_SOURCE\_MEMBER}},
\]
not as an open structural \texttt{FULL\_FAMILY\_MEMBER} constructor.
The analogous statement holds for an actual Dirac-lepton source member.
\end{theorem}

\begin{corollary}[Spectra, relative frames and weak CP are dependent readouts]
\label{p6:cor:v578-family-readouts}
Once the corresponding actual member is fixed, its spectra and relative
left-eigenframes are deterministic invariant readouts.  In particular,
\[
 H_u=U_u\Lambda_uU_u^\dagger,\qquad
 H_d=U_d\Lambda_dU_d^\dagger
\]
give the CKM rephasing class of $U_u^\dagger U_d$, and polynomial invariants
of $(H_u,H_d)$ give weak-CP readouts.  The Dirac-lepton analogue gives the
Dirac-PMNS relative frame.  These are not additional occurrence selectors.
\end{corollary}

\begin{remark}[Structural terminal compiler is not the G0 member]
\label{p6:rem:v578-terminal-compiler-firewall}
The charged raw terminal-word compiler and its rank-$2,4,6$ primitive grammar
remain structural outputs.  Their existence does not by itself constitute a
cross-sort relation identifying those words with the actual normalized
fermion-source coordinate.  Co-realization may preserve a pre-generated
relation inside an occurring composite; it may not invent that relation after
evaluation.  Majorana occurrence and pole/RG/continuum certificates retain
their separate scopes.
\end{remark}

\noindent\textbf{Status retyping:}
\[
 \boxed{\texttt{FULL\_FAMILY\_MEMBER}\ \to\
        \texttt{G0\_FERMION\_SOURCE\_MEMBER}},
\]
with \texttt{CKM\_RELATIVE\_FRAME}, \texttt{PMNS\_DIRAC\_RELATIVE\_FRAME}
and \texttt{WEAK\_CP\_MEMBER} retained as dependent readouts.
% END FERMION-SOURCE-BOUNDARY v5.78


% BEGIN TRANSITIVE-RG-POLE-SPLIT v5.79
\subsection{charged-lepton incidence evaluation correction}
\label{p6:sec:v595-charged-lepton-incidence-correction}
\noindent\textbf{Current scope.} The broad open status recorded in this historical correction is superseded in part by the main-text v5.96 tau primal--dual closure.  Its continuous counterfamily remains useful after exchanging the singlet role from $\mu$ to $\tau$; the one in-plane orientation remains open.

\begin{proposition}[Hierarchy block membership does not determine a role arrow]
\label{p6:prop:v595-block-membership-not-role-arrow}
The hierarchy/family bridge fixes the operator blocks
$\mathscr I_{32}$, $\mathscr I_{36}$ and $\mathscr I_{40}$.  It does not by
itself select a particular element of one of these blocks.  Consequently the
notation $A_{s,r}^{\rm hier}$ in Release v5.89 is a conditional typed role
arrow: its use in
\[
 \iota_{s,r}=\Phi_{\rm HG}(A_{s,r}^{\rm hier})
\]
is valid after the role arrow has been generated, but does not constitute its
generation.  The mathematical terminal support and mass coefficient remain
fully defined independently of this correction.
\end{proposition}

\begin{corollary}[Charged-lepton normalized evaluation remains physical data]
\label{p6:cor:v595-we-formal-open-evaluation}
Release v5.94 correctly defines
\[
 W_e=E_e(E_e^\dagger E_e)^{-1/2}
\]
for a supplied source-owned evaluation $E_e$.  The present source chain does
not evaluate the three charged-lepton role images in $G_L$.  Therefore
\[
 \boxed{
 \texttt{LEPTON\_ROLE\_TO\_FAMILY\_INCIDENCE}
 =\texttt{OPEN\_PHYSICAL\_EVALUATION}.}
\]
The v5.94 existence/normalization construction is retained as a formal
readout; it is not a numerical current-incidence closure.  This correction
does not re-adjudicate the separately established quark-current results.
\end{corollary}

\begin{definition}[Structural exceptional flag and orientation moment]
\label{p6:def:v595-lepton-orientation-moment}
Retain the generated rooted family decomposition
\[
 Q_0=\frac13\mathbf1\mathbf1^\dagger,\qquad
 Q_\perp=I-Q_0,
\]
and the two $C_3$ character projectors
\[
 E_\pm=\frac12\left(Q_\perp\pm\frac{C_{2,G}}{\sqrt3}\right).
\]
This is a structural flag; it is not by itself the charged-lepton mass/current
frame.  For an evaluated electron projector $P_e$, define the dimensionless
owner readout
\[
 \boxed{\zeta_e:=\frac1{\sqrt3}\operatorname{Tr}(P_eC_{2,G}).}
\]
Set
\[
 \boxed{
 \texttt{CHARGED\_LEPTON\_ORIENTATION\_MOMENT}
 =\texttt{OPEN\_SOURCE\_EVALUATION}.}
\]
\end{definition}

\begin{theorem}[Continuous in-plane counterfamily]
\label{p6:thm:v595-lepton-inplane-counterfamily}
Assume only the stronger partial placement
\[
 P_\mu=Q_0,\qquad P_e+P_\tau=Q_\perp.
\]
Choose normalized character vectors $e_\pm$ with projectors $E_\pm$.  Every
orthonormal charged-lepton flag compatible with these two equations contains
a representative
\[
 v_e=\cos\theta\,e_+ +e^{i\phi}\sin\theta\,e_-,
\]
\[
 v_\tau=-e^{-i\phi}\sin\theta\,e_+ +\cos\theta\,e_-,
\qquad
 v_\mu=\frac1{\sqrt3}(1,1,1)^T,
\]
up to role-line phases.  All charged terminal mass coefficients are unchanged,
while
\[
 \boxed{\zeta_e=\cos(2\theta).}
\]
Hence terminal magnitudes, the singlet/plane split and the existence of the
structural flag do not fix the charged-lepton CP orientation.
\end{theorem}

\begin{proof}
The displayed pair is the general $U(2)$ rotation of an orthonormal basis of
$Q_\perp$ after quotienting independent role-line phases.  Since
$C_{2,G}e_\pm=\pm\sqrt3e_\pm$, direct substitution gives
$\operatorname{Tr}(P_eC_{2,G})=\sqrt3(\cos^2\theta-\sin^2\theta)$.
\end{proof}

\begin{remark}[Structural flag firewall]
\label{p6:rem:v595-structural-flag-firewall}
Assigning $(e,\mu,\tau)$ to $(E_+,Q_0,E_-)$ would set $|\zeta_e|=1$, but this
assignment is not licensed by the positive/inverse terminal transfer words
alone.  It requires a source-owned role-to-family incidence theorem.
\end{remark}
% END CHARGED-LEPTON-INCIDENCE-EVALUATION-CORRECTION v5.95



% BEGIN INLINE SOURCE: Paper5B-ChargedTransferDiagnostics-Appendix-v1_0.tex
% Framework v2.168: integrated finite charged-source derivations.
\subsection{Charged source evaluation lineage and comparison scope}
\label{p6:app:v2189-electron-execution-lineage}

This appendix records execution bookkeeping and comparison cautions for
Section~\ref{p6:sec:v2189-electron-word-moments}; it supplies no premise
hidden from that positive chain. Its formulas retain the same licensed
canonical chart and actual-owner typing. The conclusion is a finite list
of contractions to evaluate, not a new charged source-member selector:
the rooted assembly remains closed by
Corollary~\ref{p5:cor:v397-charged-preterminal-assembly-closure}.

\subsubsection{Primitive evaluation, directed owner arrows and source synthesis}
\label{p6:app:v2189-directed-arrow-lineage}
Paper I-C's primitive owner maps are explicit. In their actual types,
\[
 \lambda_nN_n(z)=zI_{E_n},\qquad
 \epsilon_{E_n}\lambda_n=N_n^+.
\]
For $n=4$, the normalized central owner direction $s_4=I_4/2$ obeys
$\epsilon_{E_4}(s_4)=1/2$, as in
Theorem~\ref{p1b:thm:v2168-primitive-owner-support}.
These evaluate primitive transfers and retain their owner support.
The retraction has type $\operatorname{End}(E_4)\to\mathbb C$; it is
not itself a directed mixed-block arrow $Q_\perp G\leftarrow Q_0G$.
A scalar terminal split likewise does not identify a physical incidence
operator with a normalized current arrow.

The rooted $q=5$ operator frame and its Gram normalization are supplied by
Theorem~\ref{p1b:thm:v137-q5-operator-trimer}; the unitary $J_{\rm HG}$
and $*$-isomorphism
\[
 \Phi_{\rm HG}(X)=J_{\rm HG}^\dagger XJ_{\rm HG}
\]
are supplied by
Theorem~\ref{p1b:thm:v138-hierarchy-family-block-isomorphism}.
They transport an existing hierarchy arrow into the common generated
family host. The owner-resolved packet of
Definition~\ref{p6:def:v588-owner-resolved-charged-packet} retains that
arrow; Proposition~\ref{p6:prop:v589-mass-incidence-sibling-export}
explicitly distinguishes its evaluation from block membership.
The signed map $\chi$ gives the channel representative
$W_r=\chi^{-1}(\iota_r)$, and
Theorem~\ref{p6:thm:v590-charged-channel-representatives} supplies the
linear synthesis $R=\sum_r\varepsilon_r\iota_r$.
An inverse presentation map evaluates a supplied input; it does not
replace a missing row-entry evaluation by a new physical choice.

For the already occurring actual hierarchy source arrows, the output
coordinates of the main theorem are precisely
\[
 \boxed{\begin{aligned}
 A&=L^\dagger\Phi_{\rm HG}(A_e^{\rm hier})W,\\
 b&=L^\dagger\Phi_{\rm HG}(A_\mu^{\rm hier})c,\\
 d&=W^\dagger\Phi_{\rm HG}(A_\tau^{\rm hier})^\dagger c.
 \end{aligned}}
\]
For example, the first electron scalar contribution to the marked mass
second moment is the actual contraction
\[
 g_e^\dagger\Phi_{\rm HG}
       (A_e^{\rm hier}(A_e^{\rm hier})^\dagger)g_e,
\]
with its muon counterpart also retained. The two actual return scalars
in \eqref{p6:eq:v2189-mass-return-defect} test the smaller marked cyclic
claim directly. These are evaluation outputs of a declared rooted
source, not adjustable parameters or an arbitrary family of sources.
The packet and host formulas above do not print these numerical
contractions. Their execution remains separate from the closed charged
source-value realization of
Theorem~\ref{p5:thm:v397-charged-preterminal-srcval-realization}.
No native impossibility or distinct-actual-actions counterexample is
inferred from that execution boundary.

\subsubsection{Finite source coefficients and spectral-weight comparisons}
\label{p6:app:v2189-electron-source-comparisons}
The unchanged terminal formulas are
\[
 \varepsilon_e=\frac{E_{40}}{1+\sqrt5}(24/25)^{-3/2},\qquad
 \varepsilon_\mu=\frac{E_{36}}2(24/25),\qquad
 \varepsilon_\tau=E_{36}(2\sqrt6+\sqrt{12})(24/25).
\]
Replaying only the already printed $E_{36},E_{40}$ ledger gives,
to its input precision,
\[
 (\varepsilon_e,\varepsilon_\mu,\varepsilon_\tau)
 \simeq(0.5112237095,105.3978706892,1762.9018808308)\ {\mathrm{MeV}}.
\]
These are fixed finite source coefficients, not newly calculated poles.
In particular no source theorem licenses setting the actual incidence
contractions in \eqref{p6:eq:v2189-mass-second-moment} to one merely
to reproduce these coefficients as pole values.

For the complete normalized stationary CAR source, the actual source
Liouvillian has a positive finite spectral measure
$G(z)=\sum_\lambda w_\lambda/(z-\lambda)$ with
\[
 w_\lambda\ge0,\qquad \sum_\lambda w_\lambda=1,\qquad
 \sum_\lambda w_\lambda\lambda=\mu_1,\qquad
 \sum_\lambda w_\lambda\lambda^2=\mu_2.
\]
Consequently $\mu_2-\mu_1^2\ge0$ and, for $t>0$,
\[
 \sum_{|\lambda-\mu_1|\ge t}w_\lambda
       \le\frac{\mu_2-\mu_1^2}{t^2}.
\]
These are weighted source-frequency centroid and tail identities, not
individual physical mass determinations. They follow from the stationary
positive pairing and the spectral theorem; the bound follows by summing
$(\lambda-\mu_1)^2\ge t^2$ on the indicated subset. Actual complete
$\mu_1,\mu_2$ still consume the full inherited coefficients and density.

\subsubsection{Constraint, state and mixed-return discipline}
\label{p6:app:v2189-electron-state-scope}
The complete-word Hessian uses all actual finite CAR slots and their even
normal word. If an actual physical source compression changes CAR,
compressed multiplication and all occupation/constraint boundary returns
must be evaluated in that chart. A zero classical Grassmann background
does not evaluate the quantum weighted occupation contractions.
Likewise the direct constant mass--weak cancellation does not cancel
quantum scalar/Yukawa or higher mixed coefficient terms, later pair
returns, or an inherited interacting density.

The free filled-negative projection in
\eqref{p4:eq:v2169-b11-l9} follows from an actual invertible purely
off-diagonal free $h_R$; both chiral compressions are $I/2$, independently
of terminal-return normalization. Its marked weak leakage is nevertheless
nonzero on the licensed free coordinate unit, as evaluated in
\eqref{p4:eq:v2169-b11-l12}. A complete stationary support must retain
actual cancellations or the interacting dressed state. Neither an empty
Fock reset nor an optional thermalization premise replaces that support.
The fresh neutral Fock unit is a branch-origin construction; by itself it
does not certify a complete interacting ground density.

% END INLINE SOURCE: Paper5B-ElectronExecutionLineage-Appendix-v1_0.tex




% BEGIN GEOMETRY-FIRST CONSEQUENCES v2.201
\setlength{\emergencystretch}{2em}
\subsection{Rooted negative-Hopf hierarchy blocks and the exact positive return fibre}
\label{p6:sec:v229-rooted-charged-converse}

\begin{theorem}[Actual rooted-isometry transport of the negative-Hopf relation]
\label{p6:thm:v229-rooted-negative-hopf-hierarchy-blocks}
Use the pre-existing rooted hierarchy/family isometry of
Theorem~\ref{p1b:thm:v138-hierarchy-family-isometry},
\[
 J_{\rm HG}=|u\rangle\langle c|+\Psi_uQ_\perp,
 \qquad J_{\rm HG}Q_0J_{\rm HG}^\dagger=Q_4,
 \quad J_{\rm HG}Q_\perp J_{\rm HG}^\dagger=Q_5,
\]
not the aligned-coordinate reflection $2Q_0-I$ as a substitute for this
map.  With $R_-=\operatorname{diag}(1,i,-i)$,
$\widehat R_-=J_{\rm HG}R_-J_{\rm HG}^\dagger$,
$\widetilde B_r=J_{\rm HG}B_rJ_{\rm HG}^\dagger$ and
$s=2\sqrt2/3$, the exact hierarchy-space compressions are
\[
 \boxed{Q_5\widehat R_-Q_4=s\widetilde B_\mu,\qquad
 Q_4\widehat R_-Q_5=s\widetilde B_\tau,\qquad
 Q_5\widehat R_-Q_5=J_{\rm HG}F_{-1/3}J_{\rm HG}^\dagger.}
\]
For the orthonormal pairs $(g_e,g_\mu)$ and $(h_e,h_\tau)$, let
$F_\gamma=g_eh_e^\dagger+\gamma g_\mu h_\tau^\dagger$.  Then
\[
 \boxed{P_eF_\gamma=F_\gamma\quad\Longleftrightarrow\quad\gamma=0.}
\]
Consequently an actual physical electron source \emph{whose left role
support is certified to be $P_e$} selects $B_e=g_eh_e^\dagger$ within
this family.  The geometric transport alone does not assert that the
actual charged action consumes these arrows.
\end{theorem}
\begin{proof}
The two transported projectors satisfy $Q_5J_{\rm HG}=J_{\rm HG}Q_\perp$
and $Q_4J_{\rm HG}=J_{\rm HG}Q_0$.  Therefore the existing relations
$Q_\perp R_-Q_0=sB_\mu$ and $Q_0R_-Q_\perp=sB_\tau$
transport without a sign change.  The plane compression transports
likewise.  Finally $P_eg_e=g_e$, $P_eg_\mu=0$, and the rank-one term
$g_\mu h_\tau^\dagger$ is nonzero, yielding the iff claim.
\end{proof}


% BEGIN FRAMEWORK v2.239: EXACT HOPF REDUCING PAIR
\begin{proposition}[Exact rooted negative-Hopf $1\oplus2$ reducing normal form]
\label{p6:prop:v239-rooted-hopf-one-plus-two}
Use precisely the rooted vectors and unitary $R_-$ of
Proposition~\ref{p6:prop:v2168-angular-hierarchy}. In the left frame
$U=(g_e,g_\mu,c)$ and the \emph{reordered} unit right frame
$V^\sharp=(h_e,h_\tau,c)$, one has
\begin{equation}\label{p6:eq:v239-hopf-normal-form}
 \boxed{U^\dagger R_-V^\sharp=
 \begin{pmatrix}
 1&0&0\\
 0&-1/3&2\sqrt2/3\\
 0&2\sqrt2/3&1/3
 \end{pmatrix}=1\oplus\mathcal R_2,\qquad \mathcal R_2^2=I_2.}
\end{equation}
In particular $R_-h_e=g_e$ and $R_-^\dagger g_e=h_e$ are exact
geometric reducing-singular-pair identities. The independent
$P_e$-supported rank-one arrow $B_e=g_eh_e^\dagger$ and the full
rank-two plane compression $F_{-1/3}$ are different operator claims:
the latter retains a $g_\mu h_\tau^\dagger$ component which does not
affect the electron reducing ray. The scalar $Q_0R_-Q_0=1/3$
and unnormalized cross amplitudes $2\sqrt2/3$ of the \emph{complete}
$R_-$ are also retained; the full $R_-$ is not thereby identified
with the three-corner charged source whose $Q_0$--$Q_0$ corner is zero.
\end{proposition}
\begin{proof}
Substitute the printed $g_e,g_\mu,c,h_e,h_\tau$ and
$R_-=\operatorname{diag}(1,i,-i)$ and multiply. The first row and
column follow from $R_-h_e=g_e$ and unitarity; the plane block is
exactly the displayed reflection and squares to $I_2$.
The right-frame order $V^\sharp=(h_e,h_\tau,c)$ is deliberately
distinguished from $(h_e,c,h_\tau)$ used in the unit-arrow ledger;
only a presentation permutation has been performed.
\end{proof}

% BEGIN FRAMEWORK v2.240: CANONICAL POLAR HOPF ARROWS
\begin{theorem}[Canonical normalized three-role arrows from a single rooted unitary]
\label{p6:thm:v240-canonical-polar-hopf-triple}
Let $W\in U(3)$ and let $c$ be a fixed unit singlet already present on
its carrier. Assume $s^2=1-|c^\dagger Wc|^2>0$. Define, using only
these generated inputs,
\[
 P_0=cc^\dagger,\quad Q=I-P_0,\quad
 P_\mu=s^{-2}QWP_0W^\dagger Q,\quad P_e=Q-P_\mu,
\]
\begin{equation}\label{p6:eq:v240-polar-hopf-arrows}
 \boxed{B_e^{\rm pol}=P_eWQ,\qquad
 B_\mu^{\rm pol}=s^{-1}QWP_0,\qquad
 B_\tau^{\rm pol}=s^{-1}P_0WQ.}
\end{equation}
Then $P_e,P_\mu,P_0$ are orthogonal rank-one projections summing to
$I$, and
\begin{equation}\label{p6:eq:v240-hopf-full-left-return}
 \boxed{B_r^{\rm pol}(B_s^{\rm pol})^\dagger
        =\delta_{rs}P_r,\qquad
 T_{\rm pol}T_{\rm pol}^\dagger=e^2P_e+m^2P_\mu+t^2P_0,}
\end{equation}
where $T_{\rm pol}=eB_e^{\rm pol}+mB_\mu^{\rm pol}
+tB_\tau^{\rm pol}$ and $(P_e,P_\mu,P_\tau)=(P_e,P_\mu,P_0)$.
No right-frame representative, mass fit or new selector is needed.
For the existing rooted negative-Hopf parent
$W=R_-=\operatorname{diag}(1,i,-i)$ and
$c=(1,1,1)^T/\sqrt3$, one has
\[
 \boxed{s=2\sqrt2/3,\quad B_e^{\rm pol}=B_e,\quad
 B_\mu^{\rm pol}=B_\mu,\quad B_\tau^{\rm pol}=B_\tau.}
\]
Thus the free tagged left resolvents on this \emph{explicitly defined}
operator family are $z/(z^2-e^2)$, $z/(z^2-m^2)$ and
$z/(z^2-t^2)$ in the respective role lines. The result does not
assert that the physical source Hessian has already consumed
$T_{\rm pol}$.
\end{theorem}
\begin{proof}
The vector $QWc$ has norm $s$, so $P_\mu$ is a rank-one
projection and $P_e$ its rank-one complement in $Q$. In particular
$P_eWP_0=0$. Unitarity gives
\[
 P_eWQW^\dagger P_e=P_e(I-WP_0W^\dagger)P_e=P_e,
\]
$QWP_0W^\dagger Q=s^2P_\mu$ and
$P_0WQW^\dagger P_0=s^2P_0$. Thus each arrow has the asserted
left return. The mixed returns vanish because $QP_0=0$ and
$P_eWP_0=0$, including
$P_eWQW^\dagger P_0=P_e(I-WP_0W^\dagger)P_0=0$.
Expansion proves the weighted sum. For $R_-$ use the existing rooted
vectors and $c^\dagger R_-c=1/3$; substituting gives the stated
$B_e,B_\mu,B_\tau$.
\end{proof}

% BEGIN FRAMEWORK INTEGRATION v3.024: SHARP NEAREST ROOTED HOPF SOURCE
\begin{theorem}[Sharp nearest zero-defect representatives of the canonical rooted Hopf triple]
\label{p6:thm:v3024-rooted-hopf-nearest-source}
Use the fixed source/target rooted chart of
Theorem~\ref{p6:thm:v240-canonical-polar-hopf-triple} and
the hierarchy-corner defect theorem of Paper~V-A.
Let \(W\in U(3)\), \(P_0=cc^\dagger\),
\(Q=I-P_0\), \(a=c^\dagger Wc\) and
\(s=(1-|a|^2)^{1/2}\), with \(0<|a|<1\).
Let \(B_e^{\rm pol},B_\mu^{\rm pol},B_\tau^{\rm pol}\) be the
already defined polar arrows.  Set
\[
 T=B_e^{\rm pol}+sB_\mu^{\rm pol}+sB_\tau^{\rm pol},
 \qquad U=B_e^{\rm pol}+B_\mu^{\rm pol}+B_\tau^{\rm pol}.
\]
For the dimensionless rooted hierarchy defect on this chart,
\[
 \boxed{\min_{\substack{R\ne0\\\Delta_{\rm hier}(R)=0}}
 \|W-R\|_{\rm HS}^2=2|a|^2,\quad
 \operatorname*{argmin}_{\substack{R\ne0\\\Delta_{\rm hier}(R)=0}}
 \|W-R\|_{\rm HS}^2=\{T\},}
\]
\[
 \boxed{\min_{\substack{V\in U(3)\\\Delta_{\rm hier}(V)=0}}
 \|W-V\|_{\rm HS}^2=4(1-s),\quad
 \operatorname*{argmin}_{\substack{V\in U(3)\\\Delta_{\rm hier}(V)=0}}
 \|W-V\|_{\rm HS}^2=\{U\}.}
\]
Both minimizers have nonzero three-arrow blocks.  Uniqueness is for the
\emph{fixed rooted chart} and is not a physical action selection rule.
For the existing negative-Hopf parent
\(W=R_-=\operatorname{diag}(1,i,-i)\),
\(c=(1,1,1)^\mathsf T/\sqrt3\), one has
\[
 a=\frac13,\qquad s=\frac{2\sqrt2}{3},\qquad
 \|W-T\|_{\rm HS}^2=\frac29,\qquad
 \|W-U\|_{\rm HS}^2=4-\frac{8\sqrt2}{3}.
\]
\end{theorem}
\begin{proof}
Write the four orthogonal root corners
\(D=P_0WP_0=aP_0\), \(A=QWP_0\), \(B=P_0WQ\),
\(C=QWQ\).  By unitarity and the definition of \(P_\mu\),
\[
 CC^\dagger=Q-QWP_0W^\dagger Q
  =P_e+|a|^2P_\mu.
\]
Thus the two singular values of \(C\) on the augmentation are
\(1,|a|\).  Their strict gap makes
\(B_e^{\rm pol}=P_eWQ\) the unique best rank-one
Hilbert--Schmidt approximation of \(C\), with squared residual
\(|a|^2\).
If \(\Delta_{\rm hier}(R)=0\), the source defect theorem forces
\(P_0RP_0=0\) and \(\operatorname{rank}(QRQ)\le1\).
Hence the squared distance is at least \(|a|^2\) from the singlet
corner and another \(|a|^2\) from the augmentation corner.
Equality fixes both cross blocks to those of \(W\) and the
augmentation corner to \(B_e^{\rm pol}\), uniquely producing
\(T\).  By the polar-arrow source and target orthogonality,
\(\Delta_{\rm hier}(T)=0\).

If \(R=V\) is also unitary, its vanishing singlet corner forces
the two crossing blocks to have Hilbert--Schmidt norm one, whereas
those of \(W\) have norm \(s\).  Each crossing therefore costs
at least \((1-s)^2\), and equality aligns it with the corresponding
normalized polar arrow.  The singlet and rank-one augmentation costs
remain \(|a|^2\) each.  Thus
\[
 \|W-V\|_{\rm HS}^2\ge2|a|^2+2(1-s)^2=4(1-s).
\]
The unweighted sum \(U\) is unitary by the full left-return identity
of Theorem~\ref{p6:thm:v240-canonical-polar-hopf-triple}; it has
zero defect and attains equality.  Strict singular-value separation
and equality in the two crossing Cauchy--Schwarz inequalities establish
uniqueness.
\end{proof}

\begin{corollary}[Commuting zero-defect negative-Hopf control]
\label{p6:cor:v3024-commuting-hopf-control}
In the same fixed rooted chart for
\(W=\operatorname{diag}(1,i,-i)\) and
\(c=(1,1,1)^\mathsf T/\sqrt3\), one has
\[
 \boxed{\min_{\substack{V\in U(3),\ [V,W]=0\\
                   \Delta_{\rm hier}(V)=0}}
     \|W-V\|_{\rm HS}^2=4-2\sqrt3
     >4-\frac{8\sqrt2}{3}.}
\]
Thus the existing rooted polar representative improves the
Hilbert--Schmidt comparison over every commuting zero-defect unitary;
commuting zero-defect representatives nevertheless exist.
Neither inequality requires noncommuting physical source dynamics.
\end{corollary}
\begin{proof}
The three eigenvalues of \(W\) are distinct, so a commuting unitary
\(V\) is diagonal with phases \(z_1,z_2,z_3\).
The vanishing singlet corner requires
\(z_1+z_2+z_3=0\); thus the phases are, up to common phase and
permutation, the three cubic roots
\(1,\omega,\omega^2\), \(\omega=e^{2\pi i/3}\).
For such a unitary, zero defect follows from the rank-one singlet
and unitary block relations.  Maximizing
\(\Re\Tr(W^\dagger V)\) over the six permutations and one common
phase gives \(1+\sqrt3\).  The identity
\(\|W-V\|_{\rm HS}^2=6-2\Re\Tr(W^\dagger V)\)
gives the first minimum.  The strict inequality follows from
\(8\sqrt2/3>2\sqrt3\).
\end{proof}
% END FRAMEWORK INTEGRATION v3.024


\begin{proposition}[Unreduced negative-Hopf plane and nominal heavy-role obstruction]
\label{p6:prop:v240-unreduced-hopf-spectrum}
In the native frames of
Proposition~\ref{p6:prop:v239-rooted-hopf-one-plus-two}, replace only
the electron arrow by the \emph{unreduced} plane compression
$F_{-1/3}=QR_-Q$. The resulting family word is
\[
 T_{\rm unred}=\begin{pmatrix}
 e&0&0\\0&-e/3&m\\0&t&0
 \end{pmatrix}.
\]
Its electron squared singular value is exactly $e^2$, whereas the
remaining values $\lambda_\pm$ satisfy
\[
 \boxed{\lambda_++\lambda_-=m^2+t^2+e^2/9,\qquad
        \lambda_+\lambda_-=m^2t^2.}
\]
For $e>0$ both cannot simultaneously equal $m^2,t^2$.
The geometric role projection $P_eR_-Q$ removes the additional
$g_\mu h_\tau^\dagger$ return without changing the electron ray.
\end{proposition}
\begin{proof}
Use the $1\oplus2$ normal form and the normalized cross-corners.
The lower $2\times2$ matrix is
$\left(\begin{smallmatrix}-e/3&m\\t&0\end{smallmatrix}\right)$;
its squared singular-value trace and determinant are the displayed
quantities. The electron row and column reduce separately.
\end{proof}
\begin{remark}[Source-vs-relation origin certificate]
\label{p6:rem:v240-hopf-physical-origin-firewall}
The geometric arrows in
\eqref{p6:eq:v240-polar-hopf-arrows} are fixed equivariantly by the
same $R_-,P_0$; they do not choose a point of the existing positive
character-return fibre. For a \emph{physical} charged claim one must
still establish that the actual mixed source packet
$\iota_r=\Phi_{\rm HG}(A_r^{\rm hier})$ has the required return,
not merely that the normalized relation arrows exist.
Paper~III-F, Proposition~\ref{p3f:prop:v240-action-source-gram-test},
states the minimal action test.
\end{remark}
% END FRAMEWORK v2.240: CANONICAL POLAR HOPF ARROWS

% BEGIN FRAMEWORK v2.242: GEOMETRIC PHASE AND COMPARISON SCOPE
\begin{remark}[Signed phase realization and finite nonlinear comparison]
\label{p6:rem:v242-geometric-provider-scope}
The rooted \(R_\pm\) used in this section have the signed real-vector
realization of Corollary~\ref{p1:cor:v242-joint-phase-vector-lift}.
This is compatible with the canonical polar Hopf arrows above and does
not assign a new value to their physical source coefficient.
The complete-channel value
\(\rho^\dagger\rho=\operatorname{diag}(1,5,5)\) is a Gram on its
stated support, not a replacement for the terminal rank-five owner.
Likewise Corollary~\ref{p1b:cor:v242-family-cubic-comparison}
calculates an explicit coefficient-mode cubic whose raw root-register
exclusion is \(-C_H/25\). It cannot be relabelled a physical Yukawa
coupling or used to infer a mass multiplier without the actual field
embedding and action comparison. The printed charged coefficients and
the source-vs-relation origin certificate remain unchanged.
\end{remark}
% END FRAMEWORK v2.242: GEOMETRIC PHASE AND COMPARISON SCOPE


% END FRAMEWORK v2.239: EXACT HOPF REDUCING PAIR

\begin{theorem}[Necessity and sufficiency of the positive opposite-odd return fibre]
\label{p6:thm:v229-opposite-odd-positive-converse}
Fix the rank-one arrows $(B_e,B_\mu,B_\tau)$ and the commuting positive
right dressings
\[
 C_e=Q_0+a_+E_++a_-E_-,\qquad
 C_\tau=Q_0+d_+E_++d_-E_-,\qquad a_\pm,d_\pm>0.
\]
Set $\iota_e=B_eC_e$, $\iota_\mu=B_\mu$ and
$\iota_\tau=B_\tau C_\tau$.  The complete orthonormal three-role return
\[
 \iota_r\iota_s^\dagger=\delta_{rs}P_r\quad\hbox{for all }r,s
\]
holds \emph{if and only if} there exist real $A,B$ satisfying
\[
 \boxed{C_e=Q_0+AQ_\perp+BC_2,\qquad
 C_\tau=Q_0+AQ_\perp-BC_2,}
\]
\[
 \boxed{A^2-3AB+3B^2=1,\qquad A>\sqrt3|B|.}
\]
The entire positive arc has the same mass-only left return
$TT^\dagger=\sum_r\varepsilon_r^2P_r$ for the corresponding weighted
role rows, without selecting a point $(A,B)$.  Membership of the
\emph{actual} preterminal source word in this fibre is a separate
certificate.
\end{theorem}
\begin{proof}
Write $w_+=(2-\sqrt3)/4$ and $w_-=(2+\sqrt3)/4$.
The two diagonal norms and the vanishing electron--tau cross return are
respectively equivalent to
\[
 w_+a_+^2+w_-a_-^2=1,\qquad
 w_-d_+^2+w_+d_-^2=1,\qquad a_+d_+=a_-d_-=:q>0.
\]
The last equation gives $d_+=q/a_+$, $d_-=q/a_-$.  Substitution in the
second norm, using the first, gives $q^2/(a_+^2a_-^2)=1$.
Strict positivity fixes $q=a_+a_-$ and hence $d_+=a_-$,
$d_-=a_+$.  With $A=(a_++a_-)/2$ and
$B=(a_+-a_-)/(2\sqrt3)$, the first norm is exactly
$A^2-3AB+3B^2=1$ and positivity is $A>\sqrt3|B|$.
Conversely these relations imply all three norm/orthogonality tests;
this includes the sufficient branch of
Proposition~\ref{p1b:prop:v228-opposite-odd-fibre}.
\end{proof}

% BEGIN FRAMEWORK v2.267: POSITIVE FIBRE AS A RIGHT-UNITARY ORBIT
\begin{corollary}[Complete explicit right-unitary positive-fibre orbit]
\label{p6:cor:v267-positive-right-unitary-orbit}
For the \emph{geometric} positive opposite-odd fibre of
Theorem~\ref{p6:thm:v229-opposite-odd-positive-converse}, put
\[
 A=\cos\theta+\sqrt3\sin\theta,\qquad
 B=\frac{2}{\sqrt3}\sin\theta,\qquad
 -\frac\pi{12}<\theta<\frac{5\pi}{12}.
\]
Let $V_0=(h_e,h_\tau,c)$ and $V_\theta=(C_eh_e,C_\tau h_\tau,c)$.
In this ordered right frame, with $u=e^{-i\pi/4}$, one has
\[
 \boxed{V_\theta=V_0
 \begin{pmatrix}
 \cos\theta&u\sin\theta&0\\
 -\bar u\sin\theta&\cos\theta&0\\
 0&0&1
 \end{pmatrix}.}
\]
Writing $V_\theta^\sharp=(C_eh_e,c,C_\tau h_\tau)$,
$U_L=(g_e,g_\mu,c)$, and
$T_\theta=U_L\operatorname{diag}(e,m,t)(V_\theta^\sharp)^\dagger$,
\[
 \boxed{T_\theta=T_0U_R(\theta),\qquad
 U_R(\theta)=V_0^\sharp(V_\theta^\sharp)^\dagger
 =e^{\theta J_R},\quad J_R^\dagger=-J_R.}
\]
Thus $T_\theta T_\theta^\dagger=T_0T_0^\dagger$ on the whole
positive arc.  For the finite chiral Hamiltonians
$h_\theta=\bigl(\begin{smallmatrix}0&T_\theta\\
 T_\theta^\dagger&0\end{smallmatrix}\bigr)$ and every finite
spectral function $f$,
\[
 \boxed{P_Lf(h_\theta)P_L=P_Lf(h_0)P_L.}
\]
Accordingly no positive-fibre member selector is needed for
\emph{free left-only observables} of this geometric packet.
\end{corollary}
\begin{proof}
The displayed trigonometric substitution solves the quadratic fibre
equation identically; $A>\sqrt3|B|$ is precisely the open interval
shown. In the existing rooted right basis, direct use of the defined
$C_2=-i(X-X^\dagger)$ and the dressings gives the printed two-plane
matrix, which is unitary and has determinant one.  Reorder its columns
to obtain $T_\theta=T_0U_R(\theta)$.  The right unitary is a one-parameter
subgroup because the two-plane rotation is exponential.
With $W_\theta=\operatorname{diag}(I_L,U_R(\theta)^\dagger)$,
$h_\theta=W_\theta h_0W_\theta^\dagger$ and
$P_LW_\theta=P_L$.  Covariance of finite functional calculus proves
the left compression identity.
\end{proof}
\begin{remark}[Right sensitivity is a distinct possible selector]
\label{p6:rem:v267-positive-orbit-right-mark}
If a \emph{fixed actual} right inverse kernel is split on the two
rotated right slots, $K_R^{-1}=\operatorname{diag}(u_0,w_0,c_0)$
with $u_0\ne w_0$ in that rooted basis, the electron Schur return is
$e^2(u_0\cos^2\theta+w_0\sin^2\theta)$.  A same-action equality
to $e^2u_0$ forces $\theta=0$ on the positive arc.  Alternatively,
commutation $[K_R,J_R]=0$ makes every member invisible to the whole
left Schur response on an invertible common chart.  Neither branch
constructs the missing \emph{physical} hierarchy arrow: its origin
certificate remains Definition~\ref{p5:def:v228-hierarchy-origin}.
\end{remark}
% END FRAMEWORK v2.267: POSITIVE FIBRE AS A RIGHT-UNITARY ORBIT

\begin{remark}[Canonical comparison-loop test and actual owner assignments]
\label{p6:rem:v229-canonical-loop-test}
For a loop that multiplies the three quadratic source/kinetic blocks by
$(\eta_H^{k_L},\eta_H^{k_R},\eta_H^{k_D})$, canonical normalization
changes the response by
\[
 \boxed{D_{\rm can}\longmapsto
 \eta_H^{\,k_D-(k_L+k_R)/2}D_{\rm can}.}
\]
It is invisible to this normalized response precisely when
$k_D=(k_L+k_R)/2$.  A source-only extra loop is not removable by
\textup{GP3} while keeping its changed physical response.  The nine
current exterior/Hodge multipliers remain exact geometric candidates;
the actual chirality, exterior-grade and loop-word assignments require
the preterminal mixed action derivative of
Definition~\ref{p2a:def:v229-action-incidence-certificate}.
\end{remark}
% END FRAMEWORK v2.229: ROOTED CHARGED CONVERSES AND COMPARISON SCOPE


\section{Flavor-coordinate dictionaries and current invariants}
\label{p5b:sec:organized-flavor}
\subsection{Complete generic flavor readouts and source-selection equations}
\label{p6:sec:v199-flavor-readouts}
Use the canonically normalized left Grams on the common generation carrier,
as in Theorem~\ref{p5:thm:v199-weak-current-moments}.  This section concerns
positive simple spectra and the generic three-family Dirac quotient.  Any
normalization constraint on a larger multi-sector tuple must still be
imposed; the readout coordinates do not add independent physical parameters.

\begin{theorem}[Reconstruction and admissibility of the generic Dirac orbit]
\label{p6:thm:v199-flavor-reconstruction}
The six spectral moments \(a_r=\operatorname{tr}A^r\),
\(b_r=\operatorname{tr}B^r\), \(r=1,2,3\), and four mixed moments
\(M_{11},M_{12},M_{21},M_{22}\), where
\(M_{pq}=\operatorname{tr}(A^pB^q)\), recover the pair up to its generic
CP-conjugate ambiguity.  More explicitly, order the two spectra
\(a_i^{\rm eig},b_j^{\rm eig}\), put
\((S_A)_{pi}=(a_i^{\rm eig})^p\),
\((S_B)_{qj}=(b_j^{\rm eig})^q\) for \(p,q=0,1,2\), and fill the
remaining entries of \(M\) by \(M_{00}=3\), \(M_{p0}=a_p\),
\(M_{0q}=b_q\).  Then
\begin{equation}
 W=(|V_{ij}|^2)=S_A^{-1}MS_B^{-\mathsf T}.
 \label{p6:eq:v199-vandermonde-reconstruction}
\end{equation}
Let \(r_j=W_{1j}W_{2j}\).  The squared Dirac quartet is
\begin{equation}
 J^2=\tfrac14\left(2\sum_{i<j}r_ir_j-\sum_i r_i^2\right)
     =r_1r_2-\tfrac14(r_3-r_1-r_2)^2.
 \label{p6:eq:v199-cp-triangle}
\end{equation}
A nonnegative \(3\times3\) matrix with unit row and column sums is
unistochastic precisely when this triangle expression is nonnegative.
On the strictly positive, positive-area locus, the two orientations are
distinguished by the sign of
\(\chi=\operatorname{tr}[A,B]^3/(6i)\), whose magnitude obeys
\begin{equation}
 \chi^2=\Delta_A^2\Delta_B^2J^2,
 \qquad \Delta_A=(a_1^{\rm eig}-a_2^{\rm eig})
 (a_2^{\rm eig}-a_3^{\rm eig})(a_3^{\rm eig}-a_1^{\rm eig}),
 \label{p6:eq:v199-cp-syzygy}
\end{equation}
and similarly for \(\Delta_B\).
\end{theorem}
\begin{proof}
Apply Paper~1B, Theorem~\ref{p1b:thm:rank-three}, to the positive
simple-spectrum flavor pair already constructed on the common generation
carrier. Its Newton and mixed-moment reconstruction gives
\eqref{p6:eq:v199-vandermonde-reconstruction}; its triangle criterion
and discriminant identity give \eqref{p6:eq:v199-cp-triangle} and
\eqref{p6:eq:v199-cp-syzygy}. On the stated generic positive-area locus
the two orientations are distinguished once the local signed quartet
convention is fixed, as in Paper~IV,
Proposition~\ref{p4:prop:rank-three-jarlskog-identity}, and
\cite{Jarlskog1985}. The framework's construction supplies the source
pair and its flavor assignment.
\end{proof}

The CP magnitude is thus constrained, rather than an eleventh independently
selectable real coordinate.  Degenerate, zero-entry and CP-conserving strata
need their own quotient charts.  The same construction applies to the
charged-lepton/Dirac-neutrino pair; it is not a Majorana classification.

\begin{proposition}[Positive mixed resolvent and determinant-source probe]
\label{p6:prop:v199-mixed-resolvent}
For \(z,w>0\), put \(R_A=(zI+A)^{-1}\), \(R_B=(wI+B)^{-1}\).
Then
\begin{equation}
 \mathcal G(z,w)=\operatorname{tr}(R_AR_B)
 =\sum_{i,j}\frac{W_{ij}}{(z+a_i^{\rm eig})(w+b_j^{\rm eig})}
 \label{p6:eq:v199-mixed-stieltjes}
\end{equation}
is the double Stieltjes transform of the positive measure
\(\sum W_{ij}\delta_{a_i^{\rm eig}}\otimes\delta_{b_j^{\rm eig}}\),
whose two marginals put one unit at each eigenvalue.  In particular
\[
 (-1)^{m+n}\partial_z^m\partial_w^n\mathcal G
 =m!n!\sum_{i,j}\frac{W_{ij}}
 {(z+a_i^{\rm eig})^{m+1}(w+b_j^{\rm eig})^{n+1}}>0.
\]
The weak-current identification of the two generation copies permits the
block probe
\[
 \mathcal K=\begin{pmatrix}zI+A&\eta I\\\bar\eta I&wI+B\end{pmatrix}.
\]
For \(t=|\eta|^2\) with
\(t\lambda_{\max}(R_A^{1/2}R_BR_A^{1/2})<1\), it is positive and
\begin{equation}
 -\log\frac{\det\mathcal K}{\det(zI+A)\det(wI+B)}
 =\sum_{k\ge1}\frac{t^k}{k}\operatorname{tr}[(R_AR_B)^k].
 \label{p6:eq:v199-determinant-probe}
\end{equation}
\end{proposition}
\begin{proof}
The mixed resolvent is Paper~1B, Theorem~\ref{p1b:thm:rank-three},
for the already constructed flavor pair. Differentiating its finite
positive sum gives the stated complete monotonicity; unit row and
column sums give the two marginals. The weak-current identification
supplies the identity off-diagonal blocks of the present probe.
Schur factorization then yields
$\det\mathcal K/(\det(zI+A)\det(wI+B))=\det(I-tR_BR_A)$.
The product is similar to a positive Hermitian matrix, so the logarithmic
series converges on the displayed positive domain and gives
\eqref{p6:eq:v199-determinant-probe}.
\end{proof}
The first source derivative is \(\mathcal G\), and the large-\(z,w\)
expansion recovers every mixed moment.  This is a probe determinant.  The
physical fermion determinant still requires the actual Dirac source,
Clifford factors, seagull terms, regulator and action/state weight.


The full logarithm of this simple probe has only commuting angular
stationary pairs, as proved in
Proposition~\ref{p6:prop:v203-full-probe-stationarity}. It therefore
cannot alone select generic mixing. The actual source-owned determinant
must be evaluated with its licensed Dirac/current, Clifford, seagull,
regulator and normalization data.

\begin{proposition}[Derivative compiler for an admitted determinant term]
\label{p6:prop:v203-native-determinant-derivatives}
On an invertible smooth finite matrix branch with a local logarithm,
an admitted term $F=-\log\det D(c)$ obeys
\[
 \partial_aF=-\operatorname{Tr}(D^{-1}D_a),\qquad
 \partial_a\partial_bF=
 \operatorname{Tr}(D^{-1}D_bD^{-1}D_a)-\operatorname{Tr}(D^{-1}D_{ab}).
\]
Here $c$ are actual licensed source coordinates and subscripts denote
their derivatives; the action fixes signs, multiplicities and real/phase
prescriptions of all such terms.
\end{proposition}
\begin{proof}
Apply Paper~1B,
Proposition~\ref{p1b:result:determinant-derivatives-and-relative-frame-stationarity},
on the displayed smooth invertible branch with its local logarithm.
The identities apply to the actual source coordinates before taking
the action's fixed real or phase prescription, retaining its signs and
multiplicities.
\end{proof}
The angular Hessian calculation below is retained as a historical/diagnostic selection test.  Framework v2.86 closes the finite rooted transport selector upstream by the Paper-V-A $C_3$ source-cell theorem; the Hessian problem below concerns only relative flavor-frame/member dynamics.
Its positive definiteness would remove local angular flat directions on
that chart. Global minima, CP-state occurrence and pole/RG transport remain
distinct. The derivatives do not make Yukawa parameters into new dynamical
scalars: their variational role must already belong to the source action.

\begin{proposition}[Invariant coordinates and conditional selection Hessian]
\label{p6:prop:v199-invariant-selection-equations}
At fixed simple spectra, the four mixed moments are local angular coordinates
on each generic CP sheet.  If the existing primitive-trimer source variables
are admitted as variational variables of the same action and its effective
function is \(S(c)=\Phi(I(c))\), with
\[
 I=(a_1,a_2,a_3,b_1,b_2,b_3,M_{11},M_{12},M_{21},M_{22}),
\]
the invertible Jacobian \(J_{\rm tri}=\partial I/\partial c\) of
Theorem~\ref{p6:thm:v548-trimer-local-submersion} gives
\begin{equation}
 \nabla_c S=J_{\rm tri}^{\mathsf T}\nabla_I\Phi,
 \qquad \nabla_cS=0\ \Longleftrightarrow\ \nabla_I\Phi=0.
 \label{p6:eq:v199-invariant-stationarity}
\end{equation}
At such a critical point,
\begin{equation}
 \operatorname{Hess}_c S
 =J_{\rm tri}^{\mathsf T}(\operatorname{Hess}_I\Phi)J_{\rm tri}.
 \label{p6:eq:v199-invariant-stability}
\end{equation}
\end{proposition}
\begin{proof}
Use the upper-left \(2\times2\) submatrix of \(W\) as affine coordinates
on the doubly stochastic plane.  Then
\(\delta M=X_A\delta W X_B^{\mathsf T}\), where
\((X_A)_{pi}=(a_i^{\rm eig})^p-(a_3^{\rm eig})^p\), \(p,i=1,2\),
and likewise for \(B\).  Its Jacobian has absolute determinant
\(\Delta_A^2\Delta_B^2>0\).
Writing \(a=s_{12}^2,b=s_{23}^2,c_0=s_{13}^2\), the Jacobian of
\((W_{13},W_{12},W_{23},W_{21})\) with respect to
\((c_0,a,b,\delta)\) is \(-2(1-c_0)J\): use
\(W_{13}=c_0\), \(W_{12}=a(1-c_0)\), \(W_{23}=b(1-c_0)\), and
\(W_{21}=a(1-b)+(1-a)bc_0+
2\sqrt{a(1-a)b(1-b)c_0}\cos\delta\).
It is nonzero on the stated locus.  The remaining assertions are the chain
rule and its second derivative, with the gradient term zero at criticality.
\end{proof}

For an implemented invariant action its matrix gradients are explicitly
\begin{align}
 P_A&=\sum_{r=1}^3r\Phi_{a_r}A^{r-1}
 +\sum_{p,q=1}^2\Phi_{pq}\sum_{k=0}^{p-1}A^kB^qA^{p-1-k},\nonumber\\
 P_B&=\sum_{r=1}^3r\Phi_{b_r}B^{r-1}
 +\sum_{p,q=1}^2\Phi_{pq}\sum_{k=0}^{q-1}B^kA^pB^{q-1-k}.
 \label{p6:eq:v199-matrix-gradients}
\end{align}
Contract these with the actual trimer tangent matrices to obtain its ten
equations.  For linear Gram coordinates with basis \(T_\alpha\), they are
\(\operatorname{tr}(T_\alpha P_A)=
\operatorname{tr}(T_\alpha P_B)=0\).
Relative-frame stationarity is \([B,P_B]=0\), and common-basis invariance
gives \([A,P_A]+[B,P_B]=0\).

The remaining selection target is an action-owned \(\Phi\) with a
positive-definite Hessian at its physical source member.  Merely possessing
the invariant operators fixes no coefficients.  Appendix~\ref{p6:app:v199-flavor-selection}
proves that the four affine mixed invariants, even with one additional
\(\operatorname{tr}(ABAB)\) term, cannot isolate a nondegenerate generic
angular minimum, and that CP-invariant data cannot distinguish conjugate
orientations.  The finite rooted transport selector is closed upstream.  The downstream physical pole/RG map, CP-state occurrence and relative flavor-frame/member dynamics retain their separate scopes.  No Yukawa parameter is promoted to a new dynamical field here.

% BEGIN M8 PREPARED FLAVOR OPERATOR HANDOFF
\paragraph{Determined low-mode contractions on the spatial anchor.}
\label{p6:par:m8td-flavor-handoff}
On the round spatial doublet realization, Paper~III-F,
Theorem~\ref{p3f:thm:m8td-low-mode}, now fixes the zero projector and
the complementary inverse. An actual spin-scalar internal field contracts
to $(2V_R)^{-1}\int\operatorname{tr}_E M\,d\mu$, while the actual full
Clifford/odd-field member is evaluated by its own $PVP$ and the exact
Feshbach expression \eqref{p3f:eq:m8td-feshbach}. This supplies specified
operator contractions once that member and representation are supplied;
it does not identify an arbitrary probe determinant with the physical
flavor action. The normalized source tuple, full spectator/chiral
multiplicities, state and pole/RG map in the preceding selection problem
remain its inputs. In particular, the spatial doublet zero alone neither
counts physical generations nor predicts a physical massless fermion.
% END M8 PREPARED FLAVOR OPERATOR HANDOFF

\subsection{Exact finite dynamics and the irreversible boundary}
\label{p6:sec:no-width}

\begin{theorem}[Finite-Hermitian no-width theorem]
\label{p6:thm:finite-hermitian-no-width}
Let \(H=H^*\) act on a finite-dimensional Hilbert space.  Then
\[
H=\sum_nE_nP_n,\qquad E_n\in\RR,
\qquad
(z-H)^{-1}=\sum_n\frac{P_n}{z-E_n}.
\]
Every exact resolvent pole is real.  For a normalized state \(\psi\),
\[
A_\psi(t)
=
\langle\psi,e^{-iHt}\psi\rangle
=
\sum_np_ne^{-iE_nt}
\]
is a finite almost-periodic sum and cannot equal
\[
e^{-iE_rt-\Gamma t/2}
\]
for every \(t\ge0\) with \(\Gamma>0\).  Therefore
\[
\boxed{\Gamma_{\rm exact\ finite}=0.}
\]
\end{theorem}

\begin{proof}
The finite spectral theorem gives the resolvent and survival amplitude.
Finite sums of phases have recurrence subsequences returning arbitrarily
close to their initial value, whereas a positive-width exponential tends to
zero.
\end{proof}

\begin{corollary}[Required rate data]
\label{p6:cor:required-rate-data}
A nonzero width requires either a continuum boundary value
\[
\Sigma(E+i0)
=
\Delta(E)-\frac i2\Gamma(E),
\qquad
\Gamma(E)=2\pi V^\dagger\delta(E-H_{\rm ch})V,
\]
or a controlled open-system limit with jump operators \(L_k\) and
\[
\Gamma
=
\sum_k\langle r|L_k^\dagger L_k|r\rangle .
\]
Channel projectors, interaction normalization, state normalization, and the
limit domain are indispensable.  Rescaling \(V\mapsto gV\) preserves the
finite carrier while changing the weak-coupling rate by \(g^2\).
\end{corollary}

\begin{proposition}[Meaning of \(1/(192\pi)\)]
\label{p6:prop:one-real-scalar-interface}
Assume the canonical Einstein-frame Starobinsky scalaron normalization
\(\phi T^\mu{}_\mu/(\sqrt6\,\overline M_{\rm Pl})\), and a Jordan-frame real
scalar with action
\[
-\frac12\int\sqrt{-g}\,
\left((\nabla\chi)^2+\xi R\chi^2\right)d^4x .
\]
For interpretation as an \(S\)-matrix partial width, also assume the applicable
\textnormal{(SM1)--(SM3)} channel construction, \textnormal{(SM5)} parent-pole
continuation, and phase-space/channel-normalization parts of
\textnormal{(SM7)} in
Definition~\ref{p12:def:claim-indexed-smat-package}, with
\textnormal{(SM4)} when the final-state normalization is obtained by LSZ and a
route from Definition~\ref{p12:def:claim-indexed-infrared-package} whenever the
massless limit has a soft, collinear, or long-range obstruction.
In the downstream continuum and massless-\(\chi\) limit, its scalaron partial
width is
\[
\boxed{
\Gamma_{\phi\to\chi\chi}
=
(1-6\xi)^2
\frac{M_\phi^3}{192\pi\overline M_{\rm Pl}^{\,2}}.
}
\]
For \(N_s\) equivalent real components this is multiplied by \(N_s\).
Thus \(1/(192\pi)\) is the minimally coupled one-real-scalar coefficient,
not a universal Standard-Model total width.  Four minimally coupled Higgs
components give four times the result, while the massless conformal choice
\(\xi=1/6\) removes the tree scalar channel.
\end{proposition}

\begin{proof}
On the declared asymptotic domain, the on-shell trace-coupling amplitude is
\[
|\mathcal M|
=
|1-6\xi|\frac{M_\phi^2}{\sqrt6\,\overline M_{\rm Pl}}.
\]
The identical-particle formula
\[
\Gamma
=
\frac1{2M_\phi}\frac1{2!}
\int d\Phi_2\,|\mathcal M|^2,
\qquad
\int d\Phi_2=\frac1{8\pi},
\]
gives the displayed partial width.  Independent real components add
incoherently.
\end{proof}

\begin{remark}[Paper--V-B boundary]
Proposition~\ref{p6:prop:one-real-scalar-interface} is recorded to identify
the downstream coefficient and prevent a false finite derivation.  It is not
used in any Paper--V-B mass theorem.  Scalaron matching belongs to Paper~VI-A;
interactions and continuum limits belong to their QFT owner; the reheating
history is a later cosmological evaluation.
\end{remark}

\subsection{native role orbit, ordering closure, and sharp neutrino no-go boundaries}
\label{p6:sec:v568-d4-neutrino-closure-refinement}

\noindent\textbf{Framework v2.208 branch scope.}
The neutral inverse/direct/closed-Gram entries in this retained section
belong to the historical $D_4$ response constructor. They are obsolete
on the harmonic rank branch of \ref{p6:sec:v2208-neutral-rank-word}.
Their algebraic theorems neither force nor constrain that new mass word.
The charged results and branch-independent geometric/source rules keep
their stated domains.


Release v5.67 introduced the preferred conditional shape.  The present
release removes three artificial ambiguities: the numerical values of the
native inverse/direct/Gram role orbit, its ordering for \(x>1\), and any
suggestion that the dimensionless \(D_4\) response by itself could fix either
the absolute mass scale or the PMNS frame.

\begin{proposition}[One-pass native response orbit]
\label{p6:prop:v568-native-response-orbit}
Let
\[
 R=xI_G,\qquad x>0,
\]
be the positive central response on the generated rank-three family carrier.
The already licensed one-pass operations direct response, Moore--Penrose
inverse, and closed Gram return have positive scalar magnitudes
\[
 \boxed{x,\qquad x^{-1},\qquad x^2,}
\]
respectively:
\[
 R=xI_G,\qquad R^+=x^{-1}I_G,\qquad R^\dagger R=x^2I_G.
\]
Thus the unordered role multiset \(\{x^{-1},x,x^2\}\) requires no additional
continuous coefficient.  Physical use of all three roles remains part of the
terminal source-consumption branch.
\end{proposition}

\begin{corollary}[Normal ordering fixes the role assignment]
\label{p6:cor:v568-normal-order-role-assignment}
For the generated value \(x=192^{1/3}>1\),
\[
 \boxed{x^{-1}<x<x^2.}
\]
Hence, conditional on occurrence of the three-role branch and normal mass
ordering, the assignment is unique:
\[
 \boxed{m_1\leftrightarrow x^{-1},\qquad
        m_2\leftrightarrow x,\qquad
        m_3\leftrightarrow x^2.}
\]
No additional permutation selector remains on that branch.
\end{corollary}

\begin{proposition}[Absolute-scale counterfamily]
\label{p6:prop:v568-absolute-scale-counterfamily}
The \(D_4\) frame determinant and the entire scale-free neutrino shape do not
determine the positive dimensionful normalization \(\mu_{\nu,56}\).  For
every \(\lambda>0\),
\[
 M_{\nu,+}^{(\lambda)}
 :=\lambda\mu_{\nu,56}T_{\nu,D_4}
\]
has the same \(D_4\) determinant owner, the same dimensionless response
parameter \(x\), the same ordered mass ratios and the same squared-gap
quotient as the unscaled operator, while every absolute mass and squared gap
changes by \(\lambda\) or \(\lambda^2\).  Therefore no theorem using only the
v5.67/v5.68 dimensionless \(D_4\) response can derive
\(\mu_{\nu,56}\).  Any identification with \(E_{56}\), a Higgs/VEV factor or
a pole normalization requires an additional dimensionful owner map.
\end{proposition}

\begin{proof}
Multiplication by \(\lambda\) leaves the dimensionless operator
\(T_{\nu,D_4}\), hence all ratios of its eigenvalues and of differences of
squared eigenvalues, unchanged.  It rescales the spectrum linearly and its
Gram spectrum quadratically.  The geometric input \(J_{D_4}^\dagger
J_{D_4}=192\) contains no dimensionful quantity that can distinguish members
of this counterfamily.
\end{proof}

\begin{proposition}[Central-response eigenframe no-go]
\label{p6:prop:v568-central-response-eigenframe-nogo}
The generated central response \(R_{D_4}^{\rm fam}=xI_G\) contains no neutral
mass-eigenframe information.  For every \(U\in U(G)\),
\[
 UR_{D_4}^{\rm fam}U^\dagger=R_{D_4}^{\rm fam}.
\]
Consequently the \(D_4\) determinant theorem cannot determine the three
rank-one mass projectors or the PMNS frame.  Those require a noncentral
same-action family operator or an independently derived relative-frame
transport.  Identifying the structural trimer basis with all mass bases would
instead force trivial relative mixing and is therefore not a permitted
shortcut.
\end{proposition}

\begin{remark}[Contracted live neutral frontier]
\label{p6:rem:v568-contracted-neutral-frontier}
After Papers I-A v5.60 and V-A v3.74, the live neutral obligations on the
preferred branch are:
(i) the source-consumption relation
\(\mathsf{N56}_{D_4}^{\rm consume}\), which is a component condition on the
same G15 source member rather than a new independent occurrence cluster;
(ii) one dimensionful normalization owner for \(\mu_{\nu,56}\); and
(iii) one noncentral relative-eigenframe/PMNS theorem.  The determinant count,
central cube-root lift, role magnitudes and normal-order permutation are no
longer live selector data.
\end{remark}
% END D4-NEUTRINO-CLOSURE-REFINEMENT v5.68


% BEGIN D4-TRIALITY-NEUTRINO-CONSUMPTION v5.69
\subsection{triality-parent consumption of the preferred terminal-neutral shape}
\label{p6:sec:v569-d4-triality-neutrino-consumption}

\noindent\textbf{Framework v2.208 branch scope.}
The neutral inverse/direct/closed-Gram entries in this retained section
belong to the historical $D_4$ response constructor. They are obsolete
on the harmonic rank branch of \ref{p6:sec:v2208-neutral-rank-word}.
Their algebraic theorems neither force nor constrain that new mass word.
The charged results and branch-independent geometric/source rules keep
their stated domains.


Paper~V-A v3.75 proves that, once a nonzero neutral mass/source coordinate
actually occurs, the generated \(D_4\) triality parent is not an optional
same-action port.  Claim-complete source ownership and GP2 exactness force its
consumption.  Therefore the v5.67--v5.68 \(D_4\) shape is no longer conditional
on an independent relation \(\mathsf{N56}_{D_4}^{\rm consume}\).

\begin{theorem}[Preferred finite neutral shape after source actuality]
\label{p6:thm:v569-preferred-neutral-shape-after-actuality}
Let a nonzero terminal-neutral source member occur in the already generated
Yukawa--Dirac source role.  Put
\[
 x:=192^{1/3}.
\]
Then the generated inverse/direct/Gram response orbit on the preferred
\(D_4\) branch is consumed in the same source owner, and normal ordering fixes
\[
 \boxed{m_1:m_2:m_3=x^{-1}:x:x^2}
\]
up to the still independent common dimensionful factor \(\mu_{\nu,56}\).
Equivalently,
\[
 M_{\nu,+}=\mu_{\nu,56}
 \left(x^{-1}P_1+xP_2+x^2P_3\right)
\]
for whatever noncentral rank-one projectors are supplied by the later neutral
eigenframe theorem.
\end{theorem}

\begin{proof}
Paper~I-A v5.60 generates \(R_{D_4}^{\rm fam}=xI_G\); the direct,
Moore--Penrose inverse and Gram operations give \(x,x^{-1},x^2\).  Paper~V-A
v3.75 removes independent source-consumption as a theoretical selector once
the nonzero neutral source is actual.  The v5.68 ordering theorem then assigns
the increasing orbit uniquely to \((m_1,m_2,m_3)\).  The scale and the
noncentral projectors are excluded from this conclusion by the v5.68 scale and
eigenframe no-go theorems.
\end{proof}

\begin{remark}[Updated neutral frontier]
\label{p6:rem:v569-updated-neutral-frontier}
On the preferred branch the finite shape now has only the ordinary G0
actuality antecedent for a nonzero neutral source; there is no separate
\(\mathsf{N56}_{D_4}^{\rm consume}\) assumption.  The remaining neutral
frontier is exactly:
(i) one dimensionful normalization/pole owner for \(\mu_{\nu,56}\); and
(ii) one noncentral relative-eigenframe theorem for PMNS.  The historical
\((1/\sqrt8,\sqrt5,13)\) radical remains legacy comparison only.
\end{remark}
% END D4-TRIALITY-NEUTRINO-CONSUMPTION v5.69


% BEGIN D4-TERMINAL-ROLE-SCALE-REFINEMENT v5.70
\subsection{synchronization to the rooted \texorpdfstring{$C_3$}{C3} source-cell closure}
\label{p6:sec:v575-g15-c3-sync}
The v5.73 statement of a surviving one-member gate
\(U_Q=X_G\) is superseded by Paper~V-A v3.81.  On a nonzero occurring odd
fermion-source branch, the same rooted triality/source owner has the frozen
nonredundant transport presentation
\[
 (I_G,C_{1,G},C_{2,G}),
\]
and Paper~II-A generated-constituent co-realization retains the complete cell.
Equivalently,
\[
 \boxed{\mathscr T_{C_3}^{\rm src}=\{I_G,X_G,X_G^{-1}\}.}
\]
Rooted inversion fixes \(I_G\) and exchanges \(X_G\) with \(X_G^{-1}\).
Thus no spectrum calculation in this paper is permitted to reintroduce a
matter-versus-antimatter orientation selector.  The two nontrivial
orientations are conjugate presentation members of one realized source cell;
\(I_G\) is the conjugation-fixed/neutral transport member.  Gauge neutrality
and Majorana self-conjugacy remain separately typed core questions.

Accordingly
\[
 \boxed{\begin{gathered}
 \texttt{G15\_C3\_SOURCE\_CELL}\\
 =\texttt{CLOSED\_DEPENDENT\_ON\_G0\_NONZERO}\\
 \texttt{\_FERMION\_SOURCE\_ACTUALITY}.
 \end{gathered}}
\]
The remaining spectrum obligations are the numerical physical family-orbit
members, optional neutral/Majorana lift data where claimed, and the
claim-scoped pole/RG analytic certificates.  CKM, Dirac-PMNS, and weak-CP
coordinates are dependent readouts once the corresponding orbit member is
fixed.  None reopens the finite transport
selector.
% END G15-C3-SYNC v5.75

% BEGIN FERMION-SOURCE-BOUNDARY v5.78
\paragraph{Historical antecedent and effective rooted specialization (v2.109).}
The following v5.78--v5.79 source-member notation records its historical
conditional antecedent.  On the active rooted branch the primitive source
occurs by V-A v3.92; a physical noncentral family and its terminal matching
were subject to the former same-action SVD candidate assembly. These
historical paragraphs neither reopen primitive occurrence nor a charged
assembly gap now closed by Paper~V-A Release v3.97.

\subsection{transitive running-operator/pole split}
\label{p6:sec:v579-transitive-rg-pole-split}

The v5.78 G0 retyping and Remark~\ref{p6:rem:v236-qrg-pole-split} imply a
strict premise-minimization rule for all downstream consumers.

\begin{theorem}[Running source operator does not consume the pole certificate]
\label{p6:thm:v579-running-source-no-pole-premise}
Fix an actual nonzero G0 fermion-source member and hence its presentation-free
quark source orbit \(\mathcal O_q\).  On any renormalization interval for which
the declared QRG/threshold matching map exists, the running quark
mass/source operator \(M_q(\mu)\) is a deterministic transport of that same
microscopic member.  Its construction does not require an isolated physical
pole, a pole sheet, or a residue normalization.

If a downstream claim is instead a pole/mass-shell claim, it additionally
consumes the renormalized inverse-response continuation together with the
sheet/domain, isolated-zero and residue hypotheses stated for G16.
\end{theorem}
\begin{proof}
By Theorem~\ref{p6:thm:v578-family-member-g0-retyping}, the physical family
member is already the occurring G0 source member rather than a second
structural selector.  Remark~\ref{p6:rem:v236-qrg-pole-split} states that QRG
transports an already owned operator between renormalization coordinates and
thresholds, while the pole conditions are additional conditions on the
continued two-point/inverse-response object.  Therefore the latter cannot be a
premise of a claim whose target is only the running operator.
\end{proof}

\begin{corollary}[Minimal downstream G16 dependency]
\label{p6:cor:v579-minimal-g16-dependency}
For a QCD or chiral consumer whose target is \(M_q(\mu)\), the live G16 input is
only the RG/threshold matching-and-running certificate on the consumed scale
interval.  The G16 pole certificate branches off only when a pole observable is
claimed.  No downstream consumer may reintroduce a separate
\texttt{FULL\_FAMILY\_MEMBER} or quark-mass constructor.
\end{corollary}

\begin{remark}[Supersession of historical live-package wording]
\label{p6:rem:v579-historical-live-package-superseded}
Earlier sections that described a free G15 family member, a G5Q-loop member, or
a full pole package as part of every downstream mass handoff are retained as
historical derivation provenance.  Their live dependency meaning is superseded
by Theorem~\ref{p6:thm:v578-family-member-g0-retyping} and
Corollary~\ref{p6:cor:v579-minimal-g16-dependency}: actuality is G0, the owner
family is generated, running is RG transport, and pole analyticity is
claim-scoped.
\end{remark}
% END TRANSITIVE-RG-POLE-SPLIT v5.79

% BEGIN SAFEGUARD-G5Q-RETIREMENT v5.80

\section{Neutral response geometry and conditional mass branch}
\label{p5b:sec:organized-neutral}
\subsection{Exclusive light-neutrino operator}
\label{p6:sec:exclusive-neutrino}

\begin{proposition}[Former inverse-shell spectrum]
\label{p6:prop:inverse-shell-neutrino}
At the retained source order, the inverse-shell branch has mass ratios
\[
1:8.56985590055:53.5380616385
\]
and adjacent squared-gap quotient
\[
R_{\rm shell}=38.5531187706\ldots .
\]
Its auxiliary two-variable expression \(q(a,b)\) was never defined, so the
additional claimed \(56\)-line closure is incomplete.
\end{proposition}

\begin{proposition}[Radical and affine spectra]
\label{p6:prop:radical-affine-neutrino}
The radical branch has
\[
(m_1,m_2,m_3)
=
\sqrt{\etaH}E_{56}
\left(\frac1{\sqrt8},\sqrt5,13\right)
\]
and squared gaps proportional to
\[
39:1312:1351,
\qquad
R_{\rm rad}=\frac{1312}{39}.
\]
By contrast, for
\[
\DfamI=\diag(2/3,-1/3,-1/3),
\qquad
\DfamII=\diag(0,1,-1),
\]
the affine operator
\[
\alpha I+\beta\DfamI+\frac{\beta}{56}\DfamII
\]
has pairwise eigenvalue gaps proportional to
\[
2:55:57,
\qquad
R_{\rm aff}=\frac{55}{2}.
\]
\end{proposition}

\begin{theorem}[Neutrino incompatibility]
\label{p6:thm:neutrino-incompatibility}
The inverse-shell, radical, and affine-\(1/56\) constructions are not
normalizations of one nondegenerate spectrum when each positive
mass-squared spectrum is placed in increasing order.
\end{theorem}

\begin{proof}
The three adjacent-gap quotients
\[
38.5531187706\ldots,\qquad
\frac{1312}{39},\qquad
\frac{55}{2}
\]
are pairwise unequal.  Scaling preserves them, adding a multiple of the
identity preserves squared gaps, and relabeling only permutes one fixed
multiset before the canonical increasing order is restored.
\end{proof}

\begin{definition}[Light-neutrino realization family]
\label{p6:def:light-neutrino-realization-family}
Before the flavor package \(\mathsf F\) is frozen, the inverse-shell,
radical, and affine constructions above are inequivalent competing neutral
operator realizations; Theorem~\ref{p6:thm:neutrino-incompatibility} proves
that they cannot be regarded as normalizations of one spectrum.  Releases
through v5.66 retained the radical operator below.  Release v5.67 supersedes
that numerical selection by the conditional Hurwitz--\(D_4\) terminal-
determinant branch of Section~\ref{p6:sec:v567-d4-terminal-neutrino}.  The
radical operator is retained here only as a legacy comparison target.
\end{definition}

\begin{definition}[Legacy radical operator (pre-v5.67 comparison)]
\label{p6:def:selected-radical-neutrino}
Fix a real orthonormal neutral basis
\((e_1,e_2,e_3)\) and its real symmetric rank-one projectors
\[
P_i=|e_i\rangle\langle e_i|,
\qquad
P_iP_j=\delta_{ij}P_i,
\qquad
P_1+P_2+P_3=I.
\]
The exclusive light-neutrino mass-squared operator and its positive root are
\begin{equation}
\label{p6:eq:selected-neutrino-square}
H_\nu^{(56)}
=
\etaH E_{56}^2
\left(\frac18P_1+5P_2+169P_3\right),
\end{equation}
\begin{equation}
\label{p6:eq:selected-positive-neutrino}
\boxed{
M_{\nu,+}
=
\sqrt{\etaH}E_{56}
\left(\frac1{\sqrt8}P_1+\sqrt5P_2+13P_3\right).
}
\end{equation}
The inverse-shell neutrino spectrum and the common \(1/56\) light-neutrino
slope are not part of the reconstructed model.
\end{definition}

\begin{proposition}[Affine encoding of the selected square]
\label{p6:prop:radical-affine-encoding}
Identify \((e_1,e_2,e_3)\) with the ordered coordinate basis in which
\(\DfamI,\DfamII\) have the displayed diagonal forms.  Then
\[
H_\nu^{(56)}
=
\etaH E_{56}^2
\left(
\frac{1393}{24}I
-\frac{695}{8}\DfamI
-82\DfamII
\right).
\]
This is the unique element of the three-dimensional span
\[
\operatorname{span}\{I,\DfamI,\DfamII\}
\]
whose ordered eigenvalues are
\[
\etaH E_{56}^2\left(\frac18,5,169\right).
\]
\end{proposition}

\begin{proof}
The three diagonal coefficient vectors are linearly independent.  Solving
the three scalar equations gives the displayed coefficients; direct
substitution returns \(1/8,5,169\).
\end{proof}

\begin{corollary}[Legacy conditional neutral numerical packet]
\label{p6:cor:neutral-numerical-packet}
On the source \(\mathrm{REG}_0\) member with the declared absolute calibration,
\(E_{56}=3.924776812825316\,{\rm meV}\),
\[
(m_1,m_2,m_3)
\simeq
(1.359583,\ 8.598755,\ 49.991243)\ {\rm meV},
\qquad
\sum_i m_i\simeq59.949581\ {\rm meV},
\]
and
\[
\begin{aligned}
\Delta m_{21}^2&\simeq7.209013\times10^{-5}\ {\rm eV}^2,\\
\Delta m_{31}^2&\simeq2.497276\times10^{-3}\ {\rm eV}^2,\\
\Delta m_{32}^2&\simeq2.425186\times10^{-3}\ {\rm eV}^2.
\end{aligned}
\]
These values are conditional on \(\mathsf{SLOT}\), \(\mathsf{TMF}\),
\(\CMPexp\), and \(\Uroleabs\); the \(K=56\) anchor and energy ratio do not
require \(\CLKrem\).
\end{corollary}

\begin{theorem}[External Takagi sign classification]
\label{p6:thm:majorana-signs}
If a real symmetric Majorana operator has
\[
O^TMO=\diag(\lambda_1,\lambda_2,\lambda_3),
\]
then its Takagi masses are \(|\lambda_i|\).  A negative eigenvalue requires a
column factor \(i\) and therefore contributes relative CP parity \(\pi\).
Realness alone implies only
\[
\alpha_{21},\alpha_{31}\in\{0,\pi\}.
\]
This is a theorem inside the external Majorana/Takagi comparison dictionary; it does not define the framework-native N56 terminal masses or current incidence.
\end{theorem}

\begin{proof}
For \(\lambda_i<0\), multiplication of the Takagi column by \(i\) changes
\(\lambda_i\) to \(i^2\lambda_i=|\lambda_i|\) under congruence.  This
Majorana sign cannot be removed by a continuous mass-eigenstate rephasing.
\end{proof}

\begin{corollary}[Why the selected branch has zero relative parity]
\label{p6:cor:selected-majorana-parity}
Equation~\eqref{p6:eq:selected-positive-neutrino} defines the mass operator
itself, not merely its square, and all three coefficients are positive.
Hence the retained positive neutral operator has the explicit all-positive parity branch.
The conclusion follows from positivity, not from realness.
\end{corollary}

\subsection{N56 neutral energy/incidence sibling closure}
\label{p6:sec:v593-n56-neutral-incidence-closure}

\noindent\textbf{Framework v2.208 branch scope.}
The neutral inverse/direct/closed-Gram entries in this retained section
belong to the historical $D_4$ response constructor. They are obsolete
on the harmonic rank branch of \ref{p6:sec:v2208-neutral-rank-word}.
Their algebraic theorems neither force nor constrain that new mass word.
The charged results and branch-independent geometric/source rules keep
their stated domains.


\begin{theorem}[Complete neutral terminal energy/incidence sibling packet]
\label{p6:thm:v593-neutral-energy-incidence-packet}
Import the rooted neutral response/transport sibling theorem of Paper~V-A
Release v3.100.  With
\[
 \mathcal E_{\rm inv}=E_{56}x^{-1},\qquad
 \mathcal E_{\rm dir}=E_{56}x,\qquad
 \mathcal E_{\rm Gram}=E_{56}x^2,
 \qquad x^3=192,
\]
and the canonical neutral role evaluation $W_\nu$, define
\[
 \iota_{\nu,r}:=W_\nu\Pi_r^{\rm role},
 \qquad r\in\{\mathrm{inv},\mathrm{dir},\mathrm{Gram}\}.
\]
Then the complete neutral terminal packet is
\[
 \boxed{
 \mathfrak P_{\nu,56}
 =\bigl((\mathcal E_r,\iota_{\nu,r})\bigr)_r,}
\]
represented in the rooted transport presentation by
\[
 \boxed{
 (E_{56}x^{-1},X_G^{-1}),\qquad
 (E_{56}x,X_G),\qquad
 (E_{56}x^2,I_G).}
\]
The first entry of each pair is the already generated terminal role energy and
the second is its preterminal current incidence sibling.  Neither is inferred
from the spectrum of the other.
\end{theorem}

\begin{corollary}[Neutral current evaluation closes before any PMNS dictionary]
\label{p6:cor:v593-neutral-current-evaluation-closed}
The neutral terminal-role evaluation
\[
 \boxed{W_\nu:\mathscr L_{56}\xrightarrow{\cong}G_L}
\]
is generated up to the rooted $D_3$ presentation quotient.  No independent
$J_\nu$, Takagi frame, Majorana mass eigenbasis or family-unitary selector is
required to define the neutral current incidence.  Accordingly
\[
 \boxed{
 \texttt{NEUTRAL\_CURRENT\_INCIDENCE}
 =\texttt{CLOSED\_DEPENDENT\_ON\_ROOTED\_D4/C3\_SIBLING}.}
\]
The terminal-energy to physical pole/rest-mass identification remains the
separately typed downstream bridge already carried by the framework.
\end{corollary}

\begin{remark}[Compatibility with the v5.87 role-space firewall]
The v5.87 statement that the central operator $xI_G$ contains no family-frame
information remains valid.  Release v5.93 does not extract $W_\nu$ from
$xI_G$.  It retains $X_G$ as a sibling output of the same rooted source before
scalarization and carries the common inverse/direct/closed-Gram operation
labels through both outputs.
\end{remark}
% END N56-NEUTRAL-INCIDENCE-CLOSURE v5.93

% BEGIN CHARGED-CURRENT-INCIDENCE-CLOSURE v5.94
\subsection{K56 geometry and the conditional harmonic rank response word}
\label{p6:sec:v2208-neutral-rank-word}
\noindent\textbf{Current branch scope.}
The old inverse/direct/closed-Gram word $x^{-1},x,x^2$ is obsolete on
the branch addressed here. Its existing $D_4/C_3$ theorems are retained
as historical statements within their declared constructors. They impose
no compatibility gate on the new harmonic rank proposal. No physical
closure of that replacement word is claimed in this release.

The native K56 anchor is the scalar Clifford endpoint
$P_0=\mathbb E_{\rm IR}$ of \ref{p3:cor:derived-infrared-projector}.
Paper I-C now proves its rooted geometric family bridge and the harmonic
boundary map in \ref{p1b:sec:v2208-neutral-boundary-geometry}.
The bridge has a full nine-dimensional operator corner. Its endpoint
pullback $Q_c$ is one geometric operator element, while the spectator
extension of the Clifford core map retains every family coefficient
by \ref{p5:thm:v2208-terminal-tensor-extension}.
Thus the scalar anchor alone determines neither three finite mass
coefficients nor a rank-one family mass.

\begin{proposition}[Typed collective rank response magnitudes]
\label{p6:prop:v2208-collective-rank-word}
For an $n$-dimensional owner $E$ with its ordinary Hilbert metric, equip
$\End(E)$ with the unnormalized Hilbert--Schmidt metric and set
\[
 \iota_E:\mathbb C\longrightarrow\End(E),\qquad
 \iota_E(z)=zI_E.
\]
Then
\[
 \iota_E^\dagger\iota_E=nI_{\mathbb C},\qquad
 \iota_E^+(A)=\frac{\Tr A}{n},\qquad
 |\iota_E|=\sqrt n,\quad |\iota_E^+|=1/\sqrt n,
\]
where the final magnitudes are the single nonzero singular values.
For the owners
$E_{163}=\ker B$ and $E_5=B^\dagger\ker\varepsilon$ of
\ref{p1b:thm:v2208-boundary-projectors}, the typed response packet
\[
 (\iota_{E_{163}}^+,\iota_{E_5},\iota_{E_{163}})
\]
therefore has collective magnitudes
\[
 \boxed{1/\sqrt{163},\quad\sqrt5,\quad\sqrt{163}.}
\]
\end{proposition}
\begin{proof}
The adjoint of $z\mapsto zI_E$ is $A\mapsto\Tr A$ for the stated metric,
and $\Tr I_E=n$. Its positive Gram is $nI$ and its Moore--Penrose
inverse is $(\iota_E^\dagger\iota_E)^{-1}\iota_E^\dagger=\Tr/n$.
Their singular magnitudes follow. Substitute the ranks proved in Paper I-C.
\end{proof}
These ports have their displayed operator domains; they are not an
automatic field compression into $G$ or a newly forced physical role packet.
The order inverse-kernel/direct-augmentation/direct-kernel specifies the
mathematical proposal. Its physical selection is not established by this definition.


% BEGIN FRAMEWORK v2.239: NEUTRAL COMPLETE-OWNER RETURN
\begin{corollary}[Exact nonsinglet/augmentation cross response and the K56 action test]
\label{p6:cor:v239-neutral-green-port-test}
On the native harmonic boundary of
Theorem~\ref{p1b:thm:v239-nonsinglet-green-schur}, the complete
nonsinglet owner $E_{168}=E_{163}\oplus E_5$ and the augmentation
share the nonzero canonical Green cross return
\[
 \boxed{\Tr(P_{\rm ns}K_{12}^{-1}P_{\rm aug})=\frac5{168},
 \qquad \Tr(QK_{12}^{-1}P_{\rm aug})=0.}
\]
The trace-scalar inverse on the full nonsinglet owner evaluates its
augmentation projector to
$\iota_{E_{168}}^+(P_{\rm aug}|_{E_{168}})=5/168$;
the restricted augmentation on $E_{163}$ vanishes. The two owners
therefore cannot be interchanged inside a common preterminal
constructor consuming their cross return. The original conditional
$163,5,1/163$ rank-magnitude word remains valid on its stated
\emph{reduced} owner; the $168$ Green return is not a replacement
physical mass word or a derivation of a neutrino splitting ratio.
The remaining K56 certificate is to determine whether its \emph{actual}
action consumes the unreduced Green packet and, if it reduces it,
where the augmentation Schur elimination occurs relative to the
physical terminal/pole map.
\end{corollary}
\begin{proof}
The two traces and reduction are
\eqref{p1b:eq:v239-nonsinglet-green-gram}. The trace-scalar inverse
is $\Tr(\cdot)/168$ on $E_{168}$ by
Proposition~\ref{p6:prop:v2208-collective-rank-word};
$Q P_{\rm aug}=0$ gives the zero kernel restriction.
The conclusion about the ordering is
Theorem~\ref{p2a:thm:v3102-preterminal-relation-preservation}.
\end{proof}
% END FRAMEWORK v2.239: NEUTRAL COMPLETE-OWNER RETURN

\begin{corollary}[Conditional terminal row-energy coordinates]
\label{p6:cor:v2208-conditional-rank-energies}
If the independently licensed new K56 constructor requires precisely this
harmonic input and collective response packet as its terminal roles, then
the existing row-energy insertion gives
\[
 \boxed{E_{56}\{1/\sqrt{163},\sqrt5,\sqrt{163}\}.}
\]
The source handoff of \ref{p5:thm:v2208-neutral-boundary-export} applies
once that required mass/incidence packet and insertion are generated.
\end{corollary}
\begin{proof}
Apply \ref{p6:thm:v570-e56-terminal-row-energy-scale} to the dimensionless
role magnitudes under the stated constructor hypothesis. The conditional
export theorem then supplies inherited source occurrence in its stated sort.
\end{proof}
% BEGIN FRAMEWORK v2.243: NEUTRAL ROOTED BOUNDARY OWNER CHECK
\begin{proposition}[Canonical harmonic owner and the complete five-port geometric return]
\label{p6:prop:v243-neutral-icosian-owner-test}
The first nonconstant $2I$-invariant degree-$12$ harmonic
$h_{12}$ of Theorem~\ref{p1b:thm:v243-icosian-h12-origin} is a
canonical nonzero \emph{geometric} seed for the full
$169$-dimensional round harmonic representation.  Under the actual
rooted six-point boundary coisometry it satisfies
\[
 \boxed{P_{\rm aug}h_{12}=0,\qquad
 \|P_{\rm sing}h_{12}\|^2=\frac{13}{3},\qquad
 \|Q_{163}h_{12}\|^2=\frac{26}{3}.}
\]
The complete five-dimensional boundary augmentation nonetheless has
the explicit isometric geometric provider
$\mathcal I_5:G_\perp\oplus\mathbb C^3\xrightarrow{\cong}E_5$
from \eqref{p1b:eq:v243-five-augmentation-isometry}.  Thus a
neutral physical source word which first scalarizes to $h_{12}$
alone cannot subsequently produce a nonzero $E_5$ boundary
response by the unchanged evaluation map.  If the physical K56
construction needs that channel, it must retain a licensed
unsummed six-port source or an independently proved action-owned
incidence before such scalarization.  The $163\oplus5\oplus1$
geometric packet is complete; this proposition makes no
physical K56 occurrence, mass, normalization or pole assertion.
\end{proposition}
\begin{proof}
Combine Theorems~\ref{p1b:thm:v243-icosian-h12-origin} and
\ref{p1b:thm:v243-neutral-five-isometry}.  The group-averaged
source is constant $13$ on the six rooted Hurwitz points, so its
normalized evaluation is a singlet.  The exact even two-dimensional
family and odd three-dimensional quaternion providers are instead
orthogonal and exhaustive in the mean-zero boundary image.
\end{proof}
% END FRAMEWORK v2.243: NEUTRAL ROOTED BOUNDARY OWNER CHECK

The actual K56 identification with $\mathcal H_{12}$, the required new
collective word and its incidence packet remain
\texttt{OPEN\_PHYSICAL\_ORIGIN}. Co-consumption retains required outputs
of a licensed constructor; it does not create that constructor from the
available ranks. The geometric K56 bridge is not the neutral physical
role-to-family realization. No K16 pairing or production calculation,
new binding factor, obsolete-word comparison, or later SVD/Takagi
definition of terminal masses is introduced as a prerequisite.
The common $E_{56}$ row scale is retained, and observed pole/rest claims
remain within their separately typed downstream readout domain.
% END NEUTRAL-RANK-WORD-SCOPE v2.208

% BEGIN LITERATURE-VERIFIED SOURCE AND HARMONIC TESTS v2.210
\subsection{conditional Hurwitz--\texorpdfstring{\(D_4\)}{D4} terminal-determinant neutrino branch}
\label{p6:sec:v567-d4-terminal-neutrino}

\noindent\textbf{Framework v2.208 branch scope.}
The neutral inverse/direct/closed-Gram entries in this retained section
belong to the historical $D_4$ response constructor. They are obsolete
on the harmonic rank branch of \ref{p6:sec:v2208-neutral-rank-word}.
Their algebraic theorems neither force nor constrain that new mass word.
The charged results and branch-independent geometric/source rules keep
their stated domains.


This section is controlling for the preferred light-neutrino shape branch.
It does not alter the charged compiler of Release v5.66.  The branch is
conditional on the Paper-V-A occurrence gate
\(\mathsf{N56}_{D_4}^{\rm occ}\); absent that occurrence certificate the
construction remains mathematical capacity rather than a physical spectrum.

\begin{definition}[Generated terminal response parameter]
\label{p6:def:v567-d4-terminal-x}
Conditional on \(\mathsf{N56}_{D_4}^{\rm occ}\), define
\[
 \boxed{x:=192^{1/3}>1.}
\]
The number \(192\) is not fitted in Paper~V-B: it is the closed oriented-frame
determinant return imported from Paper~I-A.
\end{definition}

\begin{definition}[Minimal inverse/direct/Gram neutral branch]
\label{p6:def:v567-inverse-direct-gram-branch}
Let \((P_1,P_2,P_3)\) be the orthogonal rank-one projectors of the neutral
mass basis.  The minimal three-role orbit of one positive response uses the
already available inverse, direct, and closed Gram constructors.  Define its
positive shape operator by
\[
 \boxed{
 T_{\nu,D_4}
 =x^{-1}P_1+xP_2+x^2P_3 .}
\]
For an as-yet unproved positive dimensionful normalization \(\mu_{\nu,56}\),
set
\[
 \boxed{
 M_{\nu,+}^{D_4}=\mu_{\nu,56}T_{\nu,D_4},\qquad
 H_\nu^{D_4}=\mu_{\nu,56}^{\,2}
 \bigl(x^{-2}P_1+x^2P_2+x^4P_3\bigr).}
\]
The association of the three ordered mass channels with the inverse, direct,
and Gram roles is part of the terminal occurrence branch; it introduces no
continuous coefficient.  The neutral eigenframe itself remains downstream.
\end{definition}

\begin{theorem}[Scale-free \(D_4\) neutrino shape]
\label{p6:thm:v567-d4-scale-free-neutrino-shape}
On the branch of Definition~\ref{p6:def:v567-inverse-direct-gram-branch}, the
ordered mass ratios are
\[
 \boxed{m_1:m_2:m_3=x^{-1}:x:x^2,
 \qquad x^3=192.}
\]
The overall normalization cancels from the squared-gap quotient, giving
\[
 \boxed{
 \mathcal R_{31/21}^{D_4}
 :=\frac{\Delta m_{31}^2}{\Delta m_{21}^2}
 =\frac{x^4-x^{-2}}{x^2-x^{-2}}
 =\frac{x^6-1}{x^4-1}
 =33.3105115561\ldots .}
\]
Consequently
\[
 \frac{\Delta m_{32}^2}{\Delta m_{21}^2}
 =\mathcal R_{31/21}^{D_4}-1
 =32.3105115561\ldots .
\]
These are shape predictions of the conditional branch; they do not determine
\(\mu_{\nu,56}\).
\end{theorem}

\begin{proof}
The eigenvalues of \(M_{\nu,+}^{D_4}\) are the three displayed positive
coefficients times \(\mu_{\nu,56}\).  Therefore
\[
 \Delta m_{21}^2=\mu_{\nu,56}^{\,2}(x^2-x^{-2}),\qquad
 \Delta m_{31}^2=\mu_{\nu,56}^{\,2}(x^4-x^{-2}).
\]
Cancel the common scale and multiply numerator and denominator by \(x^2\).
Substitution of \(x^3=192\) gives the numerical value.
\end{proof}

\begin{corollary}[One-scale reconstruction after an external absolute input]
\label{p6:cor:v567-one-scale-reconstruction}
Once any one independent absolute neutrino datum fixes the positive scale
\(\mu_{\nu,56}\), the other two masses are fixed by the generated shape.  For
example, a supplied value of \(\Delta m_{21}^2\) gives
\[
 \mu_{\nu,56}
 =\sqrt{\frac{\Delta m_{21}^2}{x^2-x^{-2}}},
\]
followed by
\[
 (m_1,m_2,m_3)
 =\mu_{\nu,56}(x^{-1},x,x^2).
\]
This is a reconstruction formula, not an internal derivation of the absolute
scale.
\end{corollary}

\begin{remark}[Status of the pre-v5.67 radical branch]
\label{p6:rem:v567-legacy-radical-status}
\leavevmode\par\noindent
The earlier coefficients
\((1/\sqrt8,\sqrt5,13)\) remain useful historical diagnostics, but on the
preferred v5.67 branch they are not independent live selector data.  Their
former numerical packet is not used to normalize the new branch.  The live
neutral obligations are the same-action gate
\(\mathsf{N56}_{D_4}^{\rm occ}\), the absolute scale
\(\mu_{\nu,56}\), and the neutral mass-eigenframe/PMNS occurrence.
\end{remark}
% END D4-TERMINAL-NEUTRINO v5.67

% BEGIN D4-NEUTRINO-CLOSURE-REFINEMENT v5.68
\subsection{terminal shell insertion, depth-one role closure, and rooted-frame narrowing}
\label{p6:sec:v570-d4-terminal-role-scale-refinement}

\noindent\textbf{Framework v2.208 branch scope.}
The neutral inverse/direct/closed-Gram entries in this retained section
belong to the historical $D_4$ response constructor. They are obsolete
on the harmonic rank branch of \ref{p6:sec:v2208-neutral-rank-word}.
Their algebraic theorems neither force nor constrain that new mass word.
The charged results and branch-independent geometric/source rules keep
their stated domains.


Release v5.69 correctly removes an independent occurrence gate for the
consumed \(D_4\) parent, but its statement that the complete inverse/direct/
Gram orbit is thereby physically consumed is stronger than the presently
proved constructor ancestry.  This release separates mathematical response
closure from simultaneous terminal-role consumption and synchronizes the
already derived \(K=56\) energy anchor.

\begin{theorem}[Complete nontrivial depth-one response closure]
\label{p6:thm:v570-depth-one-response-closure}
Let
\[
 R=xI_G,\qquad x=192^{1/3}>1,
\]
be the generated positive rooted-family response.  From the seed response
itself, the nontrivial depth-one positive readouts licensed by the existing
boundary/direct/inverse/closed-return grammar are
\[
 \boxed{R^+=x^{-1}I_G,\qquad R=xI_G,\qquad R^\dagger R=x^2I_G.}
\]
Thus the complete nontrivial first-generation scalar magnitude set is
\[
 \boxed{\{x^{-1},x,x^2\}.}
\]
Values such as \(x^{-2}\) or higher powers require an additional constructor
generation, while positive codiagonals require extra biproduct/phase data.
\end{theorem}

\begin{definition}[Terminal three-role completion gate]
\label{p6:def:v570-terminal-three-role-gate}
The remaining finite neutral-role statement is
\[
 \boxed{\begin{gathered}
 \mathsf{N56}_{\rm role}:\ \text{the claim-complete terminal neutral}\\
 \text{constructor consumes all three members}\\
 \{R^+,R,R^\dagger R\}\ \text{as its three one-pass channels}.
 \end{gathered}}
\]
The mathematical existence and values of the three readouts are derived by
Theorem~\ref{p6:thm:v570-depth-one-response-closure}; their simultaneous
physical use is not inferred merely from occurrence of the parent response.
\end{definition}

\begin{theorem}[Native \(E_{56}\) terminal row-energy scale]
\label{p6:thm:v570-e56-terminal-row-energy-scale}
On any terminal neutral branch whose role index is \(K=56\), the framework's
terminal row-energy functor already supplies the unique dimensionful shell
insertion \(E_{56}\).  Therefore a dimensionless neutral response
\(T_\nu\) has premise-minimal terminal row-energy operator
\[
 \boxed{\mathcal E_\nu^{\rm term}=E_{56}\,T_\nu.}
\]
There is no independent neutrino-specific dimensionful owner
\(\mu_{\nu,56}\) at this row-energy level unless the same source word
generates an additional dimensionless scalar normalization
\(\kappa_{56}\), in which case \(\mu_{\nu,56}=\kappa_{56}E_{56}\).  In the
absence of such a generated port, no-insertion gives \(\kappa_{56}=1\).
This theorem does not identify terminal row energy with a downstream physical
pole/rest-mass readout.
\end{theorem}

\begin{corollary}[Conditional terminal spectrum on the three-role branch]
\label{p6:cor:v570-e56-three-role-spectrum}
If \(\mathsf{N56}_{\rm role}\) holds, then
\[
 \boxed{\mathcal E_{\nu,1}:\mathcal E_{\nu,2}:\mathcal E_{\nu,3}
 =x^{-1}:x:x^2,\qquad x^3=192,}
\]
and, for normal ordering,
\[
 \boxed{\mathcal E_\nu^{\rm term}
 =E_{56}\bigl(x^{-1}P_1+xP_2+x^2P_3\bigr).}
\]
The scale-free gap quotient of v5.67 is unchanged.  Any comparison with pole
or rest masses remains under the universal downstream readout bridge rather
than a neutrino-specific scale selector.
\end{corollary}

\begin{proposition}[Rooted axial source narrows the neutral frame residual]
\label{p6:prop:v570-rooted-axial-neutral-frame-narrowing}
On the rooted axial terminal source word of Paper~I-A v5.62, the triality
contraction identifies the two complementary trimers by \(I_G\).  Hence the
rank-one rooted coordinate projectors \(P_a,P_b,P_c\) are generated endpoint
projectors of that source word.  An arbitrary independently supplied
\(U(3)\)-rotated projector triple is therefore not premise-minimal on this
branch.  The remaining neutral assignment is a finite matching of the three
response roles \(R^+,R,R^\dagger R\) to \(P_a,P_b,P_c\), modulo the already
generated cyclic/orientation presentation symmetries.
\end{proposition}

\begin{remark}[Supersession of the v5.68 scale no-go]
\label{p6:rem:v570-scale-nogo-scope}
The v5.68 scaling counterfamily remains correct for the \emph{dimensionless
\(D_4\) response considered alone}.  It does not imply a second independent
neutrino energy scale once the same terminal word already carries the derived
\(E_{56}\) shell anchor.  The residual scale issue is only the universal
terminal-row-energy to physical-pole/readout bridge, together with any
explicitly generated dimensionless normalization port.
\end{remark}

\begin{remark}[Updated neutral frontier]
\label{p6:rem:v570-updated-neutral-frontier}
The preferred neutral branch now retains: (i) the finite
\(\mathsf{N56}_{\rm role}\) simultaneous-role gate; (ii) a finite
role-to-rooted-trimer matching on the axial source branch; and (iii) the
relative neutral/charged frame theorem required for PMNS-type readouts.
\texttt{MU\_NU56} is retired as an independent neutrino-specific dimensionful
assumption at terminal-row-energy level.
\end{remark}
% END D4-TERMINAL-ROLE-SCALE-REFINEMENT v5.70


% BEGIN G15-Q5-SYNC v5.73

\section{Physical pole, clock and source-comparison interface}
\label{p5b:sec:organized-physical}
\subsection{Finite pole matrices and the vertex boundary}
\label{p6:sec:pole-matrices}

\begin{theorem}[Complete coherent-boundary family]
\label{p6:thm:coherent-boundary-family}
Let \(\{|\beta_a\rangle\}_{a=1}^R\) be orthonormal boundary modes, with a
transitive permutation action, and let the source line
\(\CC|B\rangle\) be fixed.  Every equivariant coherent map has the form
\[
C_b|B\rangle=b\sum_{a=1}^R|\beta_a\rangle,
\qquad b\in\CC.
\]
Its only nonzero singular value is
\[
\sigma(C_b)=|b|\sqrt R.
\]
For
\[
\mathbb K_b=
\begin{pmatrix}
N&C_b^*\\
C_b&NI_R
\end{pmatrix},
\]
\[
\Spec(\mathbb K_b)
=
\{N+|b|\sqrt R,\ N-|b|\sqrt R,\ N^{(R-1)}\}.
\]
\end{theorem}

\begin{proof}
Transitivity makes the invariant boundary line the span of
\(\sum_a|\beta_a\rangle\), but leaves the coefficient \(b\) free.
Orthogonality gives
\[
C_b^*C_b=|b|^2R\,|B\rangle\langle B|.
\]
Only the normalized coherent vector couples to \(|B\rangle\), leaving a
two-dimensional block with off-diagonal entry \(|b|\sqrt R\).
\end{proof}

\begin{corollary}[Status of \(1/120\)]
\label{p6:cor:one-over-120}
The former pole rule is the conditional choice
\[
b_B=\frac1{120},
\qquad
(R_W,R_Z,R_H)=(1,7,1),
\qquad
\varepsilon_B=+1.
\]
Equal phase fixes the coherent direction, not the magnitude, ranks, or upper
sign.  These data form the finite matrix subpackage
\(\mathsf{PV}_{\rm mat}\).
\end{corollary}

\begin{definition}[Conditional pole-coordinate packet]
\label{p6:def:conditional-pole-packet}
Under \(\mathsf{PV}_{\rm mat}\), the retained comparison coordinates are
\[
\begin{aligned}
(M_W^{\rm mat})^2
&=\etaH E_{32}^2\left(7+\frac1{120}\right),\\
(M_Z^{\rm mat})^2
&=\etaH E_{32}^2\left(9+\frac{\sqrt7}{120}\right),\\
(M_H^{\rm mat})^2
&=\etaH E_{32}^2\left(17+\frac1{120}\right).
\end{aligned}
\]
They are exact eigenvalues of the declared finite matrices.  Calling them
physical poles additionally requires an interaction self-energy,
renormalization prescription, and a proof selecting the displayed branch.
\end{definition}

\begin{theorem}[Mass matrices do not fix
\(\sin^2\theta_{\rm eff}^{\ell}\)]
\label{p6:thm:effective-angle-no-go}
The matrix packet may define the mass-ratio coordinate
\[
s_{\rm mat}^2
=1-\frac{(M_W^{\rm mat})^2}{(M_Z^{\rm mat})^2}.
\]
It cannot by itself determine the effective leptonic angle defined through
\[
\frac{g_V^\ell}{g_A^\ell}
=1-4|Q_\ell|\sin^2\theta_{\rm eff}^{\ell}.
\]
In particular, adding \(1/120\) to \(s_{\rm mat}^2\) is a vertex bridge, not
an eigenvalue theorem.
\end{theorem}

\begin{proof}
Keep the two finite mass matrices fixed and vary a finite, gauge-compatible
correction to the renormalized \(Z\ell\ell\) vertex.  The ratio
\(g_V^\ell/g_A^\ell\) changes while both matrix eigenvalues remain fixed.
Relating the two requires the interaction action, photon separation,
self-energies, Ward identities, and a renormalization scheme.
\end{proof}

\subsection{Quark comparison requires an RG bridge}
\label{p6:sec:qrg}

The continuum comparison language follows the standard running-mass,
threshold-matching, and reporting-point conventions summarized by the
Particle Data Group \pvcite{PDG2024}.

\begin{theorem}[Mixed reporting points are not one spectrum]
\label{p6:thm:mixed-quark-reporting}
The tuple
\[
\bigl(
m_u(2\,{\rm GeV}),m_d(2\,{\rm GeV}),m_s(2\,{\rm GeV}),
m_c(m_c),m_b(m_b),m_t(\mu_{\rm EW})
\bigr)
\]
contains renormalized coordinates at different scales, and the top entry may
also use a different short-distance convention.  Ratios formed directly
between such entries are not common-scheme observables.
\end{theorem}

\begin{proof}
Renormalized quark masses have a nonzero anomalous dimension.  Changing scale
or finite renormalization changes the coordinates without changing the
underlying finite carrier.  A ratio is transport invariant only after both
entries are placed in the same scheme, active-flavor theory, and scale.
\end{proof}

\begin{theorem}[Exact common-slice transport]
\label{p6:thm:exact-qrg-transport}
Fix a continuum scheme \(S\).  On an interval with \(n_f\) active flavors and a continuous coupling
trajectory, assume \(\beta_S^{(n_f)}(a)\neq0\) on the coupling
image and \(\gamma_{m,S}^{(n_f)}/\beta_S^{(n_f)}\) integrable there.
Only on this change-of-variable domain let
\[
\mu\frac{d\alpha_s}{d\mu}
=
\beta_S^{(n_f)}(\alpha_s),
\qquad
\mu\frac{dm_q^S}{d\mu}
=
-\gamma_{m,S}^{(n_f)}(\alpha_s)m_q^S .
\]
Then
\[
\boxed{
m_q^S(\mu_2)
=
m_q^S(\mu_1)
\exp\left[
-\int_{\alpha_s(\mu_1)}^{\alpha_s(\mu_2)}
\frac{\gamma_{m,S}^{(n_f)}(a)}
{\beta_S^{(n_f)}(a)}\,da
\right].
}
\]
At a heavy threshold \(\mu_h\), adjacent effective theories require fixed
matching maps
\[
\alpha_s^{(n_f-1)}(\mu_h)
=
\zeta_\alpha\alpha_s^{(n_f)}(\mu_h),
\qquad
m_q^{S,(n_f-1)}(\mu_h)
=
\zeta_m m_q^{S,(n_f)}(\mu_h).
\]
\end{theorem}

\begin{proof}
On the stated nonzero-beta coupling path, divide the mass RGE by the
coupling RGE and integrate; threshold matching is the finite decoupling
map needed to identify amplitudes in adjacent effective theories.
If the coupling is stationary at a beta zero, the coupling-variable integral
need not exist even though the original mass RGE still evolves. On such a
path one must use the always-typed time-scale expression
\[
 m_q^S(\mu_2)=m_q^S(\mu_1)
 \exp\left[-\int_{\log\mu_1}^{\log\mu_2}
 \gamma_{m,S}^{(n_f)}(\alpha_s(e^t))\,dt\right],
\]
under its ordinary solution/integrability conditions.
\end{proof}

\begin{definition}[Quark comparison package]
\label{p6:def:qrg-package}
\(\mathsf{QRG}\) consists of:
\begin{enumerate}[label=(\roman*),leftmargin=2em]
\item a native mass-renormalization condition for each finite spurion;
\item a finite-to-continuum conversion into one named scheme \(S\), including
the start scale of every converted mass;
\item \(\beta_S\) and \(\gamma_{m,S}\) to a stated order, together with one
coupling boundary value \(\alpha_s^S(\mu_0)\) or the equivalent
\(\Lambda_{\rm QCD}^S\) datum;
\item active-flavor intervals, threshold scales, and
\(\zeta_\alpha,\zeta_m\);
\item one common target scale; and
\item a separate short-distance or pole conversion for the top coordinate.
\end{enumerate}
For a fixed package, Theorem~\ref{p6:thm:exact-qrg-transport} gives a unique
common-slice table to the declared truncation order.  Paper~V-B supplies the
dictionary, not a completed high-order transport.
\end{definition}

\begin{proposition}[Split packet is not yet a quark mass operator]
\label{p6:prop:split-packet-not-operator}
The finite data
\[
(6=3+3,\ 16=7+9,\ \kappa_u,\kappa_d)
\]
specify a carrier and a coefficient ledger.  Writing
\[
D_u=\diag(m_u,m_c,m_t),
\qquad
D_d=\diag(m_d,m_s,m_b)
\]
after the six numbers have been assigned does not derive those eigenvalues.
A native theorem would require an operator-valued map
\[
\mathcal S_q^{\rm split}
\longmapsto(H_u,H_d)
\longmapsto(\Spec H_u,\Spec H_d).
\]
Until such a map is derived, \(D_u,D_d\) are components of \(\mathsf F\).
\end{proposition}

\begin{corollary}[Withdrawal of the \(15:4\) comparison]
\label{p6:cor:withdraw-fifteen-four}
The conditional coordinates
\[
m_b=4.207890\,{\rm GeV},
\qquad
m_c=1.264774\,{\rm GeV}
\]
give
\[
\frac{m_b}{m_c}\simeq3.326988\ne\frac{15}{4}.
\]
They are moreover quoted at different native reporting points.  The former
\(15:4\) assertion is removed rather than treated as a mass theorem.
\end{corollary}

\subsection{Canonical physical-source tests and the harmonic gap ratio}
\label{p6:sec:v2210-lit-physical-source-test}
The charged mass/incidence compiler remains closed.  The native coefficients
are not redefined through a later SVD.  The following results test the
physical response of its already occurring source and make the new harmonic
proposal testable without reviving its obsolete \(D_4\) mass word.

\begin{theorem}[Canonical kinetic mass response and source covariance]
\label{p6:thm:v2210-lit-kinetic-source}
On an actual positive quadratic chiral kinetic chart let
\(Z_L>0,Z_R>0\) be the generated left/right metrics and let \(R\) be the
already compiled chiral source.  Set
\[
 M=Z_L^{-1/2}RZ_R^{-1/2},\qquad
 h_M=\begin{pmatrix}0&M\\M^\dagger&0\end{pmatrix},\qquad B=MM^\dagger.
\]
For an assigned canonically normalized left source \(v\), the mass-only
resolvent response is, on its resolvent domain,
\[
 \boxed{\langle(v,0),(z-h_M)^{-1}(v,0)\rangle
 =z\langle v,(z^2-B)^{-1}v\rangle.}
\]
Under arbitrary invertible chiral coordinate changes \(S_L,S_R\), transform
\[
 Z_s'=S_s^\dagger Z_sS_s,\qquad R'=S_L^\dagger RS_R.
\]
Then \(M'=U_L^\dagger MU_R\), with
\[
 U_s=Z_s^{1/2}S_s(Z_s')^{-1/2}
\]
unitary.  With \(v'=U_L^\dagger v\), the displayed response and all of its
visible spectral weights are unchanged.  A common positive action factor
\((Z_L,R,Z_R)\mapsto c(Z_L,R,Z_R)\) also leaves \(M\) unchanged.
\end{theorem}
\begin{proof}
Since \(h_M^2=\operatorname{diag}(B,M^\dagger M)\), multiply
\((z-h_M)^{-1}\) by \(z+h_M\) to obtain the left block
\(z(z^2-B)^{-1}\).  Direct multiplication gives
\(U_s^\dagger U_s=I\) and the stated transformation of \(M\).
Thus \(B'=U_L^\dagger BU_L\) and source contraction gives the same scalar
response.  The common \(c\) cancels between the two square roots and \(R\).
\end{proof}

\begin{corollary}[A tagged free coefficient is visible precisely on its source]
\label{p6:cor:v2210-lit-tagged-defect}
For \(\varepsilon>0\) and a unit source \(v\), the complete mass-only
response above has only the two poles \(\pm\varepsilon\) if and only if
\[
 \boxed{(B-\varepsilon^2I)v=0.}
\]
Both poles then have residue \(1/2\).  More generally, put
\(\mathcal D_v=\|(B-\varepsilon^2I)v\|^2\).  For every
\(\Delta>0\), the source spectral measure obeys
\[
 \nu_v\{\lambda:|\lambda-\varepsilon^2|\ge\Delta\}
 \le\min\{1,\mathcal D_v/\Delta^2\}.
\]
These are physical-response tests of a fixed coefficient, not its mass
genesis.  The actual interacting response additionally retains every
action-owned return, constraint and source term.
\end{corollary}
\begin{proof}
The positive spectral measure of \(B\) gives a sum of terms
\(z\nu_v(\{\lambda\})/(z^2-\lambda)\); positive weights cannot cancel
distinct poles.  Only \(\lambda=\varepsilon^2\) is visible exactly when
\((B-\varepsilon^2I)v=0\).  Then
\(z/(z^2-\varepsilon^2)=\tfrac12[(z-\varepsilon)^{-1}
+(z+\varepsilon)^{-1}]\).  Finally
\(\mathcal D_v=\int(\lambda-\varepsilon^2)^2d\nu_v(\lambda)\)
implies the bound by positivity.  This also recovers the existing electron
defect bound on its stated source.
\end{proof}

\begin{proposition}[Source-visible interacting Schur root]
\label{p6:prop:v2210-lit-schur-source}
Let \(H=H^\dagger\) be the actual finite physical Hamiltonian on a certified
constraint/source domain, let \(P+Q=I\), and let the observed source
\(s\ne0\) lie in \(P\mathcal H\).  For real \(\lambda\) such that
\(\lambda-QHQ\) is invertible, set
\[
 S_P(\lambda)=\lambda-PHP
 -PHQ(\lambda-QHQ)^{-1}QHP.
\]
Then \(\lambda\) is visible in the source resolvent
\(\langle s,(z-H)^{-1}s\rangle\) if and only if some
\(u\in\ker S_P(\lambda)\) has \(\langle s,u\rangle\ne0\).
\end{proposition}
\begin{proof}
Block elimination gives the bijection
\[
 u\longmapsto
 \psi=u+(\lambda-QHQ)^{-1}QHPu
\]
between the Schur kernel and the full \(\lambda\)-eigenspace.
Since \(s=Ps\), \(\langle s,\psi\rangle=\langle s,u\rangle\).
The source resolvent has residue
\(\|\mathbf1_{\{\lambda\}}(H)s\|^2\), which is nonzero precisely when
that eigenspace is not orthogonal to \(s\).  This finite elimination is the
corresponding sharp instance of Feshbach--Schur reduction
\cite{LitSMFeshbach2007}; no smooth-map analytic domain is needed for these
finite invertible blocks.
\end{proof}
If the complement is singular, the full eigenspace/source test must be used
instead.  A determinant zero without source visibility is insufficient.
For a state-dependent many-body correlator use the existing complete marked
support polynomial and its transition weights, rather than substituting
this one-source resolvent for that differently typed observable.  Named rest,
pole-continuation, charged infrared and limiting identifications remain
claim-scoped.  The actual hierarchy arrows and completed roots are still
unevaluated; none is chosen by forcing \(S_P(\varepsilon_r)=0\).

\begin{theorem}[Exact harmonic row-gap ratio on its physical phase domain]
\label{p6:thm:v2210-lit-harmonic-ratio}
Suppose the independently licensed harmonic constructor of
\ref{p6:sec:v2208-neutral-rank-word} supplies the complete mass/incidence
packet with row magnitudes
\((n^{-1/2},\sqrt k,\sqrt n)\), \(1/n<k<n\), and one common
positive energy scale \(E\).  If the actually consumed rest/pole readout is
\(m=(E/\sqrt n,E\sqrt k,E\sqrt n)\) and propagation has the common
quadratic mass phase, then
\[
 \boxed{\frac{\Delta m_{21}^2}{\Delta m_{31}^2}
 =\frac{kn-1}{n^2-1},\qquad
 \frac{\Delta m_{21}^2}{\Delta m_{32}^2}
 =\frac{kn-1}{n(n-k)}.}
\]
For the proposed \((n,k)=(163,5)\), the first ratio is
\[
 \boxed{\frac{407}{13284}=0.03063836194\ldots.}
\]
It is invariant under a common rescaling of \(E\).  It is not an
unconditional propagation theorem or a derivation of the harmonic
constructor.
\end{theorem}
\begin{proof}
Squaring the three entries gives
\(E^2/n,kE^2,nE^2\).  Their successive differences yield
\(\Delta m_{21}^2=E^2(kn-1)/n\),
\(\Delta m_{31}^2=E^2(n^2-1)/n\) and
\(\Delta m_{32}^2=E^2(n-k)\).  Cancel \(E^2\) and substitute the ranks.
The additional phase/rest assumptions are exactly what identifies those
differences with measured oscillation frequencies.
\end{proof}

% BEGIN FRAMEWORK v2.245: NEUTRAL ACTUAL THREE-INVARIANT TEST
\begin{proposition}[Three invariant tests for the conditional harmonic neutral spectrum]
\label{p6:prop:v245-neutral-three-invariant-action-test}
Let $H_\nu\succeq0$ be the \emph{actual} canonically normalized
$3\times3$ neutral mass-squared operator obtained from the complete
source/current/kinetic action word, and let $E_{56}>0$ be its
independently generated common scale. Put $A_\nu=H_\nu/E_{56}^2$
and $e_2(A)=((\Tr A)^2-\Tr A^2)/2$. The action-derived squared
spectrum is exactly $\{1/163,5,163\}$ if and only if
\[
 \boxed{\Tr A_\nu=168+\frac1{163},\quad
 e_2(A_\nu)=816+\frac5{163},\quad
 \det A_\nu=5.}
\]
On this typed action domain these three equalities certify the
\emph{mass-only} rank word and its ratio $407/13284$. They do
not select $H_\nu$ from the harmonic projection ranks, nor identify
its eigenvalues with interacting physical poles or propagation phases.
The unreduced $168$-owner cross return of
Corollary~\ref{p6:cor:v239-neutral-green-port-test} must be included
before any justified reduction to the $163$-owner.
\end{proposition}
\begin{proof}
The three conditions are precisely the coefficients of
$\det(xI-A_\nu)=(x-1/163)(x-5)(x-163)$ by the
characteristic polynomial of a Hermitian three-dimensional operator.
Hermitian spectral calculus proves the two directions. The row ratio
is then the already established conditional identity in
Theorem~\ref{p6:thm:v2210-lit-harmonic-ratio}.
\end{proof}
% END FRAMEWORK v2.245: NEUTRAL ACTUAL THREE-INVARIANT TEST

\begin{remark}[Named experimental test, not a selected neutral constructor]
\label{p6:rem:v2210-lit-harmonic-observation}
At the printed \(E_{56}=3.924776812825316\,\mathrm{meV}\), the conditional
entries are
\[
 m=(0.307412245,
       8.776067750,
       50.108195976)\,\mathrm{meV},
\]
and the conditional gaps are
\[
 \Delta m_{21}^2=7.692486286\times10^{-5}\,\mathrm{eV}^2,
 \qquad
 \Delta m_{31}^2=2.510736802\times10^{-3}\,\mathrm{eV}^2.
\]
The named JUNO preprint submitted on 29 September 2026 reports
\(\Delta m_{21}^2=(7.388\pm0.078)\times10^{-5}\,\mathrm{eV}^2\) and
\(\Delta m_{31}^2=2.509^{+0.027}_{-0.025}\times10^{-3}\,\mathrm{eV}^2\)
for its normal-ordering fit \cite{LitSMJUNO2026}.  The solar difference is
about \(3.90\) times its quoted experimental error under the direct
identification above, with no theoretical-scale error or joint likelihood
included.  This is a conditional tension rather than a full model exclusion.
Common scale matching cannot change the exact ratio; an actual different
phase/readout map must be derived forward and tested jointly.  The older
\(D_4\) mass word, a new fitted angle or a K16 pairing condition is not an
antecedent of the new harmonic physical-origin calculation.
\end{remark}

\paragraph{Current local obligation ledger.}
The source-covariant response tests above narrow charged G16 to the actual
kinetic/source/constraint chart and completed visible root calculation.  The
quark comparison additionally requires the native-to-named-scheme conversion
and start scale already stated in \ref{p6:sec:qrg}; high-order running and
decoupling algorithms \cite{LitSMQuarkRunning2017} operate only after that
conversion.  They do not choose it.  The harmonic neutral origin retains
\texttt{OPEN\_PHYSICAL\_ORIGIN} independently of its conditional ratio
test.  Its three-role incidence and physical detector map are still required;
the old neutral incidence theorem is not automatically transferred to a
constructor which replaces its mass word.

All charged support/compiler statements are preserved, including
\[
 24=8+16,\quad16=7+9,\quad
 24=6+6+12,\quad12=4+4+4,\quad6=5+1,\quad5=3+2,
\]
the bottom unital \(+1\), Higgs \(17=16+1\), and the already generated
charged \(24/25\) polar degrees.  These are finite ancestry statements;
their physical source normalization is tested by the same-action canonical
map, not by an independently fitted mass multiplier.
% END LITERATURE-VERIFIED SOURCE AND HARMONIC TESTS v2.210



\subsection{Carrier-four gravitational coefficient port (framework v2.219)}
\label{p5b:sec:v2219-next-carrier-port}
The physical ultraviolet Laurent coefficient is indexed by the first
\emph{constituted} nonzero active carrier depth:
\[
 \operatorname{spec}\mathcal Q_{\rm act}=\{0,4,5,7\},\quad
 P_2=P_3=0,\quad P_4\ne0,\quad
 \boxed{j_{\rm dress}=4,\quad R_{B,4}=C_2=E_{32}^{\,2}.}
\]
This exports a carrier \emph{selection rule}, not an additional
Standard-Model particle. $E_{32}$ remains the existing Hurwitz $2T$
energy anchor. The later positive-part freeze $C_4$ is not a new
finite carrier at a formal $E_{64}$.  The gravitational response coupling
from this existing grade-four projector to the scalaron/particle Schur
vertex is evaluated, if it occurs, by the same parent action; equality
of the squared anchors does not itself determine a matrix element.


% BEGIN FRAMEWORK v2.224: particle response and geometric source obligations
\subsection{Particle response after the geometric source comparison}
\label{p5b:sec:v224-source-response-status}
Paper~V-A, Section~\ref{p5:sec:v224-geometric-source-comparison}, links
the established particle carrier to the rooted antipodal icosian
comparison and the retained primitive readout. These strengthen the
geometric organization without changing any terminal mass formula.
The hierarchy--family isomorphism and charged assembly
\[
 R_s^{\rm ch}=\sum_r\varepsilon_{s,r}\,\iota_{s,r}
\]
retain their declared native source arrows and domains; the label
$R_s^{\rm ch}$ here is a charged-source symbol, not the gravity cap
curvature scale $R_s$. Physical response tests consume the actual arrows,
kinetic and constraint chart, source-visible roots and residues, and
propagation or detector maps. These downstream tests do not redefine
terminal mass genesis as a subsequent matrix diagonalization.

The exact comparison-frame factor $24/25$ and its exterior descendants
must retain their physical-owner qualifications. No new assignment of
the nine grades or their chiral source contractions has been evaluated
in this integration. The established electron-return test remains a
concrete forward calculation; no additional observed particle pole is
declared closed by the cap source identity.

The finite leptonic quartet remains distinct from an observed mixing
invariant, with the phase and detector maps already typed in Paper~V-C.
A proposed harmonic neutral mass word still needs its own constructor,
three-role incidence and physical response. An incidence theorem for
an earlier $D_4$ mass word is not inherited merely by replacing that word.
Neither the 48-dimensional comparison nor the gravity source norm supplies
those maps. This update preserves the existing charged and neutral
predictions on their original premises and introduces no new physical
mass or mixing closure.
% END FRAMEWORK v2.224: particle response and geometric source obligations


% BEGIN FRAMEWORK v2.228: CHARGED ORIGIN AND CHARACTER TESTS
\subsection[Mass-source arrow availability and free-response nonuniqueness]{Mass-source arrow availability and\\ free-response nonuniqueness}
\label{p6:app:v2191-electron-arrow-limits}
This appendix supplies supplementary mathematical diagnostics for the
actual mass-source execution. Its comparison operators do not constitute
alternative occurring actions. The current physical source remains the
owner-resolved charged output of Paper V-A, and the current-relation sort
remains distinct from its physical mass-source value.

\begin{proposition}[Rooted parent-algebra availability]
\label{p6:prop:v2191-parent-algebra}
For $c=(1,1,1)^T/\sqrt3$, $Q_0=cc^\dagger$ and
$R_-=\operatorname{diag}(1,i,-i)$, the unital mathematical algebra generated
by $Q_0,R_-$ is $M_3(\mathbb C)$. In contrast the reduct
$\alpha Q_0+\beta Q_\perp$ has zero mixed corners. Moreover
\[
 \|[Q_0,R_-]\|_{\rm HS}^2=16/9.
\]
\end{proposition}
\begin{proof}
The three distinct eigenvalues of $R_-$ give each coordinate projector
$E_j$ as a polynomial in $R_-$. Then
$3E_jQ_0E_k=|j\rangle\langle k|$, so every matrix unit is in the
generated algebra. The mixed reduct corners vanish by $Q_0Q_\perp=0$.
For the commutator, unitarity and rank one give
$\|[Q_0,R_-]\|_{\rm HS}^2=2(1-|c^\dagger R_-c|^2)$, and
$c^\dagger R_-c=1/3$.
\end{proof}

These facts establish host availability and a nonzero current-relation
block; neither fact supplies a retained typed physical constructor word
for the occurring mass-source arrows. Mere block membership is also not
such a word. The source occurrence already licensed by Paper V-A is not
reopened by this distinction.

\begin{proposition}[Free left response does not reconstruct right arrows]
\label{p6:prop:v2191-right-arrow-nonuniqueness}
Let $R$ be a full-rank mathematical source with the three hierarchy corners
of Theorem~\ref{p3f:thm:v2168-hierarchy-source-return}. For every unitary
$W_R$ commuting with $Q_0$ and equal to the identity on $Q_0G$,
$R'=RW_R$ has the same corner types and
\[
 R'R'^\dagger=RR^\dagger.
\]
Thus all free left-current responses and poles agree. A nonidentity
$W_R$ changes $R$, and generally its directed right arrows.
\end{proposition}
\begin{proof}
Right multiplication by the unitary preserves the left return exactly.
Commutation with $Q_0$ preserves the two right support planes and hence
the three corner types. Since $R$ is invertible, $RW_R=R$ implies
$W_R=I$. The free left response depends only on $RR^\dagger$ by
Theorem~\ref{p3f:thm:v2168-left-return-test}.
\end{proof}

This comparison prevents reconstructing a unique physical source word
from free masses alone. It does not authorize varying the actual source.
An equivalent physical frame change must transform every consumed source,
kinetic, current and interaction block together. Conversely, the
electron-only defect tests a weaker observable than right-arrow equality.

Historical fourth-order mass readouts, the contemporary native terminal
tuple and path/common-column inverses keep their distinct existing types
and scale domains. Numerical agreement in one entry cannot derive a
universal factor or an interacting physical pole shift. The original
conversation proposal is preserved as evidence; its historical scale,
right-arrow rival and interaction-shift comparisons are not adopted as
new physical predictions. No $1/\sqrt6$ or $2\,\mathrm{GeV}$ electron-input
correction is introduced by the present integration.

% END INLINE SOURCE: Paper5B-ElectronArrowLimits-Appendix-v1_0.tex


\subsection{finite rows are not interacting poles}
\label{p6:sec:v164-finite-row-pole-boundary}

\begin{proposition}[Lookup/resolvent type boundary]
\label{p6:prop:v164-finite-row-not-pole}
A row energy or singular value of the finite spectrum tables is a finite
operator readout on its declared carrier.  An interacting pole instead
belongs to the meromorphic continuation of
\[
 G_R(z)=\bigl(K_R(z)-\Sigma_R(z)\bigr)^{-1},
 \qquad K_R=\delta^2S_R.
\]
Changing the generated self-energy \(\Sigma_R\) preserves the same finite row
while moving or dissolving the pole.  Hence no general theorem can equate a
finite row with a physical pole.  The tables may seed a quadratic comparison;
the pole position and residue must be computed from the same-action two-point
operator.
\end{proposition}

\begin{corollary}[Physical pole is not an independent mass selector]
\label{p6:cor:v577-pole-owner-direction-reversal}
Let two admissible renormalized presentations of the same two-point owner be
related on one common analytic pole chart by
\[
 \Gamma_{R'}^{(2)}(p)=A(p)\Gamma_R^{(2)}(p)B(p),
\]
where \(A,B\) are analytic and invertible on that chart.  Then
\[
 \det\Gamma_{R'}^{(2)}
 =\det A\,\det B\,\det\Gamma_R^{(2)},
\]
so the isolated zero set, with multiplicity on the fixed sheet, is unchanged.
Accordingly G16 contains no second physical mass selector: scheme choice is a
representation/calibration coordinate, while analytic continuation, sheet and
domain control, pole isolation, residue normalization, matching, and continuum
passage are claim-scoped analytic certificates for the same renormalized
two-point owner.
\end{corollary}

\begin{proof}
Analytic invertibility implies \(\det A\neq0\) and \(\det B\neq0\) on the
chart, so multiplication by these factors cannot create, remove, or move an
isolated zero of \(\det\Gamma_R^{(2)}\).  Residue representatives transform
covariantly and require the separately declared channel normalization.
\end{proof}

\subsection{Geometry-first consequences and physical scope (framework v2.201)}
\label{p6:sec:v2201-geometry-first-consequences}
\paragraph{Geometry ancestry.}
Hurwitz root and signed triality $\to$ hierarchy/family isometry $\to$ charged sibling packet and actual source arrows $\to$ tagged-pole tests.
\paragraph{Status.}
These are consequences on the domains stated below.
A mathematical certificate or an executed restriction retains every
physical source, state, clock and readout hypothesis of its inputs.

\subsubsection{D05: Complete normal form of the charged free-pole certificate}
\label{p6:sec:v2201-derived-d05}
\noindent\textit{Status: proved result on stated domain.}

Write $e=\varepsilon_e$, $m=\varepsilon_\mu$, $t=\varepsilon_\tau$ and use
the native hierarchy coordinates
\[
 \widetilde R=\begin{pmatrix}eA&mb\\td^\dagger&0\end{pmatrix},\quad
 A\in M_2(\mathbb C),\quad b,d\in\mathbb C^2.
\]
The full tagged certificate is
\[
 \widetilde R\widetilde R^\dagger=\operatorname{diag}(e^2,m^2,t^2).
\]
It holds if and only if there are orthonormal $v,d\in\mathbb C^2$ and
$S\in U(2)$ with
\[
 \boxed{A=\binom{S_{11}}{(m/e)S_{21}}v^\dagger,\qquad
 b=\binom{(e/m)S_{12}}{S_{22}},\qquad \|d\|=1.}
\]
\begin{proof}
The cross corner gives $Ad=0$, and the bottom corner gives
$\|d\|=1$. Thus $A=av^\dagger$ for a unit vector $v\perp d$.
The top corner becomes $XX^\dagger=D^2$, with
$X=(ea,mb)$ and $D=\operatorname{diag}(e,m)$. Hence $X=DS$ for a unitary
$S$, giving the formula. Conversely direct multiplication gives all three
corners. Positivity of $D^2$ forces $A$ to have rank one and $b$ to have an
independent left image. The parameters have the usual phase redundancy.
\end{proof}
This classifies \emph{evaluation outputs that pass the test}; it does not
select any of them as the actual physical arrows. Four support-forced
cross returns were already zero in the source. The surviving
electron--tau return is now explicitly $Ad$.
\par\noindent\textit{Inherited inputs:} \ref{p1b:thm:v138-hierarchy-family-block-isomorphism}, \ref{p6:thm:v590-charged-channel-representatives}, \ref{p5:thm:v397-charged-preterminal-srcval-realization}, \ref{p3f:thm:v2168-hierarchy-source-return}.

\subsubsection{D06: Unit source arrows would force left-current alignment}
\label{p6:sec:v2201-derived-d06}
\noindent\textit{Status: proved result on stated domain.}

Assume the \hyperref[p6:sec:v2201-derived-d05]{D05} certificate and the additional independently established
normalization $\|A\|_{\rm HS}=1$, with $e\ne m$. Then
\[
 \|A\|_{\rm HS}^2
 =1+\left(\frac{m^2}{e^2}-1\right)|S_{21}|^2
\]
forces $S_{21}=S_{12}=0$. Consequently
\[
 A=e^{i\theta}e_1v^\dagger,\qquad b=e^{i\varphi}e_2,
 \qquad Ad=0,\quad v\perp d.
\]
Thus the normalized electron and muon images would be exactly their
generated current lines, up to phases. A right-plane unitary can remain;
the left free-pole test never needed one preferred right frame.
\begin{proof}
Insert the first column of the unitary $S$ into the norm formula.
Unitarity gives $|S_{11}|^2=1-|S_{21}|^2$ and makes $S_{12}$ vanish with
$S_{21}$.
\end{proof}
The missed opportunity is to prove this normalization on the \emph{actual}
hierarchy arrows. The explicitly available unit angular arrows do satisfy
it, but the source text expressly distinguishes them from those arrows.
No numerical electron defect is asserted without that evaluation.
\par\noindent\textit{Inherited inputs:} \ref{p6:prop:v2168-angular-hierarchy}, \ref{p3f:cor:v2191-action-corner-inverse}, \ref{p6:cor:v2191-electron-defect}.

\subsubsection{A2: a hierarchy-protection bound on the actual electron arrows}
\label{p6:sec:v2201-derived-a2}
\noindent\textit{Status: proved corollary on stated domain.}

On the existing mass-only source chart, let
\(e=\varepsilon_e\), \(m=\varepsilon_\mu\), \(t=\varepsilon_\tau>0\).
The left return has
\[
B_{11}=e^2\|A_{1\bullet}\|^2+m^2|b_1|^2,
\qquad B_{31}=et\overline{(Ad)_1}.
\]
If the exact electron-only two-pole certificate \(Be_1=e^2e_1\) holds, then
\[
\boxed{|b_1|\le e/m,\qquad \|A_{1\bullet}\|\le1,\qquad(Ad)_1=0.}
\]
For the current REG0 coefficients,
\[
\frac em=\frac{2}{1+\sqrt5}\,
\alpha_{\rm hier}^{-1/8}\left(\frac{24}{25}\right)^{-5/2}
=0.0048504178\ldots.
\]
\textbf{Proof.} Positivity of the two terms in \(B_{11}=e^2\) bounds each
separately. The third component of \(Be_1=e^2e_1\) gives the exact tau-return
cancellation since \(et\ne0\). Substitution of the already compiled
coefficients gives the ratio. These are necessary conditions, not a
numerical evaluation or a freely chosen small family angle.
\par\noindent\textit{Inherited inputs:} \ref{p6:thm:v2189-hierarchy-mass-moments}, \ref{p6:cor:v2191-electron-defect}.

This exposes what hierarchy protection would require from the actual
source arrow. If \(b_1\) is structurally zero, that fact should be proved
on its parent relation. If it is nonzero, its bound and the two mixed-return
contractions are stringent execution targets.

\subsubsection{A3: a quantitative source-visible spectral bound}
\label{p6:sec:v2201-derived-a3}
\noindent\textit{Status: proved corollary on stated domain.}

The paper already defines
\(\mathcal D_e=\|(B-e^2I)e_1\|^2\). Let \(\nu_B\) be the normalized
spectral measure of \(B\) on \(e_1\). For every squared-mass window
\(\Delta>0\),
\[
\boxed{\nu_B\{\lambda:|\lambda-e^2|\ge\Delta\}
 \le\min\{1,\mathcal D_e/\Delta^2\}.}
\]
For the mass-only chiral Hamiltonian the corresponding source weight outside
\(|z^2-e^2|<\Delta\) obeys the same bound.

\textbf{Proof.} The finite spectral theorem gives
\(\mathcal D_e=\int(\lambda-e^2)^2d\nu_B(\lambda)\). On the indicated set
the integrand is at least \(\Delta^2\). The chiral block has
\(h_R^2\) left return \(B\), so its source measure pushes forward to
\(\nu_B\) under \(z\mapsto z^2\).
\par\noindent\textit{Inherited inputs:} \ref{p6:cor:v2191-electron-defect}, \ref{p6:cor:v2191-mass-cofactor}.

An exact zero defect proves the two-pole statement. A small nonzero defect
now gives a controlled source-weight statement instead of an all-or-nothing
closure label. It does not establish the interacting pole or its physical
rest interpretation; those still use the complete source word and pairing.

% END GEOMETRY-FIRST CONSEQUENCES v2.201

% BEGIN FRAMEWORK v2.226: KINEMATIC CLOSURES
\subsection{Native coefficient completeness and physical pole ownership}
\label{p5b:sec:v226-pole-coefficients}
The coordinate-only algebra, its separate chirality projectors and its
complete finite coefficient tests propagate to any actual intrinsic
momentum corner used by the particle comparison. Equalities are statements
about those full actual corner matrices, rather than their dimensions.

\begin{corollary}[A supplied native pole matrix has complete finite tests]
\label{p5b:cor:v226-pole-coefficients}
For a supplied complete native kinetic/pole matrix with its physical source
and clock dictionary, the 137 dual coefficients determine that matrix and
all admitted source derivatives. On a common positive kinetic normalization,
coefficient errors give the spectral bounds of
Proposition~\ref{p2d:prop:v226-coefficient-spectral-error}.
\end{corollary}
\begin{proof}
Apply Theorem~\ref{p1b:thm:v226-dual-coefficients} and its source-jet
corollary to the supplied matrix, followed by the cited kinetic-pair bound.
If the pole is a reduced response, first retain the complete Schur source,
contacts and residues of Theorem~\ref{p3f:thm:v226-response-reconstruction}.
\end{proof}
The fixed upstream terminal coefficients, mass formulas and comparison-owner
assignments are preserved. A coordinate-polynomial presentation does not
supply the still-unexhibited complete electron member or its source-cyclic
quotient, occupation embeddings, response roots, detector map or mixing
invariants. The actual required matrices must precede these finite tests.

% END FRAMEWORK v2.226: KINEMATIC CLOSURES

% BEGIN FRAMEWORK v2.227: DOWNWARD CLOSURES
\subsection{Complete-source pole comparison without a unique dark completion}
\label{p5b:sec:v227-pole-comparison}
\begin{corollary}[A complete finite response fixes the same-source pole]
\label{p5b:cor:v227-pole-comparison}
For two licensed Hermitian parents with one actual retained/source frame,
equal complete response matrices at the upper-half-plane samples of
Theorem~\ref{p1b:thm:v227-finite-response} give the same admitted retained
Green response and hence the same isolated same-source poles and residues.
A different uncoupled dark extension is immaterial to this claim.
\end{corollary}
\begin{proof}
The finite tests give equality as meromorphic matrix functions. Invert
on the common regular domain and take the same contour residues.
The minimal visible environment is fixed up to unitary equivalence by
Theorem~\ref{p1b:thm:response-minimal-visible-realization}.
\end{proof}
This augments the existing common-source cofactor and moment routes. It
requires actual source normalization, support, clock and source-port data;
a zero overlap still removes a pole from that source. The presently
uncomputed complete charged external/dual/remainder returns therefore
remain execution obligations. No new electron or other particle mass is
inferred from the finite sampling bound, and odd-source poles do not
inherit the neutral central-offset quotient.

% END FRAMEWORK v2.227: DOWNWARD CLOSURES


% BEGIN FRAMEWORK v2.229: ROOTED CHARGED CONVERSES AND COMPARISON SCOPE

\paragraph{Preserved extended mass-source obstruction.}
The complete rooted candidate and its offset-only obstruction,
with all original positive and negative certificates, is given in
Paper~V-D, Section~\ref{p5d:sec:D}. Its negative result
limits physical pole identification and is not superseded by a
numerical source-table match.

\section{Exact results, admissible comparisons and remaining ports}
\label{p5b:sec:organized-results}
\subsection{Integrated closure theorem}
\label{p6:sec:integrated-theorem}

\begin{theorem}[Full Paper--V-B reconstruction]
\label{p6:thm:full-reconstruction}
The revised manuscript has the following exhaustive status.
\begin{enumerate}[label=(\roman*),leftmargin=2em]
\item From \(\Ubase\), the corrected finite core \(\CoreSix\) contains:
the six structural and eighteen family projectors; the rank-five coset
no-go; the \(A_2\) trace identities; the fact that the nearest \(2I\) shell
generates \(2I\); the five conjugate binary tetrahedral subgroups; the
twenty-five-cell incidence and singular spectrum; the classification of
equivariant terminal and coherent-boundary maps; the Hurwitz walk
rank-fifteen projector; the binary-octahedral normalizer and quotient-odd
rank-three carrier; and every finite matrix identity whose hypotheses lie
entirely in \(\Ubase\).

\item From \(\Ubase\) and the compatible structural/readout subtuple
\(b_{\rm str}\), enlarged to \(\Uabs\) for absolute coordinates, the paper
constructs one terminal lookup operator and self-adjoint dilation plus one
canonical scalar readout coordinate.  Node magnitudes follow from the one
atomicity clause by the first-minimum theorem; comparison factorization through
\(n_{\etaH}\) and the invariant scalar embedding/state are theorems rather than
extra members of this subtuple.

\item From \(\Urole=\Ubase\), a compatible full member of
\(\BridgeRaw\), the anchored \(\mathrm{REG}_0\) charged-root
member (or a separately root-certified \(\mathrm{REG}_\sigma\) member),
\((\mathrm{CRP4})\), and \(\mathrm{CAL}_{E_0}\) wherever a named absolute
number is shown, the paper constructs one internally consistent conditional
physical-role model: one retained-order charged branch, one \emph{historical comparison}
positive radical neutrino operator, one standard CKM representative, one
standard PMNS representative, and one affine heavy-family operator.
The current K56 harmonic rank proposal has no certified replacement
physical neutral operator; none of the comparison matrices selects it.  A Majorana operator
occurrence and \(\mathsf{SLOT}_{\rm phys}\) remain separately explicit when
their corresponding claims are made.

\item The unrestricted equivariant terminal comparison class has a
nontrivial singular-value modulus. Non-singleton structural or physical-role
\emph{admissible} quotients require the additional compatible membership
and role-extension certificates of
Theorem~\ref{p6:thm:positive-dimensional-obstruction}. Additional
clock-to-role placement, scalar-coordinate, family-eigenframe, heavy-family,
REG finite-part, and coherent-boundary counterfamilies remain, but they are not
premises of the non-singleton theorem.  The interpolation function itself is
already fixed by the derived shell law.

\item Physical poles, effective vertices, QCD common-slice comparisons,
nonzero widths, scalaron matching, reheating histories, and thermal
heavy-scale placement require their downstream owners.  None is used as a
premise of the finite Paper--V-B construction.
\end{enumerate}
\end{theorem}

\begin{proof}
Part~\ref{p6:part:exact-core} proves the structural, trace, group, incidence,
projector, and representation-domain statements.  Part~\ref{p6:part:conditional-model}
proves the terminal Hom-lattice minimum, comparison-output quotient, canonical
scalar embedding/state and field-rescaling quotient, local no-go theorems,
selected charged and neutral operators, and exact matrix representatives.
Part~\ref{p6:part:interfaces}
proves the heavy-carrier, scheme-covariance, QRG-dictionary, no-width, and
sufficiency statements.  The positive-dimensional counterfamily is
Theorem~\ref{p6:thm:positive-dimensional-obstruction}; conditional
single-valuedness is
Theorem~\ref{p6:thm:conditional-finite-sufficiency}.
\end{proof}

\begin{proposition}[Historical coherent comparison member]
\label{p6:prop:selected-coherent-branch}
On \(\Urole\), the following explicit member of the historical optional atlas
is internally compatible.  It is a comparison witness, not the
premise-minimal physical branch.  Source coordinates are preserved only where
they remain mutually compatible.
\begin{enumerate}[leftmargin=2em]
\item Use the G0-occurring conjugate \(2T\subset2I\) supplied by the derived
\(q=4\) active realization, together with atomic first-minimum node transfers,
common-domain columns, positive codiagonals, and aligned path phases.
\item Use the orthogonal cell-label carrier with
\(\etaH=24/25\) and the declared minimized comparison-exponent ledger.
\item Use the derived active anchors and the selected residual clock data
\[
\begin{aligned}
2T&\mapsto K=32,& C_6&\mapsto K=40,\\
\mathbb R\mathbf1&\mapsto K=56,& E_{36}&=\sqrt{E_{32}E_{40}},\\
2O/2T&\mapsto K=16&&\text{(historical optional association)}.
\end{aligned}
\]
while making no role assignments at \(K=8,24\); the absence of a native
active carrier at \(K=48\) is already a theorem.
\item Use the printed historical scalar comparison coordinate
\[
\kappa_{\rm print}=2\pi\sqrt{\frac{24}{15}},
\]
while scalar stationary/effective-action branches remain claim-scoped in Paper~V-C and are not required for native electroweak reduction.
\item Select
\[
\Delta_L=0.125314767527813\ldots
\]
as the isolated root of the explicit \([4]\)-model, and use only the
inverse-shell map of Definition~\ref{p6:def:charged-absolute-map} for its
retained-order charged coordinates.  These are the values disclosed in
Corollary~\ref{p6:cor:actual-charged-output}, not the diagnostic table rows.
\item Use
\[
M_{\nu,+}
=
\sqrt{\etaH}E_{56}
\diag(1/\sqrt8,\sqrt5,13)
\]
and remove the two incompatible light-neutrino mass claims.
\item Use the exact triangle-angle CKM branch
\[
\begin{aligned}
s_{13}&=0.003794424824531,&
\delta_{\rm std}&=64.2021006467^\circ,\\
\rho&=0.176003116174467,&
\eta_q&=0.364114000362982,
\end{aligned}
\]
the retained PMNS packet, and the positive-Majorana operator basis of
Proposition~\ref{p6:prop:positive-majorana-frame}.
\item Retain the heavy odd carrier exactly and the affine
\((E_N,\Delta_M,\widehat D_N)\) operator conditionally, with
\(E_N,\Delta_M\) symbolic.
\end{enumerate}
\end{proposition}

\begin{remark}[Downstream comparison appendices]
\label{p6:rem:selected-downstream-comparisons}
The historical optional comparison member is not enlarged by the following
instructions.  The finite pole packet may use \(b=1/120\), ranks
\((1,7,1)\), and the upper finite branch as
\(\mathsf{PV}_{\rm mat}\), but not as a vertex or physical-pole theorem.
Every quark comparison must first receive \(\mathsf{QRG}\), while scalar
potential, widths, scalaron matching, reheating, and thermal heavy placement
remain with the owners in Table~\ref{p6:tab:ownership}.
\end{remark}

\begin{corollary}[Nonempty historical comparison locus]
\label{p6:cor:nonempty-compatible-locus}
The compatible historical comparison locus is nonempty.
\end{corollary}

\begin{proof}
Every transfer and matrix in
Proposition~\ref{p6:prop:selected-coherent-branch} is explicitly defined on a
finite-dimensional Hilbert space.  Common-domain columns, tensor comparison
maps, self-adjoint doubling, positive spectral roots, and unitary
representatives exist by direct finite-dimensional construction.  The
charged branch is the explicitly defined \([4]\)-polynomial representative
with its isolated root interval.  Therefore the finite/raw tuple is
consistent.  This proves existence, not upstream selection or convergence
of an unknown exact charged completion and not physical selection.
\end{proof}

\subsection{Claim-by-claim source disposition}
\label{p6:sec:source-disposition}

\begin{longtable}{@{}L{0.29\textwidth}L{0.19\textwidth}
L{0.44\textwidth}@{}}
\caption{Disposition of the principal source claims.}
\label{p6:tab:source-disposition}\\
\toprule
Former claim or construction & Revision status & Integrated replacement \\
\midrule
\endfirsthead
\toprule
Former claim or construction & Revision status & Integrated replacement \\
\midrule
\endhead
\bottomrule
\endlastfoot

\(A_2\) normalized-trace Gram values \(2/3,-1/3\)
& corrected
& Those are ordinary-trace values; normalized trace gives \(2/9,-1/9\).
The defect eigenvalues themselves remain \(2/3,-1/3,-1/3\).\\

Nearest \(2I\) shell ``becomes \(2T\)''
& withdrawn
& The shell generates \(2I\).  The derived \(q=4\) active realization
supplies one G0-occurring conjugate \(2T\), or all five conjugates may be
treated invariantly.\\

Diagonal \(A_5\) transitivity on twenty-five cells
& corrected
& Diagonal action has \(5+20\) orbits; the left--right
\(A_5\times A_5\) action is transitive.\\

\(24/25\) as an intrinsic incidence amplitude
& retyped
& It is self-exclusion in the orthogonal label carrier.  Vertex incidence
instead gives \(4/5\).\\

Direct-sum singular value \(\sqrt{\sum_i|S_i|}\)
& corrected
& It belongs to a common-domain column map, not a direct sum with independent
domains.\\

\(7|9\) exchange-odd quark operator
& withdrawn
& The retained object is a weighted signed contrast with an explicit
exchange-even part.\\

Finite split rank implies exact truncation
& withdrawn
& Finite rank gives recurrence/resummation.  Exact termination needs a new
nilpotent or exterior grammar.\\

\(1+1+1+2=5\) proves particle completeness
& withdrawn
& The Paper--V-A matter carrier has a true six/eighteen-projector exhaustion;
the rank-five coset packet has only the canonical \(1\oplus4\) split without
extra selectors.\\

Terminal radicals are fixed by carrier size
& premise-reduced
& On the atomic integral-transfer branch, the unsigned node radical is the
first minimum of the invariant Hom-lattice.  Atomicity remains in
\(\mathsf{TMF}\); a chosen magnitude and node positivity do not.\\

\(K=32,40,56\) are selected clock anchors
& derived/retyped
& The active realization proves \((A^{\rm even},2T)\mapsto32\),
\([C_6\curvearrowright E_3]\mapsto40\), and
scalar reclosure \(\mapsto56\).  Particle names and mass operations remain
\(\mathsf{SLOT}+\mathsf{TMF}\) data.\\

\(K=36\) protected scale / interface role
& split status
& The scale \(E_{36}=\sqrt{E_{32}E_{40}}\) is derived from the continuous shell
law.  Assigning the interface packet to that coordinate is \(\CLKrem\) data.\\

\(K=16\) normalizer placement
& conditional ancestry
& The quotient and its canonical real rank-two multiplicity are exact; \(K=16\) still requires the typed \(\mathsf{ANC}_O\) map into the protected hierarchy owner.  Once that map is proved, the exponent/normalization is inherited rather than selected.\\

Three comparison copies are forced
& eliminated for current outputs
& Every native evaluation factors through one exponent
\(n_{\etaH}\in\tfrac12\ZZ\).  A three-factor word is an unidentifiable lift
unless a future factor-sensitive observable is supplied.\\

Rank-fifteen response fixes \(v\)
& retyped
& \(15/24\) is a maximally mixed trace, while the invariant state response is
one.  The invariant embedding projector/state and canonical field coordinate
are derived; the printed scalar coefficient is the remaining
\(\mathsf H_{\rm comp}\) comparison coordinate, not a live scalar selector.\\

Paper--V-B derivation of EWSB
& moved downstream
& Paper~V-C must derive the scalar potential and select a nonzero global
minimum independently.\\

Unique charged shell/Koide solution
& conditional
& Two transverse positive roots are certified for \(\mathrm{REG}_0\)
(and persist locally by transversality), not for all
\(\mathrm{REG}_\sigma\). The explicit \([4]\)-model selects the
first isolated interval; convergence of an unknown exact completion remains
open.\\

Three simultaneous light-neutrino spectra
& replaced
& The former positive radical \(56\)-line operator is a historical
comparison member; the active K56 harmonic-rank proposal has no certified
physical neutral-operator replacement.\\

Reality fixes Majorana phases to zero
& corrected
& Reality permits \(0\) or \(\pi\); positivity of the selected operator fixes
the all-positive branch.\\

Two CKM constructions represent one orbit
& corrected
& One exact standard convention and one triangle-angle branch are used.
The former inequivalent block is removed.\\

CKM/PMNS matrices predicted by carrier support
& conditional
& Explicit operator-first representatives realize supplied frames; arbitrary
relative eigenframes remain possible upstream.\\

Parity fixes CP-defect orders and coefficients
& withdrawn
& Orders must be calculated from explicit sector-operator series.\\

\(1/120\) and upper sign are unique pole data
& conditional
& They are one member of the coherent-boundary family
\(\mathsf{PV}_{\rm mat}\).\\

Mass matrix fixes \(\sin^2\theta_{\rm eff}^\ell\)
& withdrawn
& The effective angle is a vertex/form-factor observable requiring
interaction and renormalization data.\\

``No free coefficient'' proves predictivity
& replaced
& Predictivity is constancy of an observable on the admissible physical
operator quotient.\\

Heavy quotient fixes a heavy threshold
& corrected
& The quotient fixes an odd rank-three carrier.  \(x_N=1\) is projector
normalization; \(E_N,\Delta_M\) remain bridge data.\\

Mixed-scale quark table is one precision spectrum
& conditional
& Native coordinates require \(\mathsf{QRG}\) before common-slice
comparison; the \(15:4\) claim is removed.\\

Finite dynamics yields a universal scalaron width
& withdrawn
& Exact finite width is zero.  \(1/(192\pi)\) is one minimally coupled real
scalar partial width in a downstream continuum interaction theory.\\

Scalaron mass and reheating are Paper--V-B outputs
& moved downstream
& Paper~VI-A owns \(R+R^2\) matching and the conditional FLRW/reheating
evaluation.\\
\end{longtable}

\subsection{Conclusion}
\label{p6:sec:conclusion}

Paper~V-B now has a rigorous theorem boundary.  Its finite geometry is not
discarded: the corrected family plane, binary polyhedral groups, cell
incidence, terminal-map classifications, projector exhaustions, finite
matrix spectra, and quotient-odd carrier all survive in exact form.  What
does not survive is the inference that carrier symmetry alone fixes every
normalization, residual clock assignment, particle slot, family matrix, pole
residue, scheme, or rate.  Conversely, the active-realization update now
derives \((A^{\rm even},2T)\mapsto32\),
\([C_6\curvearrowright E_3]\mapsto40\), scalar reclosure
\(\mapsto56\), and the heavy/light finite split grammars; these are no longer
members of the bridge.  Scalar reclosure is the derived Haar projector, so an
independent infrared-extraction premise is absent as well.

The present repair further removes four redundant coordinates or coordinate
families.  Physical
representation names are licensed on \(\Urole=\Ubase\) by the upstream
structural theorem rather than a second residual; on the residual atomic integral branch,
transitive node radicals are the first minima of invariant Hom-lattices and
need no chosen magnitude or node positivity; every native comparison output
factors through one half-integer exponent; and the multiplicity-one invariant
walk line fixes the scalar embedding projector/state while field rescaling
removes an independent positive kinetic coefficient.  None of these
reductions selects a mass, mixing orbit, or dynamical vacuum.

For the energy hierarchy, the natural physical cone and the independent
contact/filling calculation now derive the positive fourth-density degree.
By themselves they do not close C13: the exact natural positive
$\lambda$-family retains every character weight.  The imported CRP4
occurrence supplies the missing common parent over the pullback base, so
Theorem~\ref{p6:thm:crp4-c13-discharge} derives \(\lambda_{ED}=1\) and closes
C13.  The counterfamily remains the minimality witness; a bare line
isomorphism or transported connection is still not counted as closure.

The reconstruction now separates exact finite owner/readout machinery from
nonhorizontal comparison and continuum interfaces.  The generic orbit
calculations remain useful for counting and for constructing invariant
readouts, but they no longer imply an unresolved family-member selector: the
primitive and charged noncentral source packets occur on the rooted branch.
The charged mass/incidence synthesis and native current evaluation are
generated descendants. Historical directed SVD matching remains only a
conditional comparison calculation; it is not a live charged family selector.

The synchronized logical statement is
\[
\boxed{
\Ubase\vdash\CoreSix,
\qquad
\Urole+(\mathrm{CRP4}_{\rm rooted}^{\rm der})+
\texttt{G16}_{\rm claim}
+(\mathrm{MOMENERGY}_{0}^{\rm der})
\vdash\mathfrak P_{5B}^{\rm current}
}
\]
with \(\mathrm{REG}_{0}\) fixed on the occurring anchored branch; the affine
\(\mathrm{REG}_{\sigma}\) family is retained only in the unanchored
comparison/scheme atlas, and pole certification is added only when required.  Hence
\[
\boxed{\begin{gathered}
N_{\rm exact\ core\ selector\ families}=0,\\
\texttt{CHARGED\_PRETERMINAL\_ASSEMBLY}=\texttt{CLOSED},\\
N_{\rm live\ Paper\text{-}V\text{-}B\ local\ physical\ selector\ families}=0,\\
N_{\rm nonhorizontal\ numerical\ scheme\ boundaries}^{\rm anchored}=0,\\
N_{\rm comparison\ affine\ scheme\ families}=1.
\end{gathered}}
\]
The former G5Q-LOOP item is not a second selector.  Its weighted-quiver,
spectral and orbit formulas are readout machinery attached to the inherited
actual member.  Scalar stationary/effective-action branches are likewise
claim-scoped and cannot re-enter native fermion-mass ownership.  Successful later spectrum/flavor or
observable branches may not be cited backwards to reselect those upstream
objects.  This is a complete conditional reconstruction with an exact finite
core, not a coefficient-free numerical fit and not an unconditional
parameter-free Standard-Model derivation.


% BEGIN TERMINAL-SUPPORT-OPERATORIZATION v5.85

% BEGIN FRAMEWORK v2.246: JOINT-GRAM ORIENTATION FIBRE
\begin{proposition}[Complete orientation fibre at fixed separate current Grams]
\label{p6:prop:v246-current-gram-orientation-fibre}
Fix $Z_u,Z_d>0$ on two three-dimensional role spaces. The cross Grams
$X$ for which
\[
 \begin{pmatrix}Z_u&X\\X^\dagger&Z_d\end{pmatrix}\ge0,
 \qquad Z_d-X^\dagger Z_u^{-1}X=0
\]
are exactly
\[
 \boxed{X=Z_u^{1/2}CZ_d^{1/2},\qquad C\in U(3).}
\]
Every such block Gram has rank three and is realized in a common
three-dimensional receiver by
$E_u=Z_u^{1/2}$, $E_d=CZ_d^{1/2}$. Thus the structural receiver,
positivity, separate normalization Grams and zero Schur defect leave
the full unitary current undetermined until the actual cross Gram is
calculated. For $Z_u=Z_d=I$, both $C=I$ and the Fourier current
$C_{jk}=3^{-1/2}\omega^{jk}$ ($j,k=0,1,2$, $\omega=e^{2\pi i/3}$)
satisfy all these conditions, while their quartet invariants include
$J=0$ and $|J|=1/(6\sqrt3)$, respectively.
\end{proposition}
\begin{proof}
Set $C=Z_u^{-1/2}XZ_d^{-1/2}$. The Schur equation is
$C^\dagger C=I$; since $C$ is square, it is unitary. Conversely
substitution gives the zero Schur complement and the stated Gram
factorization, whose first three columns have full rank. The identity
current has zero off-diagonal quartets. For the Fourier current the
quartet with rows and columns $0,1$ is $\omega/9$ or its conjugate,
depending on the chosen quartet convention; its imaginary magnitude is
$\sqrt3/18=1/(6\sqrt3)$.
\end{proof}
\begin{remark}[Algebraic fibre is not a new physical selector]
This proposition characterizes the information not contained in the
structural current certificates. It does not license freely varying
$E_u,E_d$ in the frozen rooted action. Their source-owned evaluation
selects the actual $X$; the point of the fibre is that separate Grams
and receiver rank cannot substitute for that calculation.
\end{remark}
% END FRAMEWORK v2.246: JOINT-GRAM ORIENTATION FIBRE


% Paper V-B, after p6:prop:v245-neutral-three-invariant-action-test.
\begin{proposition}[Two-channel Green traces do not identify a physical kinetic owner]
\label{p6:prop:v246-neutral-green-trace-nonuniqueness}
Let $Q=Q_{163}$, $P=P_{\rm aug}$ and $P_{\rm ns}=Q+P$ act on
$E_{168}$. There is a nonconstant family of positive, $S_3$-equivariant
operators $K_\epsilon$ whose inverse has exactly the two-channel trace
Gram of \eqref{p1b:eq:v239-nonsinglet-green-gram}.
Thus that Gram and its scalar Schur residual do not identify an
unspecified action-owned kinetic operator.
\end{proposition}
\begin{proof}
The $S_3$-fixed subspace of $E_{163}$ has dimension $28$ by
Proposition~\ref{p1b:prop:v240-k56-orthogonal-green-firewall}.
Choose two orthonormal fixed vectors $v,w$, and set
$T=vv^\dagger-ww^\dagger$. Then $T=T^\dagger$, $QT=TQ=T$,
$PT=TP=0$, $\Tr T=0$, and $T$ is equivariant. For
$|\epsilon|<1/168$, define
\[
 G_\epsilon=I_{E_{168}}/168+\epsilon T\succ0,
 \qquad K_\epsilon=G_\epsilon^{-1}.
\]
For $(P_1,P_2)=(P_{\rm ns},P)$, cyclicity of trace and the displayed
identities give
\[
 \bigl(\Tr(P_aG_\epsilon P_b)\bigr)_{a,b=1}^2
 =\begin{pmatrix}1&5/168\\5/168&5/168\end{pmatrix}.
\]
Its scalar Schur residual is $163/168$ for every $\epsilon$.
The two selected eigenvalues of $K_\epsilon$ are
$(1/168+\epsilon)^{-1}$ and $(1/168-\epsilon)^{-1}$,
so the operators differ for nonzero $\epsilon$.
\end{proof}
This is an insufficiency theorem for identifying an action from the
listed response data. It does not alter the defined geometric round
operator $K_{12}=168I$, nor assert that the comparison family is a
licensed physical constructor.



% BEGIN FRAMEWORK v2.321: PETERSEN SPECTRAL VISIBILITY AND ICOSIAN-SHEET COMPARISON
\section{Rooted Petersen sector visibility and terminal-spectrum boundary}
\label{p6:sec:v321-petersen-sector-visibility}

The induced edge space in Paper~V-A admits four pairwise orthogonal,
$A_5$-equivariant projectors $\Pi_4,\Pi_5,\Pi_3,\Pi_{3'}$ of ranks
$4,5,3,3$. For a marked tetrahedral core $H_i$, let $P_i$ project onto its
three internal oriented-edge directions, not onto a newly constituted
physical family or a level-one momentum multiplet.
All statements here are conditional on the occurring $2I$ support; the
generic $2T$ or $2O$ branch does not supply this fifteen-edge packet.

\begin{theorem}[Universal local-triplet sector weights]
\label{p6:thm:v321-local-petersen-sector-weights}
For all five cores, exact projector compression gives
\[
 \boxed{\begin{aligned}
 P_i\Pi_4P_i&=\frac4{15}P_i,&
 P_i\Pi_5P_i&=\frac13P_i,\\
 P_i\Pi_3P_i&=\frac15P_i,&
 P_i\Pi_{3'}P_i&=\frac15P_i.
 \end{aligned}}
\]
In particular $P_i(\Pi_3-\Pi_{3'})P_i=0$. For an auxiliary normalized
state supported in $V_i$, the corresponding geometric projective measurement
has the universal four-outcome probabilities
\[
 \boxed{(p_4,p_5,p_3,p_{3'})=(4/15,1/3,1/5,1/5).}
\]
These are exact finite sector weights, not physical hadron branching ratios.
\end{theorem}
\begin{proof}
The root stabilizer $A_{4,i}$ acts irreducibly on $V_i$, and each $A_5$
projector commutes with it. Thus by real Schur's lemma each compression
$P_i\Pi_rP_i=c_rP_i$. The five mutually orthogonal $P_i$ sum to the
identity and are permuted transitively by $A_5$. Hence
$\sum_i\operatorname{Tr}(P_i\Pi_r)=\operatorname{rank}\Pi_r$ and
$3c_r=\operatorname{rank}\Pi_r/5$, yielding
$c_r=\operatorname{rank}\Pi_r/15$. The stated values and probability law
follow immediately. An independent exact-coordinate verification appears
in Paper~V-D's appended certificate.
\end{proof}

\begin{corollary}[Mirror-odd one-step quadratic source no-go]
\label{p6:cor:v321-local-mirror-odd-no-go}
For every self-adjoint $A_5$-invariant comparison operator
$B=b_4\Pi_4+b_5\Pi_5+b_3\Pi_3+b_{3'}\Pi_{3'}$,
\[
 \boxed{P_iBP_i=\frac{4b_4+5b_5+3b_3+3b_{3'}}{15}P_i.}
\]
The antisymmetric mirror coefficient $b_3-b_{3'}$ cannot be detected by a
single quadratic expectation supported entirely in the local triple.
External-channel coherence, a source-marked operator, or higher-order
interaction is required to distinguish that coefficient. The theorem does
not prohibit such physical interactions and supplies no term in the current
physical mass operator.
\end{corollary}

\begin{proposition}[Conjugation-sheet refinement of the rooted icosian complement]
\label{p6:prop:v321-petersen-order-sheets}
On the projective quotient $A_5=2I/\{\pm1\}$, fix
$A_4=2T/\{\pm1\}$. Under \emph{conjugation by this $A_4$}, the
48-element complement $A_5\setminus A_4$ splits into four regular
$A_4$ orbits of size twelve:
\[
 \boxed{48=12_{\rm ord\,2}+12_{\rm ord\,3}
       +12_{\rm ord\,5,+}+12_{\rm ord\,5,-}.}
\]
The twelve Petersen edges not omitting the rooted label are canonically
identified, as an $A_4$-set, with the order-two orbit by
$(ab\mid cd)\mapsto(ab)(cd)$. This conjugation action is distinct from
the \emph{left-multiplication} representation used in Paper~I-A's original
four-regular-orbit comparison. The two actions have equal regular
characters but no equality of their preferred bases or physical source
intertwiners follows.
\end{proposition}
\begin{proof}
In $A_5$ there are fifteen involutions, twenty elements of order three,
and two twelve-element classes of order five. The rooted $A_4$ contains
three involutions and eight elements of order three, and no element of
order five. Thus the complement has the displayed order counts. The
centralizer in $A_4$ of an element from any of the four complement subsets
is trivial: a commuting nonidentity power or Klein-four partner cannot
fix the root label. Every conjugation orbit therefore has size $|A_4|=12$.
A Petersen edge is a partition of four labels into two pairs, hence defines
one double transposition; its omission of the rooted label is equivalent
to that double transposition fixing the rooted label. Conjugation preserves
this correspondence.
\end{proof}

\begin{remark}[Terminal mass and source firewall]
The geometric order-sheets, the fixed $4/15,1/3,1/5,1/5$ probabilities,
and the earlier two-/three-core polar amplitudes do not select a terminal
clock or source-visible pole. In particular, $E_{32},E_{36},E_{40}$ and the
current rooted charged and neutral spectrum remain unchanged. A hadronic
mass calculation requires an action-owned color-singlet contraction and its
physical source/receiver residues; a geometric eigenvalue is insufficient.
\end{remark}
% END FRAMEWORK v2.321




% BEGIN FRAMEWORK v2.338: COMPLETE PAIR-INCIDENCE SPECTRUM
\section{Complete Petersen-pair coefficient carriers and marked responses}
\label{p6:sec:v338-pair-incidence}
This section treats finite matrix coefficients of one given $2I$.
Its normalized eigenvalues are incidence strengths, not particle
rest energies.

\begin{theorem}[The rank-25 increment and full center-even closure]
\label{p6:thm:v338-pair-incidence-full-rank}
Let $B_5:\mathbb C^{25}\to\mathbb C^{120}$ be the five-core
permutation coefficient map, and $B_u:\mathbb C^{100}\to
\mathbb C^{120}$ the unsigned two-subset coefficient map:
\[(B_u)_{g,(T,S)}=\mathbf1_{\{\pi(g)S=T\}}.\]
With $P_1,P_4,P_5$ the irreducible projectors of the ten-vertex
Petersen permutation module $\mathbf1\oplus\mathbf4\oplus\mathbf5$,
\[
 B_u^\dagger B_u=120P_1\otimes P_1
              +30P_4\otimes P_4+24P_5\otimes P_5.
\]
Hence $\operatorname{rank}B_5=17$, $\operatorname{rank}B_u=42$,
$\operatorname{im}B_5\subset\operatorname{im}B_u$, and the
orthogonal rank-25 increment has projector
\[
 \Pi_{25}=\frac1{24}B_u(P_5\otimes P_5)B_u^\dagger.
\]
Its character kernel is
\[
 (\Pi_{25})_{g,h}=\frac1{24}\chi_5(\pi(g)^{-1}\pi(h)).
\]
The signed oriented-pair coefficient map $B_o$ has rank 34 and
spectrum $\{40^{(18)},30^{(16)},0^{(86)}\}$ on the full
$120$-carrier.  The two images share the 16-dimensional
four-irrep coefficient block; their sum is exactly all sixty
center-even functions on $2I$.  Their unweighted Gram sum obeys
\[
 K_+=B_uB_u^\dagger+B_oB_o^\dagger
 =120\Pi_1+60\Pi_4+24\Pi_5+40(\Pi_3+\Pi_{3'}),
\]
with full spectrum $\{120^1,60^{16},40^{18},24^{25},0^{60}\}$.
\end{theorem}
\begin{proof}
The unsigned ten-pair representation is $1\oplus4\oplus5$,
while the signed ten-pair representation is $4\oplus3\oplus3'$.
Finite-group matrix-coefficient orthogonality yields the Gram
scalars $|2I|/d$, giving $120,30,24$ and $30,40$, with
degeneracies $d^2$.  The point--pair incidence matrix has
$D^\dagger D=3I_5+J_5$, so the five-core image lies inside
the unsigned pair image.  The projector follows by polar
normalization, and the character kernel by the central idempotent.
Shared irreps account for the 16D intersection; their union carries
all five even irreducible coefficient sectors.  The sum is the
specified counting-normalized positive operator, not a physical
action choice.
\end{proof}

\begin{proposition}[Center-odd completion by spinorial coefficients]
\label{p6:prop:v338-center-odd-completion}
The defining $2I\subset SU(2)$ spinor tensored with the ordinary
pair module has the representation content
\[
 \mathbf2\otimes(\mathbf1\oplus\mathbf4\oplus\mathbf5)
 =\mathbf2\oplus\mathbf2'\oplus\mathbf4_s\oplus2\mathbf6.
\]
Its matrix coefficients span all $2^2+2^2+4^2+6^2=60$
center-odd group functions.  Together with the even coefficient
carriers this gives the exact function-space decomposition
\[
 \mathbb C^{120}=\mathcal H_{17}\oplus\mathcal H_{25}
 \oplus\mathcal H_{18}\oplus\mathcal H_{60}^-,
 \qquad 120=17+25+18+60.
\]
\end{proposition}
\begin{proof}
The stated decompositions are the finite binary-icosahedral
McKay tensor-product rules.  Matrix coefficients of inequivalent
irreducible representations are orthogonal; each irrep of degree
$d$ contributes $d^2$ functions.  The even side is the preceding
theorem, and the odd tensor product contains every center-odd
irreducible type.  This is a representation-theoretic completion,
not an identification with a physical Hilbert space.
\end{proof}

\begin{theorem}[Exact three-stratum response of the new 25-sector]
\label{p6:thm:v338-compressed-occupancy-spectrum}
For one marked pair $S$, let $D_s,D_d,D_o$ be the diagonal
indicators on the center-even carrier of the relation classes of
sizes $6,18,36$ (equivalently $12,36,72$ on $2I$).  Restrict
$R_r=\Pi_{25}D_r\Pi_{25}$ to $\operatorname{im}\Pi_{25}$. Then
\[
 \begin{array}{c|l}
 \text{relation}&\operatorname{Spec}(R_r)\\\hline
 \mathrm{same}&(1/2)^5,0^{20}\\
 \mathrm{disjoint}&(3/4)^4,(2/3)^4,(5/12)^4,(1/6)^1,0^{12}\\
 \mathrm{overlap}&1^8,(3/4)^4,(1/3)^9,(1/4)^4
 \end{array}
\]
The responses sum to $I_{25}$ and their commutator obeys
\[
 \operatorname{SingVal}([R_s,R_d])
   =\{(1/12)^{(8)},0^{(17)}\}.
\]
A vector in the eight-dimensional unit-eigenspace of $R_o$
is actually supported entirely on the overlap stratum; these are
exact source-localized modes in the finite carrier.
\end{theorem}
\begin{proof}
Use the central character kernel from the preceding theorem.
After combining equal central lifts, the integer compression kernel
has four values $5,1,-1,0$ by the relative $A_5$ cycle class.
The characteristic polynomials of the restricted principal
submatrices of that kernel factor over $\mathbb Q$; dividing the
nonzero roots by twelve yields the displayed spectra.  The exact
commutator obeys a cubic polynomial with nonzero singular value
$1/12$, and has rank eight.  If $\Pi_{25}v=v$ and
$\langle v,D_o v\rangle=\|v\|^2$, positivity implies $D_ov=v$.
All identities have rational/integer certificates.
\end{proof}

\begin{theorem}[Ten protected eight-planes reconstruct Petersen adjacency]
\label{p6:thm:v338-eight-plane-fusion}
Let $W_S\subset\operatorname{im}\Pi_{25}$ be the eight-dimensional
unit-eigenspace above, and let $Q_S$ be its orthogonal projector.
Then the fusion-frame operator $T=\sum_{|S|=2}Q_S$ has
\[
 \operatorname{Spec}(T|_{\mathcal H_{25}})
 =\{0^{(1)},2^{(10)},(5/2)^{(4)},(10/3)^{(6)},(15/2)^{(4)}\}.
\]
In particular its span has dimension 24 and the one excluded
invariant direction is the class character $\chi_5$.  The
$10\times10$ overlap matrix $\Gamma_{S,T}=\operatorname{tr}(Q_SQ_T)$
recovers Petersen adjacency exactly:
\[
 18\Gamma=89I_{10}+A_{\mathrm P}+55J_{10}.
\]
Adjacent vertices have three common subspace dimensions;
nonadjacent distinct vertices have two.  Therefore the original
Petersen quadrupole projector is also recovered from $\Gamma$:
\[
 P_{\mathrm{Hopf},5}=6\Gamma-29I_{10}-\tfrac{37}{2}J_{10}.
\]
\end{theorem}
\begin{proof}
The explicitly normalized projector onto each eight-plane is
obtained by the character kernel restricted to the overlap stratum.
Its integer numerator has square equal to $36$ times itself,
so $Q_S$ is exactly idempotent.  The group average of the ten
projectors is $A_5$-equivariant.  Character block reduction gives
the five displayed rational eigenvalues with their multiplicities;
the trace is $80=10\cdot8$.  Their Gram traces depend only on whether
$S=T$, $S\cap T=\varnothing$ or $|S\cap T|=1$, giving the
matrix identity.  Substitution into the earlier Petersen rank-five
idempotent yields the last equation.  Exact finite certificates
check every pair of markings and the characteristic polynomial.
\end{proof}

\begin{remark}[The complete symmetry-preserving response is not a hadron spectrum]
The five compressed $A_5$ conjugacy-class multipliers generate a
nine-dimensional complex algebra, isomorphic to
$\mathbb C^5\oplus M_2(\mathbb C)$, on the 25-dimensional
coefficient sector when group inversion is respected.  A Hermitian
operator from this \emph{unmarked} algebra has at most seven distinct
eigenlevels, generically with multiplicities $1,3,3,4,4,5,5$.
Its finite ranks and the number $25$ do not count 25 particles.
All occupancies and couplings remain geometric unless a same-action
physical source and rest-energy readout are proved.
\end{remark}
% END FRAMEWORK v2.338


% BEGIN FRAMEWORK v2.339: 15-MODE PETERSEN OPERATOR VISIBILITY
\begin{corollary}[The Petersen matrix modes are visible in the old rank-seventeen carrier]
\label{p6:cor:v339-petersen-15-incidence-visibility}
In the five-core coefficient carrier $\mathcal H_{17}=\operatorname{im}B_5$
of Theorem~\ref{p6:thm:v338-pair-incidence-full-rank},
\[
 \mathcal H_{17}=\mathbb C\mathbf1\oplus
 \mathbb C\chi_4\oplus \operatorname{im}\Pi_{15},
 \qquad 17=1+1+15.
\]
The last term is the complexification of the explicit real
oriented-edge operator image
$\mathfrak{sl}(\mathbf1^\perp\subset\mathbb R^5)$ of
Theorem~\ref{p1:thm:v339-petersen-sl4-isometry}.
It has orthogonal character projector
\[
 \Pi_{15}(g,h)=\frac{4\chi_4(\pi(g)^{-1}\pi(h))
 -\chi_4(\pi(g))\chi_4(\pi(h))}{120}
\]
whose central, double-transposition, three-cycle and five-cycle
source strata contribute $0,4,5,6$ to the total trace $15$.
It is orthogonal to the \emph{different} rank-twenty-five
increment $\operatorname{im}\Pi_{25}$ of the preceding theorem.
None of these finite projector leverages supplies a hadron mass,
a Hamiltonian, a marked scalar coupling or a measurement probability
without an independently owned physical action and readout.
\end{corollary}
% END FRAMEWORK v2.339


% BEGIN FRAMEWORK v2.340: FINITE MASS-SOURCE NONIDENTIFICATION
\section{Distinct 16- and 25-blocks in the terminal-source carrier}
\label{p5b:sec:v340-ranktwo-terminal-firewall}
\begin{corollary}[Central character sectors do not determine a pole]
\label{p5b:cor:v340-25-16-terminal-firewall}
In the centre-even group-function carrier induced by a marked
$2I$, the two orthogonal central blocks identified by the
square--hexagon and hexagon--decagon differences are
\[
 \operatorname{im}\mathscr P_5\cong\mathbf5\boxtimes\mathbf5^*,
 \quad\dim=25,\qquad
 \operatorname{im}\mathscr P_4\cong\mathbf4\boxtimes\mathbf4^*,
 \quad\dim=16,\qquad \mathscr P_5\mathscr P_4=0.
\]
The original transporter functions occupy
$\mathbf1\oplus\operatorname{im}\mathscr P_4$ of rank $17$,
where the latter splits under diagonal conjugation into the
one scalar and the traceless fifteen-mode Petersen matrix part.
No numerical correspondence between the $25$ register labels,
the $25$ shell-kernel modes, the $15$ real edge cochains, and
any physical multiplet follows. An actual terminal mass or hadron
pole requires the specified same-owner action/Hessian, its
normalization and an observable coupling; none is supplied by the
central character idempotents alone.
\end{corollary}
\begin{proof}
The ranks and orthogonality follow from
Theorem~\ref{p1:thm:v340-incidence-character-projectors} and
Proposition~\ref{p1:prop:v340-transporter-hexagon-intersection}.
The traceless fifteen-mode splitting is
Proposition~\ref{p1:prop:v339-incidence-rank15-contrast}.
A character projector specifies a representation subspace and not
the spectrum of a separately realized dynamical generator.
\end{proof}
% END FRAMEWORK v2.340

\printbibliography

% BEGIN FRAMEWORK v2.341: 10-WALK-CHANNEL NONIDENTIFICATION
\begin{remark}[Ten nearest-shell kernel channels are not ten axes]
\label{p5b:rem:v341-kernel-petersen-nonidentification}
Although the $2I$ nearest-shell kernel has dimension $25$ and
its common kernel with the selected $2T$ shell has dimension
$15$, the remaining ten-dimensional module is $V_5\oplus V_5$
under the regular action commuting with left convolution.
It is not the Petersen vertex module
$\mathbf1\oplus\mathbf4\oplus\mathbf5$, nor is it the
ten image channels of the symmetric transporter indicator map.
The transporter coefficients live in the trivial and
four-dimensional matrix-coefficient blocks and are orthogonal to
the nearest-shell five-dimensional coefficient block.
This representation mismatch prevents extracting a new hadron
multiplicity, terminal pole or mass rule from $25-15=10$.
Any such physical inference still requires a realized action,
its Hessian/normalization and a measured observable.  See
Proposition~\ref{p1:prop:v341-walk-kernel-no-petersen-isomorphism}.
\end{remark}
% END FRAMEWORK v2.341


% BEGIN FRAMEWORK v2.342: CLASS-TOMOGRAPHY TERMINAL FIREWALL
\begin{corollary}[Class-frame completeness is not a terminal mass determination]
\label{p5b:cor:v342-class-tomography-no-mass}
Inside the existing rank-$15$ traceless coefficient component
of the $17$-dimensional five-core transporter image, the
three-cycle conjugacy class alone supports a complete algebraic
frame with sector eigenvalues $(2/3,1/15,1/3,1/3)$ on the
$\mathbf4\oplus\mathbf5\oplus\mathbf3\oplus\mathbf3'$
operator modes.  The two five-cycle classes carry the exact
$\varphi^6$ polarization between the two three-dimensional
skew sectors but vanish on the four-dimensional one.
The three-cycle frame reconstructs finite matrices, not a
physical rest-energy operator or interaction kernel.  In
particular the eigenvalues $a,b$, the frame condition ratio
$10$, and the minimal three-generator count are \emph{not}
hadron masses, mass ratios, coupling constants or a replacement
running law.  A terminal pole claim still requires the
same-owner action Hessian, normalizations, spectral source and
an observable whose response distinguishes the relevant
class-filtered modes.
\end{corollary}
\begin{proof}
Apply Theorem~\ref{p1:thm:v342-class-frame-three-cycle-tomography}
to the rank-fifteen subspace of
Corollary~\ref{p6:cor:v339-petersen-15-incidence-visibility};
its operator metric is the exact finite Hilbert--Schmidt metric.
The necessary independent terminal-source conditions are those
stated in Corollary~\ref{p5b:cor:v340-25-16-terminal-firewall}.
\end{proof}
% END FRAMEWORK v2.342


% BEGIN FRAMEWORK v2.343: RECIPROCAL CASIMIR TERMINAL FIREWALL
\begin{corollary}[Signed reciprocal multiplicities do not determine masses]
\label{p5b:cor:v343-reciprocal-casimir-no-mass}
The exact selected carrier and quadratic operator of
Theorems~\ref{p1b:thm:v343-canonical-selected15}
and~\ref{p1b:thm:v343-s5-twist-so6} furnish the finite
multiplicities $105=15+90$, $15=6+9$, and the
\emph{dimensionless} auxiliary Casimir spectrum
\[
 \operatorname{Spec}\Delta_{\mathrm{rec}}
   =\{2^{(6)},4^{(9)},0^{(90)}\}.
\]
They do not select a physical mass operator or its energy scale,
a hadron rest-mass ratio $1:2$, a new $SO(6)$ gauge sector,
chirality breaking, a terminal pole residue or a replacement
for QCD running.  Any terminal claim must first produce a
same-owner action/response operator and a measurable readout
that distinguishes the signed selected sectors, beyond the
conditional effect algebra in
Corollary~\ref{p2e:cor:v343-signed-correlation-export}.
\end{corollary}
\begin{proof}
The abstract $\mathfrak{so}(6)$ adjoint Casimir normalization
and representation multiplicities are the finite identities
in Paper I-C.  Multiplying this operator by an arbitrary
positive scale would change every dimensionful eigenvalue
without altering the mathematical theorems.  The independent
GP2 source, effect, and terminal-pole requirements remain those
of Corollary~\ref{p5b:cor:v340-25-16-terminal-firewall}.
\end{proof}
% END FRAMEWORK v2.343


% BEGIN FRAMEWORK v2.344: TWISTED QUINTET TERMINAL SCOPE
\begin{corollary}[A full geometric quintet does not fix terminal hadron masses]
\label{p5b:cor:v344-twisted-quintet-no-mass}
The character-twisted rank-$25$ coefficient construction of
Theorems~\ref{p1:thm:v344-induced-hurwitz-quintet}
and~\ref{p1:thm:v344-twisted-25-gram} is an exact comparison
module and repairs the untwisted rank-$17$ deficiency.
It cannot by itself label a physical $s$-quark pair, prescribe
normal/inverse radical constituent factors, or assign a terminal
pole correction $-E_{36}/25$.  Even an observed nonzero
finite-flavor curvature
$M_{qqQ}-2M_{qsQ}+M_{ssQ}$ requires an
actual same-action flavor-resolved Hamiltonian or correlator;
identical energy shifts to all three members cancel.
\end{corollary}
\begin{proof}
The twisted quintet belongs to the purely mathematical
$2T\subset2I$ comparison bank, whereas the terminal physical
mass-source and the physical representation occurrence are
independent constitutional inputs.  The rank and transfer
identities fix dimensionless numbers only.  No physical flavor
selector, two-source coupling amplitude or pole resolvent is
specified by them, so distinct terminal spectra remain compatible
with all the displayed finite identities.
\end{proof}
% END FRAMEWORK v2.344


% BEGIN FRAMEWORK v2.348: TERMINAL OBSERVABLE AND LABEL TYPE BOUNDARIES
\clearpage
\section{Terminal observables: incidence descent and neutral-pole bridge}
\label{p5b:sec:v348-terminal-observable-tests}
\begin{corollary}[The label-space exclusion is not a vertex observable by descent]
\label{p5b:cor:v348-label-descent-firewall}
Let $B:\mathbb C^{25}\to\mathbb C^{120}$ denote the untwisted
$2T\subset2I$ cell incidence map and let
$Q_H=I_{25}-|e_H\rangle\langle e_H|$ be single-cell label exclusion.
Although $\langle\Omega_{25},Q_H\Omega_{25}\rangle=24/25$,
there is no linear $T$ on the vertex carrier satisfying $TB=BQ_H$.
Hence any physical consumption of the $24/25$ label coefficient must
identify a typed label-space source and its same-action readout, or
construct another explicitly licensed descendant; incidence alone is
insufficient.
\end{corollary}
\begin{proof}
Proposition~\ref{p6:prop:v179-label-exclusion-nondescent} gives
$k=P_{\ker B}e_H\in\ker B$ and
$BQ_Hk=-(8/25)Be_H\ne0$.  If $TB=BQ_H$, then
$0=TBk=BQ_Hk$, a contradiction.  The expectation $24/25$ is
an exact fact in the declared orthogonal label host, but does not
supply the forbidden vertex intertwiner.
\end{proof}

\begin{proposition}[What the $K=56$ role-gap ratio does and does not predict]
\label{p5b:prop:v348-n56-physical-gap-bridge}
Take the generated neutral role-energy triple of
Corollary~\ref{p6:cor:v587-terminal-role-energy-spectrum},
$E_{56}(x^{-1},x,x^2)$ with $x^3=192$.  Suppose a further
same-owner physical pole theorem establishes a \emph{common positive
scale} $s$ and a compatible normal assignment
$(m_1,m_2,m_3)=s(x^{-1},x,x^2)$.  Then
\[
 \boxed{\frac{m_3^2-m_1^2}{m_2^2-m_1^2}
   =\frac{x^6-1}{x^4-1}
   =x^2+\frac1{x^2+1}
   =33.3105115561\ldots .}
\]
The finite role-energy spectrum alone does \emph{not} prove this
physical mass-splitting relation: the monotone alternative
$m_i=s\,r_i^2$, where $(r_1,r_2,r_3)=(x^{-1},x,x^2)$,
preserves the ordered geometric roles but yields instead
$(x^{12}-1)/(x^8-1)$.  Thus the mass-pole/readout bridge is an
independent mathematical and physical obligation, not a free
neutrino-specific normalization parameter.
\end{proposition}
\begin{proof}
With the common scale the numerator equals
$s^2(x^4-x^{-2})$ and the denominator equals
$s^2(x^2-x^{-2})$.  Multiplication by $x^2$ gives the first
quotient; factoring $x^6-1$ yields the displayed second form.
For the alternative $r_i^2$ assignment the corresponding square
differences are $s^2(x^8-x^{-4})$ and
$s^2(x^4-x^{-4})$, producing $(x^{12}-1)/(x^8-1)$.
Both assignments preserve the geometric role input, establishing
non-implication without an action-selected pole map.
\end{proof}
% END FRAMEWORK v2.348


% BEGIN FRAMEWORK v2.365: TYPED HURWITZ MASS-GRAMMAR INTERFACE
\subsection{Hurwitz source partitions are not terminal mass additions}
\label{p5b:subsec:v365-hurwitz-terminal-scope}
The rooted Hurwitz heavy support $24=8+16=8+(7+9)$, the
chamber-relative dual $24=6+12+6$, the $2T/Q_8$ coset partition
$24=8+8+8$, and the six-point light-root refinement
$6=1+5=1+(3+2)$ label different owner spaces and geometric
operations.  Their simultaneous validity does not provide a
vertex-to-vertex equality of physical Hamiltonians, a hadron
mass-addition prescription, or an alternative charged $[4]$ map.
The exact terminal coefficient grammar and, in particular,
$E_{36}^{2}=E_{32}E_{40}$ and the established $W,Z,t,b,\tau,
\mu,c,s,e,u,d$ readout conventions are unchanged.  New finite
96-face source identities are located in
Theorem~\ref{p1:thm:v365-hurwitz-face-laplacian}; their physical
observable use requires the same-action map stated in Paper VII.
% END FRAMEWORK v2.365


% BEGIN FRAMEWORK v2.366: CANONICAL DOUBLE-ROOT HURWITZ COMPARISON
\begin{remark}[Source-matched root--facet refinement without mass alteration]
\label{p5b:rem:v366-hurwitz-chamber-refinement}
In the exactly marked transporter 24-cell of
Theorem~\ref{p1:thm:v366-dual-hurwitz-refinement}, a dual chamber
normal selects $6_++12_0+6_-$, while a separate root selects its
nearest $8$ and complementary $7+9$.  Their intersection counts
are $(2+2t,4,2-2t)$ for the rooted nearest eight and
$(4-2t,8,4+2t)$ for its complementary 16, with
$t\in\{-1,0,1\}$ the root--facet incidence.
The local $Q_8\sqcup aQ_8\sqcup a^2Q_8$ is another $8+8+8$
\emph{group coset} partition, not the nearest-eight radial layer.
These are simultaneous exact support comparisons; the original
charged terminal, tau and $E_{32}/E_{36}/E_{40}$ readouts retain
exactly their previous normalization and scope.  No hadron mass or
decay width is inferred from this intersection table.
\end{remark}
% END FRAMEWORK v2.366


% BEGIN PAPER V-B v5.200: FULL 2O ELEMENTARY-PARTICLE CARRIER ROLE
\section{Complete octahedral elementary-particle response: the working realization}
\label{p5b:sec:v5200-full-octahedral-particle-role}
\begin{definition}[Working complete $2O$ particle-response assignment]
\label{p5b:def:v5200-full-octahedral-particle-role}
Retain the already derived round momentum carrier and its local
Hurwitz core $H=2T$.  For the working \emph{physical interpretation}
of elementary-particle response, use the \emph{complete} octahedral
normalizer $\mathsf C_{\rm part}^{\rm work}=2O=N_{\operatorname{Sp}(1)}(H)$.
Its intrinsic finite support grammar is
\[
 \boxed{2O=\underbrace{8+16}_{2T,\;E_{32}\text{ anchor}}
          \;\sqcup\;
          \underbrace{6+12+6}_{2O\setminus2T,\;E_{36}\text{ response}}.}
\]
The attached sixfold interface is $C_6(u)\subset H$ together with
its six-point normalizer coset $B_6(u)\subset2O\setminus H$;
$E_3(u)=\mathbb R1\oplus\operatorname{span}B_6(u)$ is the previously
constituted $E_{40}$ active Clifford carrier.  The sixfold attachment
is a structure inside the same $48$ elements, not another particle
set and not a replacement by $2I$.  The displayed energy names are
\emph{source-owner anchors}, not one mass eigenvalue per vertex.
\end{definition}
\begin{remark}[Physical claim frontier and the legacy antipodal map]
\label{p5b:rem:v5200-full-octahedral-physical-gate}
The mathematical normalizer, quartic response, complementary split
and sixfold maximizers are established by
the Born-even quartic theorem in Paper~I-B,
the sixfold-interface proposition in Paper~I-B
and the balanced-source selector theorem in Paper~III-C.
The physical equality ``elementary-particle response carrier $=2O$''
remains the \emph{working realization} of
the typed particle-response definition of Paper~II-A, not a
new consequence of G0--GP3 without a licensed source-occurrence
certificate.  In particular the historical conditional antipodal
$X_+,X_0,X_-$ clock-to-role assignment and the present full-coset
$E_{36}$ source support are differently typed interpretations;
no numerical clock or previous physical source member is
silently overwritten.  The native carrier grading has energies
$\{E_0,E_{32},E_{40},E_{56}\}$ and
$E_{36}^2=E_{32}E_{40}$ is its balanced \emph{transition} weight,
not an extra active eigenvalue.  The prior $u,d$ light-stage
constructors and charged terminal coefficients remain unchanged.
\end{remark}
% END PAPER V-B v5.200

\end{document}
